\documentclass[11pt,a4paper,oneside,openany]{book}
\usepackage{fontspec}
\setmainfont{Noto Serif CJK SC}[BoldFont={Noto Serif CJK SC Bold}]
\setsansfont{Noto Sans CJK SC}[BoldFont={Noto Sans CJK SC Bold}]
\setmonofont{DejaVu Sans Mono}[Scale=0.85]
\XeTeXlinebreaklocale "zh"
\XeTeXlinebreakskip=0pt plus 1pt minus .1pt
\usepackage[margin=21mm,top=22mm,bottom=23mm,headheight=16pt,headsep=7mm,footskip=11mm]{geometry}
\usepackage{amsmath,amssymb,mathtools,bm}
\usepackage{booktabs,longtable,tabularx,array,multirow}
\usepackage{graphicx,float,caption,subcaption}
\usepackage{enumitem,listings,etoolbox}
\usepackage[most]{tcolorbox}
\usepackage{xcolor,setspace,titlesec,fancyhdr,needspace}
\definecolor{ink}{HTML}{152E45}
\definecolor{teal}{HTML}{087F8C}
\definecolor{rust}{HTML}{A24734}
\definecolor{softblue}{HTML}{F2F6FA}
\definecolor{softteal}{HTML}{F0F8F8}
\definecolor{softwarm}{HTML}{FBF5F0}
\usepackage[unicode,colorlinks=true,linkcolor=ink,citecolor=teal,urlcolor=teal,bookmarksnumbered=true]{hyperref}
\usepackage{bookmark}
\hypersetup{pdftitle={机械臂控制与仿真详细教材},pdfauthor={整合编写},pdfsubject={数学推导 数值算例 图解 七轴柔顺控制},pdfkeywords={机器人,机械臂,控制,辨识,仿真,阻抗,零空间}}
\renewcommand{\contentsname}{目录}
\renewcommand{\listfigurename}{图目录}
\renewcommand{\listtablename}{表目录}
\renewcommand{\figurename}{图}
\renewcommand{\tablename}{表}
\renewcommand{\bibname}{参考资料}
\renewcommand{\chaptername}{章}
\renewcommand{\appendixname}{附录}
\setcounter{tocdepth}{1}
\setcounter{secnumdepth}{2}
\makeatletter
\renewcommand{\@pnumwidth}{2.5em}
\renewcommand{\@tocrmarg}{3.5em}
\renewcommand*\l@section{\@dottedtocline{1}{1.5em}{3.4em}}
\renewcommand*\l@figure{\@dottedtocline{1}{1.5em}{3.4em}}
\makeatother
\setstretch{1.25}
\setlength{\parindent}{2em}
\setlength{\parskip}{3pt plus 1pt}
\setlength{\emergencystretch}{2em}
\renewcommand{\arraystretch}{1.35}
\setlength{\tabcolsep}{4pt}
\setlist{itemsep=3pt,topsep=5pt,parsep=1pt,leftmargin=2em}
\setlist[description]{font=\bfseries}
\newcolumntype{Y}{>{\raggedright\arraybackslash}X}
\titleformat{\chapter}[display]{\sffamily\bfseries\color{black}}{\large 第\thechapter 章}{7pt}{\LARGE}
\titlespacing*{\chapter}{0pt}{8pt}{20pt}
\titleformat{\section}{\large\sffamily\bfseries\color{ink}}{\thesection}{.7em}{\phantomsection}
\titleformat{\subsection}{\normalsize\sffamily\bfseries}{\thesubsection}{.65em}{\phantomsection}
\titlespacing*{\section}{0pt}{16pt}{7pt}
\titlespacing*{\subsection}{0pt}{11pt}{5pt}
\pagestyle{fancy}
\fancyhf{}
\fancyhead[L]{\small\sffamily\nouppercase{\leftmark}}
\fancyfoot[C]{\small\thepage}
\renewcommand{\headrulewidth}{.35pt}
\renewcommand{\chaptermark}[1]{\markboth{第\thechapter 章\quad #1}{}}
\fancypagestyle{plain}{\fancyhf{}\fancyfoot[C]{\small\thepage}\renewcommand{\headrulewidth}{0pt}}
\captionsetup{font=small,labelfont=bf,skip=7pt,justification=raggedright,singlelinecheck=false}
\tcbset{enhanced,breakable,boxrule=.5pt,arc=1mm,left=8pt,right=8pt,top=7pt,bottom=7pt,before skip=10pt,after skip=10pt,fonttitle=\sffamily\bfseries,coltitle=ink}
\newtcolorbox{examplebox}[1][]{colback=softblue,colframe=ink!40,title={#1},colbacktitle=ink!7}
\newtcolorbox{conceptbox}[1][]{colback=softteal,colframe=teal!45,title={#1},colbacktitle=teal!10}
\newtcolorbox{pitfallbox}[1][]{colback=softwarm,colframe=rust!45,title={#1},colbacktitle=rust!9}
\lstset{basicstyle=\ttfamily\small,breaklines=true,breakatwhitespace=false,columns=fullflexible,keepspaces=true,showstringspaces=false,frame=single,rulecolor=\color{ink!25},backgroundcolor=\color{softblue},aboveskip=9pt,belowskip=9pt,framesep=7pt,tabsize=4,keywordstyle=\color{teal}\bfseries,commentstyle=\color{GREY}}
\definecolor{GREY}{HTML}{53616D}
\allowdisplaybreaks[2]
\newcommand{\R}{\mathbb R}
\newcommand{\diag}{\operatorname{diag}}
\begin{document}
\frontmatter
\begin{titlepage}
\centering
\vspace*{14mm}
{\sffamily\small 系统学习教材\par}
\vspace{13mm}
{\sffamily\bfseries\fontsize{29}{40}\selectfont 机械臂控制与仿真\par}
\vspace{3mm}
{\sffamily\bfseries\fontsize{25}{35}\selectfont 详细教材\par}
\vspace{11mm}
{\sffamily\Large 从数学推导到七轴柔顺控制\par}
\vspace{7mm}
{\large 概念解释\quad 完整数值算例\quad 公式推导\quad 精确图解\par}
\vspace{17mm}
\includegraphics[width=.86\linewidth]{figures/rt_cartesian.pdf}
\vfill
\begin{minipage}{.84\linewidth}
\normalsize
以单关节和两连杆为计算主线，连接七轴机械臂的参数辨识、低阻抗控制、零空间、示教回放与相机云台。

本书根据两份上传教材重新整合编写，统一符号与假设，补充中间计算、算法边界和可复现图形。
\end{minipage}
\vspace{13mm}

{\small 整合与扩展版\quad 2026年11月23日\par}
\end{titlepage}

\chapter*{本书的目的与阅读方法}
\addcontentsline{toc}{chapter}{本书的目的与阅读方法}

这本书面向已经接触机械臂软件、希望把公式和工程现象连起来的读者。学习目标是：看到机械臂的位置、速度、力矩、外力和误差时，知道这些量分别来自哪里，能把一条控制公式拆成可解释、可计算、可检验的步骤。

两份原稿包含很广的知识，但数学、控制和工程内容反复交叉，部分独立公式仍显示为源码，部分图只有标题，某些结论缺少成立条件。新版本重新组织知识顺序，并对符号、数值和前提进行统一；它不是把两个 PDF 顺次拼接，也不是删减成知识点摘要。

\section*{原始材料和引用约定}

\begin{description}[leftmargin=2.6em,style=nextline]
\item[原稿 A]《机械臂与机器人控制算法 从零开始的基础巩固手册》，上传文件共90页。侧重数学基础、经典与现代控制、机器人控制和工程能力体系。
\item[原稿 B]《机械臂与仿真工程教材 数学及格也能懂版》，上传文件共129页。侧重运动学、辨识、摩擦、仿真、示教、遥操作和七轴云台。
\end{description}

每章末标注原稿对应页码，页码均指上传 PDF 的物理页序。涉及补充的一手软件文档和论文，采用参考文献编号。网页资料核对日期为2026年11月23日；软件示例应对应实际安装版本。原稿提到但本次没有提供完整内容的其它讲义或教程，不会被当成已经通读的独立来源。

\section*{怎样读一个公式}

每遇到一个公式，依次问五件事：它在解决什么问题；每个符号是什么量、在哪个坐标系、什么单位；等号两边的维度是否相同；它由哪一步物理关系或数学运算得到；它在什么前提下可用。只记最后一行，通常不足以排查真实代码。

例如 $\tau=J^\mathsf TF$ 并不只是“雅可比转置乘力”。必须知道 $J$ 映射哪个参考点的速度，$F$ 是哪一方作用于哪一方，力和速度是否在同一个坐标系，以及六维向量采用哪种行顺序。正确的单位、坐标与功率关系，可以比盲调参数更早发现错误。

\section*{本书每章的组织方式}

正文先说明问题和直觉，再引入概念与符号，随后给出推导与逐步计算，图形用来显示公式的结果或边界，最后用计算题和解答检查掌握程度。解答紧跟本章，便于自学；建议先遮住解答手算，再逐步对照。

全书包括23章正文、4个附录、57幅图、114个框内算例，以及83道附逐步解答的章末计算练习；正文中还包含没有单独入框的数值演算。目录和图目录可点击跳转，文末列出原稿与补充的一手参考资料。

所有新生成的曲线均来自明确方程或给定数据。示意几何和合成仿真会在图注中说明，不会写成真机结果。章节中出现的教学刚度、阻尼、采样周期和质量，只用于解释计算过程；迁移到 NERO 需要代入实际模型、接口和测量。

\section*{知识地图}

\begin{longtable}{p{24mm}p{44mm}p{86mm}}
\toprule 章节 & 主要问题 & 学完后应能亲手完成的计算\\\midrule
\endhead
第1--2章 & 数学工具 & 解线性方程，解释秩与条件数，求导积分，解简单动态方程，理解变换与离散化\\
第3--5章 & 位姿与运动学 & 坐标转换，旋转误差，2R正逆解，雅可比、奇异与速度零空间\\
第6--7章 & 动力学与执行器 & 从能量推导力矩，理解惯量耦合、电机、传动、柔性和接触\\
第8--10章 & 参数辨识与标定 & 建回归矩阵，算加权最小二乘，设计激励，分析摩擦、负载与手眼标定\\
第11--14章 & 控制与估计 & 设计PD和PID，分析频域，算LQR和离散模型，逐步更新滤波器与观测器\\
第15--16章 & 力控与七轴柔顺 & 推导阻抗、导纳和操作空间方程，检查零空间力矩何时不干扰主任务\\
第17--19章 & 运动与人机输入 & 解五次轨迹与S曲线，处理示教，计算DMP、ProMP、遥操作和视觉映射\\
第20--21章 & 工程和仿真 & 对齐接口与时序，估总线带宽，建立多速率仿真，区分模型真值与测量\\
第22--23章 & 约束和综合应用 & 手算MPC与任务权重，理解WBC，把整条链迁移到七轴云台\\
\bottomrule
\end{longtable}

\section*{按当前问题选择入口}

若当前最困惑的是“参数辨识到底在算什么”，先读第1章线性代数与最小二乘，再读第6、8、9章。若重点是“笛卡尔拖动和肘部都要软”，先读第3、5、6章，再读第15、16章。若要做运镜与示教，连接第17--21章及第23章。跳读时遇到尚不理解的符号，可先查附录的术语表，再回到首次推导的章节。

不按固定周数承诺学习速度。一个可靠的前进标准是：能不看答案重算本章的关键数字，并说清楚改一个参数后，图形会怎样变化。对于本书的二维或一维例子，这比单纯背诵算法名称更有帮助。

\chapter*{全书符号和三个贯穿模型}
\addcontentsline{toc}{chapter}{全书符号和三个贯穿模型}

\section*{全书默认约定}

关节角用弧度，时间用秒，长度用米，力用牛顿，转动力矩用牛米。角度虽然在 SI 量纲中可视为无量纲，工程表达仍保留 rad，以便区分角位移、线位移及对应增益。

\begin{longtable}{p{33mm}p{79mm}p{42mm}}
\toprule 符号 & 含义 & 常见单位或维度\\\midrule\endhead
$q,\dot q,\ddot q$ & 关节位置、速度、加速度列向量 & rad，rad/s，rad/s$^2$\\
$q_d$ 与 $e=q-q_d$ & 期望关节角与实际减期望误差 & rad\\
$\tau_{\rm cmd},\tau_{\rm ext}$ & 执行器关节侧命令，环境对机器人的外力矩 & N m\\
$M(q)$ & 关节空间惯量矩阵 & kg m$^2$，$n\times n$\\
$C(q,\dot q)\dot q$ & 科氏与离心合成项 & N m\\
$g(q),\tau_f$ & 重力项与左端约定的摩擦项 & N m\\
$p,R,T$ & 位置、旋转矩阵、齐次变换 & $3\times1$，$3\times3$，$4\times4$\\
$R_{AB}$ & 把B系向量分量转换到A系 & $v_A=R_{AB}v_B$\\
$V=(v^\mathsf T,\omega^\mathsf T)^\mathsf T$ & 线速度在前，角速度在后 & m/s，rad/s\\
$F=(f^\mathsf T,m^\mathsf T)^\mathsf T$ & 力在前，力矩在后；与$V$同点同轴 & N，N m\\
$J$ & 对应任务的几何或明确指定的误差雅可比 & 行数由任务决定\\
$H,\beta,y$ & 堆叠回归矩阵、参数、测量目标 & 由回归定义\\
$T_s,h$ & 控制采样周期、物理积分步长 & s\\
$K,D$ & 刚度与阻尼；按关节或任务空间说明 & 由空间和物理量决定\\
\bottomrule
\end{longtable}

统一的动力学符号为
\begin{equation*}
M(q)\ddot q+C(q,\dot q)\dot q+g(q)+\tau_f
=\tau_{\rm cmd}+\tau_{\rm ext}.
\end{equation*}
因此，采用 $e=q-q_d$ 的恢复反馈力矩一般带负号。个别小节为了连接标准优化或经典控制推导采用参考减输出，会在当地明确定义，不能跨小节直接移植符号。

\section*{模型一 单关节}

惯量 $J_r=0.1\,\mathrm{kg\,m^2}$，黏性系数 $b=0.2\,\mathrm{N\,m\,s/rad}$，基础方程 $J_r\ddot q+b\dot q=\tau$。它用于把控制、滤波和离散化算清楚。$J_r$ 是单轴转动惯量，与几何雅可比 $J(q)$ 不同。

\section*{模型二 竖直平面两连杆}

$l_1=l_2=0.5\,\mathrm m$，$m_1=m_2=1\,\mathrm{kg}$，$c_1=c_2=0.25\,\mathrm m$，绕运动平面法向的质心惯量 $I_1=I_2=1/48\,\mathrm{kg\,m^2}$。重力沿 $-y$，$q_1$ 从水平正 $x$ 轴起算，$q_2$ 是第二杆相对第一杆的角。少数独立算例采用其它数值时会明确重列参数。

\section*{模型三 七轴机械臂与相机任务}

以固定基座七轴机械臂作应用背景，但不虚构真实 DH、URDF 惯量、驱动器力矩限制或厂商内部算法。其主要用途是解释六维任务、冗余、零空间、工具负载和世界到相机的误差链。物理图和低维算例用于说明原理，不替代 NERO 实际模型。

\clearpage
\tableofcontents
\clearpage
\listoffigures
\clearpage

\mainmatter
\chapter{从数与单位到矩阵和最小二乘}\label{ch:math}
\section{先知道每一个数描述什么}
机械臂程序里，一列七个数可能是角度，也可能是电流、力矩或编码器计数。它们在内存中都是浮点数，物理意义却完全不同。本章从能够实际计算的最小单元开始：先把数的含义、单位和形状说清楚，再学习矩阵为什么能够描述运动、误差与辨识。

\subsection{标量 向量 矩阵及其维度}
标量（scalar）是一个数，例如一根连杆的长度 $l=0.5\,\mathrm m$。向量（vector）把一组有顺序的数排成一列，例如两关节位置
\begin{equation}
 q=\begin{bmatrix}q_1\\q_2\end{bmatrix}
 =\begin{bmatrix}0\\\pi/2\end{bmatrix}\mathrm{rad}\in\mathbb R^2.
\end{equation}
$\mathbb R^2$ 读作“二维实数空间”，意思是有两个实数分量。这里的下标是关节编号。向量不是无序的数袋：交换两个分量，会变成另一种机械臂姿态。矩阵（matrix）是一张有行和列的数表。$A\in\mathbb R^{m\times n}$ 表示 $m$ 行、$n$ 列；它把一个 $n$ 维输入变为 $m$ 维输出。

“维度”有两种常见意思，必须分开。数组维度描述有几个分量；物理量纲描述长度、时间、质量怎样组合。例如三维位置 $p$ 的数组维度为 3，三个分量的单位都为米。三维角速度 $\omega$ 也有 3 个分量，但单位为弧度每秒。两者不能因为形状相同就相加。

弧度（radian，rad）定义为弧长除以半径。整圆的弧长为 $2\pi r$，故 $360^\circ=2\pi\,\mathrm{rad}$。数学上弧度可视为无量纲，工程记录里仍应保留 rad，以防与度混淆。
\begin{examplebox}[算例 角度速度和力矩单位]
$30^\circ=30\pi/180=\pi/6\approx0.523599\,\mathrm{rad}$。若关节在 $0.2\,\mathrm s$ 内增加 $3^\circ$，平均速度是
\[
 \bar\omega=\frac{3\pi/180}{0.2}=0.261799\,\mathrm{rad/s}.
\]
长度 $0.5\,\mathrm m$ 的杆端受到与杆垂直的 $4\,\mathrm N$ 力，力矩为 $\tau=lF=2\,\mathrm{N\,m}$。如果把 $0.5\,\mathrm m$ 误写成数值 500，再继续按米计算，结果会错误地放大一千倍。单位错误不会自动引发程序异常。
\end{examplebox}

\subsection{等号 近似号 变化量和导数}
$=$ 表示在给定假设下严格相等，$\approx$ 表示近似。$\Delta q$ 常表示一次有限变化，$\delta q$ 常表示用于线性化的小变化；两者不是新的物理量。$\dot q=dq/dt$ 是角速度，$\ddot q=d^2q/dt^2$ 是角加速度。上标 $\mathsf T$ 表示转置，不是时间；$q^\mathsf T$ 把列向量改成行向量。上标 $-1$ 表示逆运算，只有在相应逆存在时才可以写。

\section{矩阵乘法是一组加权求和}
设
\begin{equation}
 A=\begin{bmatrix}2&1\\-1&3\end{bmatrix},\qquad
 x=\begin{bmatrix}4\\2\end{bmatrix}.
\end{equation}
矩阵乘向量是每一行分别与向量做乘积再求和：
\begin{equation}
 Ax=\begin{bmatrix}2\cdot4+1\cdot2\\-1\cdot4+3\cdot2\end{bmatrix}
 =\begin{bmatrix}10\\2\end{bmatrix}.
\end{equation}
它还可解释为“按输入系数组合矩阵的各列”：$Ax=4(2,-1)^\mathsf T+2(1,3)^\mathsf T$。这个看法随后会解释为什么某些末端方向无法到达，以及为什么辨识数据不足时无法区分某些参数。

矩阵乘法 $AB$ 的规则相同：$A$ 的每一行分别与 $B$ 的每一列相乘再求和。若 $A$ 为 $m\times n$，$B$ 必须有 $n$ 行。结果为 $m\times p$，其中 $p$ 是 $B$ 的列数。检查维度是检查公式的第一步。
\begin{examplebox}[算例 为什么矩阵的先后顺序不能交换]
取
\[
 A=\begin{bmatrix}1&1\\0&1\end{bmatrix},\quad
 B=\begin{bmatrix}2&0\\0&1\end{bmatrix}.
\]
直接计算得到
\[
 AB=\begin{bmatrix}2&1\\0&1\end{bmatrix},\qquad
 BA=\begin{bmatrix}2&2\\0&1\end{bmatrix}.
\]
对 $x=(0,1)^\mathsf T$，$ABx=(1,1)^\mathsf T$，而 $BAx=(2,1)^\mathsf T$。$B$ 先沿水平方向拉伸，$A$ 再把高度加入水平坐标，与先做后一步再拉伸不同。三维旋转的不可交换性也是这个道理。
\end{examplebox}

单位矩阵 $I$ 的对角元素为 1，其余为 0，满足 $Ix=x$。转置满足
\begin{equation}
 (AB)^\mathsf T=B^\mathsf TA^\mathsf T.
\end{equation}
顺序倒过来，可以从第 $i,j$ 元素验证：$[(AB)^\mathsf T]_{ij}=\sum_k A_{jk}B_{ki}$，恰好等于 $[B^\mathsf TA^\mathsf T]_{ij}$。这个顺序规则在后面的功率推导和最小二乘求导中会反复使用。

\section{长度 夹角和加权误差}
向量内积 $x^\mathsf Ty=\sum_i x_i y_i$ 是一个数。欧氏范数（Euclidean norm）也叫二范数：
\begin{equation}
 \|x\|_2=\sqrt{x^\mathsf Tx}=\sqrt{\sum_i x_i^2}.
\end{equation}
例如位置误差 $e_p=(0.003,-0.004,0)^\mathsf T\,\mathrm m$ 的长度为 $0.005\,\mathrm m$，即 $5\,\mathrm{mm}$。两向量夹角 $\gamma$ 满足 $x^\mathsf Ty=\|x\|\|y\|\cos\gamma$；内积为零称为正交（orthogonal），二维、三维中就是垂直。

当误差同时包含米和弧度时，直接计算 $\sqrt{e_x^2+e_y^2+e_z^2+e_{rx}^2+e_{ry}^2+e_{rz}^2}$ 会把两种单位混在一起。必须先声明尺度。例如允许位置误差 $p_s=0.01\,\mathrm m$、角度误差 $\theta_s=0.05\,\mathrm{rad}$，定义无量纲误差
\begin{equation}
 \tilde e=\begin{bmatrix}e_p/p_s\\e_R/\theta_s\end{bmatrix},\qquad
 \|\tilde e\|_2^2=e_p^\mathsf T p_s^{-2}I e_p+e_R^\mathsf T\theta_s^{-2}I e_R.
\end{equation}
若只有 $e_x=0.005\,\mathrm m$ 和 $e_{rz}=0.025\,\mathrm{rad}$，两者各占允许尺度的一半，故 $\|\tilde e\|=\sqrt{0.5^2+0.5^2}=0.7071$。这种缩放没有改变实际误差，却明确了“哪种误差更重要”。后面的 IK 权重、辨识条件数、控制代价都要面对相同问题。

\section{线性方程 逆矩阵 秩与零空间}
\subsection{从两个方程理解逆矩阵}
考虑 $Ax=b$，其中
\begin{equation}
 \begin{bmatrix}2&1\\1&3\end{bmatrix}
 \begin{bmatrix}x_1\\x_2\end{bmatrix}
 =\begin{bmatrix}5\\5\end{bmatrix}.
\end{equation}
第一行是 $2x_1+x_2=5$，第二行是 $x_1+3x_2=5$。从第一行得 $x_2=5-2x_1$，代入第二行得 $x_1+15-6x_1=5$，故 $x_1=2$，$x_2=1$。

对 $A=\begin{bmatrix}a&b\\c&d\end{bmatrix}$，行列式（determinant）是 $ad-bc$。若它不为零，则
\begin{equation}
 A^{-1}=\frac{1}{ad-bc}\begin{bmatrix}d&-b\\-c&a\end{bmatrix}.
\end{equation}
上例 $\det A=5$，所以 $A^{-1}=\frac15\begin{bmatrix}3&-1\\-1&2\end{bmatrix}$，乘以 $b$ 同样得到 $(2,1)^\mathsf T$。实际程序需要的是解 $Ax=b$，通常直接调用线性求解器，不必先显式构造 $A^{-1}$。

\subsection{秩描述独立信息有几份}
矩阵的列空间（column space）是各列线性组合能产生的所有输出；秩（rank）是其中独立方向的个数。若
\begin{equation}
 A=\begin{bmatrix}1&2\\2&4\end{bmatrix},
\end{equation}
第二列是第一列的两倍，因此两个输入只有一个独立的输出方向，秩为 1。$b=(3,6)^\mathsf T$ 可由 $x=(3,0)^\mathsf T$、$(1,1)^\mathsf T$ 等无穷多组输入生成；$b=(3,7)^\mathsf T$ 则不在列空间中，无论如何都无法精确生成。

零空间（null space）是所有被 $A$ 映射为零的输入：
\begin{equation}
 \mathcal N(A)=\{x:Ax=0\}.
\end{equation}
上例 $A(-2,1)^\mathsf T=0$。只要在任一解上加 $\alpha(-2,1)^\mathsf T$，输出就不变。对 $m\times n$ 矩阵，秩与零空间维数满足
\begin{equation}
 \operatorname{rank}(A)+\dim\mathcal N(A)=n.
\end{equation}
七轴机械臂在六维任务雅可比行满秩时，$n=7$、秩为 6，于是速度零空间是一维。这里有条件：任务减少为三维位置时，零空间通常为四维；任务雅可比奇异时也会改变。不能把“七轴只有一个零空间方向”当成任何情形都成立的规则。

\section{特征值说明反复作用会怎样}
如果非零向量 $v$ 满足 $Av=\lambda v$，它就是 $A$ 的特征向量（eigenvector），$\lambda$ 是相应特征值（eigenvalue）。沿这个方向输入，输出仍在同一直线上，只改变长度或正负方向。求特征值就是让
\begin{equation}
 (A-\lambda I)v=0
\end{equation}
有非零解，因而必须有 $\det(A-\lambda I)=0$。
\begin{examplebox}[算例 手算两个特征值和两个特征方向]
取 $A=\begin{bmatrix}2&1\\1&2\end{bmatrix}$，则
\[
 \det(A-\lambda I)=(2-\lambda)^2-1
 =\lambda^2-4\lambda+3=(\lambda-3)(\lambda-1).
\]
所以 $\lambda_1=3$、$\lambda_2=1$。代回方程，得到单位特征向量
\[
 v_1=\frac{1}{\sqrt2}\begin{bmatrix}1\\1\end{bmatrix},\qquad
 v_2=\frac{1}{\sqrt2}\begin{bmatrix}-1\\1\end{bmatrix}.
\]
输入 $x=(1,0)^\mathsf T=(v_1-v_2)/\sqrt2$，输出为 $(3v_1-v_2)/\sqrt2=(2,1)^\mathsf T$。反复作用 $k$ 次时，第一个分量被乘 $3^k$，第二个仍乘 1。这就是用特征值判断离散系统会放大还是衰减的来源。
\end{examplebox}

若矩阵有一组线性独立的特征向量，可写 $A=V\Lambda V^{-1}$，于是 $A^k=V\Lambda^kV^{-1}$。并非所有方阵都能这样对角化。例如 $\begin{bmatrix}1&1\\0&1\end{bmatrix}$ 只有一个独立特征向量。一般情形可用 Jordan 结构或直接矩阵指数；本书不把“任何矩阵都能拆成独立特征方向”作为前提。

对实对称矩阵，总可以选择一组相互正交的实特征向量。若其全部特征值为正，就称为对称正定（symmetric positive definite，SPD），等价于
\begin{equation}
 x^\mathsf TAx>0\quad\text{对每一个非零 }x.
\end{equation}
例如上例 $x^\mathsf TAx=2x_1^2+2x_1x_2+2x_2^2=(x_1+x_2)^2+x_1^2+x_2^2>0$。把它乘 $1/2$ 可得到类似弹簧势能的二次型。允许等于零则称半正定（positive semidefinite）。对称 $2\times2$ 矩阵 $\begin{bmatrix}a&b\\b&d\end{bmatrix}$ 正定的便捷判据是 $a>0$ 且 $ad-b^2>0$。

\begin{figure}[htbp]
\centering\includegraphics[width=0.88\linewidth]{figures/mk_eigen_energy.pdf}
\caption{左图中两条特征方向经过矩阵后不转向，长度分别乘 3 和 1。右图是二次型等高线；同样的数值沿不同方向达到时，向量长度不同。}
\end{figure}

\section{奇异值分解把任意线性映射分成三步}
奇异值分解（singular value decomposition，SVD）为
\begin{equation}
 A=U\Sigma V^\mathsf T.
\end{equation}
$V^\mathsf T$ 把输入换到一组相互正交的坐标；$\Sigma$ 沿各坐标方向拉伸；$U$ 再把结果转到输出坐标。正交矩阵可能包含反射，因此这里的“换坐标”比笼统称作“旋转”更准确。$\Sigma$ 的非负对角元素为奇异值，按 $\sigma_1\ge\sigma_2\ge\cdots\ge0$ 排列。SVD 适用于长方形矩阵，也不要求存在完整的特征向量组。

\begin{examplebox}[算例 从两连杆雅可比手算奇异值]
取本书贯穿的两连杆在 $q=(0,\pi/2)^\mathsf T$ 时的位置雅可比
\[
 J=\begin{bmatrix}-0.5&-0.5\\0.5&0\end{bmatrix}.
\]
先计算
\[
 J^\mathsf TJ=\begin{bmatrix}0.5&0.25\\0.25&0.25\end{bmatrix}.
\]
其特征方程为 $\lambda^2-0.75\lambda+0.0625=0$，根为 $0.654508$ 和 $0.095492$。奇异值是这些根的平方根，故 $\sigma_1=0.809017$、$\sigma_2=0.309017$。可选
\[
 v_1=\begin{bmatrix}0.850651\\0.525731\end{bmatrix},\quad
 v_2=\begin{bmatrix}0.525731\\-0.850651\end{bmatrix},
\]
再由 $u_i=Jv_i/\sigma_i$ 得
\[
 u_1=\begin{bmatrix}-0.850651\\0.525731\end{bmatrix},\quad
 u_2=\begin{bmatrix}0.525731\\0.850651\end{bmatrix}.
\]
输入速度长度为 $1\,\mathrm{rad/s}$ 时，输出速度在方向 $u_1$ 上最大可达 $0.809017\,\mathrm{m/s}$，在方向 $u_2$ 上最大为 $0.309017\,\mathrm{m/s}$。这两个值描述当前姿态下的几何换算，不是电机真实速度上限。
\end{examplebox}

\begin{figure}[htbp]
\centering\includegraphics[width=0.9\linewidth]{figures/mk_svd.pdf}
\caption{单位关节速度圆经过雅可比变成末端速度椭圆。椭圆主轴方向由 $U$ 给出，半轴长度为奇异值。所有曲线均由同一个数值矩阵计算。}
\end{figure}

由 $A^\mathsf TA=V\Sigma^\mathsf T\Sigma V^\mathsf T$ 可知，奇异值的平方是 $A^\mathsf TA$ 的特征值。手算可用这条关系理解概念，但数值程序应直接求 SVD，避免先构造 $A^\mathsf TA$ 增加病态程度。对称正定矩阵的奇异值等于特征值；对称不定矩阵的奇异值等于特征值绝对值。例如 $\operatorname{diag}(2,-1)$ 的特征值是 2 和 $-1$，奇异值却是 2 和 1。

\section{条件数是误差敏感程度 不是统一的及格线}
对可逆方阵，二范数条件数为
\begin{equation}
 \kappa_2(A)=\|A\|_2\|A^{-1}\|_2=\frac{\sigma_{\max}}{\sigma_{\min}}.
\end{equation}
当 $A$ 固定，仅右端从 $b$ 变成 $b+\delta b$，解相应变成 $x+\delta x$ 时，$A\delta x=\delta b$。于是
\[
 \|\delta x\|\le\|A^{-1}\|\,\|\delta b\|,
 \qquad \|b\|=\|Ax\|\le\|A\|\,\|x\|.
\]
将二者相乘整理得到
\begin{equation}
 \frac{\|\delta x\|}{\|x\|}
 \le\kappa_2(A)\frac{\|\delta b\|}{\|b\|}.
\end{equation}
这是上界，实际误差可能小得多。矩阵本身也有误差、或者问题是带残差的最小二乘时，误差界还需额外条件，不能直接套用这个最简单版本。
\begin{examplebox}[算例 数据只变一点 解却变很多]
令
\[
 A=\begin{bmatrix}1&1\\1&1.01\end{bmatrix},\qquad b=\begin{bmatrix}2\\2.01\end{bmatrix}.
\]
原方程相减得到 $0.01x_2=0.01$，故 $x=(1,1)^\mathsf T$。仅把第二个测量改为 2.02，相减变成 $0.01x_2=0.02$，于是 $x=(0,2)^\mathsf T$。
其条件数约为 $402.008$。测量的相对二范数变化为 $0.01/\sqrt{2^2+2.01^2}\approx0.3527\%$，解的相对二范数变化却为 $\sqrt{(-1)^2+1^2}/\sqrt{1^2+1^2}=100\%$，没有违反约 $141.8\%$ 的上界。
\end{examplebox}

改变单位会改变矩阵列或行的尺度，因此会改变条件数以及最小奇异值。把米改成毫米，原始几何雅可比的线速度行就放大一千倍。对位置与姿态混合的雅可比，应先指定特征长度或权重；对参数辨识矩阵，应同时说明噪声白化方式、参数尺度和基参数选取。条件数 40 只在这些口径固定后才有比较意义。

\section{最小二乘为何是一个投影}
\subsection{有噪声时无法同时满足每一个方程}
测量三组数据 $(t,y)=(1,1),(2,2),(3,2)$，希望用直线 $y=at+b$ 描述。三条方程而只有两个未知数，通常没有精确解。写成
\begin{equation}
 H\beta\approx y,\qquad
 H=\begin{bmatrix}1&1\\2&1\\3&1\end{bmatrix},\quad
 \beta=\begin{bmatrix}a\\b\end{bmatrix},\quad
 y=\begin{bmatrix}1\\2\\2\end{bmatrix}.
\end{equation}
残差（residual）$r=y-H\beta$ 是测量减模型。最小二乘（least squares）找使残差平方和最小的参数：
\begin{equation}
 E(\beta)=r^\mathsf Tr=(y-H\beta)^\mathsf T(y-H\beta).
\end{equation}
展开为 $y^\mathsf Ty-2\beta^\mathsf TH^\mathsf Ty+\beta^\mathsf TH^\mathsf TH\beta$。对每个参数求偏导后合成梯度：
\begin{equation}
 \nabla_\beta E=-2H^\mathsf Ty+2H^\mathsf TH\beta.
\end{equation}
令梯度为零便得正规方程（normal equations）$H^\mathsf TH\hat\beta=H^\mathsf Ty$。若 $H$ 列满秩，解唯一；否则需要伪逆、先验约束或改进实验。
\begin{examplebox}[算例 完整求出直线和每一个残差]
本例
\[
 H^\mathsf TH=\begin{bmatrix}14&6\\6&3\end{bmatrix},\qquad
 H^\mathsf Ty=\begin{bmatrix}11\\5\end{bmatrix}.
\]
由 $6a+3b=5$ 得 $b=5/3-2a$。代入第一行：$14a+10-12a=11$，故 $a=1/2$，$b=2/3$。预测为
\[
 \hat y=(7/6,5/3,13/6)^\mathsf T,
 \qquad r=(-1/6,1/3,-1/6)^\mathsf T.
\]
平方和为 $1/36+1/9+1/36=1/6$。并且
\[
 H^\mathsf Tr=\begin{bmatrix}-1/6+2/3-3/6\\-1/6+1/3-1/6\end{bmatrix}=0.
\]
残差与 $H$ 的每一列都正交，所以预测 $H\hat\beta$ 正是 $y$ 向 $H$ 列空间的正交投影。这个性质比“背下一个求逆公式”更有价值。
\end{examplebox}

\begin{figure}[htbp]
\centering\includegraphics[width=0.88\linewidth]{figures/mk_least_squares.pdf}
\caption{三个测量点不能被同一直线精确穿过。最小二乘使竖直残差的平方和最小；右图列出同一解的三个残差，它们的和为零。}
\end{figure}

\subsection{伪逆 加权和正则化分别解决什么}
SVD 令 $H=U\Sigma V^\mathsf T$。把非零奇异值倒数放入 $\Sigma^+$、零奇异值仍置零，得到 Moore--Penrose 伪逆 $H^+=V\Sigma^+U^\mathsf T$。$\hat\beta=H^+y$ 是最小二乘解中欧氏范数最小的一组。对可逆方阵，伪逆就是逆。伪逆不会创造缺失的数据：秩不足时，它只是选择一个数学上明确的解。

若第 $i$ 个测量的噪声标准差为 $s_i$，加权最小二乘采用 $\sum_i r_i^2/s_i^2$。实现时把 $H$ 与 $y$ 的第 $i$ 行都除以 $s_i$；这是“白化”。目标函数里的权重为 $1/s_i^2$，实际乘到方程行上的因子却是 $1/s_i$，两者不要混淆。

正则化（regularization）则额外惩罚参数离参考值 $\beta_0$ 太远：
\begin{equation}
 \min_\beta\ \|H\beta-y\|^2+\lambda\|\beta-\beta_0\|^2,
 \qquad\lambda>0.
\end{equation}
它稳定不可靠方向，但会引入偏差。若 $\beta_0$ 不正确，较大的 $\lambda$ 可能把错误固定下来。数值求解优先使用 QR 或 SVD；即使公式写有逆矩阵，也不代表实现时要直接求逆。本章先建立这些概念，后续辨识章节会把它们用于力矩数据，并讨论噪声、参数物理一致性和验证轨迹。


\subsection{QR 分解怎样避开正规方程的条件数平方}
QR 分解把 $H$ 写成 $QR$，其中 $Q$ 的列正交归一，$R$ 为上三角矩阵。对于列满秩的长矩阵，经济型 $Q$ 是 $N\times p$，并不是方阵；$Q^\mathsf TQ=I$，但 $QQ^\mathsf T$ 通常不是单位阵，而是列空间投影。因此精确的残差分解应为
\begin{equation}
 \|H\beta-y\|^2
 =\|R\beta-Q^\mathsf Ty\|^2
 +\|(I-QQ^\mathsf T)y\|^2.
\end{equation}
第二项与参数无关，故最小化第一项即可。若省掉第二项而声称两侧范数总相等，只对相应完整正交变换的写法才成立。

\begin{examplebox}[算例 对同一组三点数据手算 QR]
第一列 $h_1=(1,2,3)^\mathsf T$ 的长度为 $\sqrt{14}$，所以 $q_1=h_1/\sqrt{14}$。第二列 $h_2=(1,1,1)^\mathsf T$ 在 $q_1$ 上的投影系数为 $r_{12}=q_1^\mathsf Th_2=6/\sqrt{14}$。去掉这一投影：
\[
 u_2=h_2-q_1r_{12}=\frac17(4,1,-2)^\mathsf T.
\]
其长度为 $\sqrt{21}/7$，归一化得 $q_2=(4,1,-2)^\mathsf T/\sqrt{21}$。因此
\[
 R=\begin{bmatrix}\sqrt{14}&6/\sqrt{14}\\0&\sqrt{21}/7\end{bmatrix},
 \qquad Q^\mathsf Ty=\begin{bmatrix}11/\sqrt{14}\\2/\sqrt{21}\end{bmatrix}.
\]
从最后一行回代：$(\sqrt{21}/7)b=2/\sqrt{21}$，故 $b=2/3$。第一行乘 $\sqrt{14}$ 得 $14a+6b=11$，故 $a=1/2$。全程没有构造 $H^\mathsf TH$。这里用的正交化步骤适合手算说明，实际高精度软件常用更稳定的 Householder 反射实现 QR。
\end{examplebox}

对称正定矩阵还可用 Cholesky 分解 $A=LL^\mathsf T$，$L$ 为下三角矩阵。它相当于把二次型写成 $x^\mathsf TAx=\|L^\mathsf Tx\|^2$，既显示能量非负，也能高效解方程。
\begin{examplebox}[算例 三角回代代替显式求逆]
取 $A=\begin{bmatrix}4&2\\2&3\end{bmatrix}$，设 $L=\begin{bmatrix}l_{11}&0\\l_{21}&l_{22}\end{bmatrix}$。比较 $LL^\mathsf T$ 的元素：$l_{11}^2=4$，故 $l_{11}=2$；$l_{11}l_{21}=2$，故 $l_{21}=1$；$l_{21}^2+l_{22}^2=3$，故 $l_{22}=\sqrt2$。
求 $Ax=(6,5)^\mathsf T$，先解 $Lz=b$：$z_1=3$、$z_2=(5-3)/\sqrt2=\sqrt2$。再解 $L^\mathsf Tx=z$：$x_2=1$、$x_1=(3-1)/2=1$。这两次三角求解比先算完整逆矩阵更直接。
\end{examplebox}

\subsection{梯度 Hessian 和局部二次近似}
梯度（gradient）把一个标量函数对每个参数的偏导排成列向量。对 $E(x_1,x_2)=x_1^2+2x_1x_2+3x_2^2$，
\[
 \nabla E=\begin{bmatrix}2x_1+2x_2\\2x_1+6x_2\end{bmatrix}.
\]
它指向局部增加最快的方向，沿负梯度走可在足够小的步长下降低函数值。Hessian 矩阵（海森矩阵）再对梯度求偏导，描述曲率：$\nabla^2E=\begin{bmatrix}2&2\\2&6\end{bmatrix}$。二阶展开是
\begin{equation}
 E(x+\delta x)\approx E(x)+\nabla E(x)^\mathsf T\delta x
 +\tfrac12\delta x^\mathsf T\nabla^2E(x)\delta x.
\end{equation}
对这个二次函数，展开其实严格相等；对一般非线性函数才是局部近似。正规方程中的 $H^\mathsf TH$ 与最小二乘目标的 Hessian 直接有关：取目标 $\frac12\|H\beta-y\|^2$，Hessian 恰为 $H^\mathsf TH$。列不独立时有平坦方向，意味着某些参数变化不改变预测；这正是辨识不可区分参数的数学表现。

\section{均值 方差 协方差为观测器做准备}
随机变量（random variable）描述重复测量时可能变化的数。期望 $\mathbb E[x]$ 是理想重复实验的平均；有限 $N$ 次测量的样本均值为 $\bar x=N^{-1}\sum_i x_i$。方差衡量分散程度，单位是原量单位的平方；标准差为方差平方根，恢复原单位。估计未知总体方差常用
\begin{equation}
 s^2=\frac{1}{N-1}\sum_{i=1}^N(x_i-\bar x)^2.
\end{equation}
\begin{examplebox}[算例 位置测量的均值和标准差]
四次读数为 $9.9,10.1,10.0,10.0\,\mathrm{mm}$。均值 $\bar x=10.0\,\mathrm{mm}$，偏差平方和为 $0.01+0.01=0.02\,\mathrm{mm}^2$，所以无偏样本方差为 $0.02/3=0.006667\,\mathrm{mm}^2$，标准差为 $0.08165\,\mathrm{mm}$。若计算的是这四个数本身的平均平方偏差，分母用 4，得到 $0.005\,\mathrm{mm}^2$；两者用途不同，不是互相矛盾。
\end{examplebox}

多维随机量的协方差矩阵定义为 $P=\mathbb E[(x-\mu)(x-\mu)^\mathsf T]$。对角元素是每一维的方差，非对角元素描述两维偏差是否倾向一起变化。例如 $P=\begin{bmatrix}4&1\\1&1\end{bmatrix}\,\mathrm{mm}^2$ 表示第一维标准差为 $2\,\mathrm{mm}$、第二维为 $1\,\mathrm{mm}$，相关系数为 $1/(2\cdot1)=0.5$。误差椭圆倾斜，不能只用两个独立误差条完全描述。

高斯分布（Gaussian distribution）由均值和协方差确定，但现实噪声不一定高斯。白噪声表示不同采样时刻不相关且具有约定的恒定功率谱；“白”不等于“高斯”，不相关也不普遍等于统计独立。编码器量化、摩擦切换、温漂都可能产生有结构的误差。学习这些词是为了识别模型假设，而不是给所有残差一律贴上“白高斯噪声”的标签。


\subsection{不确定度怎样通过一个矩阵传到输出}
设随机误差 $\delta x$ 均值为零，协方差为 $P_x$；线性输出误差为 $\delta y=A\delta x$。直接代定义：
\begin{equation}
 P_y=\mathbb E[(A\delta x)(A\delta x)^\mathsf T]
 =A\mathbb E[\delta x\delta x^\mathsf T]A^\mathsf T
 =AP_xA^\mathsf T.
\end{equation}
这条式子是卡尔曼滤波协方差预测、误差传播和标定置信度计算的共同基础。若输入误差相关，必须保留协方差的非对角项。
例如两关节测量误差互相独立，标准差都为 $0.001\,\mathrm{rad}$，那么 $P_q=10^{-6}I$。在本章标准姿态，由局部误差 $\delta p\approx J\delta q$，有
\[
 P_p\approx10^{-6}JJ^\mathsf T
 =10^{-6}\begin{bmatrix}0.5&-0.25\\-0.25&0.25\end{bmatrix}\mathrm{m^2}.
\]
位置标准差分别为 $0.707107\,\mathrm{mm}$ 和 $0.5\,\mathrm{mm}$，相关系数为 $-0.25/\sqrt{0.5\times0.25}=-0.707107$。两个关节噪声独立，末端两轴误差却相关；是几何映射产生了这种相关性。这个结果是小噪声一阶传播近似，并未包含几何尺寸和零偏的不确定度。

\section{计算练习与逐步解答}
\subsection*{题目}
\begin{enumerate}
\item 已知 $A=\begin{bmatrix}3&1\\1&2\end{bmatrix}$、$b=(7,4)^\mathsf T$，求 $x$，计算 $\det A$，并判断 $A$ 是否正定。
\item 矩阵 $B=\begin{bmatrix}1&0&1\\0&1&1\end{bmatrix}$ 将三维输入变成二维输出。求一个非零零空间向量；对输出 $(1,2)^\mathsf T$ 写出所有解，并求最小范数解。
\item 用常数 $c$ 拟合测量 $1,2,3$。先求普通最小二乘，再在噪声标准差为 $1,1,2$ 时求加权最小二乘。
\item 对 $D=\operatorname{diag}(2,0.02)$，求奇异值、条件数、$D^\mathsf TD$ 的条件数。若右端 $b=(2,0.02)^\mathsf T$ 的第二分量增加 0.002，求参数变化。
\end{enumerate}
\subsection*{解答}
\begin{enumerate}
\item $\det A=3\cdot2-1=5$。$A^{-1}=\frac15\begin{bmatrix}2&-1\\-1&3\end{bmatrix}$，所以 $x=\frac15(14-4,-7+12)^\mathsf T=(2,1)^\mathsf T$。$A$ 对称、第一顺序主子式为 $3>0$、行列式为 $5>0$，故正定。代回得 $(3\cdot2+1,2+2)^\mathsf T=(7,4)^\mathsf T$。
\item 方程为 $x_1+x_3=1$、$x_2+x_3=2$。令 $x_3=t$，则 $x=(1-t,2-t,t)^\mathsf T$；零空间方向为 $(-1,-1,1)^\mathsf T$。范数平方为 $(1-t)^2+(2-t)^2+t^2=5-6t+3t^2$，导数为 $-6+6t$，在 $t=1$ 最小，故最小范数解为 $(0,1,1)^\mathsf T$。
\item 普通目标为 $(1-c)^2+(2-c)^2+(3-c)^2$，导数 $6c-12=0$，得 $c=2$。加权目标为 $(1-c)^2+(2-c)^2+(3-c)^2/4$，导数为 $2[(c-1)+(c-2)+(c-3)/4]=0$，得 $c=(1+2+3/4)/(1+1+1/4)=5/3$。第三个较不可靠的测量对结果影响变小。
\item 奇异值为 $2,0.02$，$\kappa_2(D)=100$。$D^\mathsf TD=\operatorname{diag}(4,0.0004)$，条件数为 $10\,000=100^2$。原解为 $(1,1)^\mathsf T$，新解为 $(1,1.1)^\mathsf T$。第二个参数变化 0.1；小奇异值方向的绝对噪声被放大了 $1/0.02=50$ 倍。
\end{enumerate}

\section*{本章与原稿的对应}
本章整合 [A，第2--5页] 与 [B，第3、8--17页] 的线性代数、最小二乘和概率内容。补全了单位与尺度、秩与零空间、逐项求导和数值计算；修正了对称不定矩阵的 SVD、一般矩阵可对角化及条件数适用范围的表述。

\chapter{微积分 微分方程和信号变换}\label{ch:calculus}
\section{导数把位置变化变成速度}
设关节角 $q(t)$ 随时间变化。时间区间 $[t,t+h]$ 上的平均速度为 $[q(t+h)-q(t)]/h$。令区间越来越短，若极限存在，就得到导数
\begin{equation}
 \dot q(t)=\lim_{h\to0}\frac{q(t+h)-q(t)}{h}.
\end{equation}
导数描述“这一瞬间的变化率”，不是当前位置除以当前时间。例如 $q(t)=t^2$ 时，$q(2)=4$，但瞬时速度为 $\dot q(2)=4$，而 $q(2)/2=2$ 只是从零时刻起的平均速度。

\begin{examplebox}[算例 从轨迹逐次求导并检查单位]
令
\[
 q(t)=0.2+0.3t^2-0.1t^3\quad(0\le t\le2\,\mathrm s),
\]
系数包含使结果为 rad 的相应单位。逐项求导得到
\[
 \dot q=0.6t-0.3t^2,
 \qquad\ddot q=0.6-0.6t,
 \qquad\dddot q=-0.6.
\]
在 $t=1\,\mathrm s$ 时，角度 $q=0.4\,\mathrm{rad}$，速度 $0.3\,\mathrm{rad/s}$，加速度为零，加加速度（jerk）为 $-0.6\,\mathrm{rad/s^3}$。加速度为零只表示该瞬间速度不再增加，不表示关节停止。
\end{examplebox}

积分是累积变化。若 $\dot q(t)=v(t)$，则
\begin{equation}
 q(t)=q(0)+\int_0^t v(\xi)\,d\xi.
\end{equation}
$\xi$ 是积分内部的时间变量，用来避免与上限 $t$ 混淆。若关节以 $0.2\,\mathrm{rad/s}$ 运动 $0.5\,\mathrm s$，角度增加 $0.1\,\mathrm{rad}$。力矩对角度的积分是做功，而不是力矩对时间积分：$W=\int\tau\,dq$，单位为焦耳；后者对应角动量变化。

\section{偏导数 链式法则与线性化}
\subsection{一维泰勒展开是局部直线近似}
在 $q_0$ 附近，光滑函数满足
\begin{equation}
 f(q_0+\delta q)=f(q_0)+f'(q_0)\delta q
 +\frac12 f''(q_0)(\delta q)^2+\cdots.
\end{equation}
保留一阶项称为线性化（linearization）。它只在小邻域里有效，忽略项通常为二阶小量。三角函数求导按弧度进行；如果输入数值用度，需要额外乘 $\pi/180$。
\begin{examplebox}[算例 用导数估计连杆端点的小位移]
单杆长度 $l=0.5\,\mathrm m$，横坐标 $x(q)=l\cos q$。在 $q_0=60^\circ=\pi/3$ 附近，关节增加 $\delta q=0.02\,\mathrm{rad}$：
\[
 x(q_0)=0.25,\quad x'(q_0)=-0.5\sin(\pi/3)=-0.433013.
\]
故线性估计为 $x\approx0.25-0.433013\times0.02=0.241340\,\mathrm m$。精确值为 $0.5\cos(\pi/3+0.02)=0.241290\,\mathrm m$，相差约 $0.0494\,\mathrm{mm}$。二阶项为 $\frac12(-0.25)(0.02)^2=-0.00005\,\mathrm m$，正好解释主要误差。
\end{examplebox}

\begin{figure}[htbp]
\centering\includegraphics[width=0.86\linewidth]{figures/mk_linearization.pdf}
\caption{连杆横坐标曲线及在 $60^\circ$ 处的切线。放大图显示线性化误差随关节增量增大；“雅可比乘增量”是局部关系。}
\end{figure}

\subsection{多变量时把每个输入的影响分开算}
对 $f(q_1,q_2)$，偏导数 $\partial f/\partial q_1$ 表示暂时固定 $q_2$，只改变 $q_1$ 的影响。若两者都随时间变化，则链式法则给出
\begin{equation}
 \frac{df}{dt}=\frac{\partial f}{\partial q_1}\dot q_1
 +\frac{\partial f}{\partial q_2}\dot q_2.
\end{equation}
将多个输出的偏导数排成矩阵，便得到雅可比。对位置 $p=f(q)$，
\begin{equation}
 \dot p=J_p(q)\dot q,\qquad
 \ddot p=J_p(q)\ddot q+\dot J_p(q,\dot q)\dot q.
\end{equation}
第二式来自乘积求导。$\dot J_p\dot q$ 不是装饰项：即使关节速度不变，端点沿圆弧运动仍有向心加速度。例如单杆 $p=(l\cos q,l\sin q)^\mathsf T$，$\ddot q=0$ 时，$\ddot p=-l\dot q^2(\cos q,\sin q)^\mathsf T$。

\subsection{差分为何放大测量噪声}
程序看不到无限短的 $h$，常用 $\hat{\dot q}_k=(q_k-q_{k-1})/T_s$，其中 $T_s$ 是采样周期。若两次位置噪声独立、方差均为 $\sigma_q^2$，差分噪声方差为
\begin{equation}
 \sigma_{\dot q}^2=\frac{2\sigma_q^2}{T_s^2}.
\end{equation}
若 $\sigma_q=10^{-4}\,\mathrm{rad}$，$T_s=0.005\,\mathrm s$，速度噪声标准差为 $\sqrt2\times10^{-4}/0.005=0.02828\,\mathrm{rad/s}$。采样改为 1 ms，但仍直接做这类差分，标准差会上升至 $0.14142\,\mathrm{rad/s}$。更快采样有益，但必须与合适的估计、滤波和传感器特性配合；不能把“采得快”直接等同于“数值导数更干净”。

\section{一阶微分方程完整解一次}
常微分方程（ordinary differential equation，ODE）将未知函数与其时间导数联系起来。方程的阶数是出现的最高阶导数。给定初值，是为了从一族可能轨迹中确定实际那一条。

考虑单关节在忽略重力等项后的速度方程
\begin{equation}
 J_r\dot\omega+b\omega=\tau_0,
 \quad J_r=0.1\,\mathrm{kg\,m^2},\quad
 b=0.2\,\mathrm{N\,m\,s/rad}.
\end{equation}
这里 $J_r$ 是转动惯量，避免与雅可比矩阵 $J$ 混淆。恒力矩取 $\tau_0=0.4\,\mathrm{N\,m}$，两边除以 $J_r$ 得 $\dot\omega+2\omega=4$。乘积分因子 $e^{2t}$：
\begin{equation}
 \frac{d}{dt}(e^{2t}\omega)=4e^{2t}.
\end{equation}
积分得 $e^{2t}\omega=2e^{2t}+C$，即 $\omega=2+Ce^{-2t}$。初始静止 $\omega(0)=0$ 给出 $C=-2$，故
\begin{equation}
 \omega(t)=2(1-e^{-2t})\,\mathrm{rad/s}.
\end{equation}
平衡速度 $\omega_\infty=\tau_0/b=2\,\mathrm{rad/s}$；时间常数 $T_c=J_r/b=0.5\,\mathrm s$。$t=T_c$ 时达到最终变化量的 $1-e^{-1}=63.21\%$，$t=3T_c$ 时达到 $95.02\%$，$t=5T_c$ 时达到 $99.33\%$。
\begin{examplebox}[算例 速度已经接近稳定 角度仍然持续变化]
在 $t=1\,\mathrm s$ 时，$\omega=2(1-e^{-2})=1.729329\,\mathrm{rad/s}$。再积分，且 $q(0)=0$：
\[
 q(t)=\int_0^t2(1-e^{-2\xi})\,d\xi
 =2t-(1-e^{-2t}).
\]
于是 $q(1)=1.135335\,\mathrm{rad}$。速度趋于常数不意味着位置趋于常数。本例仅研究速度阻尼，未设置回到某个位置的弹簧或位置反馈。
\end{examplebox}

一般一阶系统 $T_c\dot y+y=Ku$ 对常值输入 $u_0$ 的解是 $y(t)=Ku_0+[y(0)-Ku_0]e^{-t/T_c}$，前提 $T_c>0$。这同一个数学结构可以描述电流建立、低通滤波、速度阻尼或某些热过程，具体单位由应用决定。

\section{二阶微分方程解释振荡和阻尼}
\subsection{从特征根得到响应形状}
质量弹簧阻尼模型为
\begin{equation}
 m\ddot x+c\dot x+kx=F(t).
\end{equation}
$m$ 为质量，单位 kg；$c$ 为黏性阻尼，单位 $\mathrm{N\,s/m}$；$k$ 为刚度，单位 $\mathrm{N/m}$。先令外力为零，猜解 $x=e^{st}$。因 $\dot x=se^{st}$、$\ddot x=s^2e^{st}$，除去非零因子 $e^{st}$ 后得到特征方程
\begin{equation}
 ms^2+cs+k=0.
\end{equation}
定义无阻尼自然角频率与阻尼比
\begin{equation}
 \omega_n=\sqrt{k/m},\qquad
 \zeta=\frac{c}{2\sqrt{km}},
\end{equation}
便得到 $s_{1,2}=-\zeta\omega_n\pm\omega_n\sqrt{\zeta^2-1}$。

若 $0<\zeta<1$，根为共轭复数 $-\zeta\omega_n\pm j\omega_d$，其中 $\omega_d=\omega_n\sqrt{1-\zeta^2}$，解是指数衰减乘正弦或余弦。$\zeta=1$ 时是重实根，齐次解为 $(C_1+C_2t)e^{-\omega_nt}$；$\zeta>1$ 时为两个不同的负实根，解是两个衰减指数之和。若 $c=0$ 且 $k,m>0$，模型永远振荡；若阻尼为负，反馈可能向系统注入能量。

“临界阻尼最快而不超调”需要限定比较对象：固定 $\omega_n$、标准二阶无零点传递函数、零初值阶跃跟踪，临界阻尼是常用的快速单调选择。任意初始速度、额外零点或饱和都会改变超调，不能据一个 $\zeta$ 就承诺全部运动不越界。

\subsection{把阶跃响应完整算成数字}
\begin{examplebox}[算例 弹簧刚度相同 阻尼改变超调]
取 $m=1\,\mathrm{kg}$、$k=100\,\mathrm{N/m}$、$c=4\,\mathrm{N\,s/m}$，初始位置速度都为零，施加 $F_0=1\,\mathrm N$ 恒力。平衡位置 $x_\infty=F_0/k=0.01\,\mathrm m$，$\omega_n=10\,\mathrm{rad/s}$，$\zeta=0.2$，$\omega_d=\sqrt{96}=9.797959\,\mathrm{rad/s}$。解为
\[
 x(t)=0.01\left[1-e^{-2t}\left(\cos(\sqrt{96}t)
 +\frac{2}{\sqrt{96}}\sin(\sqrt{96}t)\right)\right].
\]
求导后正弦为零的第一个峰值出现在 $t_p=\pi/\omega_d=0.320637\,\mathrm s$。代回得相对超调
\[
 M_p=e^{-\zeta\pi/\sqrt{1-\zeta^2}}=0.526621,
\]
所以峰值为 $0.0152662\,\mathrm m$，超过平衡位置约 $5.2662\,\mathrm{mm}$。若改为临界阻尼 $c=2\sqrt{km}=20$，响应变成 $x(t)=0.01[1-(1+10t)e^{-10t}]$。
\end{examplebox}

\begin{figure}[htbp]
\centering\includegraphics[width=0.9\linewidth]{figures/mk_second_order.pdf}
\caption{相同质量和刚度下的归一化阶跃响应。阻尼太小会振荡，过大则产生慢模态。图中的四条曲线由各自完整微分方程计算。}
\end{figure}

若要求标准欠阻尼二阶阶跃超调不超过 $5\%$，由 $e^{-\zeta\pi/\sqrt{1-\zeta^2}}\le0.05$，取对数并平方得
\begin{equation}
 \zeta\ge\frac{-\ln0.05}{\sqrt{\pi^2+(\ln0.05)^2}}
 \approx0.690107.
\end{equation}
本例至少需要 $c=2\zeta\sqrt{km}\approx13.8021\,\mathrm{N\,s/m}$。实际关节还包含采样、力矩环、时延和柔性，故这是模型设计值，不是可以直接搬到真机的保证。

\section{从二阶方程到状态方程}
令 $x_1=q$、$x_2=\dot q$，则任何给出 $\ddot q=f(q,\dot q,u)$ 的系统都可以写成一阶方程组
\begin{equation}
 \frac{d}{dt}\begin{bmatrix}x_1\\x_2\end{bmatrix}
 =\begin{bmatrix}x_2\\f(x_1,x_2,u)\end{bmatrix}.
\end{equation}
状态（state）是预测未来所需的一组当前量。在给定输入及模型时，只知道位置通常不够：同一位置可能静止，也可能高速通过，因此还需要速度。

对于 $J_r\ddot q+b\dot q+kq=\tau$，状态方程为
\begin{equation}
 \dot x=Ax+B\tau,\quad
 A=\begin{bmatrix}0&1\\-k/J_r&-b/J_r\end{bmatrix},\quad
 B=\begin{bmatrix}0\\1/J_r\end{bmatrix}.
\end{equation}
线性系统中的“线性”指状态和输入以一次形式出现，且不互相相乘；并不要求轨迹在时间图上是直线。$\dot x=Ax$ 的解是 $x(t)=e^{At}x(0)$，矩阵指数定义为 $I+At+(At)^2/2!+\cdots$，通常不等于对每个矩阵元素分别取指数。


\begin{examplebox}[算例 用矩阵指数精确推进一个双积分模型]
若输入 $u$ 直接是加速度，$\ddot q=u$，则 $A=\begin{bmatrix}0&1\\0&0\end{bmatrix}$、$B=(0,1)^\mathsf T$。因为 $A^2=0$，矩阵指数级数在一次项后全部为零：
\[
 e^{AT_s}=I+AT_s=\begin{bmatrix}1&T_s\\0&1\end{bmatrix}.
\]
在一段内输入保持为 $u_k$，输入积分为
\[
 B_d=\int_0^{T_s}e^{A\xi}B\,d\xi
 =\int_0^{T_s}\begin{bmatrix}\xi\\1\end{bmatrix}d\xi
 =\begin{bmatrix}T_s^2/2\\T_s\end{bmatrix}.
\]
因此 $q_{k+1}=q_k+T_s\dot q_k+T_s^2u_k/2$，$\dot q_{k+1}=\dot q_k+T_su_k$。在 $T_s=0.005\,\mathrm s$、静止起步、$u=2\,\mathrm{rad/s^2}$ 时，第一步角度增加 $0.000025\,\mathrm{rad}$，速度增加 $0.01\,\mathrm{rad/s}$。这个结果与匀加速运动公式完全一致，说明矩阵指数只是把熟悉的积分组织成通用形式。
\end{examplebox}

\subsection{非线性模型在平衡点附近怎样变成线性模型}
设 $\dot x=f(x,u)$，平衡点满足 $f(x_0,u_0)=0$。用 $\delta x=x-x_0$、$\delta u=u-u_0$ 表示小扰动，一阶展开得到
\begin{equation}
 \delta\dot x=A\delta x+B\delta u,
 \quad A=\left.\frac{\partial f}{\partial x}\right|_{x_0,u_0},
 \quad B=\left.\frac{\partial f}{\partial u}\right|_{x_0,u_0}.
\end{equation}
单摆从竖直向下方向量角时，$\ddot\theta=-(g/l)\sin\theta$。在下垂点 $\theta_0=0$，$\sin\theta\approx\theta$，所以 $A=\begin{bmatrix}0&1\\-g/l&0\end{bmatrix}$。在倒立点 $\theta_0=\pi$，令 $\delta\theta=\theta-\pi$，有 $\sin(\pi+\delta\theta)\approx-\delta\theta$，故加速度变成 $+(g/l)\delta\theta$。只改变工作点，局部稳定性就完全不同。

\section{复数不是额外的物理维度}
复数 $z=a+jb$ 中 $j^2=-1$。它可写为 $z=re^{j\phi}$，其中 $r=\sqrt{a^2+b^2}$，$\phi=\operatorname{atan2}(b,a)$。欧拉公式
\begin{equation}
 e^{j\phi}=\cos\phi+j\sin\phi
\end{equation}
把旋转与正弦波联系起来。两个复数相乘，模相乘、相角相加。写成复数通常是计算技巧；实际位置、电压或力仍可取复表达式的实部。

\begin{examplebox}[算例 一阶低通的幅值衰减和相位]
传递函数 $G(s)=1/(s+1)$，输入角频率为 $1\,\mathrm{rad/s}$ 时，令 $s=j$：
\[
 G(j)=\frac1{1+j}=\frac{1-j}{(1+j)(1-j)}=\frac{1-j}{2}.
\]
故模为 $\sqrt{(1/2)^2+(-1/2)^2}=1/\sqrt2$，相角为 $-\pi/4$。若输入为 $\cos t$，稳定后的输出为 $0.707107\cos(t-\pi/4)$。分贝幅值为 $20\log_{10}(0.707107)=-3.0103\,\mathrm{dB}$。
\end{examplebox}

复数根 $s=\sigma+j\omega$ 对应 $e^{st}=e^{\sigma t}(\cos\omega t+j\sin\omega t)$。实部 $\sigma$ 决定增长或衰减，虚部 $\omega$ 决定振荡角频率。$\omega$ 的单位为 rad/s，普通频率 $f=\omega/(2\pi)$ 的单位为 Hz。$10\,\mathrm{rad/s}$ 约为 $1.59155\,\mathrm{Hz}$，不是 $10\,\mathrm{Hz}$。

\section{拉普拉斯变换把导数变成代数}
对 $t\ge0$ 的信号，单边拉普拉斯变换定义为
\begin{equation}
 X(s)=\mathcal L\{x(t)\}=\int_0^\infty x(t)e^{-st}\,dt.
\end{equation}
它在积分收敛的 $s$ 范围内成立。不是所有信号在任意 $s$ 上都有相同变换。例如 $x(t)=e^{-at}$ 时，积分为 $1/(s+a)$，要求 $\operatorname{Re}(s+a)>0$。

对导数做分部积分：
\begin{align}
 \mathcal L\{\dot x\}
 &=\left[x(t)e^{-st}\right]_0^\infty
   +s\int_0^\infty x(t)e^{-st}\,dt\\
 &=sX(s)-x(0),
\end{align}
其中无穷远边界项为零需要相应收敛条件。再应用一次得 $\mathcal L\{\ddot x\}=s^2X-sx(0)-\dot x(0)$。初值没有消失，而是变成代数项。

\begin{examplebox}[算例 用拉普拉斯逐步解一个微分方程]
求 $\dot y+2y=u$、$y(0)=0$ 在单位阶跃 $u(t)=1$ 下的响应。因 $\mathcal L\{1\}=1/s$，有
\[
 (s+2)Y(s)=1/s,\qquad
 Y(s)=\frac1{s(s+2)}.
\]
设 $Y=A/s+B/(s+2)$，两边乘 $s(s+2)$ 得 $1=A(s+2)+Bs$。比较常数和 $s$ 的系数，$2A=1$、$A+B=0$，所以 $A=1/2$、$B=-1/2$。逆变换得到
\[
 y(t)=\tfrac12(1-e^{-2t}).
\]
代入原方程验算：$\dot y=e^{-2t}$，$2y=1-e^{-2t}$，相加恰为 1；初值也为零。
\end{examplebox}

传递函数（transfer function）定义为零初值下输出与输入变换之比 $G(s)=Y(s)/U(s)$。上例为 $1/(s+2)$。分母的根称为极点（pole），分子的根称为零点（zero）。约分后传递函数只反映从指定输入到指定输出可见的动态；被消去或不可观的内部模态不能因此判定安全。现代控制章节会用可控性和可观性进一步说明。

终值定理给出 $\lim_{t\to\infty}y(t)=\lim_{s\to0}sY(s)$，需要 $sY(s)$ 的全部极点位于严格左半平面等相应条件。本例得 $1/2$。若 $y(t)=\sin t$，形式上 $\lim_{s\to0}s/(s^2+1)=0$，但时域根本没有终值；虚轴极点让定理条件不成立。

复变函数中的留数，是局部简单极点项 $c/(s-p)$ 的系数 $c$；部分分式就可看成用留数组织逆变换。辐角原理把闭曲线内的零点数与极点数联系起来，是后续 Nyquist 稳定判据的数学基础。初学阶段先能实际分解上面的有理函数，再理解这些抽象词会更稳固。

\section{傅里叶级数和频谱解释信号由什么组成}
\subsection{有限傅里叶级数可以直接求速度和加速度}
周期为 $T_0$、基角频率 $\omega_0=2\pi/T_0$ 的信号，可在相应收敛条件下展开为
\begin{equation}
 q(t)=a_0+\sum_{k=1}^{K}
 \left[a_k\cos(k\omega_0t)+b_k\sin(k\omega_0t)\right]
\end{equation}
的有限近似。$k$ 是谐波编号，$K$ 是最高谐波阶数；“七阶傅里叶”通常表示到第七个谐波，不是把 $t$ 提到七次方。常数项和系数可由正交性求出，例如
\begin{equation}
 a_0=\frac1{T_0}\int_0^{T_0}q(t)dt,\quad
 a_k=\frac2{T_0}\int_0^{T_0}q(t)\cos(k\omega_0t)dt.
\end{equation}
正交性指不同整周期正弦余弦的乘积积分为零；同频正弦或余弦平方的积分为 $T_0/2$。这让每个系数可以单独“提取”出来。

\begin{examplebox}[算例 两个谐波的速度和加速度]
令 $q(t)=0.1\sin(2\pi t)+0.02\sin(6\pi t)$，单位 rad。两个频率分别为 1 Hz 和 3 Hz。求导得
\[
 \dot q=0.2\pi\cos(2\pi t)+0.12\pi\cos(6\pi t),
\]
\[
 \ddot q=-0.4\pi^2\sin(2\pi t)-0.72\pi^2\sin(6\pi t).
\]
在 $t=0$，速度为 $0.32\pi=1.00531\,\mathrm{rad/s}$，加速度为零。3 Hz 的位置幅值只有 1 Hz 的五分之一，速度幅值却是它的 $3/5$，加速度幅值是它的 $9/5$。高频项对力矩的影响可能很大，不能只看位置曲线是否平滑。
\end{examplebox}

\begin{figure}[htbp]
\centering\includegraphics[width=0.9\linewidth]{figures/mk_fourier.pdf}
\caption{同一有限傅里叶轨迹的位置、速度和加速度。每次求导都会把第 $k$ 个谐波额外乘 $k\omega_0$，因此较小的高频位置分量也可能主导加速度。}
\end{figure}

有限谐波并不自动保证“不会激发共振”：若某个谐波、驱动非线性产生的高次频率或过渡段频率接近柔性模态，仍会激振。辨识激励必须同时限制位置、速度、加速度、力矩和频率，并处理进入、退出周期轨迹的过渡段。

\subsection{采样频谱的几个精确边界}
对 $N$ 个采样值 $x_n$，离散傅里叶变换（DFT）定义为
\begin{equation}
 X_k=\sum_{n=0}^{N-1}x_n e^{-j2\pi kn/N}.
\end{equation}
快速傅里叶变换（FFT）是计算同一组 $X_k$ 的高效算法，不是另一种物理变换。采样率为 $f_s$ 时，第 $k$ 个频点为 $f_k=kf_s/N$，频点间隔 $\Delta f=f_s/N$。$f_s=200\,\mathrm{Hz}$、记录 $10\,\mathrm s$，则 $N=2000$、间隔为 $0.1\,\mathrm{Hz}$。实际分辨两个频率的能力还受窗函数、噪声和幅值差影响，不能只靠频点间隔判断。

混叠（aliasing）是不同连续频率给出相同离散样本的现象。例如 $f_s=200\,\mathrm{Hz}$ 时，$120\,\mathrm{Hz}$ 正弦在采样点满足
\begin{equation}
 \sin(2\pi120n/200)=-\sin(2\pi80n/200).
\end{equation}
只看样本会误认为是带相位变化的 80 Hz 信号。防混叠必须在采样之前限制高频，事后低通不能唯一恢复已经混进来的成分。对于理想带限信号，采样率严格高于最高频率两倍是重建条件；控制采样率往往还要远高于期望闭环带宽以容纳相位损失。

不是整数周期的截断会产生谱泄漏（leakage）。加窗减少远处泄漏，但会加宽主峰、改变幅值标度；零填充让频谱曲线看起来更密，不增加原始记录的独立信息。用频谱分析实验时，应同时记录采样率、时长、窗函数及幅值归一化方式。

\section{离散时间和 z 变换}
计算机每隔 $T_s$ 处理一次，连续方程变成差分方程。单边 z 变换定义为
\begin{equation}
 X(z)=\sum_{k=0}^{\infty}x_kz^{-k}.
\end{equation}
在适当初始条件下，$z^{-1}$ 对应一个采样周期的延迟。若 $x_k=a^k$，由等比级数求和得 $X(z)=1/(1-az^{-1})=z/(z-a)$，收敛域是 $|z|>|a|$。连续模态 $e^{st}$ 在 $t=kT_s$ 被采样后变成 $(e^{sT_s})^k$，因此模态映射为 $z=e^{sT_s}$。连续左半平面映射到离散单位圆内。

\begin{examplebox}[算例 把连续一阶模型精确离散一次]
已知 $\dot y=-2y+u$，在每个采样区间把 $u=u_k$ 保持不变，称为零阶保持。由一阶方程解：
\[
 y_{k+1}=e^{-2T_s}y_k+\frac{1-e^{-2T_s}}2u_k.
\]
取 $T_s=0.005\,\mathrm s$，系数为 $a=0.99004983$、$b=0.004975083$。若 $y_0=0$，$u_0=1$，第一步 $y_1=b$。若错误地省略输入项中的 $T_s$ 尺度，系统增益会完全改变。
零初值下 $Y(z)/U(z)=b/(z-a)$。离散极点 $a<1$，故稳定；恒输入的稳态值为 $b/(1-a)=1/2$，与连续模型一致。
\end{examplebox}

对简单滤波程序 $y_k=0.9y_{k-1}+0.1u_k$，其传递函数是 $0.1/(1-0.9z^{-1})$。由 $0.9=e^{-aT_s}$ 可对应一个连续衰减率 $a=-\ln0.9/T_s$。在 $T_s=1\,\mathrm{ms}$ 时，$a=105.361\,\mathrm{s^{-1}}$。$a/(2\pi)$ 是等效连续极点频率；它不是无需说明即可当作该离散滤波器精确 $-3\,\mathrm{dB}$ 频率的定义。

\section{积分器为何会凭空增能量}
\subsection{显式欧拉}
显式欧拉（forward Euler）使用当前斜率：$x_{k+1}=x_k+h f(x_k)$。对 $\dot x=-ax$，递推因子为 $1-ah$，要衰减需要 $|1-ah|<1$，即 $0<ah<2$。连续系统即使稳定，过大的数值步长也会把它算成发散。

更直接的例子是无阻尼振子 $\dot x=v$、$\dot v=-\omega^2x$。显式欧拉给出
\[
 x_{k+1}=x_k+hv_k,\qquad v_{k+1}=v_k-h\omega^2x_k.
\]
定义单位质量能量 $E_k=\frac12(v_k^2+\omega^2x_k^2)$。展开两项平方，交叉项互相抵消，得到
\begin{equation}
 E_{k+1}=(1+h^2\omega^2)E_k.
\end{equation}
任何非零步长都会逐步增能。这是算法误差，不是控制器真的注入了力。

\subsection{半隐式 隐式和高阶方法}
半隐式欧拉先更新速度，再用新速度更新位置：$v_{k+1}=v_k-h\omega^2x_k$，$x_{k+1}=x_k+hv_{k+1}$。对于这个线性振子且 $0<h\omega<2$，能量在有界范围内振荡；它不严格保持原能量，更不意味着任意步长都稳定。后向欧拉用下一步斜率，需要解联立方程，常对快速稳定模态更稳，但可能带来额外数值耗散。

四阶 Runge--Kutta（RK4）在一步内计算四次斜率；对于足够光滑的方程，它的全局误差通常随 $h^4$ 缩小。碰撞、摩擦切换、刚性接触会破坏简单的光滑误差假设。仿真时间步、控制周期和传感采样周期是三个不同量，应分别注明。

\begin{figure}[htbp]
\centering\includegraphics[width=0.9\linewidth]{figures/mk_integrators.pdf}
\caption{无阻尼振子的数值积分比较。显式欧拉能量增长，后向欧拉能量衰减，半隐式欧拉在本例步长下有界振荡；解析解能量严格不变。}
\end{figure}


\subsection{四阶 Runge--Kutta 的四次斜率具体怎么算}
对于 $\dot x=f(t,x)$，RK4 使用
\begin{align}
 k_1&=f(t_k,x_k),\\
 k_2&=f(t_k+h/2,x_k+hk_1/2),\\
 k_3&=f(t_k+h/2,x_k+hk_2/2),\\
 k_4&=f(t_k+h,x_k+hk_3),\\
 x_{k+1}&=x_k+\frac h6(k_1+2k_2+2k_3+k_4).
\end{align}
四个 $k_i$ 都是变化率，不是已经乘过步长的状态增量；不同教材可能把 $h$ 吸收入 $k_i$，实现时必须选定一种。
例如 $\dot x=-2x$、$x_0=1$、$h=0.1$，依次得 $k_1=-2$、$k_2=-1.8$、$k_3=-1.82$、$k_4=-1.636$。于是 $x_1=1+(0.1/6)(-2-3.6-3.64-1.636)=0.818733333$。解析值 $e^{-0.2}=0.818730753$，误差约 $2.58\times10^{-6}$；显式欧拉仅得到 0.8。更精细的积分提高了对这个光滑模型的数值准确度，但模型参数错误、时延或真实接触未建模的问题仍然存在。

\section{变分法从最好的数走向最好的曲线}
\subsection{把求导的对象换成整条函数}
普通优化找参数 $a$ 使 $E(a)$ 最小；变分法（calculus of variations）寻找整条函数 $q(t)$，使
\begin{equation}
 \mathcal J[q]=\int_0^T L(t,q,\dot q)\,dt
\end{equation}
最小或满足驻值条件。这样的“输入函数、输出数值”的对象称为泛函（functional）。$L$ 在这里是积分中的代价密度；动力学中的拉格朗日函数也用 $L$，具体含义随问题给出。

给候选曲线一个小扰动 $q_\epsilon=q+\epsilon\eta$，其中 $\eta(0)=\eta(T)=0$ 保持端点位置不变。若 $q$ 是驻值，$d\mathcal J[q_\epsilon]/d\epsilon$ 在 $\epsilon=0$ 时为零：
\begin{align}
 0&=\int_0^T\left(\frac{\partial L}{\partial q}\eta
       +\frac{\partial L}{\partial\dot q}\dot\eta\right)dt\\
 &=\left[\frac{\partial L}{\partial\dot q}\eta\right]_0^T
  +\int_0^T\left[\frac{\partial L}{\partial q}
  -\frac d{dt}\left(\frac{\partial L}{\partial\dot q}\right)\right]\eta\,dt.
\end{align}
边界项为零，又因为内部扰动函数 $\eta$ 可以任取，方括号必须为零，得到欧拉--拉格朗日方程
\begin{equation}
 \frac d{dt}\frac{\partial L}{\partial\dot q}
 -\frac{\partial L}{\partial q}=0.
\end{equation}
它首先是必要的驻值条件，不能在所有问题中直接把每个解都称为全局最小值。

\subsection{最小加加速度为何得到五次多项式}
令代价 $\mathcal J[q]=\int_0^T(\dddot q)^2dt$，并固定起止位置、速度、加速度。此时允许扰动 $\eta,\dot\eta,\ddot\eta$ 在两端都为零。求一阶变分并连续分部积分三次：
\begin{equation}
 \delta\mathcal J=2\int_0^T\dddot q\,\dddot\eta\,dt
 =-2\int_0^T q^{(6)}\eta\,dt=0.
\end{equation}
因此 $q^{(6)}=0$，即 $q$ 至多为五次多项式。令起点为 0、终点为 1、时间归一化 $s=t/T$，且两端速度加速度均为零，则
\begin{equation}
 q(s)=a_3s^3+a_4s^4+a_5s^5.
\end{equation}
终点的三个条件给出 $a_3+a_4+a_5=1$、$3a_3+4a_4+5a_5=0$、$6a_3+12a_4+20a_5=0$。后两式消元得 $a_4=-15$、$a_5=6$、$a_3=10$，所以 $q=10s^3-15s^4+6s^5$。因为任何满足齐次边界条件的扰动都有
\[
 \mathcal J[q+\eta]-\mathcal J[q]
 =\int_0^T(\dddot\eta)^2dt\ge0,
\]
且交叉项因 $q^{(6)}=0$ 与边界条件消失，本例确实是最小值，而不只是驻值。

这并未自动加入最大速度、最大力矩、避障等约束；加入它们后，最优轨迹一般需要分段表达或数值优化。后续规划章节将把这一数学结果与实际运动约束连接起来。

\section{计算练习与逐步解答}
\subsection*{题目}
\begin{enumerate}
\item $q(t)=0.1+0.4t-0.05t^2$，求 $t=2\,\mathrm s$ 的位置、速度和加速度，以及从 0 到 2 s 的角度变化。
\item 求 $\dot y+4y=8$、$y(0)=1$ 的解、时间常数和 $t=0.5\,\mathrm s$ 的值。
\item $m=2\,\mathrm{kg}$、$k=200\,\mathrm{N/m}$。求临界阻尼，并在 $c=8$ 时求阻尼比、极点和振荡频率。
\item 对 $\dot x=-10x$ 用显式欧拉。求 $h=0.05,0.15,0.25\,\mathrm s$ 的递推因子，判断结果；再用精确离散化求 $h=0.25$ 时的因子。
\end{enumerate}
\subsection*{解答}
\begin{enumerate}
\item $\dot q=0.4-0.1t$，$\ddot q=-0.1$。在 $t=2$ 时，$q=0.1+0.8-0.2=0.7\,\mathrm{rad}$、$\dot q=0.2\,\mathrm{rad/s}$、$\ddot q=-0.1\,\mathrm{rad/s^2}$。角度变化为 $0.7-0.1=0.6\,\mathrm{rad}$，也等于 $\int_0^2(0.4-0.1t)dt$。
\item 平衡值为 $8/4=2$，齐次解为 $Ce^{-4t}$，故 $y=2+Ce^{-4t}$。由 $y(0)=1$ 得 $C=-1$，所以 $y=2-e^{-4t}$。时间常数为 $1/4=0.25\,\mathrm s$；$y(0.5)=2-e^{-2}=1.864665$。
\item $\omega_n=\sqrt{200/2}=10\,\mathrm{rad/s}$；$c_{\rm crit}=2\sqrt{mk}=40\,\mathrm{N\,s/m}$。$c=8$ 时 $\zeta=8/40=0.2$，极点为 $-2\pm j\sqrt{96}$，振荡角频率为 $9.797959\,\mathrm{rad/s}$，频率为 $9.797959/(2\pi)=1.55939\,\mathrm{Hz}$。
\item 递推因子 $a_E=1-10h$ 分别为 $0.5,-0.5,-1.5$。前者单调衰减，中者交替变号衰减，后者交替发散。精确离散因子为 $e^{-10\times0.25}=e^{-2.5}=0.082085$，仍稳定。发散来自积分算法，不来自原连续模型。
\end{enumerate}

\section*{本章与原稿的对应}
本章整合 [A，第3--5页] 与 [B，第17--24页]，补充了链式法则、差分噪声、初值和收敛条件、积分器能量推导、变分法边界项。拉普拉斯是积分变换，拉格朗日是能量建模与变分框架中的名称，两者不能混用。

\chapter{三维位姿 旋转和李群}\label{ch:pose}
\section{同一个点为什么需要坐标系标签}
位置向量必须回答“相对谁的原点、沿谁的坐标轴量”。一台机械臂常同时使用世界系、基座系、法兰系、工具系和相机系。设 $p^B$ 是点 $P$ 在坐标系 $B$ 中的坐标，$R_{AB}$ 把 $B$ 中表达的向量转换为 $A$ 中表达，$p_{AB}$ 是 $B$ 原点在 $A$ 中的位置，则
\begin{equation}
 p^A=R_{AB}p^B+p_{AB}.
\end{equation}
上式描述同一个物理点的两种坐标，并不是让物体真的移动。若把式子用作主动旋转物体，也可以得到相同数值形式，但必须清楚所描述的对象。实际代码中保留 \texttt{R\_base\_tool}、\texttt{p\_base\_tool} 这样的方向信息，比一个无标签的 \texttt{R} 更容易排错。

\begin{examplebox}[算例 把工具上的点换到基座坐标]
工具原点在基座中的位置为 $p_{AB}=(0.3,0.2,0.1)^\mathsf T\,\mathrm m$，工具绕基座 $z$ 轴转了 $90^\circ$，故
\[
 R_{AB}=\begin{bmatrix}0&-1&0\\1&0&0\\0&0&1\end{bmatrix}.
\]
工具上的点 $p^B=(0.1,0,0.2)^\mathsf T\,\mathrm m$，先旋转得到 $(0,0.1,0.2)^\mathsf T$，再平移得到 $p^A=(0.3,0.3,0.3)^\mathsf T\,\mathrm m$。若只是工具 $x$ 轴方向向量 $d^B=(1,0,0)^\mathsf T$，换系结果为 $d^A=R_{AB}d^B=(0,1,0)^\mathsf T$，不能加原点平移。
\end{examplebox}

\begin{figure}[htbp]
\centering\includegraphics[width=0.62\linewidth]{figures/mk_frames.pdf}
\caption{二维投影中，同一点在基座系和工具系有不同坐标。位置变换包含原点平移，方向向量只改变分量表达。图中点取 $z$ 分量为零以便看清平面关系。}
\end{figure}

\section{旋转矩阵的列就是新坐标轴}
$R_{AB}$ 的三列分别是 $B$ 的 $x,y,z$ 单位轴在 $A$ 中的表达。每根轴长度为 1，轴间相互垂直，所以
\begin{equation}
 R^\mathsf TR=I,\qquad R^{-1}=R^\mathsf T.
\end{equation}
此外要求 $\det R=+1$，排除镜像反射。满足这些条件的所有 $3\times3$ 矩阵构成 $SO(3)$，读作“三维特殊正交群”。“特殊”指行列式为正一，“正交”指列向量正交归一。

三个基本旋转矩阵为
\begin{align}
 R_x(\alpha)&=\begin{bmatrix}1&0&0\\0&\cos\alpha&-\sin\alpha\\0&\sin\alpha&\cos\alpha\end{bmatrix},\\
 R_y(\beta)&=\begin{bmatrix}\cos\beta&0&\sin\beta\\0&1&0\\-\sin\beta&0&\cos\beta\end{bmatrix},\\
 R_z(\gamma)&=\begin{bmatrix}\cos\gamma&-\sin\gamma&0\\\sin\gamma&\cos\gamma&0\\0&0&1\end{bmatrix}.
\end{align}
正角度采用右手定则：右手拇指指向正轴，其余手指弯曲方向为正转向。

\subsection{左乘和右乘必须连同轴的归属解释}
若 $R$ 表示物体当前朝向，绕固定参考系的轴再转 $R_f$，更新为 $R_{\rm new}=R_fR$。绕物体当前自身轴转 $R_b$，更新为 $R_{\rm new}=RR_b$。这两种说法不能省掉“固定”和“自身”。
\begin{examplebox}[算例 两次九十度旋转确实不交换]
从向量 $a=(0,1,0)^\mathsf T$ 出发，先绕固定 $x$ 轴转 $90^\circ$，得到 $(0,0,1)^\mathsf T$；再绕固定 $z$ 轴转 $90^\circ$，仍是 $(0,0,1)^\mathsf T$，对应 $R_zR_xa$。
反过来，先绕固定 $z$ 轴得 $(-1,0,0)^\mathsf T$，再绕固定 $x$ 轴不改变它，结果为 $(-1,0,0)^\mathsf T$，对应 $R_xR_za$。因此 $R_zR_x\ne R_xR_z$。
\end{examplebox}

数值积分后若 $\tilde R^\mathsf T\tilde R$ 偏离 $I$，可以用 SVD $\tilde R=U\Sigma V^\mathsf T$ 投影回旋转集合：
\begin{equation}
 R=U\operatorname{diag}(1,1,\det(UV^\mathsf T))V^\mathsf T.
\end{equation}
最后一项用于排除反射。直接使用 $UV^\mathsf T$ 只在其行列式已经为正时合适。更好的日常做法是用保持旋转结构的增量更新，减少这种修正需求。

\section{齐次变换让平移和旋转一起连乘}
把点扩展成 $\bar p=(p_x,p_y,p_z,1)^\mathsf T$，则
\begin{equation}
 T_{AB}=\begin{bmatrix}R_{AB}&p_{AB}\\0&1\end{bmatrix},
 \qquad\bar p^A=T_{AB}\bar p^B.
\end{equation}
下方的 0 代表 $1\times3$ 零行。这样的刚体变换构成 $SE(3)$，即三维特殊欧氏群。连乘满足
\begin{equation}
 T_{AC}=T_{AB}T_{BC}
 =\begin{bmatrix}R_{AB}R_{BC}&p_{AB}+R_{AB}p_{BC}\\0&1\end{bmatrix}.
\end{equation}
这说明后一段平移要先换到前一段坐标，才能相加。不能把不同坐标系中的两个位置向量直接加起来。

\subsection{齐次变换的逆怎样推出来}
从 $p^A=Rp^B+p$ 移项并左乘 $R^\mathsf T$，得 $p^B=R^\mathsf Tp^A-R^\mathsf Tp$，所以
\begin{equation}
 T^{-1}=\begin{bmatrix}R^\mathsf T&-R^\mathsf Tp\\0&1\end{bmatrix}.
\end{equation}
反变换的平移一般不是单纯 $-p$。对前面的 $90^\circ$ 旋转和位置，逆变换平移为
\[
 -R^\mathsf Tp=(-0.2,0.3,-0.1)^\mathsf T.
\]
代入点 $p^A=(0.3,0.3,0.3)^\mathsf T$，先得 $R^\mathsf Tp^A=(0.3,-0.3,0.3)^\mathsf T$，再加逆平移便回到 $p^B=(0.1,0,0.2)^\mathsf T$。

若工具在法兰坐标中的安装位置向量为 $(0.1,0,0)^\mathsf T$，法兰绕基座 $z$ 轴转 $90^\circ$ 后，这个安装偏移会在基座中指向 $+y$。工具中心点（tool center point，TCP）与法兰点因此有不同的位置及线速度雅可比；设置 TCP 是几何问题，设置工具质量和质心是动力学问题，两者不可互相替代。

\section{欧拉角的三次旋转和万向节锁}
本书采用明确的 ZYX 约定：
\begin{equation}
 R=R_z(\psi)R_y(\theta)R_x(\phi),
\end{equation}
$\phi$ 为横滚 roll，$\theta$ 为俯仰 pitch，$\psi$ 为偏航 yaw。它既可读为绕固定轴依次 roll、pitch、yaw，也可读为绕随动自身轴依次 yaw、pitch、roll。只写“先绕 z 再绕 y 再绕 x”而不说明固定或随动，会产生歧义。

在 $\cos\theta\ne0$ 且采用 $\theta\in[-\pi/2,\pi/2]$ 的分支时，可由矩阵反算
\begin{equation}
 \theta=\arcsin(-R_{31}),\quad
 \psi=\operatorname{atan2}(R_{21},R_{11}),\quad
 \phi=\operatorname{atan2}(R_{32},R_{33}).
\end{equation}
$\operatorname{atan2}(y,x)$ 同时利用两个分量判断象限；$\arctan(y/x)$ 不能正确区分 $x,y$ 都变号的情形。实际计算还应将 $\arcsin$ 的输入截断到 $[-1,1]$，避免舍入误差造成无效值。

\subsection{角速度不是三个欧拉角导数的直接拼接}
将各次转动轴转换到参考系，可得空间角速度
\begin{equation}
 \omega^A=E(\theta,\psi)
 \begin{bmatrix}\dot\phi\\\dot\theta\\\dot\psi\end{bmatrix},
 \quad
 E=\begin{bmatrix}
 \cos\psi\cos\theta&-\sin\psi&0\\
 \sin\psi\cos\theta&\cos\psi&0\\
 -\sin\theta&0&1
 \end{bmatrix}.
\end{equation}
第一列是 roll 轴在参考系中的方向，第二列是 pitch 轴，第三列是 yaw 轴。$\det E=\cos\theta$；当 $\theta=\pm\pi/2$，第一与第三列共线，两个角速度参数只能共同描述一个实际转动方向。这是万向节锁（gimbal lock），属于参数化奇异；刚体本身仍有三个物理转动自由度。

\begin{examplebox}[算例 欧拉角变化很大 实际姿态变化却很小]
纯偏航从 $179^\circ$ 变为 $-179^\circ$，直接相减得到 $-358^\circ$。实际相对旋转为
\[
 R_z(-179^\circ)R_z(179^\circ)^\mathsf T
 =R_z(-358^\circ)=R_z(2^\circ).
\]
最短转角只有 $2^\circ=0.0349066\,\mathrm{rad}$。这解释为什么显示角度看起来跳变时，机械臂不一定真的瞬间转了一圈；也说明不能直接用未处理的欧拉角差构造姿态反馈。
\end{examplebox}

\section{轴角 反对称矩阵和罗德里格斯公式}
\subsection{把叉乘写成矩阵乘法}
对 $a=(a_x,a_y,a_z)^\mathsf T$，定义 hat 映射
\begin{equation}
 \widehat a=[a]_\times=
 \begin{bmatrix}0&-a_z&a_y\\a_z&0&-a_x\\-a_y&a_x&0\end{bmatrix}.
\end{equation}
它满足 $\widehat a\,b=a\times b$，其中叉乘方向按右手定则、长度为 $\|a\|\|b\|\sin\gamma$。反操作 vee 映射把此反对称矩阵还原为向量：$(\widehat a)^\vee=a$。这只是同一信息的两种排列，不是新的物理运算。

绕单位轴 $u$ 转角 $\theta$ 的轴角（axis-angle）可合成旋转向量 $\varphi=\theta u$，其长度等于转角。令 $U=\widehat u$。由向量三重积 $u\times(u\times a)=u(u^\mathsf Ta)-a$，有 $U^2=uu^\mathsf T-I$；再乘 $U$ 得 $U^3=-U$。矩阵指数的级数便可按奇偶次幂合并：
\begin{align}
 e^{\theta U}
 &=I+\theta U+\frac{\theta^2}{2!}U^2
       -\frac{\theta^3}{3!}U-\frac{\theta^4}{4!}U^2+\cdots\\
 &=I+\sin\theta\,U+(1-\cos\theta)U^2.
\end{align}
这就是罗德里格斯（Rodrigues）公式。它与矩阵指数相等，不是经验近似。

\begin{examplebox}[算例 同一个一百二十度旋转的矩阵形式]
取 $u=(1,1,1)^\mathsf T/\sqrt3$、$\theta=120^\circ=2\pi/3$。因此 $\sin\theta=\sqrt3/2$、$1-\cos\theta=3/2$，且
\[
 U=\frac1{\sqrt3}\begin{bmatrix}0&-1&1\\1&0&-1\\-1&1&0\end{bmatrix},
 \qquad U^2=\frac13\begin{bmatrix}-2&1&1\\1&-2&1\\1&1&-2\end{bmatrix}.
\]
代入得到
\[
 R=\begin{bmatrix}0&0&1\\1&0&0\\0&1&0\end{bmatrix}.
\]
它把原 $x$ 方向转到 $y$，把 $y$ 转到 $z$，把 $z$ 转到 $x$；轴 $u$ 自己不动。旋转向量三个分量都为 $2\pi/(3\sqrt3)=1.209200\,\mathrm{rad}$。
\end{examplebox}

\begin{figure}[htbp]
\centering\includegraphics[width=0.62\linewidth]{figures/mk_axis_angle.pdf}
\caption{绕 $(1,1,1)/\sqrt3$ 旋转 $120^\circ$ 的精确几何关系。原坐标轴按 $x\to y\to z\to x$ 循环，灰色轴向保持不变。}
\end{figure}

\subsection{从旋转矩阵求回旋转向量}
对 $0<\theta<\pi$，可由
\begin{equation}
 \theta=\arccos\frac{\operatorname{tr}R-1}{2},\qquad
 \varphi=\frac{\theta}{2\sin\theta}(R-R^\mathsf T)^\vee
\end{equation}
求出主值旋转向量。迹 $\operatorname{tr}R$ 是对角线元素之和。上例迹为零，所以 $\theta=\arccos(-1/2)=120^\circ$。当 $\theta\to0$，$\theta/\sin\theta\to1$，应使用级数或稳定算法；当 $\theta\to\pi$，反对称部分也趋于零，需要从对角元素、特征方向或四元数稳定提轴。$\theta=\pi$ 时轴 $u$ 与 $-u$ 给出同一旋转，主值在这里不能全局连续。

\section{四元数为何有四个数却只描述三个自由度}
为了不与关节向量 $q$ 混淆，四元数记为 $\mathsf q=(w,v)$，本章固定标量在前，$v=(x,y,z)^\mathsf T$。单位四元数定义为
\begin{equation}
 \mathsf q=\left(\cos\frac\theta2,\ u\sin\frac\theta2\right),
 \qquad w^2+v^\mathsf Tv=1.
\end{equation}
四个数受一个单位长度约束，所以连续自由度仍为三。$\mathsf q$ 与 $-\mathsf q$ 表示同一物理旋转；这是双覆盖，不是两个不同的姿态。

Hamilton 乘法定义为
\begin{equation}
 (w_1,v_1)\otimes(w_2,v_2)
 =(w_1w_2-v_1^\mathsf Tv_2,
 w_1v_2+w_2v_1+v_1\times v_2).
\end{equation}
单位四元数的逆为共轭 $\mathsf q^*=(w,-v)$。向量 $a$ 经旋转后的结果，是 $\mathsf q\otimes(0,a)\otimes\mathsf q^*$ 的向量部分。展开乘法得旋转矩阵
\begin{equation}
 R=\begin{bmatrix}
 1-2(y^2+z^2)&2(xy-wz)&2(xz+wy)\\
 2(xy+wz)&1-2(x^2+z^2)&2(yz-wx)\\
 2(xz-wy)&2(yz+wx)&1-2(x^2+y^2)
 \end{bmatrix}.
\end{equation}
\begin{examplebox}[算例 用四元数重算上一节的旋转]
$\theta=120^\circ$、$u=(1,1,1)/\sqrt3$ 时，$w=\cos60^\circ=0.5$，三个向量分量均为 $\sin60^\circ/\sqrt3=0.5$。所以 $\mathsf q=(0.5,0.5,0.5,0.5)$。矩阵对角元素为 $1-2(0.25+0.25)=0$；$R_{21}=2(0.25+0.25)=1$，$R_{12}=2(0.25-0.25)=0$。九项代完，恰好回到前面的循环置换矩阵。
\end{examplebox}

传输姿态前必须声明顺序、乘法约定、参考系和主动/被动含义。四个数都在 $[-1,1]$ 内、范数也为 1，并不能证明顺序正确。一个已知 $90^\circ$ 绕单轴旋转的用例，比随机四元数更容易检验接口。

\subsection{插值与角速度积分}
两个单位四元数点积为负时，先把其中一个取负，使插值使用较短路径。令 $\Omega=\arccos(\mathsf q_0^\mathsf T\mathsf q_1)$，球面线性插值（spherical linear interpolation，SLERP）为
\begin{equation}
 \mathsf q(s)=\frac{\sin[(1-s)\Omega]}{\sin\Omega}\mathsf q_0
 +\frac{\sin(s\Omega)}{\sin\Omega}\mathsf q_1,
 \qquad 0\le s\le1.
\end{equation}
这里 $\Omega$ 是四维单位球面上的夹角，物理最短转角为 $2\Omega$。当 $\Omega$ 很小时，使用归一化线性插值可避免除以小数。

从单位姿态到 $R_z(90^\circ)$，$s=1/2$ 给出 $R_z(45^\circ)$，四元数为 $(\cos22.5^\circ,0,0,\sin22.5^\circ)\approx(0.923880,0,0,0.382683)$。若本体系角速度为 $\omega^B$，时间步为 $h$，令 $\delta=\|\omega^B\|h$，则右乘增量
\begin{equation}
 \mathsf q_{k+1}=\mathsf q_k\otimes
 \left(\cos\frac\delta2,
 \frac{\omega^B}{\|\omega^B\|}\sin\frac\delta2\right).
\end{equation}
$\omega^B=0$ 时增量取单位四元数。若角速度在参考系表达，则要左乘相应增量。数值误差可通过归一化控制，但归一化不能修正坐标系或乘法顺序错误。

\section{李群和李代数是在结构内做小变化}
群（group）表示一种组合运算具有封闭性、单位元、逆和结合律；旋转连续地组合仍是旋转，单位旋转是不动，每个旋转都可转回去。李群（Lie group）还具有光滑结构，因此能讨论局部小变化。李代数（Lie algebra）是单位元处的切空间；$SO(3)$ 的李代数由 $3\times3$ 反对称矩阵组成。

为何是反对称矩阵？对 $R(t)R(t)^\mathsf T=I$ 求导：
\begin{equation}
 \dot R R^\mathsf T+R\dot R^\mathsf T=0.
\end{equation}
因此 $\dot R R^\mathsf T$ 与自己的转置互为相反数，可写成 $\widehat\omega^A$。得到
\begin{equation}
 \dot R=\widehat\omega^A R=R\widehat\omega^B,
 \qquad\omega^A=R\omega^B.
\end{equation}
指数映射 $\operatorname{Exp}(\widehat\varphi)$ 把小变化或有限旋转向量送到旋转矩阵；对数映射 $\operatorname{Log}(R)$ 做局部逆运算。有限旋转向量也能被当作普通三维数字存储，但直接相加并不等于组合两个有限旋转。

Baker--Campbell--Hausdorff（BCH）展开的前几项说明差别：
\begin{equation}
 \operatorname{Log}\left(e^{\widehat a}e^{\widehat b}\right)^\vee
 =a+b+\tfrac12(a\times b)+\text{更高阶项}.
\end{equation}
若 $a=(0.1,0,0)^\mathsf T$、$b=(0,0.1,0)^\mathsf T$，相加只得到 $(0.1,0.1,0)$，加入二阶项得 $(0.1,0.1,0.005)$。多出来的 $z$ 分量反映旋转顺序，而非传感器噪声。

\subsection{姿态误差同时包含方向 符号和坐标系}
设当前工具朝向 $R$、期望朝向 $R_d$ 均将各自局部分量转换到同一基座系。为与“实际减期望”的位置误差一致，定义基座表达的旋转误差
\begin{equation}
 e_R^A=\operatorname{Log}(RR_d^\mathsf T)^\vee,
 \qquad
 e_R^D=\operatorname{Log}(R_d^\mathsf TR)^\vee,
 \qquad e_R^A=R_d e_R^D.
\end{equation}
第二个式子可由 $R=R_d\operatorname{Exp}(\widehat e_R^D)$ 理解，它以期望坐标为右扰动的参考。若改用 $\operatorname{Log}(R^\mathsf TR_d)^\vee$，则表示从当前朝向走向期望朝向的增量，符号反转。只写一个三维误差而省略约定，会使后面的恢复力矩方向难以核对。

例如 $R=R_z(10^\circ)$、$R_d=I$，则 $e_R^A=(0,0,0.174533)^\mathsf T$。恢复方向是负 $z$ 转动。对于一般大姿态误差，$\dot e_R$ 并不严格等于 $\omega-\omega_d$；中间还存在对数映射的雅可比。小误差时该雅可比接近单位阵，这也是很多工程控制式的局部近似基础。

对固定 $R_d$，若 $R R_d^\mathsf T=\operatorname{Exp}(\widehat e)$，有 $\omega^A=J_l(e)\dot e$。$J_l$ 称为 $SO(3)$ 左雅可比：
\begin{equation}
 J_l(e)=I+\frac{1-\cos\theta}{\theta^2}\widehat e
 +\frac{\theta-\sin\theta}{\theta^3}\widehat e^{\,2},\qquad\theta=\|e\|.
\end{equation}
这里的 $J_l$ 是旋转参数映射的雅可比，不是机械臂关节到末端的雅可比。理解它们都来自“局部导数”，同时保留不同下标，是避免混淆的关键。

\section{位姿旋量和伴随变换}
\subsection{六维数的前后顺序和作用点必须明确}
本书六维速度统一采用线速度在前、角速度在后，即 $V=(v^\mathsf T,\omega^\mathsf T)^\mathsf T$。但仅有顺序还不够：线速度是哪个点的？由 $T=\begin{bmatrix}R&p\\0&1\end{bmatrix}$，空间旋量与本体旋量分别定义为
\begin{equation}
 \widehat V_s=\dot TT^{-1},\qquad
 \widehat V_b=T^{-1}\dot T.
\end{equation}
直接相乘得到
\begin{equation}
 V_s=\begin{bmatrix}\dot p-\omega^A\times p\\\omega^A\end{bmatrix},
 \qquad
 V_b=\begin{bmatrix}R^\mathsf T\dot p\\R^\mathsf T\omega^A\end{bmatrix}.
\end{equation}
空间旋量的前三项一般不是工具原点的速度 $\dot p$。若一个接口返回基座方向表达、但以工具原点为参考点的几何速度 $(\dot p,\omega^A)$，它与上面的空间旋量不同。许多“WORLD”与“工具原点对齐”接口的差异正发生在这里。

两种旋量用伴随矩阵（adjoint matrix）转换：
\begin{equation}
 V_s=\operatorname{Ad}_T V_b,
 \qquad
 \operatorname{Ad}_T=\begin{bmatrix}R&\widehat pR\\0&R\end{bmatrix}.
\end{equation}
如果只是同一个物理点的线速度和角速度换表达坐标，不移动参考点，则只需块对角旋转 $\operatorname{diag}(R,R)$；不要额外加入 $\widehat pR$。

\subsection{六维指数中的平移不是简单照抄前三项}
对积分旋量 $\xi=(\rho,\varphi)$，hat 矩阵定义为 $\widehat\xi=\begin{bmatrix}\widehat\varphi&\rho\\0&0\end{bmatrix}$。其指数是
\begin{equation}
 \operatorname{Exp}(\widehat\xi)
 =\begin{bmatrix}\operatorname{Exp}(\widehat\varphi)&J_l(\varphi)\rho\\0&1\end{bmatrix}.
\end{equation}
$\rho$ 是旋量坐标，不总等于最终位置 $p$。因为在转动过程中，局部平移方向也在改变。
\begin{examplebox}[算例 一边前进一边转弯的一秒运动]
从单位位姿开始，以恒定本体速度 $v^B=(1,0,0)^\mathsf T\,\mathrm{m/s}$、角速度 $\omega^B=(0,0,\pi/2)^\mathsf T\,\mathrm{rad/s}$ 运动 1 s。此时 $\rho=(1,0,0)^\mathsf T$、$\varphi=(0,0,\pi/2)^\mathsf T$。指数给出的位移为
\[
 p=\begin{bmatrix}\sin(\pi/2)/(\pi/2)\\(1-\cos(\pi/2))/(\pi/2)\\0\end{bmatrix}
 =\begin{bmatrix}2/\pi\\2/\pi\\0\end{bmatrix}\mathrm m.
\]
物体走了半径 $2/\pi\,\mathrm m$ 的四分之一圆弧，最后朝向转过 $90^\circ$。如果把平移直接写成 $(1,0,0)$，就把转弯错误地算成了直走。
\end{examplebox}

\begin{figure}[htbp]
\centering\includegraphics[width=0.62\linewidth]{figures/mk_screw.pdf}
\caption{恒定本体速度与恒定角速度形成圆弧，箭头显示局部前进轴。位姿指数同时积分了旋转与平移的耦合。}
\end{figure}

\subsection{力旋量用功率不变决定变换规则}
力旋量记为 $F=(f,m)$，前面是力、后面是关于指定参考点的力矩。瞬时功率为 $P=F^\mathsf TV=f^\mathsf Tv+m^\mathsf T\omega$。若 $V_A=\operatorname{Ad}_T V_B$，要求两种表达的功率相等，便有
\begin{equation}
 F_A=\operatorname{Ad}_T^{-\mathsf T}F_B.
\end{equation}
展开为 $f_A=Rf_B$，$m_A=p\times Rf_B+Rm_B$。例如作用点距参考原点 $p=(0.3,0,0)^\mathsf T\,\mathrm m$，力为 $f=(0,10,0)^\mathsf T\,\mathrm N$，该点自身力矩为零，则关于参考原点的力矩为 $(0,0,3)^\mathsf T\,\mathrm{N\,m}$。换参考点时漏掉这一项，会让负载与外力估计出现姿态相关误差。

\section{小扰动导数和姿态平均}
若旋转后的点为 $b=Ra$，对右扰动 $R'=R\operatorname{Exp}(\widehat{\delta\varphi})$，一阶展开
\[
 b'\approx R(I+\widehat{\delta\varphi})a
 =Ra-R\widehat a\,\delta\varphi.
\]
所以 $\partial b/\partial\delta\varphi=-R\widehat a$。对左扰动则得到 $-\widehat{Ra}$。标定求解、姿态估计和数值 IK 里使用的导数必须与扰动定义一致。

多个相近姿态的平均，也可先在当前估计 $\bar R$ 的切空间中求平均误差 $\bar e=N^{-1}\sum_i\operatorname{Log}(\bar R^\mathsf TR_i)^\vee$，再更新 $\bar R\leftarrow\bar R\operatorname{Exp}(\widehat{\bar e})$，重复到小误差。姿态分散接近半圈时，均值可能不唯一或依赖初始化；“先平均四元数再归一化”若不处理符号，也会把相同姿态的 $\mathsf q$ 与 $-\mathsf q$ 平均成零。

\section{计算练习与逐步解答}
\subsection*{题目}
\begin{enumerate}
\item $R_{AB}=R_z(90^\circ)$、$p_{AB}=(1,0,0)^\mathsf T\,\mathrm m$，点 $p^B=(0,2,0)^\mathsf T\,\mathrm m$。求 $p^A$ 与逆变换平移。
\item 绕单位轴 $(0,1,0)^\mathsf T$ 转 $60^\circ$，求标量在前的单位四元数与旋转矩阵，检验它对 $(1,0,0)^\mathsf T$ 的作用。
\item 当前 $R=R_z(-170^\circ)$、期望 $R_d=R_z(170^\circ)$。按本书实际减期望的定义求主值姿态误差；若以恒定 $10^\circ/\mathrm s$ 沿最短方向回到期望，需多久？
\item 一作用点相对基座原点为 $(0,0.2,0)^\mathsf T\,\mathrm m$，施加力 $(5,0,0)^\mathsf T\,\mathrm N$，作用点自身无力矩。求关于基座原点的力矩。
\end{enumerate}
\subsection*{解答}
\begin{enumerate}
\item $Rp^B=(-2,0,0)^\mathsf T$，加平移得到 $p^A=(-1,0,0)^\mathsf T\,\mathrm m$。逆平移为 $-R^\mathsf Tp_{AB}=-(0,-1,0)^\mathsf T=(0,1,0)^\mathsf T\,\mathrm m$。使用逆变换可恢复原点坐标。
\item $\mathsf q=(\cos30^\circ,0,\sin30^\circ,0)=(\sqrt3/2,0,1/2,0)$。$R_y(60^\circ)=\begin{bmatrix}1/2&0&\sqrt3/2\\0&1&0\\-\sqrt3/2&0&1/2\end{bmatrix}$。输入 $(1,0,0)^\mathsf T$ 得 $(1/2,0,-\sqrt3/2)^\mathsf T$，长度仍为 1。
\item $RR_d^\mathsf T=R_z(-340^\circ)=R_z(20^\circ)$，主值误差为 $(0,0,\pi/9)^\mathsf T$，即正 $20^\circ$。回到期望应沿负 $z$ 转 $20^\circ$，用时 $20/10=2\,\mathrm s$；这与直接相减得到的 $-340^\circ$ 不同。
\item 叉乘 $p\times f=(0,0,-0.2\cdot5)^\mathsf T=(0,0,-1)^\mathsf T\,\mathrm{N\,m}$。力矩方向由右手定则确定，不能只写正的“力乘距离”。
\end{enumerate}

\section*{本章与原稿的对应}
本章整合 [A，第22--26页] 与 [B，第3--8、24--27页]。所有旋转约定均在正文中固定，并补充对数误差符号、旋量参考点、力旋量对偶变换和位姿指数的数字推导，避免把几何速度、空间旋量和欧拉角导数混为一谈。

\chapter{从关节角到末端位姿 再求逆运动学}\label{ch:kinematics}
\section{几何模型回答能到哪里}
运动学（kinematics）研究位置、速度、加速度之间的关系，不直接讨论产生运动所需的力。正运动学（forward kinematics，FK）从关节量求末端位姿；逆运动学（inverse kinematics，IK）反过来求关节量。动力学（dynamics）随后再回答需要多少力矩。一个位置在几何上可达，并不意味着可以以任意速度到达，更不意味着驱动器能承担所需负载。

旋转关节用 R 表示，关节量是角度；移动关节用 P 表示，关节量是距离。自由度（degree of freedom，DoF）是局部独立运动变量的个数。空间刚体最多需要三个平移量和三个旋转量；要在某个姿态附近任意改变六维末端位姿，任务雅可比至少要有六个独立行方向。拥有六个关节只是必要的维数条件之一，不保证所有位姿可达。

本章使用两连杆平面机械臂，参数固定为
\begin{equation}
 l_1=l_2=0.5\,\mathrm m,
 \quad q_1\text{ 从基座 }+x\text{ 轴逆时针量起},
 \quad q_2\text{ 相对第一杆方向量起}.
\end{equation}
第二杆相对基座的绝对方向为 $q_1+q_2$。这里两个关节足以控制平面末端的二维位置，但一般不能独立再指定末端朝向。

\section{两连杆正运动学从投影逐项写出}
第一杆端点为
\begin{equation}
 p_1=\begin{bmatrix}l_1\cos q_1\\l_1\sin q_1\end{bmatrix}.
\end{equation}
第二杆相对第一个关节点的位移为
\begin{equation}
 p_2-p_1=\begin{bmatrix}l_2\cos(q_1+q_2)\\l_2\sin(q_1+q_2)\end{bmatrix}.
\end{equation}
两段位移都已在基座表达，因此可以直接相加：
\begin{equation}\label{eq:two_r_fk}
 p(q)=\begin{bmatrix}x\\y\end{bmatrix}
 =\begin{bmatrix}
 l_1\cos q_1+l_2\cos(q_1+q_2)\\
 l_1\sin q_1+l_2\sin(q_1+q_2)
 \end{bmatrix},\qquad\theta_e=q_1+q_2.
\end{equation}
\begin{examplebox}[算例 完整算一次关节角到末端]
在 $q_1=0$、$q_2=\pi/2$ 时，第一杆端点为 $(0.5,0)$，第二杆增量为 $(0,0.5)$，故末端为 $(0.5,0.5)\,\mathrm m$，朝向为 $90^\circ$。
若 $q_1=30^\circ$、$q_2=60^\circ$，第二杆绝对朝向仍为 $90^\circ$，于是
\[
 x=0.5\frac{\sqrt3}{2}+0=0.433013\,\mathrm m,
 \qquad y=0.5\frac12+0.5=0.75\,\mathrm m.
\]
两组关节角可给出相同末端朝向，却有不同末端位置。只看末端朝向不能判断机械臂是否回到原处。
\end{examplebox}

\subsection{工作空间是关节允许范围映射出的集合}
当两个关节都可自由转动、没有障碍和限位时，末端到基座距离 $r$ 满足
\begin{equation}
 |l_1-l_2|\le r\le l_1+l_2.
\end{equation}
外边界是完全伸直，内边界是完全折叠。相等长度 $0.5\,\mathrm m$ 时，位置工作空间是半径 1 m 的闭圆盘；若长度为 $0.6$ 与 $0.4\,\mathrm m$，则是内半径 0.2 m、外半径 1 m 的圆环。加入关节限位、自碰撞、底座和工具体积后，可用工作空间会缩小，而且可能不再是简单圆环。

\begin{figure}[htbp]
\centering\includegraphics[width=0.9\linewidth]{figures/mk_workspace.pdf}
\caption{左图是两连杆的角度定义；右图中外圆是两种长度的共同外界，虚线只是不等长连杆的内界；等长连杆仍能到达虚线内的区域。图中尚未加入限位和碰撞，因此不能直接当作真机可用工作区。}
\end{figure}

\section{标准 DH 只是整理变换的一种约定}
Denavit--Hartenberg 参数简称 DH 参数，用四个数描述按规则建立的相邻坐标系。标准 DH 本书采用
\begin{equation}
 T_{i-1,i}=R_z(\theta_i)\,T_z(d_i)\,T_x(a_i)\,R_x(\alpha_i).
\end{equation}
$\theta_i$ 是绕前一坐标系 $z$ 轴的角度，$d_i$ 是沿该轴的偏移，$a_i$ 是沿连杆公垂线的长度，$\alpha_i$ 是两关节轴间的扭角。旋转关节通常以 $\theta_i=q_i+\theta_{i,0}$ 为变量，移动关节通常以 $d_i$ 为变量。按四个基本变换相乘：
\begin{equation}
 T_{i-1,i}=\begin{bmatrix}
 c_\theta&-s_\theta c_\alpha&s_\theta s_\alpha&a_ic_\theta\\
 s_\theta&c_\theta c_\alpha&-c_\theta s_\alpha&a_is_\theta\\
 0&s_\alpha&c_\alpha&d_i\\
 0&0&0&1
 \end{bmatrix},
\end{equation}
其中 $c_\theta=\cos\theta_i$、$s_\theta=\sin\theta_i$，其余同理。这些缩写只为让矩阵排版清楚，不是新增变量。

两连杆采用 $d_i=0$、$\alpha_i=0$、$a_i=l_i$、$\theta_i=q_i$，于是
\begin{equation}
 T_{01}=\begin{bmatrix}R_z(q_1)&(l_1\cos q_1,l_1\sin q_1,0)^\mathsf T\\0&1\end{bmatrix},
\end{equation}
第二段形式相同，连乘后的旋转为 $R_z(q_1+q_2)$，平移恰好为式~\eqref{eq:two_r_fk}。由此可以检验 DH 表是否与几何直觉一致。

改进 DH（modified DH，MDH）采用不同的坐标放置和变换顺序，同一机械臂的数字不能直接塞进标准 DH 的矩阵。URDF 则通常直接给出关节固定原点变换和轴方向；它不要求用 DH 参数表达。通用链条可以写为
\begin{equation}
 T_{\rm base,tool}(q)=T_{\rm base,0}
 \prod_{i=1}^n T_{i-1,i}(q_i)\,T_{n,\rm tool}.
\end{equation}
固定基座安装与固定工具安装都是链条的一部分。末端误差也可能来自这两处，而不一定是关节控制误差。

\section{指数积公式用零位关节轴描述整条链}
对一个旋转轴，设单位方向为 $\omega_i$，轴上任一点为 $r_i$，则沿本书线速度在前的约定，零位螺旋轴为
\begin{equation}
 S_i=\begin{bmatrix}-\omega_i\times r_i\\\omega_i\end{bmatrix}.
\end{equation}
因为绕固定轴转动时，空间点 $p$ 的速度是 $\omega_i\times(p-r_i)=(-\omega_i\times r_i)+\omega_i\times p$。对移动关节，取 $S_i=(v_i,0)$，其中 $v_i$ 为单位移动方向。若所有轴均在零位参考系描述，末端零位位姿为 $T_0$，正运动学为
\begin{equation}
 T(q)=e^{\widehat S_1q_1}e^{\widehat S_2q_2}\cdots e^{\widehat S_nq_n}T_0.
\end{equation}
称为指数积（product of exponentials，PoE）公式。这里使用 $T_0$ 表示零位位姿，以免与后续惯量矩阵 $M$ 混淆。

\begin{examplebox}[算例 两连杆不用 DH 也算到同一个点]
两轴都沿 $+z$。第一轴经过原点，第二轴零位经过 $(0.5,0,0)$，所以
\[
 S_1=(0,0,0,0,0,1)^\mathsf T,
 \quad S_2=(0,-0.5,0,0,0,1)^\mathsf T,
 \quad T_0=\begin{bmatrix}I&(1,0,0)^\mathsf T\\0&1\end{bmatrix}.
\]
当 $q_1=0,q_2=\pi/2$ 时，第二轴的变换绕点 $r=(0.5,0,0)$ 旋转。其旋转为 $R_z(90^\circ)$，平移为 $(I-R)r=(0.5,-0.5,0)^\mathsf T$。乘零位末端位置得到
\[
 p=R(1,0,0)^\mathsf T+(0.5,-0.5,0)^\mathsf T=(0.5,0.5,0)^\mathsf T.
\]
与几何法和 DH 法完全一致。不同方法只是模型表达不同，不应产生不同的物理末端位置。
\end{examplebox}

\section{解析逆运动学从余弦定理推起}
\subsection{先求肘关节 再求肩关节}
对式~\eqref{eq:two_r_fk} 平方相加，用 $\cos(a-b)=\cos a\cos b+\sin a\sin b$，可得
\begin{align}
 x^2+y^2
 &=l_1^2+l_2^2
   +2l_1l_2[\cos q_1\cos(q_1+q_2)+\sin q_1\sin(q_1+q_2)]\\
 &=l_1^2+l_2^2+2l_1l_2\cos q_2.
\end{align}
因此
\begin{equation}
 c_2=\cos q_2=\frac{x^2+y^2-l_1^2-l_2^2}{2l_1l_2},
 \qquad
 q_2=\operatorname{atan2}\left(\pm\sqrt{1-c_2^2},c_2\right).
\end{equation}
若 $|c_2|>1$ 且超过数值容差，就没有实数关节解。$\pm$ 给出两个肘部构型分支。

将末端向量分解成
\begin{equation}
 \begin{bmatrix}x\\y\end{bmatrix}
 =R_z^{(2)}(q_1)
 \begin{bmatrix}l_1+l_2\cos q_2\\l_2\sin q_2\end{bmatrix},
\end{equation}
其中 $R_z^{(2)}$ 是二维旋转矩阵。末端方向角记为 $\gamma$，括号向量相对第一杆的方向角记为 $\beta$，分别为
\[
 \gamma=\operatorname{atan2}(y,x),\qquad
 \beta=\operatorname{atan2}(l_2\sin q_2,l_1+l_2\cos q_2).
\]
相减 $q_1=\gamma-\beta$，便得
\begin{equation}
 q_1=\operatorname{atan2}(y,x)
 -\operatorname{atan2}(l_2\sin q_2,l_1+l_2\cos q_2).
\end{equation}
这个推导解释每一个反三角项对应哪一个几何角，不必死记符号。

\begin{examplebox}[算例 目标位置有两组关节解]
目标 $p_d=(0.5,0.5)^\mathsf T\,\mathrm m$，有 $x^2+y^2=0.5$，$c_2=(0.5-0.25-0.25)/0.5=0$。所以 $q_2=\pm90^\circ$。
正分支中，$q_1=45^\circ-\operatorname{atan2}(0.5,0.5)=0$；负分支中，$q_1=45^\circ-\operatorname{atan2}(-0.5,0.5)=90^\circ$。
两组解为 $(0^\circ,90^\circ)$ 与 $(90^\circ,-90^\circ)$。前者末端朝向 $90^\circ$，后者朝向 $0^\circ$。因此“同一末端位置”与“同一末端完整位姿”不是同一约束。
\end{examplebox}

\begin{figure}[htbp]
\centering\includegraphics[width=0.68\linewidth]{figures/mk_ik_branches.pdf}
\caption{相同目标位置的两个逆解分支。两根末端小箭头显示不同末端朝向。额外指定朝向可能排除一个分支，也可能使任务完全不可达。}
\end{figure}

\begin{examplebox}[算例 一个非整角目标逐步代数值]
目标 $(0.6,0.4)^\mathsf T\,\mathrm m$，$c_2=(0.36+0.16-0.5)/0.5=0.04$。所以 $q_2=\pm87.707557^\circ$。末端方向为 $\operatorname{atan2}(0.4,0.6)=33.690068^\circ$；正分支的第二项为 $43.853779^\circ$，故 $q_1=-10.163711^\circ$。另一组为 $(77.543846^\circ,-87.707557^\circ)$。把每组代回正运动学，都恢复 $(0.6,0.4)$，这是逆解计算必须执行的闭环核对。
\end{examplebox}

\subsection{二维位置可达不代表二维位姿可达}
仍要求位置 $(0.5,0.5)$，若额外指定朝向 $45^\circ$，则第二杆必须等于 $0.5(\cos45^\circ,\sin45^\circ)$。第一杆端点应为
\[
 (0.5,0.5)-0.5(\sqrt2/2,\sqrt2/2)
 =(0.146447,0.146447)\,\mathrm m.
\]
它到基座的距离为 $0.207107\,\mathrm m$，不等于第一杆固定长度 $0.5\,\mathrm m$，因此无解。这是任务维数过多的具体表现。把一个二维可达位置提交给带姿态约束的 IK 求解器，返回失败并不必然是算法不好。

\section{数值逆解是反复使用局部近似}
\subsection{先定义残差 再明确更新方向}
位置残差定义为 $r=p_d-p(q)$，采用“目标减当前”，仅用于本节求修正量。由一阶展开 $p(q+\Delta q)\approx p(q)+J(q)\Delta q$，希望 $J\Delta q\approx r$。在方阵可逆且远离奇异时，$\Delta q=J^{-1}r$；对非方阵可用伪逆。一次修正后必须重新计算非线性的 $p$ 和 $J$，不能把第一步的雅可比一直重复使用。

\begin{examplebox}[算例 一次牛顿修正并不直接等于最终答案]
从 $q=(0,\pi/2)$ 出发，要移到 $(0.45,0.55)\,\mathrm m$。当前 $p=(0.5,0.5)$，残差 $r=(-0.05,0.05)^\mathsf T$，雅可比为
\[
 J=\begin{bmatrix}-0.5&-0.5\\0.5&0\end{bmatrix}.
\]
解 $J\Delta q=r$：第二行给出 $\Delta q_1=0.1$，第一行给出 $\Delta q_2=0$。新关节角 $(0.1,\pi/2)$ 的真实末端为 $(0.447585,0.547419)\,\mathrm m$，仍有约 $3.535\,\mathrm{mm}$ 位置误差。局部线性近似消除了主要误差，却没有消除二阶曲率误差，必须继续迭代。
\end{examplebox}

\subsection{阻尼最小二乘从代价函数推导}
接近奇异时，伪逆中的 $1/\sigma$ 可能很大。改为求解
\begin{equation}
 \min_{\Delta q}\ \|J\Delta q-r\|_2^2+\lambda^2\|\Delta q\|_2^2,
 \qquad\lambda>0.
\end{equation}
第一项要求完成任务，第二项惩罚过大的关节修正。求导等于零得到
\begin{equation}
 (J^\mathsf TJ+\lambda^2I)\Delta q=J^\mathsf Tr.
\end{equation}
等价形式为
\begin{equation}\label{eq:dls}
 \Delta q=J^\mathsf T(JJ^\mathsf T+\lambda^2I)^{-1}r.
\end{equation}
因为 $(J^\mathsf TJ+\lambda^2I)J^\mathsf T=J^\mathsf T(JJ^\mathsf T+\lambda^2I)$，两种表达作用相同；当任务维度比关节数小，后一形式求解更小的线性方程。实际实现仍使用求解器，不显式取逆。

SVD 下，每个奇异方向的放大系数由 $1/\sigma$ 变为 $\sigma/(\sigma^2+\lambda^2)$。当 $\sigma\to0$，它趋于零，因此不会产生无穷大修正；代价是该方向的任务误差保留下来。阻尼不是让不可达方向变得可达。

对上一个数字例子取 $\lambda=0.05$，得到 $\Delta q=(0.0990195,0.0009708)^\mathsf T\,\mathrm{rad}$，更新后末端为 $(0.447639,0.546931)\,\mathrm m$，误差约 $3.872\,\mathrm{mm}$。它比普通逆的一步误差略大，换来了限制放大能力。$\lambda$ 必须与使用的任务尺度一致；若米换成毫米，不能保持所有权重数值不变就认为算法等价。

\subsection{完整迭代流程}
一个能解释的局部 IK 过程需要以下环节：
\begin{enumerate}
\item 以当前或上一时刻关节角为初值，计算当前正运动学。
\item 计算与任务定义一致的残差，例如位置用米、姿态用旋转对数，并明确两者的权重。
\item 若位置与姿态分别达到容差，则接受结果；不能只把混合单位的一个范数当作唯一标准。
\item 计算相同坐标约定的残差导数或局部几何雅可比，解 DLS 或受约束子问题。
\item 限制步长或做线搜索，确保残差实际减小，同时检查限位及碰撞。
\item 更新关节角，重新计算。超出迭代次数、步长过小却误差不降、或任务不可行时，明确报告失败。
\end{enumerate}
对于完整 $SE(3)$ 对数误差，远离目标时误差导数一般还包含对数映射雅可比。不能把任意姿态误差直接塞给任意六维雅可比而不核对。位置-only 例子不涉及这一额外旋转参数导数，因此适合先完整掌握。

\begin{figure}[htbp]
\centering\includegraphics[width=0.9\linewidth]{figures/mk_ik_convergence.pdf}
\caption{从 $(0,\pi/2)$ 求解目标 $(0.45,0.55)$ 的 DLS 迭代。左图为真实末端逐步靠近目标；右图使用对数纵轴显示每步重新计算后的残差。}
\end{figure}

上一时刻解作为初值通常促进连续性，但不保证连续：路径穿越奇异位形、碰撞约束变化、任务突然跳变或某分支碰到限位时，局部解仍可能丢失。必要时应在规划层改变路径、任务优先级或姿态约束。


\subsection{可达目标也可能让局部迭代原地不动}
一个容易误判的例子：等长两连杆从 $q=(0,0)$ 完全伸直出发，目标是 $(0.9,0)\,\mathrm m$。目标在工作圆盘内，解析 IK 有解，但当前位置为 $(1,0)$，残差 $r=(-0.1,0)$，雅可比为 $\begin{bmatrix}0&0\\1&0.5\end{bmatrix}$。因此 $J^\mathsf Tr=0$，无论普通最小二乘还是加正阻尼，局部一步都会得到 $\Delta q=0$。

这不是目标不存在，而是当前一阶模型看不见向内的二阶运动。若令 $q_1=\alpha$、$q_2=-2\alpha$，两杆绝对朝向为 $\alpha$ 与 $-\alpha$，末端为 $(\cos\alpha,0)$。取 $\cos\alpha=0.9$，有 $\alpha=25.841933^\circ$，便得到目标。这个有限动作的最初向内位移与 $\alpha^2$ 同阶，所以初始一阶梯度为零。可采用合适的非奇异初值、多起点、带几何信息的搜索或重新设计起始路径；单纯增大迭代次数不会让零修正自动变成正确方向。

\section{把限位纳入求解 而不是算完后盲目截断}
静态完整 IK 的约束 $p(q)=p_d$、姿态约束和碰撞距离一般是非线性的。所谓“QP 逆解”通常指每一步线性化后的二次规划（quadratic programming），不是整个非线性 IK 天然就是一个 QP。局部问题可写成
\begin{equation}
 \begin{aligned}
 \min_{\Delta q}\quad&\tfrac12\|J\Delta q-r\|_W^2
                  +\tfrac12\lambda^2\|\Delta q\|^2,\\
 \text{满足}\quad&q_{\min}-q\le\Delta q\le q_{\max}-q,\\
                 &-\Delta q_{\max}\le\Delta q\le\Delta q_{\max}.
 \end{aligned}
\end{equation}
$\|z\|_W^2=z^\mathsf TWz$，其中 $W$ 是对称半正定任务权重。第一组不等式保证下一步位置不越限，第二组限制步长。碰撞距离也可在当前姿态附近线性化，但需要信赖域和后验检查以控制近似误差。

若普通求解给出 $\Delta q$ 后，逐轴任意截断，原先 $J\Delta q\approx r$ 的关系可能被破坏；如果采用带约束求解，算法会在可行集合中寻找最佳折中。遇到无解，应减小任务、重新规划或停止推进，不应把任意截断结果宣称为“已到达”。

\section{选解和模型校验要对应实际任务}
多组解析逆解可按关节范围、碰撞、与当前角度距离、奇异程度及后续路径连续性筛选。关节距离要考虑周期关节的等价角，以及实际机械限位。六维任务中位置误差和姿态误差应分别报告，例如 $\|p-p_d\|$ 与 $\|\operatorname{Log}(R_d^\mathsf TR)^\vee\|$，而不是只报“综合误差 $10^{-4}$”。

一条有用的几何验证链是：给定关节角，用独立几何公式或厂家已核实的姿态求 FK；对所得末端再做 IK；对每个候选解再次做 FK；最后确认关节限位和工具坐标。比较 FK 时，必须使用相同零位、关节正方向、安装姿态和 TCP。模型与真机相差一毫米，不应先用控制积分项掩盖，因为错误可能来自几何链条。

\begin{examplebox}[算例 关节零偏如何变成端点误差]
在标准姿态 $q=(0,\pi/2)$，若第一关节有 $\delta q_1=0.01\,\mathrm{rad}$ 零偏、第二关节正确，局部位置误差为
\[
 \delta p\approx J\begin{bmatrix}0.01\\0\end{bmatrix}
 =\begin{bmatrix}-0.005\\0.005\end{bmatrix}\mathrm m,
\]
长度约为 $7.071\,\mathrm{mm}$。真实 FK 差为约 $(-5.025,4.975)\,\mathrm{mm}$，与线性估计接近。关节角看起来只差 $0.573^\circ$，末端已经出现毫米级误差。
\end{examplebox}

\section{计算练习与逐步解答}
\subsection*{题目}
\begin{enumerate}
\item 对等长两连杆，计算 $q=(45^\circ,0^\circ)$ 的末端位置与朝向。再计算 $q=(45^\circ,180^\circ)$ 的末端位置。
\item 目标为 $(1.1,0)\,\mathrm m$ 时，用解析 IK 判断是否可达。再考虑目标 $(0.5,0.5)\,\mathrm m$，关节范围为
\[
q_1\in[-45^\circ,45^\circ],\qquad q_2\in[-100^\circ,100^\circ].
\]
根据这两个限位，筛选本章的两组逆解。
\item 当前 $q=(0,\pi/2)$，目标位置微增量为 $(0.01,0.02)^\mathsf T\,\mathrm m$。用当前雅可比求无阻尼关节增量，并给出一阶重构结果。
\item 法兰位姿为 $R_z(90^\circ)$、$p_f=(0.5,0.5,0)^\mathsf T\,\mathrm m$，TCP 在法兰内为 $(0.1,0,0)^\mathsf T\,\mathrm m$。求 TCP 在基座位置，说明是否可以直接加到基座 $x$ 分量。
\end{enumerate}
\subsection*{解答}
\begin{enumerate}
\item 第一种完全伸直，$p=(\cos45^\circ,\sin45^\circ)=(\sqrt2/2,\sqrt2/2)\,\mathrm m$，朝向 $45^\circ$。第二种完全折叠，两根等长杆方向相反，位置为 $(0,0)$，朝向为 $225^\circ$。在原点位置不能唯一判断 $q_1$，这是特殊退化情况。
\item 第一目标 $c_2=(1.21-0.5)/0.5=1.42>1$，不可达。第二目标的 $(0^\circ,90^\circ)$ 在给定范围内；$(90^\circ,-90^\circ)$ 超出第一关节范围，排除。限位筛选只说明这些角度可用，还未检查碰撞与动态约束。
\item 第二行 $0.5\Delta q_1=0.02$，得 $\Delta q_1=0.04$。第一行 $-0.5(\Delta q_1+\Delta q_2)=0.01$，得 $\Delta q_2=-0.06$，单位 rad。代回 $J\Delta q=(0.01,0.02)^\mathsf T\,\mathrm m$。有限更新后真实位置仍有二阶误差，因此这只是一次局部修正。
\item 旋转后的工具偏移为 $(0,0.1,0)^\mathsf T$，故 $p_{\rm TCP}=(0.5,0.6,0)^\mathsf T\,\mathrm m$。法兰局部 $x$ 轴当前指向基座 $y$ 轴，所以不能直接把 0.1 加到基座 $x$ 分量。
\end{enumerate}

\section*{本章与原稿的对应}
本章整合 [A，第26--28页] 与 [B，第26--32页] 的正逆运动学、DH、PoE 和 IK 方法；把“IK 用 QP”限定为局部线性化求解，并补充可达任务维数、分支、限位、TCP 与零偏的完整算例。算法名称不能替代具体模型和约束的核对。

\chapter{雅可比 奇异性和七轴的运动学零空间}\label{ch:jacobian}
\section{雅可比的每一列是一种关节单独运动}
正运动学描述有限位置关系 $p=f(q)$。让一个关节微小运动、其他关节固定，端点的变化率就是关于该关节的偏导数。把所有列排在一起：
\begin{equation}
 J_p(q)=\frac{\partial p}{\partial q}
 =\begin{bmatrix}
 \partial p_x/\partial q_1&\cdots&\partial p_x/\partial q_n\\
 \partial p_y/\partial q_1&\cdots&\partial p_y/\partial q_n\\
 \partial p_z/\partial q_1&\cdots&\partial p_z/\partial q_n
 \end{bmatrix},
 \qquad\dot p=J_p(q)\dot q.
\end{equation}
位置雅可比 $J_p$ 的输出是指定点的线速度。六维几何雅可比则把同一点线速度与刚体角速度堆叠：
\begin{equation}
 \begin{bmatrix}\dot p\\\omega\end{bmatrix}
 =\begin{bmatrix}J_p\\J_\omega\end{bmatrix}\dot q
 =J_g(q)\dot q.
\end{equation}
角速度不是任意三个姿态参数的导数。若用欧拉角作为输出坐标，对应的解析雅可比需要乘入第~\ref{ch:pose}~章的姿态参数变换；不能把它与几何雅可比混用。

\section{两连杆雅可比逐项推导}
对式~\eqref{eq:two_r_fk} 求偏导。$q_1$ 同时改变第一杆和第二杆绝对朝向，$q_2$ 只改变第二杆，因此
\begin{align}
 \frac{\partial x}{\partial q_1}&=-l_1\sin q_1-l_2\sin(q_1+q_2),
 &\frac{\partial x}{\partial q_2}&=-l_2\sin(q_1+q_2),\\
 \frac{\partial y}{\partial q_1}&=l_1\cos q_1+l_2\cos(q_1+q_2),
 &\frac{\partial y}{\partial q_2}&=l_2\cos(q_1+q_2).
\end{align}
故平面位置雅可比为
\begin{equation}
 J_p=\begin{bmatrix}
 -l_1s_1-l_2s_{12}&-l_2s_{12}\\
 l_1c_1+l_2c_{12}&l_2c_{12}
 \end{bmatrix},
\end{equation}
其中 $s_1=\sin q_1$、$s_{12}=\sin(q_1+q_2)$，余弦缩写同理。平面末端角速度是 $\dot\theta_e=\dot q_1+\dot q_2$，因此若包含朝向，第三行是 $(1,1)$。这时 $J_g$ 为 $3\times2$，无法完成任意三维平面位姿速度。

\begin{examplebox}[算例 末端向上 为什么两个关节都要动]
在 $q=(0,\pi/2)$，位置雅可比为 $J=\begin{bmatrix}-0.5&-0.5\\0.5&0\end{bmatrix}$。要求 $v=(0,0.1)^\mathsf T\,\mathrm{m/s}$，第二行给出 $\dot q_1=0.2\,\mathrm{rad/s}$；第一行要求 $\dot q_1+\dot q_2=0$，所以 $\dot q_2=-0.2\,\mathrm{rad/s}$。
第一关节单独正转会让端点同时向左和向上；第二关节负转补回向左的部分，因此最终只向上。这是几何耦合，不是控制器必须逐轴独立的理由。
\end{examplebox}

\subsection{加速度为什么多出一项}
对 $\dot p=J\dot q$ 求导：
\begin{equation}
 \ddot p=J\ddot q+\dot J\dot q.
\end{equation}
在上例中，保持 $\dot q=(0.2,-0.2)^\mathsf T$、$\ddot q=0$。第二杆绝对角速度为零，第一杆却仍绕基座作圆周运动，故
\[
 \ddot p=-l_1\dot q_1^2(\cos q_1,\sin q_1)^\mathsf T
 =(-0.02,0)^\mathsf T\,\mathrm{m/s^2}.
\]
这正是 $\dot J\dot q$。如果要在有限时间内让末端沿直线匀速上移，关节速度必须随姿态更新，不能一直使用初始的两个常数速度。

\section{三维几何构造比手工求偏导更直观}
设第 $i$ 个旋转关节轴在基座中方向为单位向量 $z_i$，轴上参考点为 $p_i$，工具点为 $p_e$。刚体绕轴角速度为 $z_i\dot q_i$ 时，工具点速度为
\[
 \dot p_e=(z_i\dot q_i)\times(p_e-p_i).
\]
所以雅可比第 $i$ 列为
\begin{equation}
 J_{g,i}=\begin{bmatrix}z_i\times(p_e-p_i)\\z_i\end{bmatrix}
 \quad\text{旋转关节};\qquad
 J_{g,i}=\begin{bmatrix}z_i\\0\end{bmatrix}
 \quad\text{移动关节}.
\end{equation}
所有轴和位置必须先表达在同一坐标系。把“关节轴在局部的 $(0,0,1)$”未经变换直接代入基座点坐标，是常见错误。

若工具相对法兰有当前基座表达的偏移 $r$，则 $p_t=p_f+r$、$\dot r=\omega\times r$，因此
\begin{equation}
 J_{p,t}=J_{p,f}-\widehat rJ_\omega,
 \qquad J_{\omega,t}=J_{\omega,f}.
\end{equation}
同一刚体上不同点具有相同角速度，却一般有不同线速度。长工具尤其容易放大姿态变化造成的位置变化。

\section{雅可比转置从功率相等推出来}
力 $f$ 和力矩 $m$ 均为环境施加在工具上的量，且与几何速度使用同一坐标系、同一作用参考点，记 $F=(f,m)$。关节力矩记为 $\tau_{\rm ext}$。理想刚性无损几何传动中，任意瞬间的功率应相同：
\begin{equation}
 F^\mathsf T V=\tau_{\rm ext}^\mathsf T\dot q.
\end{equation}
代入 $V=J_g\dot q$，得到 $F^\mathsf TJ_g\dot q=\tau_{\rm ext}^\mathsf T\dot q$。因为对任意 $\dot q$ 都应成立，所以
\begin{equation}\label{eq:jac_force}
 \tau_{\rm ext}=J_g^\mathsf TF.
\end{equation}
这不是把速度式中的 $J$ 随意转置，而是由功率对偶关系决定。若要讨论机器人对环境的力，符号与环境对机器人的力相反，应先明确作用反作用的定义。

\begin{examplebox}[算例 验证末端功率与关节功率完全一致]
仍取标准两连杆姿态，末端受力 $f=(5,10)^\mathsf T\,\mathrm N$。由
\[
 \tau_{\rm ext}=\begin{bmatrix}-0.5&0.5\\-0.5&0\end{bmatrix}
 \begin{bmatrix}5\\10\end{bmatrix}
 =\begin{bmatrix}2.5\\-2.5\end{bmatrix}\mathrm{N\,m}.
\]
若关节速度为 $(0.2,-0.2)^\mathsf T\,\mathrm{rad/s}$，关节功率为 $2.5\times0.2+(-2.5)(-0.2)=1\,\mathrm W$。末端速度为 $(0,0.1)$，末端功率为 $5\times0+10\times0.1=1\,\mathrm W$。两侧一致，符号和转置方向得到核对。
\end{examplebox}

反过来用关节外力矩估计末端力，求解的是 $J^\mathsf TF\approx\tau_{\rm ext}$。只有在适当秩条件和接触假设下才有唯一解释。多点接触、模型偏差、摩擦和传感器零偏都可能产生相似关节力矩，不能把一次伪逆运算就称为“准确的六维力传感器”。

\section{奇异位形究竟失去了什么}
\subsection{两连杆的行列式逐项化简}
对平面位置雅可比，
\begin{align}
 \det J
 &=(-l_1s_1-l_2s_{12})(l_2c_{12})
  -(-l_2s_{12})(l_1c_1+l_2c_{12})\\
 &=l_1l_2(s_{12}c_1-s_1c_{12})
 =l_1l_2\sin q_2.
\end{align}
当 $q_2=0$ 或 $\pi$，秩下降。在 $q_1=q_2=0$ 的完全伸直状态，
\begin{equation}
 J=\begin{bmatrix}0&0\\1&0.5\end{bmatrix}.
\end{equation}
无论关节速度怎么选，瞬时 $x$ 方向速度都为零。沿 $y$ 的运动仍然存在，所以奇异不等于“整台机械臂完全不能动”。它表示所选任务的局部速度映射丢失独立方向。

\begin{examplebox}[算例 靠近奇异时关节速度迅速增大]
取 $q_1=0$、$q_2=0.01\,\mathrm{rad}$，要求末端 $v=(0.1,0)^\mathsf T\,\mathrm{m/s}$。求解两行线性方程得到
\[
 \dot q_1=0.2\cot q_2\approx19.9993\,\mathrm{rad/s},
 \qquad
 \dot q_2=-\frac{0.2(1+\cos q_2)}{\sin q_2}
 \approx-39.9997\,\mathrm{rad/s}.
\]
同样的 $0.1\,\mathrm{m/s}$ 目标，在 $q_2=90^\circ$ 时只需 $(0,-0.2)^\mathsf T\,\mathrm{rad/s}$。危险来自映射在指定方向上的放大，而不只是目标速度本身“大”。
\end{examplebox}

\begin{figure}[htbp]
\centering\includegraphics[width=0.9\linewidth]{figures/mk_singularity.pdf}
\caption{从屈肘到接近伸直，速度椭圆变得越来越扁。右图显示本例最小奇异值随肘角接近零而下降；某些方向所需的关节速度随之放大。}
\end{figure}

在理想刚性模型中，奇异方向的某些外力可沿结构传递而不产生相应关节广义力；这不表示真实机构能够承受无限载荷。连杆弯曲、轴承、减速器和底座都有有限强度与柔性。任务可操作性指标不能替代结构载荷判断。

\subsection{奇异值 条件数和可操作度不是同一个目标}
若 $J\in\mathbb R^{m\times n}$ 满行秩，速度单位球映射成椭球。最小奇异值 $\sigma_{\min}$ 描述最弱方向的速度换算，条件数 $\sigma_{\max}/\sigma_{\min}$ 描述各方向不均衡程度，可操作度（manipulability）
\begin{equation}
 w(q)=\sqrt{\det(JJ^\mathsf T)}=\prod_{i=1}^m\sigma_i
\end{equation}
与椭球体积成正比。它们都需要先固定单位、任务尺度及关节速度度量。

对本书等长两连杆的二维位置任务，$w=0.25|\sin q_2|$，在 $q_2=\pm90^\circ$ 时最大。但这并不等于条件数最小：$q_2=90^\circ$ 时条件数约为 2.618；$\cos q_2=-2/3$，即 $q_2\approx131.810^\circ$ 时条件数约为 1.618。前者追求椭圆面积，后者追求更接近圆形的比例。选择姿态必须根据任务需求，而不能把一个指标当作万能评分。

\section{伪逆给出最小关节速度解}
满行秩矩阵的右伪逆为
\begin{equation}
 J^+=J^\mathsf T(JJ^\mathsf T)^{-1}.
\end{equation}
它满足 $JJ^+=I_m$，所以 $\dot q_*=J^+v_d$ 能实现期望任务速度。为何是最小范数？任何其他精确解都能写成 $\dot q=\dot q_*+n$，其中 $Jn=0$。$\dot q_*$ 位于 $J$ 的行空间，$n$ 在零空间，两者正交，因此
\begin{equation}
 \|\dot q\|^2=\|\dot q_*\|^2+\|n\|^2\ge\|\dot q_*\|^2.
\end{equation}
这只是在同一瞬间、给定欧氏速度度量下的最小范数，不等于整条轨迹的最短时间、最小力矩或最小电能。

若希望某些关节尽量少动，可选正定权重 $W$，最小化 $\frac12\dot q^\mathsf TW\dot q$，约束 $J\dot q=v_d$。拉格朗日函数为 $\frac12\dot q^\mathsf TW\dot q-\lambda^\mathsf T(J\dot q-v_d)$。令导数为零得到 $W\dot q=J^\mathsf T\lambda$，代回约束后得到加权伪逆
\begin{equation}
 J_W^+=W^{-1}J^\mathsf T(JW^{-1}J^\mathsf T)^{-1}.
\end{equation}
\begin{examplebox}[算例 给第一关节更高代价]
任务是 $\dot q_1+\dot q_2=0.2$，即 $J=(1,1)$。普通最小范数解为 $(0.1,0.1)$。若 $W=\operatorname{diag}(4,1)$，则 $W^{-1}J^\mathsf T=(1/4,1)^\mathsf T$，$JW^{-1}J^\mathsf T=1.25$，所以
\[
 \dot q=(1/4,1)^\mathsf T\frac{0.2}{1.25}=(0.04,0.16)^\mathsf T.
\]
两种解都完成任务，但高代价关节动得更少。权重含义必须明确；取 $W=M$ 时关联的是瞬时动能度量，不可直接称为普遍的最小电能方案。
\end{examplebox}

\section{运动学零空间让主任务不变}
\subsection{投影矩阵的三个性质}
定义
\begin{equation}
 N=I-J^+J.
\end{equation}
由于 $JJ^+J=J$，有 $JN=0$。对 Moore--Penrose 伪逆，还满足 $N^\mathsf T=N$ 与 $N^2=N$，所以它是向 $J$ 零空间的正交投影。$N^2=N$ 表示“已经投影进去，再投一次不会改变”。完整速度解为
\begin{equation}
 \dot q=J^+v_d+N\dot q_0.
\end{equation}
$\dot q_0$ 是未经约束的次任务愿望，$N\dot q_0$ 保留其中不改变当前主任务速度的部分。若该愿望完全沿主任务方向，它可能被投影成零；这不是控制器失效，而是任务之间在当前几何状态下不兼容。

\begin{examplebox}[算例 三个输入完成两个任务的全部计算]
用简化矩阵
\[
 J=\begin{bmatrix}1&0&1\\0&1&1\end{bmatrix},\qquad v_d=(0.1,0.2)^\mathsf T.
\]
可算得
\[
 J^+=\frac13\begin{bmatrix}2&-1\\-1&2\\1&1\end{bmatrix},
 \qquad
 N=\frac13\begin{bmatrix}1&1&-1\\1&1&-1\\-1&-1&1\end{bmatrix}.
\]
主任务解为 $\dot q_*=(0,0.1,0.1)^\mathsf T$。若次任务希望第三个量降低，取 $\dot q_0=(0,0,-0.3)^\mathsf T$，投影为 $(0.1,0.1,-0.1)^\mathsf T$。合成速度为 $(0.1,0.2,0)^\mathsf T$，验算
\[
 J\dot q=(0.1+0,0.2+0)^\mathsf T=v_d.
\]
第三关节最终不动，是主任务正贡献与零空间负贡献抵消的结果，不能仅看某一项就判断最终运动。
\end{examplebox}

\subsection{七轴的一个冗余方向有明确条件}
七轴臂控制完整六维位姿且雅可比秩为 6 时，$\dim\mathcal N(J)=7-6=1$。若只控制三维位置且秩为 3，零空间通常为四维；若只控制相机视轴方向而允许绕视轴旋转，任务维数又不同。首先定义实际任务，比先说“给七轴加零空间”更重要。

为在图上看清冗余，本书用平面三关节臂控制二维位置：三关节减两约束，同样剩一维。固定末端位置而改变末端朝向，每个朝向对应一个不同的腕点，再由前两杆 IK 求解；于是得到一系列同末端位置、不同肘部形状。七轴六维任务的自运动（self-motion）具有类似直觉，但实际零空间方向由它自己的几何模型决定，不能照抄平面公式。

\begin{figure}[htbp]
\centering\includegraphics[width=0.9\linewidth]{figures/mk_nullspace.pdf}
\caption{平面三关节臂控制二维位置的冗余示例。各条臂形的末端都精确位于同一目标，关节和中间连杆可以改变；这是零空间的直观模型，不是 NERO 的结构图。}
\end{figure}

\subsection{次任务目标怎样变成一个速度愿望}
例如希望远离关节限位，定义关节中点 $q_{i,c}=(q_{i,\max}+q_{i,\min})/2$、范围 $r_i=q_{i,\max}-q_{i,\min}$，代价为
\begin{equation}
 h(q)=\frac12\sum_i\left(\frac{q_i-q_{i,c}}{r_i}\right)^2.
\end{equation}
梯度为 $[\nabla h]_i=(q_i-q_{i,c})/r_i^2$。沿负梯度走通常降低这个代价，因此可取 $\dot q_0=-k\nabla h$，再投影到零空间。投影后仍需检查实际降低速度：$\dot h=\nabla h^\mathsf TN\dot q_0=-k\nabla h^\mathsf TN\nabla h\le0$，这里利用普通正交投影 $N$ 半正定，且没有叠加主任务项的影响。主任务本身也会改变 $h$，所以总运动不保证限位代价永远下降。

避障、保持舒适姿态、扩大可操作度也可构造目标，但梯度法只是局部方法。若主任务把机器人逼向障碍而可用零空间不足，不能承诺零空间一定避得开。需要在优化问题中加入明确的不等式约束，或由规划层改变主任务路径。


\subsection{任务只控制部分方向时要重新定义雅可比}
六维末端速度可写成 $V=(v_x,v_y,v_z,\omega_x,\omega_y,\omega_z)^\mathsf T$。若任务只规定三个平移速度，用行选择矩阵 $S=[I_3\ \ 0]$ 提取，得到 $v_d=SV$、$J_{\rm task}=SJ_g$。应对这个三行矩阵计算伪逆和零空间，而不是先对六行矩阵严格约束全部姿态，再期待它释放额外姿态自由度。

以两连杆为例，在标准姿态若只规定向上速度，任务雅可比是 $J_y=(0.5,0)$。它的零空间包含第二关节单独运动，因为第二关节的瞬时运动不改变 $y$。但它会改变 $x$，所以这不是保持完整位置的零空间。说“末端不动”必须说明哪些方向被定义为主任务。

所选方向还依赖坐标系。工具绕基座 $z$ 轴转过 $90^\circ$ 时，工具 $x$ 方向就是基座 $y$ 方向。若在工具系选择一个轴，应先把同一点的线角速度转换到工具轴表达，再选择相应行。方向选择、坐标换算与参考点变化是三件不同的操作，应分别记录。

\subsection{两个任务有严格优先级时怎样串起来}
设第一任务为 $J_1\dot q=v_1$，先取 $\dot q_1=J_1^+v_1$、$N_1=I-J_1^+J_1$。所有保持第一任务的补充速度都可写成 $N_1z$。第二任务要求
\[
 J_2(\dot q_1+N_1z)=v_2,
\]
所以它在剩余空间中的方程是 $(J_2N_1)z=v_2-J_2\dot q_1$，给出
\begin{equation}
 \dot q=\dot q_1+N_1(J_2N_1)^+(v_2-J_2\dot q_1).
\end{equation}
第二任务只使用第一任务剩余的能力。如果剩余能力不足，第二任务只能做最小二乘折中，第一任务仍保持其已满足部分。这个结论假设使用一致的精确伪逆，且没有后续独立截断打破关系。

\begin{examplebox}[算例 两级任务都能满足时逐项求速度]
取 $J_1=(1,1,0)$、$v_1=0.2$，得 $\dot q_1=(0.1,0.1,0)^\mathsf T$，
\[
 N_1=\begin{bmatrix}0.5&-0.5&0\\-0.5&0.5&0\\0&0&1\end{bmatrix}.
\]
第二任务取 $J_2=(0,1,1)$、$v_2=0.3$。剩余误差为 $0.3-0.1=0.2$，剩余雅可比为 $J_2N_1=(-0.5,0.5,1)$，平方范数为 1.5。所需补充速度为 $(-0.5,0.5,1)^\mathsf T(0.2/1.5)=(-1/15,1/15,2/15)^\mathsf T$，它已经位于第一任务零空间。
最终速度为 $(1/30,1/6,2/15)^\mathsf T$。验算第一任务 $1/30+1/6=0.2$，第二任务 $1/6+2/15=0.3$。若给这两个任务随意加权求一个总平方误差，只能表达软优先级；权重再大也应与严格等式优先级区分。
\end{examplebox}

\section{阻尼伪逆的零空间并不严格为零}
工程上常把 DLS 矩阵 $J_\lambda^+=J^\mathsf T(JJ^\mathsf T+\lambda^2I)^{-1}$ 代入 $N_\lambda=I-J_\lambda^+J$。但
\begin{equation}
 JN_\lambda
 =\lambda^2(JJ^\mathsf T+\lambda^2I)^{-1}J,
\end{equation}
一般不为零。推导只需把 $JJ^\mathsf T$ 写成 $(JJ^\mathsf T+\lambda^2I)-\lambda^2I$。阻尼让逆映射平滑，也会让次任务影响主任务，称为任务泄漏；$N_\lambda$ 通常也不满足 $N_\lambda^2=N_\lambda$。

\begin{examplebox}[算例 阻尼投影后的主任务速度不是零]
取 $J=(1,1)$、$\lambda=0.5$，则 $J_\lambda^+=(1,1)^\mathsf T/2.25$。次任务愿望 $\dot q_0=(0.1,0.1)^\mathsf T$ 本来完全沿主任务方向；精确零空间投影会将它消掉。但阻尼投影得到
\[
 N_\lambda\dot q_0
 =(0.1,0.1)^\mathsf T-(1,1)^\mathsf T\frac{0.2}{2.25}
 =(0.011111,0.011111)^\mathsf T.
\]
所以主任务速度为 $J N_\lambda\dot q_0=0.022222$，并不为零。若要求严格优先级，应使用满足任务约束的分层求解，或明确允许并监测这种误差折中。
\end{examplebox}

另外，即便使用精确 $N(q)$，它也只保证当前瞬时速度不改变主任务。有限欧拉更新 $q_{k+1}=q_k+hN(q_k)\dot q_0$ 仍产生二阶位置误差；长时间累积可能漂移。每步更新雅可比、闭环纠正末端误差、使用更合适的约束积分或重投影，才能控制漂移。

\section{运动学投影与力矩投影的边界}
$JN=0$ 说的是关节速度映射为零，不是任意附加力矩都不会影响末端。力矩先经过惯量矩阵和其他动力学，再产生关节加速度。希望零空间力矩不引起任务加速度时，应检查的是相应动力学条件，例如涉及 $JM^{-1}$ 的映射。后续操作空间与阻抗章节会推导动力学一致投影。

同样，运动学 IK 只能生成几何目标，不能解释“为什么可以手推”或“为什么位置很硬”。这些表现取决于驱动接口、反馈增益、重力和摩擦补偿、延迟等。把位置目标经 IK 算得准确，并不意味着已实现力控或低阻抗。

\section{怎样核对你计算的雅可比}
对位置雅可比第 $i$ 列，可用中心差分
\begin{equation}
 J_{p,i}^{\rm num}
 =\frac{p(q+\epsilon e_i)-p(q-\epsilon e_i)}{2\epsilon}
\end{equation}
核对解析或库结果，其中 $e_i$ 是第 $i$ 个单位坐标向量。$\epsilon$ 太大时非线性误差明显，太小时舍入误差占主导；在双精度、弧度及米量级的教学例中可先试 $10^{-6}$，再比较相邻尺度的稳定性。它不是任何模型都唯一正确的常数。

旋转部分不能直接对欧拉角做差后与几何角速度列比较；应使用相对旋转对数并保持坐标约定。核对工具点、参考系、线角顺序和关节索引之后，再看数值误差。这样可以避免“公式没错，但取的是法兰雅可比而误用于相机 TCP”的问题。

\section{计算练习与逐步解答}
\subsection*{题目}
\begin{enumerate}
\item 在标准姿态 $q=(0,\pi/2)$，关节速度 $(0.3,-0.1)^\mathsf T\,\mathrm{rad/s}$ 会产生什么末端线速度和角速度？
\item 同一姿态，环境施加末端力 $(-4,6)^\mathsf T\,\mathrm N$，求关节外力矩；用上一题速度核对功率。
\item 对 $J=(1,1)$，求普通伪逆与精确零空间投影。将 $\dot q_0=(0.3,-0.1)^\mathsf T$ 投影，并检查输出。
\item 完全伸直时 $J=\begin{bmatrix}0&0\\1&0.5\end{bmatrix}$，目标速度 $v_d=(0.1,0.1)^\mathsf T$，取 $\lambda=0.1$ 求 DLS 速度，并计算真实可实现的任务速度。
\end{enumerate}
\subsection*{解答}
\begin{enumerate}
\item $v_x=-0.5\times0.3-0.5(-0.1)=-0.1\,\mathrm{m/s}$；$v_y=0.5\times0.3=0.15\,\mathrm{m/s}$。末端角速度为 $\dot q_1+\dot q_2=0.2\,\mathrm{rad/s}$，即绕平面法向正转。
\item $\tau_1=(-0.5)(-4)+0.5\times6=5\,\mathrm{N\,m}$，$\tau_2=(-0.5)(-4)=2\,\mathrm{N\,m}$。关节功率为 $5\times0.3+2(-0.1)=1.3\,\mathrm W$；末端功率为 $(-4)(-0.1)+6(0.15)=1.3\,\mathrm W$。
\item $J^+=(1/2,1/2)^\mathsf T$，$N=\frac12\begin{bmatrix}1&-1\\-1&1\end{bmatrix}$。投影结果为 $(0.2,-0.2)^\mathsf T$，$J(0.2,-0.2)^\mathsf T=0$。再投一次不改变结果。
\item $JJ^\mathsf T+\lambda^2I=\operatorname{diag}(0.01,1.26)$。先求 $(JJ^\mathsf T+\lambda^2I)^{-1}v_d=(10,0.0793651)^\mathsf T$，再左乘 $J^\mathsf T$ 得 $\dot q=(0.0793651,0.0396825)^\mathsf T$。实际任务速度为 $(0,0.0992063)^\mathsf T$；不可达的 $x$ 方向仍然为零，$y$ 方向也因阻尼略有误差。
\end{enumerate}

\section*{本章与原稿的对应}
本章整合 [A，第28--30页] 与 [B，第14--17、28--32页]，补充了功率推导、雅可比作用点、加速度项、加权速度度量、有限步长漂移与阻尼投影泄漏。七轴的主任务与零空间应先明确所控任务维度，再进入后续动力学与力控设计。

\chapter{从能量逐步建立机械臂动力学}\label{ch:dynamics}
\section{动力学究竟要回答什么}
运动学给出关节角与末端位置之间的关系。动力学进一步回答：沿着这条运动轨迹，驱动器要输出多大力矩；当人推了一下机械臂，关节又会怎样加速。前一个问题叫逆动力学（inverse dynamics），后一个叫正动力学（forward dynamics）。这里的“正”和“逆”是输入与输出的方向，不代表计算难度或正确与错误。

本章采用如下统一模型：
\begin{equation}\label{eq:full_dynamics_di}
 M(q)\ddot q+C(q,\dot q)\dot q+g(q)+\tau_f=\tau_{\mathrm{cmd}}+\tau_{\mathrm{ext}}.
\end{equation}
其中，$q\in\mathbb R^n$ 是 $n$ 个关节角组成的列向量，$\dot q$、$\ddot q$ 分别是关节角速度和角加速度；$M\in\mathbb R^{n\times n}$ 是惯量矩阵；$C\dot q$ 是速度耦合力矩；$g$ 是抵消重力所需的力矩；$\tau_f$ 是写在方程左侧的摩擦阻力矩；$\tau_{\mathrm{cmd}}$ 是驱动器实际作用于关节的力矩；$\tau_{\mathrm{ext}}$ 是环境作用于机械臂的外力矩。这里的“实际”很重要：驱动器收到的指令与执行器真正实现的力矩可能因限流、延迟、标定误差而不同。

对转动关节，角度、角速度、角加速度的单位分别写作 $\mathrm{rad}$、$\mathrm{rad/s}$、$\mathrm{rad/s^2}$，力矩用 $\mathrm{N\,m}$。惯量单位是 $\mathrm{kg\,m^2}$。弧度在量纲分析中视为无量纲，但工程接口仍必须标明弧度，避免把角度制数字直接代入。

\begin{conceptbox}[各项的直觉]
$M\ddot q$ 是“要加速就要出力”；$C\dot q$ 是“各连杆运动会互相影响”；$g$ 是“举住自身和负载的重量”；$\tau_f$ 是“克服轴承、密封、传动和接触的阻力”。低速时，速度耦合常变小，但加速度未必小：缓慢运动在急停瞬间仍可能有很大的惯性力矩。判断能否忽略某一项，应比较它的数值与允许的力矩误差。
\end{conceptbox}

给定状态和力矩，正动力学是解线性方程
\begin{equation}
 M(q)\ddot q=\tau_{\mathrm{cmd}}+\tau_{\mathrm{ext}}-C(q,\dot q)\dot q-g(q)-\tau_f.
\end{equation}
数学上可以写 $M^{-1}$，程序中应解方程，不应每个周期显式求逆。得到加速度后，数值积分器再更新速度、位置。逆动力学则直接按式\eqref{eq:full_dynamics_di}逐项计算所需力矩。

\section{一个刚体为什么需要质量 质心与惯量}
\subsection{质量决定平动惯性 转动惯量决定转动惯性}
一个质量为 $m$ 的点，以速度 $v$ 平动时，动能为 $\tfrac12 mv^2$。它绕某根轴转动，离轴距离为 $r$，则 $v=r\omega$，所以
\begin{equation}
 T=\frac12m(r\omega)^2=\frac12(mr^2)\omega^2.
\end{equation}
括号里的 $mr^2$ 就是这个点对该轴的转动惯量。物体由许多小质量组成时，把它们相加得到 $I=\int r^2\,dm$。质量离轴越远，转动惯量越大；相同质量并不意味着相同转动惯量。

质心（center of mass）是质量的加权平均位置。考虑三个质量点，其质量与横坐标分别为
\[
 (m_1,m_2,m_3)=(1,2,1)\ \mathrm{kg},\quad
 (x_1,x_2,x_3)=(0,0.1,0.4)\ \mathrm m.
\]
总质心的横坐标为
\begin{equation}
 c_x=\frac{1\times0+2\times0.1+1\times0.4}{1+2+1}=0.15\ \mathrm m.
\end{equation}
不是把三个位置简单平均，而是质量越大权重越大。

\subsection{从一个惯量数到惯量张量}
三维空间中，物体绕不同轴转动需要不同的力矩。用对称矩阵 $I_C$ 描述它绕质心的转动惯性：
\begin{equation}
 I_C=\int\left(\|r\|^2I_3-rr^{\mathsf T}\right)dm
 =\begin{bmatrix}I_{xx}&I_{xy}&I_{xz}\\I_{xy}&I_{yy}&I_{yz}\\I_{xz}&I_{yz}&I_{zz}\end{bmatrix}.
\end{equation}
$C$ 表示质心，$I_3$ 表示三阶单位矩阵，$r$ 是相对于质心的位置。按这个定义，$I_{xy}=-\int xy\,dm$；某些教材把“惯性积”定义成正的 $\int xy\,dm$，组装矩阵时再加负号，因此跨资料比较必须核对定义。

惯量矩阵的特征向量给出主惯量轴，沿这些互相垂直的方向表达时，矩阵为对角阵；三个特征值就是主惯量。矩阵对称意味着六个数即可独立决定它。加上质量一个数、质心三个数，一个刚体的惯性描述共有十个独立标量。质心坐标、惯量参考点和表达坐标系必须同时说明。只有数值而没有坐标约定的惯量不能可靠地用于控制。

在与物体固连、原点位于质心的坐标系中，牛顿欧拉方程写作
\begin{equation}
 f_C=m a_C,\qquad m_C=I_C\dot\omega+\omega\times(I_C\omega).
\end{equation}
$f_C$ 是合外力，$m_C$ 是关于质心的合外力矩，$\omega$ 是角速度；这里 $\dot\omega$ 为角加速度在同一坐标系的分量。第二项是陀螺项。若 $\omega$ 恰沿主惯量轴，$I_C\omega$ 与 $\omega$ 平行，叉乘为零；其他方向一般不为零。

\begin{examplebox}[算例 6A 陀螺力矩可以在角速度恒定时存在]
取 $I_C=\operatorname{diag}(0.01,0.02,0.03)\ \mathrm{kg\,m^2}$，$\omega=(1,2,0)^{\mathsf T}\ \mathrm{rad/s}$，并令 $\dot\omega=0$。先算
\[
 I_C\omega=(0.01,0.04,0)^{\mathsf T}.
\]
再算叉乘
\[
 \omega\times(I_C\omega)=(0,0,1\times0.04-2\times0.01)^{\mathsf T}=(0,0,0.02)^{\mathsf T}\ \mathrm{N\,m}.
\]
所以恒定角速度并不自动意味着合力矩为零。原因是角动量与角速度不平行，而物体坐标轴本身在转动。
\end{examplebox}

\subsection{平行轴定理与坐标旋转}
已知绕质心的惯量 $I_C$，质心相对于坐标原点 $O$ 的位置是 $c$，则绕 $O$ 的惯量为
\begin{equation}\label{eq:parallel_di}
 I_O=I_C+m\left((c^{\mathsf T}c)I_3-cc^{\mathsf T}\right).
\end{equation}
可从定义直接推导：把相对 $O$ 的位置写成 $r_O=c+r_C$，代入惯量积分；含 $r_C$ 的一次项因为 $\int r_Cdm=0$ 消失，剩下式\eqref{eq:parallel_di}。这也说明平行轴定理为什么必须以质心为出发点。

\begin{examplebox}[算例 6B 半米均匀杆的惯量]
一根均匀细杆质量 $1\ \mathrm{kg}$，长 $0.5\ \mathrm m$，绕垂直杆且通过中点的轴：
\[
 I_C=\frac{ml^2}{12}=\frac{0.25}{12}=0.0208333\ \mathrm{kg\,m^2}.
\]
换到端点，$c=0.25\ \mathrm m$，所以
\[
 I_O=0.0208333+1\times0.25^2=0.0833333\ \mathrm{kg\,m^2}.
\]
端点惯量是中点惯量的四倍。把这两个数弄反，会明显改变加速力矩，但静止时的重力补偿可能仍看不出错误。
\end{examplebox}

若原点不变，只把表达坐标从 $B$ 换到 $A$，则 $I_A=R_{AB}I_BR_{AB}^{\mathsf T}$。平移要用平行轴定理，旋转要用矩阵共轭，不能只把惯量矩阵的对角线交换。

\section{用同一台两连杆机械臂贯穿推导}
本章选择竖直平面内的两根均匀细杆，而不是把质量集中在杆端。几何与参数为
\[
 l_1=l_2=0.5\ \mathrm m,\quad m_1=m_2=1\ \mathrm{kg},\quad
 c_1=c_2=0.25\ \mathrm m,\quad I_1=I_2=\frac1{48}\ \mathrm{kg\,m^2}.
\]
$c_i$ 是第 $i$ 根杆从近端关节到质心的距离；$I_i$ 是绕该杆质心、垂直运动平面的转动惯量。$q_1$ 从水平正 $x$ 方向逆时针测量，$q_2$ 是第二杆相对于第一杆的转角。重力沿负 $y$，大小 $g_0=9.81\ \mathrm{m/s^2}$。

\begin{figure}[htbp]\centering
\includegraphics[width=0.82\linewidth,height=0.37\textheight,keepaspectratio]{figures/di_2r_geometry.pdf}
\caption{两连杆模型中的关节角、质心、重力和水平力臂。质量在质心处，而转动动能仍需要质心惯量。}\label{fig:di_2r_geometry}
\end{figure}

原稿中的部分算例采用“质量在杆端”的模型，所得数字与本章不同是模型不同，并非算式互相矛盾。教学时必须先确定质量分布，再谈惯量矩阵是否正确。

\subsection{第一步 写出两个质心的位置}
记 $q_{12}=q_1+q_2$，则
\begin{align}
 p_{C1}&=\begin{bmatrix}c_1\cos q_1\\c_1\sin q_1\end{bmatrix},\\
 p_{C2}&=\begin{bmatrix}l_1\cos q_1+c_2\cos q_{12}\\l_1\sin q_1+c_2\sin q_{12}\end{bmatrix}.
\end{align}
对时间求导，使用 $\frac d{dt}\cos q=-\sin q\,\dot q$，得
\begin{align}
 v_{C1}&=\begin{bmatrix}-c_1\sin q_1\dot q_1\\c_1\cos q_1\dot q_1\end{bmatrix},\\
 v_{C2}&=\begin{bmatrix}-l_1\sin q_1\dot q_1-c_2\sin q_{12}(\dot q_1+\dot q_2)\\
 l_1\cos q_1\dot q_1+c_2\cos q_{12}(\dot q_1+\dot q_2)\end{bmatrix}.
\end{align}
把横纵分量平方后相加；利用 $\sin^2q+\cos^2q=1$ 以及 $\cos a\cos b+\sin a\sin b=\cos(a-b)$，得到
\begin{align}
 \|v_{C1}\|^2&=c_1^2\dot q_1^2,\\
 \|v_{C2}\|^2&=l_1^2\dot q_1^2+c_2^2(\dot q_1+\dot q_2)^2
 +2l_1c_2\cos q_2\dot q_1(\dot q_1+\dot q_2).
\end{align}
最后一项来自两部分速度的内积，正是后来出现关节耦合的来源。

\subsection{第二步 平动与转动动能都要算}
两杆的绝对角速度分别为 $\dot q_1$、$\dot q_1+\dot q_2$。总动能是
\begin{align}
 T={}&\tfrac12m_1\|v_{C1}\|^2+\tfrac12I_1\dot q_1^2
 +\tfrac12m_2\|v_{C2}\|^2+\tfrac12I_2(\dot q_1+\dot q_2)^2.
\end{align}
为减少重复，定义三个常数
\begin{equation}
 a=I_1+I_2+m_1c_1^2+m_2(l_1^2+c_2^2),\quad
 b=m_2l_1c_2,\quad d=I_2+m_2c_2^2.
\end{equation}
把平方展开并按 $\dot q_1^2$、$\dot q_1\dot q_2$、$\dot q_2^2$ 收集：
\begin{equation}\label{eq:T_2r_di}
 T=\frac12(a+2b\cos q_2)\dot q_1^2
 +(d+b\cos q_2)\dot q_1\dot q_2+\frac12d\dot q_2^2.
\end{equation}
势能等于质量乘重力加速度乘高度：
\begin{equation}
 V=g_0\left[(m_1c_1+m_2l_1)\sin q_1+m_2c_2\sin(q_1+q_2)\right].
\end{equation}
高度零点任选，只要一致；给势能加一个常数不改变其导数，因此不改变力矩。

\subsection{第三步 对关节一完整执行欧拉拉格朗日运算}
拉格朗日函数 $L=T-V$。忽略摩擦和外部接触时，广义力矩满足
\begin{equation}
 \frac d{dt}\left(\frac{\partial L}{\partial\dot q_i}\right)
 -\frac{\partial L}{\partial q_i}=\tau_i.
\end{equation}
先对 $\dot q_1$ 求偏导：
\[
 \frac{\partial L}{\partial\dot q_1}
 =(a+2b\cos q_2)\dot q_1+(d+b\cos q_2)\dot q_2.
\]
系数含 $q_2(t)$，所以对时间求导时必须使用乘积法则：
\begin{align}
 \frac d{dt}\frac{\partial L}{\partial\dot q_1}
 ={}&(a+2b\cos q_2)\ddot q_1+(d+b\cos q_2)\ddot q_2\\
 &-2b\sin q_2\dot q_1\dot q_2-b\sin q_2\dot q_2^2.
\end{align}
动能不含 $q_1$，因此 $\partial L/\partial q_1=-\partial V/\partial q_1$。算出
\[
 \frac{\partial V}{\partial q_1}
 =g_0[(m_1c_1+m_2l_1)\cos q_1+m_2c_2\cos(q_1+q_2)].
\]
于是关节一方程为
\begin{align}
 \tau_1={}&(a+2b\cos q_2)\ddot q_1+(d+b\cos q_2)\ddot q_2\\
 &-b\sin q_2(2\dot q_1\dot q_2+\dot q_2^2)+\frac{\partial V}{\partial q_1}.
\end{align}

\subsection{第四步 对关节二求导并观察抵消项}
对 $\dot q_2$ 求偏导再求时间导数：
\begin{align}
 \frac{\partial L}{\partial\dot q_2}&=(d+b\cos q_2)\dot q_1+d\dot q_2,\\
 \frac d{dt}\frac{\partial L}{\partial\dot q_2}
 &=(d+b\cos q_2)\ddot q_1+d\ddot q_2-b\sin q_2\dot q_1\dot q_2.
\end{align}
这次动能含 $q_2$，所以
\begin{align}
 \frac{\partial L}{\partial q_2}
 &=-b\sin q_2\dot q_1^2-b\sin q_2\dot q_1\dot q_2
 -m_2c_2g_0\cos(q_1+q_2).
\end{align}
两式相减，含 $\dot q_1\dot q_2$ 的项正好抵消，得
\begin{equation}
 \tau_2=(d+b\cos q_2)\ddot q_1+d\ddot q_2
 +b\sin q_2\dot q_1^2+m_2c_2g_0\cos(q_1+q_2).
\end{equation}
因此
\begin{align}\label{eq:2rM_di}
 M&=\begin{bmatrix}a+2b\cos q_2&d+b\cos q_2\\d+b\cos q_2&d\end{bmatrix},\\
 c(q,\dot q):=C\dot q&=\begin{bmatrix}-b\sin q_2(2\dot q_1\dot q_2+\dot q_2^2)\\b\sin q_2\dot q_1^2\end{bmatrix},\\
 g&=g_0\begin{bmatrix}(m_1c_1+m_2l_1)\cos q_1+m_2c_2\cos(q_1+q_2)\\m_2c_2\cos(q_1+q_2)\end{bmatrix}.
\end{align}
用小写 $c(q,\dot q)$ 表示整个速度力矩向量，可以避免误以为 $C$ 是唯一的。满足 $C\dot q=c$ 的矩阵不唯一；稳定性证明常使用由克里斯托费尔符号构造的特定 $C$。

\begin{examplebox}[算例 6C 一次算完惯性 速度与重力力矩]
本章参数给出 $a=5/12$、$b=1/8$、$d=1/12$。令
\[
 q=(0,\pi/2)^{\mathsf T},\quad\dot q=(0.2,-0.1)^{\mathsf T},\quad
 \ddot q=(0.5,-0.2)^{\mathsf T}.
\]
此时 $\cos q_2=0$、$\sin q_2=1$，所以
\[
 M=\begin{bmatrix}0.4166667&0.0833333\\0.0833333&0.0833333\end{bmatrix}.
\]
惯性力矩逐行相乘：
\[
 M\ddot q=\begin{bmatrix}0.4166667\times0.5-0.0833333\times0.2\\0.0833333\times0.5-0.0833333\times0.2\end{bmatrix}
 =\begin{bmatrix}0.1916667\\0.025\end{bmatrix}.
\]
速度项是
\[
 c=\begin{bmatrix}-0.125(2\times0.2\times(-0.1)+(-0.1)^2)\\0.125\times0.2^2\end{bmatrix}
 =\begin{bmatrix}0.00375\\0.005\end{bmatrix}.
\]
重力项为 $g=(7.3575,0)^{\mathsf T}$，因此不计摩擦和外力时
\[
 \tau_{\mathrm{cmd}}=(7.5529167,0.0300000)^{\mathsf T}\ \mathrm{N\,m}.
\]
关节二的重力项为零，只意味着重力对它没有力臂，不意味着这个关节在运动时不用出力。
\end{examplebox}

\section{三个必须理解的结构性质}
\subsection{惯量矩阵为什么对称正定}
式\eqref{eq:T_2r_di}可以写成 $T=\tfrac12\dot q^{\mathsf T}M\dot q$。动能是各连杆平动与转动动能之和，因此正常刚体模型对任何非零有效速度都具有正动能，即 $\dot q^{\mathsf T}M\dot q>0$。这叫正定性。矩阵的交叉项来自同一个速度内积，所以 $M=M^{\mathsf T}$。

在算例姿态下，$M_{11}>0$ 且
\[
 \det M=\frac5{12}\frac1{12}-\left(\frac1{12}\right)^2=\frac1{36}>0,
\]
所以两阶矩阵正定。某个辨识结果使 $M$ 有负特征值，就意味着某个速度方向的“动能”为负，已经违反该机械模型的物理意义。

\subsection{科氏离心项参与能量交换 不能简单说它不做功}
对本模型，可选
\begin{equation}
 C=\begin{bmatrix}-h\dot q_2&-h(\dot q_1+\dot q_2)\\h\dot q_1&0\end{bmatrix},\qquad h=b\sin q_2.
\end{equation}
直接对 $M$ 求时间导数：
\[
 \dot M=\begin{bmatrix}-2h\dot q_2&-h\dot q_2\\-h\dot q_2&0\end{bmatrix}.
\]
于是
\[
 \dot M-2C=\begin{bmatrix}0&h(2\dot q_1+\dot q_2)\\-h(2\dot q_1+\dot q_2)&0\end{bmatrix}.
\]
这确实反对称，即转置等于自身负值。对任意实向量 $z$，$z^{\mathsf T}(\dot M-2C)z=0$。

不过，这并不意味着 $\dot q^{\mathsf T}C\dot q=0$。正确关系是
\begin{equation}
 \dot q^{\mathsf T}C\dot q=\frac12\dot q^{\mathsf T}\dot M\dot q.
\end{equation}
右侧通常非零，表示惯量随构型变化引起的动能交换。在算例 6C 中，$\dot q^{\mathsf T}c=0.2\times0.00375-0.1\times0.005=0.00025\ \mathrm W$，就不是零。

\subsection{从动力学严格得到能量平衡}
对总机械能 $E=T+V$ 求导：
\begin{align}
 \dot E&=\dot q^{\mathsf T}M\ddot q+\tfrac12\dot q^{\mathsf T}\dot M\dot q+\dot q^{\mathsf T}g\\
 &=\dot q^{\mathsf T}(\tau_{\mathrm{cmd}}+\tau_{\mathrm{ext}}-\tau_f)
 +\tfrac12\dot q^{\mathsf T}(\dot M-2C)\dot q\\
 &=\dot q^{\mathsf T}(\tau_{\mathrm{cmd}}+\tau_{\mathrm{ext}})-\dot q^{\mathsf T}\tau_f.
\end{align}
这条式子说明：执行器与环境输入的功率增加机械能；耗散摩擦消耗机械能。所谓机械系统的无源性，核心就是它不能无缘无故创造能量。后续讨论阻抗控制、力反馈稳定性时，将反复使用这条能量关系。

\section{外力怎样变成关节力矩}
末端受到平面外力 $f=(f_x,f_y)^{\mathsf T}$，若 $J$ 满足 $\dot p=J\dot q$，由虚功相等得到
\[
 f^{\mathsf T}\delta p=\tau_{\mathrm{ext}}^{\mathsf T}\delta q,\qquad
 \delta p=J\delta q,
\]
所以
\begin{equation}
 \tau_{\mathrm{ext}}=J^{\mathsf T}f.
\end{equation}
这是雅可比转置来自物理功率一致性的原因。三维时将速度扩展为 $[v;\omega]$、力旋量扩展为 $[f;m]$；两者必须关于同一点且在同一坐标系表达。

\begin{examplebox}[算例 6D 人向下拉末端时执行器要多出多少力矩]
在 $q=(0,\pi/2)$ 时，$J=\begin{bmatrix}-0.5&-0.5\\0.5&0\end{bmatrix}$。令环境作用力 $f=(2,-3)^{\mathsf T}\ \mathrm N$，则
\[
 \tau_{\mathrm{ext}}=\begin{bmatrix}-0.5&0.5\\-0.5&0\end{bmatrix}
 \begin{bmatrix}2\\-3\end{bmatrix}=\begin{bmatrix}-2.5\\-1\end{bmatrix}\ \mathrm{N\,m}.
\]
若仍要实现算例 6C 的运动，动力学要求
\[
 \tau_{\mathrm{cmd}}=M\ddot q+c+g-\tau_{\mathrm{ext}}
 =(10.0529167,1.0300000)^{\mathsf T}\ \mathrm{N\,m}.
\]
外力写在方程右侧，因而反算执行器力矩时要减去它。把这个符号弄反，会使控制器在受推时进一步加速。
\end{examplebox}

\section{用质心雅可比把推导推广到任意关节数}
两连杆中的求导并不是只能用于平面。对第 $i$ 个三维连杆，设其质心线速度和角速度满足
\[
 v_{Ci}=J_{vi}(q)\dot q,\qquad\omega_i=J_{\omega i}(q)\dot q.
\]
$J_{vi}$、$J_{\omega i}$ 都有三行、$n$ 列；如果全部向量在基座系表达，惯量也必须写成同一坐标系的 $I_{Ci}^B=R_{Bi}I_{Ci}^iR_{Bi}^{\mathsf T}$。该连杆动能为
\begin{align}
 T_i&=\tfrac12m_i\dot q^{\mathsf T}J_{vi}^{\mathsf T}J_{vi}\dot q
 +\tfrac12\dot q^{\mathsf T}J_{\omega i}^{\mathsf T}I_{Ci}^BJ_{\omega i}\dot q.
\end{align}
对所有连杆相加，可直接读出惯量矩阵：
\begin{equation}
 M(q)=\sum_i\left(m_iJ_{vi}^{\mathsf T}J_{vi}
 +J_{\omega i}^{\mathsf T}I_{Ci}^BJ_{\omega i}\right).
\end{equation}
每一项都把某根杆的物理惯性通过速度映射换到关节空间。这个公式同时解释了为什么同一末端负载会影响上游多个关节，以及为什么 $M$ 的非对角元素不能普遍忽略。计算某根杆的质心雅可比时，它下游的关节不会移动这根杆，因此相应列为零；这个机构树结构使递归算法能够避免无用计算。

若基座系重力加速度为常向量 $a_g^B$，势能可写为
\[
 V(q)=-\sum_i m_i(a_g^B)^{\mathsf T}p_{Ci}(q).
\]
对 $q$ 求导得到补偿重力的关节向量：
\begin{equation}
 g(q)=-\sum_i J_{vi}^{\mathsf T}m_i a_g^B.
\end{equation}
例如本章重力沿负 $y$，负号与 $a_g^B=(0,-9.81,0)$ 相乘，使水平杆的补偿力矩为正。换基座安装方向时，必须同步改变重力在基座系中的表示，而不是在某一个关节公式前随意改符号。

如果已经得到 $M(q)$，还可用克里斯托费尔系数组装一个具有所需能量性质的 $C$：
\begin{equation}
 C_{ij}=\sum_k\frac12\left(\frac{\partial M_{ij}}{\partial q_k}
 +\frac{\partial M_{ik}}{\partial q_j}
 -\frac{\partial M_{jk}}{\partial q_i}\right)\dot q_k.
\end{equation}
三个偏导项分别来自对动能执行乘积法则以及减去位置偏导。公式里的 $i,j,k$ 都是关节索引，不是三维空间坐标。把它代回 $C\dot q$，应与直接用拉格朗日求出的速度力矩一致。真实计算通常使用解析或自动微分和递归算法；粗糙地对 $M$ 做数值差分可能把差分误差注入控制器。

由这些关系可以形成一套有物理意义的实现核对：质心速度应与质心位置的差分一致；两种动能算法应一致；势能的梯度应与重力补偿一致；$\dot M-2C$ 应满足相应反对称性质。它们比只检查“机器人是否动了”更能定位错误发生在几何、惯量参考点还是求导环节。

\section{从手推两关节到计算七关节}
拉格朗日方法适合理解结构和核对小模型，但七关节完整展开会非常长。递归牛顿欧拉算法 RNEA 的思路仍是本章的物理：第一遍从基座到末端传播各杆速度、加速度；第二遍从末端到基座汇总每根杆的惯性力、重力和下游传递的力，再投影到关节轴得到力矩。它避免展开所有三角函数。

\begin{figure}[htbp]\centering
\includegraphics[width=0.89\linewidth]{figures/di_inertia_energy.pdf}
\caption{同一两连杆模型的惯量随肘角变化，以及在不同速度组合下的动能。每条曲线均由本章显式模型计算。}\label{fig:di_inertia_energy}
\end{figure}

CRBA 计算惯量矩阵，ABA 求正动力学，它们与 RNEA 是不同任务的高效算法。可以用下面三个恒等关系检查一个动力学实现：
\begin{align}
 g(q)&=\operatorname{ID}(q,0,0),\\
 c(q,\dot q)&=\operatorname{ID}(q,\dot q,0)-g(q),\\
 M(q)e_i&=\operatorname{ID}(q,0,e_i)-g(q).
\end{align}
$e_i$ 是第 $i$ 个分量为一、其他为零的单位向量；最后一式返回 $M$ 的第 $i$ 列。这些关系指同一个刚体惯性模型。若某个库在逆动力学里额外加入摩擦或其他被动力，比较时必须使用一致约定。

\section{计算练习与分步解答}
\subsection*{练习 6.1 水平伸直需要多少重力补偿}
令本章两杆 $q_1=q_2=0$，求两个关节的静态重力力矩。

\textbf{解答。}关节一看到杆一质心的水平力臂 $0.25$，杆二质心的水平力臂 $0.5+0.25=0.75$，所以 $g_1=9.81(0.25+0.75)=9.81\ \mathrm{N\,m}$。关节二只看到杆二，$g_2=9.81\times0.25=2.4525\ \mathrm{N\,m}$。代入本章矩阵公式也应得到相同结果。

\subsection*{练习 6.2 扣除重力后求加速度}
在 $q=(0,\pi/2)$、$\dot q=0$ 时，执行器比重力补偿多输出 $(0.1,0)^{\mathsf T}\ \mathrm{N\,m}$，忽略摩擦，求加速度。

\textbf{解答。}解 $\begin{bmatrix}5/12&1/12\\1/12&1/12\end{bmatrix}\ddot q=(0.1,0)^{\mathsf T}$。第二行给 $\ddot q_2=-\ddot q_1$；代入第一行，$(4/12)\ddot q_1=0.1$，所以 $\ddot q_1=0.3$、$\ddot q_2=-0.3\ \mathrm{rad/s^2}$。只给关节一附加力矩，关节二也会加速，这就是动力学耦合。

\subsection*{练习 6.3 用动能检查矩阵}
在算例 6C 的姿态和速度下，求总动能。

\textbf{解答。}先算 $M\dot q=(0.075,0.0083333)^{\mathsf T}$，再算
\[
 T=\tfrac12(0.2\times0.075-0.1\times0.0083333)=0.00708333\ \mathrm J.
\]
分别计算各质心平动与杆转动动能后相加，也应得到这个数。两种方式不一致时，常见原因是漏掉质心转动惯量或把相对角速度误当成第二杆的绝对角速度。

\subsection*{练习 6.4 位于端点的质量与均匀杆为什么不同}
若把两根杆各自的 $1\ \mathrm{kg}$ 质量全部集中在远端，几何长度不变，求 $q=(0,\pi/2)$ 下的 $M$ 和 $g$。

\textbf{解答。}此时 $c_1=c_2=0.5$，$I_1=I_2=0$。重新计算 $a=0.75$、$b=0.25$、$d=0.25$，因此 $M=\begin{bmatrix}0.75&0.25\\0.25&0.25\end{bmatrix}$，$g=(9.81,0)^{\mathsf T}$。这与原稿点质量算例一致。不能把这个矩阵与均匀杆的质心参数混搭。

\section*{本章与原稿的对应}
本章整合并展开 [A，第30--34页] 与 [B，第32--40页] 的刚体动力学内容。重新统一质量分布、外力方向和坐标约定，补足两关节欧拉拉格朗日的中间求导步骤，并修正把速度耦合项简单解释为“不做功”的表述。接触与柔性模型在下一章继续。

\chapter{电机 传动 柔性关节与接触}\label{ch:actuation}
\section{从电流到关节力矩的完整链条}
控制程序中的一个力矩数字，最终要经过电机、驱动器和传动系统作用到连杆。理解这条链，可以解释三个常见现象：电流反馈看起来很平滑，末端却有摩擦死区；模型重力已经补偿，关节仍然难拖动；调高刚度后精度改善，但系统开始振动。

对永磁电机，在磁路未饱和且采用相应电流定义时，电机力矩近似为 $\tau_m=K_t i_q$。$K_t$ 是力矩常数，单位 $\mathrm{N\,m/A}$；$i_q$ 是转子坐标系中主要产生力矩的电流。减速比定义为
\[
 N=\frac{\dot\theta_m}{\dot q},
\]
其中 $\theta_m$ 为电机角，$q$ 为输出关节角。理想无损齿轮满足功率守恒 $\tau_m\dot\theta_m=\tau_j\dot q$，于是 $\tau_j=N\tau_m$。若把传动损失压缩成一个近似效率 $\eta$，在指定驱动方向和工作区间可写
\begin{equation}
 \tau_j\approx N\eta K_ti_q=K_{\mathrm{eff}}i_q.
\end{equation}
$K_{\mathrm{eff}}$ 是关节侧等效力矩系数。它便于标定，但不应把方向相关摩擦、温度变化和回驱损失都当成一个永远不变的比例。制动、回驱与电机驱动下，简单效率公式也需要重新核对。

\begin{examplebox}[算例 7A 电流换算与模型偏差]
假设 $N=50$、$K_t=0.08\ \mathrm{N\,m/A}$、该工作点 $\eta=0.85$，反馈电流为 $2\ \mathrm A$。估算关节侧力矩
\[
 \tau_j=50\times0.08\times0.85\times2=6.8\ \mathrm{N\,m}.
\]
若软件忽略效率，会记录为 $8\ \mathrm{N\,m}$，比这个估算值高 $17.65\%$。辨识算法可能通过放大质量或惯量来“解释”这部分差异，得到力矩预测尚可但物理参数失真的结果。这里的效率仅为教学假设，不是任何实际关节的实测值。
\end{examplebox}

因此先核对反馈单位和标定链：电流是相电流还是 $q$ 轴电流、峰值还是有效值，输出是电机侧还是关节侧，是否已经乘过减速比。重复乘一次减速比会产生数量级错误。若设备把换算写成 $i\times k\times b\times c$，字母本身没有物理保证，必须查清每个系数的定义、单位及是否随工况改变。

\section{转子惯量为什么乘减速比的平方}
电机转子的动能为 $T_m=\tfrac12J_m\dot\theta_m^2$，理想刚性齿轮有 $\dot\theta_m=N\dot q$，代入得
\[
 T_m=\frac12J_mN^2\dot q^2=\frac12J_{m,\mathrm{out}}\dot q^2,
 \qquad J_{m,\mathrm{out}}=N^2J_m.
\]
这就是把转子惯量折算到关节输出端的公式。若转子速度包含多个关节速度的线性组合，更一般的写法是 $\dot\theta_m=A(q)\dot q$，由此增加的惯量矩阵为 $A^{\mathsf T}J_mA$，未必只是对角线。把 $N^2J_m$ 单独加在每个关节对角线上，是常用的独立传动近似。

\begin{examplebox}[算例 7B 一个很小的转子可以有明显的输出端惯量]
取转子惯量 $J_m=2\times10^{-5}\ \mathrm{kg\,m^2}$、减速比 $N=100$，则
\[
 J_{m,\mathrm{out}}=100^2\times2\times10^{-5}=0.2\ \mathrm{kg\,m^2}.
\]
若连杆对该关节的等效惯量是 $0.1\ \mathrm{kg\,m^2}$，总惯量约 $0.3\ \mathrm{kg\,m^2}$。实现 $2\ \mathrm{rad/s^2}$ 加速度所需惯性力矩由 $0.2$ 增至 $0.6\ \mathrm{N\,m}$。只看 CAD 连杆参数会漏掉其中 $0.4\ \mathrm{N\,m}$。
\end{examplebox}

高减速比有时让对角惯量增加，降低相对耦合程度，但它不会消除耦合，更不能推出“某类机器人必须用某个控制器”。合适方法取决于精度、带宽、摩擦、柔性和可用接口。低减速比也可以使用 PD；高减速比机械臂也可能从模型前馈中获益。

\section{磁场定向控制的必要数学}
FOC 是 field oriented control，中文为磁场定向控制。它把三相交流电流变换到随转子磁场旋转的坐标系，方便分别控制产生磁场和产生力矩的分量。

\subsection{先把三相投影到平面}
采用幅值不变的 Clarke 变换，三相电流满足 $i_a+i_b+i_c=0$ 时：
\begin{equation}
 \begin{bmatrix}i_\alpha\\i_\beta\end{bmatrix}
 =\frac23\begin{bmatrix}1&-1/2&-1/2\\0&\sqrt3/2&-\sqrt3/2\end{bmatrix}
 \begin{bmatrix}i_a\\i_b\\i_c\end{bmatrix}.
\end{equation}
然后使用电角度 $\theta_e$ 做 Park 变换：
\begin{equation}
 \begin{bmatrix}i_d\\i_q\end{bmatrix}
 =\begin{bmatrix}\cos\theta_e&\sin\theta_e\\-\sin\theta_e&\cos\theta_e\end{bmatrix}
 \begin{bmatrix}i_\alpha\\i_\beta\end{bmatrix}.
\end{equation}
若极对数为 $p$，则 $\theta_e=p\theta_m+\theta_0$。$\theta_0$ 是编码器机械零点与磁场零点间的偏置。电角度错误会把本想用于出力矩的电流分到错误方向，所以编码器角度正确并不足以证明 FOC 零位正确。

\begin{examplebox}[算例 7C 三相电流怎样变成力矩电流]
给定 $(i_a,i_b,i_c)=(2,-1,-1)\ \mathrm A$，先得 $(i_\alpha,i_\beta)=(2,0)\ \mathrm A$。若 $\theta_e=-\pi/2$，则
\[
 i_d=0\times2+(-1)\times0=0,\qquad
 i_q=-(-1)\times2+0\times0=2\ \mathrm A.
\]
若 $p=7$、永磁磁链 $\psi_f=0.012\ \mathrm{Wb}$，在表贴式永磁电机近似下
\[
 \tau_m=\frac32p\psi_f i_q=1.5\times7\times0.012\times2=0.252\ \mathrm{N\,m}.
\]
此处 $K_t=0.126\ \mathrm{N\,m/A}$，针对的是上述幅值不变变换下的 $i_q$ 定义，不能不经换算就与其他电流规范的铭牌常数比较。
\end{examplebox}

\subsection{电压方程与电流环}
永磁同步电机的一种常用 $dq$ 模型为
\begin{align}
 v_d&=R_si_d+L_d\dot i_d-\omega_eL_qi_q,\\
 v_q&=R_si_q+L_q\dot i_q+\omega_e(L_di_d+\psi_f).
\end{align}
$R_s$ 是定子电阻，$L_d,L_q$ 是两轴电感，$\omega_e$ 是电角速度。交叉项描述旋转坐标的耦合，$\omega_e\psi_f$ 是反电动势。一般力矩还包括磁阻项：
\begin{equation}
 \tau_m=\frac32p\left[\psi_fi_q+(L_d-L_q)i_di_q\right].
\end{equation}
当 $L_d\approx L_q$ 或 $i_d=0$，才退化为前面简单的 $K_t i_q$。

若用前馈补偿交叉项，且只分析局部电流环，可近似 $G(s)=1/(Ls+R_s)$。PI 控制器 $K_p+K_i/s$ 的零点放在 $-R_s/L$，并指定闭环带宽 $\omega_c$，有
\[
 K_p=L\omega_c,\qquad K_i=R_s\omega_c.
\]
推导是 $C(s)G(s)=\omega_c/s$，因此闭环 $CG/(1+CG)=\omega_c/(s+\omega_c)$。这依赖于模型准确、没有饱和以及可忽略延迟，不能把它当作真实驱动器的无条件整定公式。

\begin{examplebox}[算例 7D 电流 PI 的单位与数值]
取 $L=0.5\ \mathrm{mH}=0.0005\ \mathrm H$，$R_s=0.5\ \Omega$，希望局部一阶近似带宽为 $1000\ \mathrm{Hz}$，则 $\omega_c=2\pi\times1000=6283.19\ \mathrm{rad/s}$。
\[
 K_p=3.14159\ \mathrm{V/A},\qquad K_i=3141.59\ \mathrm{V/(A\,s)}.
\]
若采样或 PWM 更新带来延迟，这个闭环模型就会改变。实际可实现带宽必须结合开关频率、采样方式、计算延迟和相位裕度分析。
\end{examplebox}

\section{刚性齿轮假设失效后会发生什么}
\subsection{把电机与负载看作两个惯量}
选择同一个输出端角度单位，令 $\theta$ 是电机角折算到输出端后的角度，$q$ 是连杆角。两者之间用扭簧和阻尼连接：
\begin{align}
 J_m^{\prime}\ddot\theta+d_s(\dot\theta-\dot q)+k_s(\theta-q)&=\tau_m^{\prime},\\
 J_l\ddot q-d_s(\dot\theta-\dot q)-k_s(\theta-q)&=-\tau_L.
\end{align}
$J_m^{\prime}$、$\tau_m^{\prime}$ 都已折算到输出端，不能再额外乘一次 $N$。$k_s$ 单位是 $\mathrm{N\,m/rad}$，表示扭转刚度；$d_s$ 单位是 $\mathrm{N\,m\,s/rad}$；$\tau_L$ 是外部负载阻力矩。

静止时有 $\theta-q=\tau_L/k_s$。电机编码器显示到位，负载仍可因为弹性而偏离。这也是在力矩精度和输出位置精度之间，双编码器比单电机编码器多提供了什么信息。

\begin{examplebox}[算例 7E 编码器到位与末端到位]
若 $k_s=500\ \mathrm{N\,m/rad}$、负载力矩 $\tau_L=5\ \mathrm{N\,m}$，则扭转角 $0.01\ \mathrm{rad}=0.573^\circ$。若这一角误差作用在 $0.5\ \mathrm m$ 力臂上，末端小角度位移约 $0.5\times0.01=0.005\ \mathrm m$，即 $5\ \mathrm{mm}$。仅提高电机角度显示精度无法消除这个结构误差。
\end{examplebox}

\subsection{推导谐振与反谐振频率}
先令阻尼和外负载为零，采用零初始条件。将电机侧角度、负载侧角度与输入力矩的拉普拉斯变换分别记为 $\Theta(s)$、$Q(s)$、$\mathcal T(s)$，有：
\[
 (J_m^{\prime}s^2+k_s)\Theta-k_sQ=\mathcal T,\qquad
 -k_s\Theta+(J_ls^2+k_s)Q=0.
\]
由第二式 $Q=k_s\Theta/(J_ls^2+k_s)$，代回第一式，得
\begin{equation}
 \frac{s\Theta(s)}{\mathcal T(s)}
 =\frac{J_ls^2+k_s}{s[J_m^{\prime}J_ls^2+k_s(J_m^{\prime}+J_l)]}.
\end{equation}
分子为零对应电机侧速度响应的反谐振：
\[
 \omega_{ar}=\sqrt{k_s/J_l}.
\]
非零振荡极点对应谐振：
\[
 \omega_r=\sqrt{k_s(1/J_m^{\prime}+1/J_l)}.
\]
反谐振位置与测量输出有关：如果改看负载侧角速度，传递函数的分子不同，不能把电机侧的“谷”机械地套过去。

\begin{figure}[htbp]\centering
\includegraphics[width=0.90\linewidth]{figures/di_flexible_response.pdf}
\caption{双惯量系统的电机侧和负载侧速度响应。曲线加入小扭转阻尼使峰值有限；竖线标出无阻尼公式预测的特征频率。}
\end{figure}

取输出端电机惯量 $J_m^{\prime}=0.05\ \mathrm{kg\,m^2}$，负载惯量 $J_l=0.2\ \mathrm{kg\,m^2}$，扭转刚度 $k_s=5000\ \mathrm{N\,m/rad}$。分别得到
\[
 \omega_{ar}=158.114\ \mathrm{rad/s},\quad f_{ar}=25.16\ \mathrm{Hz},
\]
\[
 \omega_r=353.553\ \mathrm{rad/s},\quad f_r=56.27\ \mathrm{Hz}.
\]
这些频率由惯量和刚度共同决定。只要控制或轨迹在相关频段输入能量，就可能激起振动。

\section{接触以后 外力由约束共同决定}
\subsection{刚性接触的互补条件}
设 $\phi(q)$ 是机器人与环境的间隙，$\phi>0$ 表示分离，$\phi=0$ 表示贴合。理想无黏附单侧接触要求
\begin{equation}
 \phi\ge0,\qquad\lambda_n\ge0,\qquad\phi\lambda_n=0.
\end{equation}
$\lambda_n$ 为环境对机器人的法向压力。三个条件分别表示：不穿透、接触只能推不能拉、分离时没有接触力。这称为互补条件。它不是“力恒等于刚度乘位移”，因为理想刚性接触不允许压缩量。

动态方程增加 $J_c^{\mathsf T}\lambda$：
\[
 M\ddot q+c+g=\tau_{\mathrm{cmd}}+J_c^{\mathsf T}\lambda.
\]
若接触保持不动，还需满足 $J_c\ddot q+\dot J_c\dot q=0$。动力学与约束共同决定加速度和接触力；不能先随便给一个法向力，再期望它自动与运动一致。

\subsection{柔性接触与不能出现拉力的修正}
令穿透量 $\delta=\max(0,-\phi)$。简单弹簧阻尼模型在接触区采用
\[
 \lambda_n=\max(0,k_c\delta+d_c\dot\delta).
\]
$k_c$ 是接触刚度，$d_c$ 是接触阻尼。外层最大值防止卸载时出现非物理拉力，但切换本身可能带来不光滑。更复杂的材料模型可让阻尼随压缩量变化；无论选哪种，参数都需要与材料、接触形状和积分算法相匹配。

\begin{examplebox}[算例 7F 接触刚度决定很小位移对应多大力]
取 $k_c=20000\ \mathrm{N/m}$，压入 $0.5\ \mathrm{mm}=0.0005\ \mathrm m$，静态接触力为 $10\ \mathrm N$。如果目标是 $5\ \mathrm N$，理想静态压缩量仅为 $5/20000=0.25\ \mathrm{mm}$。这说明硬接触中，小位置误差也能产生明显力误差。
\end{examplebox}

\subsection{摩擦锥说明夹持能力}
接触切向摩擦满足 $\|\lambda_t\|\le\mu\lambda_n$。$\mu$ 是无量纲摩擦系数，静止时摩擦力可处于边界内部，恰好平衡需要抵抗的切向力；滑动时，在理想库仑模型下摩擦位于边界且反向于滑动速度。

\begin{figure}[htbp]\centering
\includegraphics[width=0.90\linewidth]{figures/di_contact_friction.pdf}
\caption{左图为单侧弹簧接触的力与间隙关系。右图为二维切向截面中的摩擦锥，锥内的接触力可满足无滑动摩擦约束。}
\end{figure}

“夹爪夹持力 $20\ \mathrm N$”可能指每个手指的法向力，也可能指一对手指的合计指标，必须先核对。若每个接触面均为 $20\ \mathrm N$、两个接触面、$\mu=0.6$，最大静态切向支撑力是 $2\times0.6\times20=24\ \mathrm N$，对应重力质量 $24/9.81=2.45\ \mathrm{kg}$。这是理想静态极限，尚未计入加速度、接触不均匀和实际摩擦变化。

\subsection{撞击用冲量计算}
极短时间内，接触力很大而持续时间很短，直接研究冲量 $\Lambda=\int\lambda\,dt$ 更方便。忽略冲击时刻其他有限力矩的积分，得到
\[
 M(\dot q^+-\dot q^-)=J_n^{\mathsf T}\Lambda,
\]
$-$、$+$ 分别表示撞击前后。静止环境下的法向恢复规律为 $v_n^+=-e_r v_n^-$，$e_r\in[0,1]$ 是恢复系数。于是
\begin{equation}
 \Lambda=-\frac{(1+e_r)v_n^-}{J_nM^{-1}J_n^{\mathsf T}}.
\end{equation}
分母的倒数 $m_{\mathrm{eff}}=(J_nM^{-1}J_n^{\mathsf T})^{-1}$ 称为该接触方向的有效质量。它取决于机械臂姿态和接触方向，并不等于所有连杆质量简单相加。

\begin{examplebox}[算例 7G 两连杆撞到水平面时的有效质量]
沿用本书两杆模型，$q=(0,\pi/2)$，法向取正 $y$，所以 $J_n=(0.5,0)$。此时
\[
 M^{-1}=\begin{bmatrix}3&-3\\-3&15\end{bmatrix},\qquad
 J_nM^{-1}J_n^{\mathsf T}=0.75,
\]
从而 $m_{\mathrm{eff}}=1.33333\ \mathrm{kg}$。若撞前法向速度 $-0.2\ \mathrm{m/s}$，$e_r=0.1$，则
\[
 \Lambda=\frac{1.1\times0.2}{0.75}=0.293333\ \mathrm{N\,s}.
\]
关节速度改变量为 $M^{-1}J_n^{\mathsf T}\Lambda=(0.44,-0.44)^{\mathsf T}\ \mathrm{rad/s}$，末端法向速度增加 $0.5\times0.44=0.22\ \mathrm{m/s}$，因此撞后速度是 $0.02\ \mathrm{m/s}$，恰好等于 $-0.1\times(-0.2)$。
\end{examplebox}

\section{持续力矩与热比峰值力矩更容易被忽略}
电机铜耗与电流平方近似成正比，所以周期工况常用电流有效值 $i_{\mathrm{rms}}=\sqrt{\frac1T\int_0^T i^2dt}$ 估计发热，而不是用平均电流。对举着相机的机械臂，静止时速度为零、机械输出功率 $\tau\dot q$ 为零，电流仍然产生铜耗。这解释了“没有动却越来越热”。

\begin{examplebox}[算例 7H 平均电流为零仍然发热]
前半周期 $i=3\ \mathrm A$，后半周期 $i=-3\ \mathrm A$，平均值为零，但 $i_{\mathrm{rms}}=3\ \mathrm A$。若等效电阻 $0.4\ \Omega$，相应电阻损耗量级为 $i_{\mathrm{rms}}^2R=3.6\ \mathrm W$。多相电机的总铜耗还需按实际电流定义和相数计算，不能把这个单等效电阻示例直接当成总电机损耗。
\end{examplebox}

减速器选型和电机选型应围绕实际工况：持载时的连续力矩、加速时的短时峰值、可允许的温升、弹性变形和传动损失。本书不会从一个型号类别直接推断可反驱性或力控精度；这些都应由实际关节测量确认。

\section{计算练习与分步解答}
\subsection*{练习 7.1 改变减速比后的惯量}
转子惯量 $10^{-5}\ \mathrm{kg\,m^2}$，减速比由 $50$ 改为 $80$，求输出端转子惯量的变化倍数。

\textbf{解答。}原来是 $50^2\times10^{-5}=0.025$，后来是 $80^2\times10^{-5}=0.064\ \mathrm{kg\,m^2}$，倍数 $0.064/0.025=2.56$。力矩放大只与 $N$ 一次方相关，惯量放大与 $N^2$ 相关，因此提高减速比不只带来更大的输出力矩。

\subsection*{练习 7.2 从双编码器估计扭转力矩}
输出端角为 $0.500\ \mathrm{rad}$，电机角折算到输出端为 $0.506\ \mathrm{rad}$，扭转刚度 $800\ \mathrm{N\,m/rad}$，忽略阻尼。求弹性传递力矩。

\textbf{解答。}角差为 $0.006\ \mathrm{rad}$，力矩 $800\times0.006=4.8\ \mathrm{N\,m}$。若两个编码器的相对零偏为 $0.001\ \mathrm{rad}$，就会产生 $0.8\ \mathrm{N\,m}$ 的估计偏置，因此双编码器方法同样需要零偏和刚度标定。

\subsection*{练习 7.3 接触自然频率}
有效质量 $2\ \mathrm{kg}$，接触刚度 $80000\ \mathrm{N/m}$，求无阻尼自然频率。

\textbf{解答。}$\omega_n=\sqrt{k/m}=\sqrt{40000}=200\ \mathrm{rad/s}$；$f_n=200/(2\pi)=31.83\ \mathrm{Hz}$，周期约 $31.4\ \mathrm{ms}$。积分步长与控制周期如何选择还取决于阻尼、积分方法和延迟，不能仅凭“采样比两倍频率快”就保证接触仿真和闭环稳定。

\subsection*{练习 7.4 静态夹持与向上加速}
两个指面各提供 $15\ \mathrm N$ 法向力，$\mu=0.5$。若物体要向上加速 $2\ \mathrm{m/s^2}$，理想最大质量是多少？

\textbf{解答。}总可用向上摩擦力为 $2\times15\times0.5=15\ \mathrm N$。牛顿方程要求 $F=m(g_0+a)$，因此 $m=15/(9.81+2)=1.270\ \mathrm{kg}$。静止时则为 $15/9.81=1.529\ \mathrm{kg}$。加速需求会直接消耗静态夹持余量。

\section*{本章与原稿的对应}
本章整合 [A，第33页与35--38页] 及 [B，第33页、第37--40页与116--120页]。补充了惯量折算的动能推导、明确坐标侧的双惯量模型、FOC 变换与力矩系数约定、有效质量和撞击的完整数值计算。通信和软件层面的执行链在工程实现章节继续。

\chapter{把动力学辨识写成可以计算的线性代数}\label{ch:identification}
\section{辨识不是重新发明机械臂的几何}
参数辨识（parameter identification）是利用测得的输入与运动，估计模型中的未知参数。URDF 或 DH 参数告诉我们连杆怎样连接、关节轴在哪里；它们使我们能够计算几何和动力学结构，但不保证每个质量、质心、转子惯量、摩擦和电流力矩系数都准确。辨识的任务是让这些数值与实际系统一致。

先区分三个问题。运动学标定修正关节零位和几何尺寸，主要解释位置误差；惯性参数辨识修正质量、质心和惯量，主要解释重力与加速力矩；摩擦辨识解释随速度、方向、温度和负载变化的阻力。三者可能互相污染，但它们不是同一个未知量。

本章以惯性与线性摩擦辨识为主线。目标不是让训练数据的误差尽量接近零，而是让没有用于拟合的新轨迹也能得到可靠力矩预测，并明确哪些参数实际可辨识。

\section{为什么对运动非线性 却对惯性参数线性}
\subsection{把参数选对 才能得到线性模型}
考虑绕水平轴转动的单杆。关节角从水平起量，质量 $m$，质心距离 $c$，质心惯量 $I_C$，则
\[
 \tau=(I_C+mc^2)\ddot q+mg_0c\cos q.
\]
若把未知量选为 $(m,c,I_C)$，式中出现 $mc^2$ 和 $mc$，并不是对这三个未知量线性。改用
\[
 I_O=I_C+mc^2,\qquad h=mc,
\]
就得到
\begin{equation}
 \tau=\begin{bmatrix}\ddot q&g_0\cos q\end{bmatrix}
 \begin{bmatrix}I_O\\h\end{bmatrix}.
\end{equation}
$h$ 叫一阶质量矩，单位 $\mathrm{kg\,m}$。运动量在前面的矩阵里，未知常数在后面的向量里。这种写法称为线性参数化，并不意味着运动方程对 $q$ 变成线性的。

对于三维刚体，用标准参数
\begin{equation}
 \pi_i=(m_i,h_{ix},h_{iy},h_{iz},I_{O,xx},I_{O,xy},I_{O,xz},I_{O,yy},I_{O,yz},I_{O,zz})^{\mathsf T}
\end{equation}
可以把刚体逆动力学写成 $\tau=Y(q,\dot q,\ddot q)\pi$。这里 $I_O$ 关于连杆坐标原点，而不是关于质心；所有惯量分量在同一连杆坐标系表达。库函数的参数顺序可能不同，调用时必须按其定义排列。

其根本原因是：刚体动能和重力势能对质量分布的积分是线性的，而欧拉拉格朗日的求导与相减也是线性运算。把质量积分归并成上述十个常数，剩余系数就仅依赖于几何和运动。

\subsection{一行数据到底包含什么}
加入库仑摩擦和黏性摩擦，对单关节可写
\begin{equation}\label{eq:single_reg_di}
 \tau_k=\underbrace{\begin{bmatrix}\ddot q_k&g_0\cos q_k&\operatorname{sgn}(\dot q_k)&\dot q_k\end{bmatrix}}_{H_k}
 \underbrace{\begin{bmatrix}I_O\\h\\F_c\\F_v\end{bmatrix}}_{\beta}+\epsilon_k.
\end{equation}
下标 $k$ 表示第 $k$ 次采样。$H_k$ 是回归行向量，$\beta$ 是待估参数，$\epsilon_k$ 是测量误差和模型遗漏的总残差。$F_c$ 单位 $\mathrm{N\,m}$，$F_v$ 单位 $\mathrm{N\,m\,s/rad}$。

例如在 $q=0$、$\dot q=0.4$、$\ddot q=1.2$ 时，$H_k=(1.2,9.81,1,0.4)$。如果 $\beta=(0.1,0.25,0.2,0.05)^{\mathsf T}$，则预测力矩为
\[
 1.2\times0.1+9.81\times0.25+1\times0.2+0.4\times0.05=2.7925\ \mathrm{N\,m}.
\]
每个数分别来自加速、重力、库仑摩擦和黏性摩擦，可以逐项核对。

七关节在一个采样时刻提供七行方程；采 $N$ 个时刻后，把所有行竖着堆叠，得到
\begin{equation}
 y=H\beta+\epsilon,\qquad y\in\mathbb R^{7N},\quad H\in\mathbb R^{7N\times p}.
\end{equation}
$p$ 为选用的参数个数。数据行很多不代表信息一定充分：把同一状态重复一万次，只会更精确地测到同一个方向的信息。

\section{基参数是可辨识的组合 不是某几根杆的原始参数}
\subsection{一个无法被任何轨迹分开的例子}
前面的单杆力矩只依赖于 $I_O$ 和 $h$。取以下两组参数：
\[
 (m,c,I_C)=(1,0.25,0.0375),\qquad(2,0.125,0.06875).
\]
它们都给出 $h=0.25$、$I_O=0.1$，因此在该单关节模型下，对任何 $q,\dot q,\ddot q$ 产生相同力矩。即使没有噪声，也不能仅凭这类关节运动数据区分这两组物理参数。这称为结构不可辨识。

基参数（base parameters）就是足以决定输入输出关系的一组独立参数组合。它们可用于力矩预测，但不一定能直接填写到某根连杆的质量和惯量字段里。

\subsection{用一个三列矩阵看清楚秩亏}
考虑
\[
 H=\begin{bmatrix}1&2&0\\2&4&1\\3&6&1\end{bmatrix}.
\]
第二列恒等于第一列的两倍。因此
\[
 H\begin{bmatrix}\beta_1\\\beta_2\\\beta_3\end{bmatrix}
 =\begin{bmatrix}1&0\\2&1\\3&1\end{bmatrix}
 \begin{bmatrix}\beta_1+2\beta_2\\\beta_3\end{bmatrix}.
\]
只能辨识 $\beta_{b1}=\beta_1+2\beta_2$ 和 $\beta_{b2}=\beta_3$。例如 $(1,2,3)$ 与 $(3,1,3)$ 的输出完全相同。矩阵有三列但秩只有二，直接计算 $(H^{\mathsf T}H)^{-1}$ 会失败。

\subsection{数值寻找基参数的方法}
在充分多样的无噪声模型状态上构造堆叠回归矩阵，用带列主元的 QR 分解选择独立列。若置换后的矩阵为 $[H_1\ H_2]$，其中 $H_2=H_1K$，则
\[
 H_1\pi_1+H_2\pi_2=H_1(\pi_1+K\pi_2)=H_1\pi_b.
\]
这里 $K$ 由线性方程求出；$\pi_b=\pi_1+K\pi_2$ 是基参数。SVD 同样可给出可辨识子空间，只是参数组合通常更稠密。

应同时检查：状态样本是否覆盖足够构型、各列是否合理缩放、阈值如何相对于最大奇异值设置。一个数值很小的奇异值可能表示结构依赖，也可能仅表示这次选取的状态没有激励到该方向。

\begin{pitfallbox}[七轴不自动等于四十三个惯性基参数]
基参数个数取决于关节类型、几何、基座约束、重力方向、是否含转子惯量以及参数化方式。某个七轴模型得到 $43$ 列，只说明那个模型和那套约定的结果；不能由“七轴”直接推断。每关节增加两列摩擦也不保证样本矩阵的秩恰好增加二：若速度只取一个非零值，$\operatorname{sgn}(\dot q)$ 与 $\dot q$ 列就成比例。
\end{pitfallbox}

\section{普通最小二乘从目标函数推到数字}
普通最小二乘 OLS 的目标是使平方残差之和最小：
\[
 J(\beta)=\|y-H\beta\|_2^2=(y-H\beta)^{\mathsf T}(y-H\beta).
\]
展开：
\[
 J=y^{\mathsf T}y-2\beta^{\mathsf T}H^{\mathsf T}y+\beta^{\mathsf T}H^{\mathsf T}H\beta.
\]
对 $\beta$ 求梯度并令其为零：
\[
 \nabla J=-2H^{\mathsf T}y+2H^{\mathsf T}H\hat\beta=0.
\]
因此正规方程是 $H^{\mathsf T}H\hat\beta=H^{\mathsf T}y$。仅当 $H$ 列满秩时，才能写出唯一解
\begin{equation}
 \hat\beta=(H^{\mathsf T}H)^{-1}H^{\mathsf T}y.
\end{equation}
这条式子解释数学，不推荐照着显式求逆实现。QR 或 SVD 可避免先构造病态程度被平方的矩阵。

\begin{examplebox}[算例 8A 四条数据完整求惯量与黏性系数]
假设已扣除重力与库仑摩擦，模型为 $\tau=J_r\ddot q+b\dot q$。数据矩阵与力矩为
\[
 H=\begin{bmatrix}1&1\\2&1\\1&2\\2&2\end{bmatrix},\qquad
 y=\begin{bmatrix}0.31\\0.39\\0.49\\0.62\end{bmatrix}.
\]
第一列是 $\mathrm{rad/s^2}$，第二列是 $\mathrm{rad/s}$，输出为 $\mathrm{N\,m}$。逐项相乘相加，得
\[
 H^{\mathsf T}H=\begin{bmatrix}10&9\\9&10\end{bmatrix},\qquad
 H^{\mathsf T}y=\begin{bmatrix}2.82\\2.92\end{bmatrix}.
\]
两阶逆矩阵的分母为 $10\times10-9\times9=19$，所以
\[
 \hat\beta=\frac1{19}\begin{bmatrix}10&-9\\-9&10\end{bmatrix}
 \begin{bmatrix}2.82\\2.92\end{bmatrix}
 =\begin{bmatrix}0.10105263\\0.20105263\end{bmatrix}.
\]
即 $\hat J_r=0.10105263\ \mathrm{kg\,m^2}$，$\hat b=0.20105263\ \mathrm{N\,m\,s/rad}$。再代回每条数据：
\[
 \hat y=(0.3021053,0.4031579,0.5031579,0.6042105)^{\mathsf T}.
\]
残差为 $r=y-\hat y=(0.0078947,-0.0131579,-0.0131579,0.0157895)^{\mathsf T}$。计算应满足 $H^{\mathsf T}r\approx0$：最优残差与回归矩阵列空间正交，但每个样本残差并不为零。
\end{examplebox}

\section{加权最小二乘的权重是噪声精度}
若各观测噪声方差不同，直接相加平方残差就会把不同可靠程度的数据当成一样可靠。设 $\epsilon$ 的协方差为 $\Sigma$，加权最小二乘 WLS 最小化
\begin{equation}
 J_W=(y-H\beta)^{\mathsf T}\Sigma^{-1}(y-H\beta).
\end{equation}
若噪声独立，$\Sigma=\operatorname{diag}(\sigma_1^2,\ldots,\sigma_m^2)$，选择
\[
 W=\operatorname{diag}(1/\sigma_1,\ldots,1/\sigma_m),
\]
则目标就是 $\|Wy-WH\beta\|^2$。也就是说，每行力矩和该行所有回归系数一起除以该行的噪声标准差。这里的 $W$ 是白化矩阵；目标函数中的权重是 $W^{\mathsf T}W=\Sigma^{-1}$。把方差与标准差混用，会产生错误的加权强度。

若噪声存在关节间或时间上的相关性，需要能满足 $W^{\mathsf T}W=\Sigma^{-1}$ 的一般白化变换，不再是简单逐行缩放。滤波后的残差通常存在时间相关性，不能自动把每个采样点当作独立观测。

\begin{examplebox}[算例 8B 两个测量值的最优加权平均]
两次测量同一个常数力矩，得到 $y_1=2$、$y_2=4\ \mathrm{N\,m}$；标准差分别为 $1$ 和 $2\ \mathrm{N\,m}$。目标函数
\[
 J=(2-\beta)^2+\frac14(4-\beta)^2.
\]
求导为 $2(\beta-2)+\tfrac12(\beta-4)=0$，得到
\[
 \hat\beta=\frac{2/1^2+4/2^2}{1/1^2+1/2^2}=2.4\ \mathrm{N\,m}.
\]
简单平均是 $3$；WLS 更相信标准差较小的第一个测量。权重来自噪声，不应仅因为某关节额定力矩大就人为降低它的重要性；若人为强调腕部精度，那是另一个设计目标，应与统计白化分开记录。
\end{examplebox}

\section{条件数 量纲与激励强度要一起看}
对列满秩矩阵，二范数条件数为 $\kappa_2(H)=\sigma_{\max}(H)/\sigma_{\min}(H)$。奇异值衡量不同参数方向在输出中的放大程度。最小奇异值很小时，某个参数组合对输出影响很弱，测量噪声容易掩盖它。

\subsection{为什么条件数会随单位改变}
例如 $H=\operatorname{diag}(1,0.01)$ 的条件数为 $100$。若第二个参数原来以米为单位，改用厘米，使相应回归列乘以 $0.01$，条件数就会进一步改变，物理实验本身却没有改变。带不同物理单位的列直接比较，数值可能主要反映单位选择。

选取固定参数尺度 $D=\operatorname{diag}(d_1,\ldots,d_p)$，定义 $\beta=D\vartheta$，则有
\begin{equation}
 H_s=WHD,\qquad y_s=Wy,\qquad y_s=H_s\vartheta.
\end{equation}
$\vartheta$ 是无量纲或相近量级的参数。计算条件数、比较轨迹时，$W$ 与 $D$ 应固定，否则每条轨迹各自重新归一化会掩盖绝对信息量的变化。

\subsection{条件数小并不等于信息多}
矩阵 $H_1=I_2$ 和 $H_2=0.001I_2$ 都有条件数一，但在相同输出噪声下，$H_2$ 的参数标准差是 $H_1$ 的一千倍。前者和后者“各方向一样均衡”，却不一样强。因此应联合观察最小奇异值、信息矩阵、参数协方差和实际力矩信号幅值。

当 $H$ 列满秩时，$H^{\mathsf T}H$ 的特征值是奇异值平方，所以
\[
 \kappa_2(H^{\mathsf T}H)=\kappa_2(H)^2.
\]
例如 $\kappa_2(H)=40$ 时，$\kappa_2(H^{\mathsf T}H)=1600$。若比较对象经过不同白化、缩放或剔除了不同列，这个数字就没有直接可比性。秩亏时最小奇异值为零，通常把条件数视为无穷大。

\begin{figure}[htbp]\centering
\includegraphics[width=0.91\linewidth]{figures/di_information_ellipses.pdf}
\caption{不同回归矩阵对应的参数不确定性椭圆。均衡程度由条件数描述，绝对大小还取决于总体激励强度。图中为合成二维线性模型。}
\end{figure}

\section{参数不确定性怎样计算 怎样正确解释}
在 $H$ 无误差、列满秩、噪声独立同方差且零均值的条件下，OLS 估计协方差为
\[
 \operatorname{Cov}(\hat\beta)=\sigma^2(H^{\mathsf T}H)^{-1}.
\]
如果 $\sigma$ 未知，可以用 $\hat\sigma^2=\|r\|^2/(m-p)$ 估计，$m$ 是观测行数，$p$ 是估计参数数目。分母扣掉 $p$，因为拟合参数消耗了自由度。模型偏差、滤波相关性或回归矩阵噪声存在时，这个公式不再给出完整不确定性。

算例 8A 中 $m=4,p=2$，残差平方和为 $0.0006578947$，所以
\[
 \hat\sigma=\sqrt{0.0006578947/2}=0.0181369\ \mathrm{N\,m}.
\]
参数协方差估计约为
\[
 \begin{bmatrix}0.00017313&-0.00015582\\-0.00015582&0.00017313\end{bmatrix}.
\]
两个参数标准差都约为 $0.013158$，但相对标准差不同，因为参数本身分别约为 $0.101$ 和 $0.201$。较大的负相关说明：惯量估高一些时，黏性系数估低一些，仍可部分保持当前数据的力矩预测。

新的运动状态给出回归行 $h_*$，参数不确定性引起的预测方差为
\begin{equation}
 \sigma_{\hat\tau,*}^2=h_*\operatorname{Cov}(\hat\beta)h_*^{\mathsf T}.
\end{equation}
如果要预测一次新的有噪声测量，还需加该测量噪声方差。只看每个参数的独立标准差而忽略参数相关性，会错误估计总预测误差。

\section{回归矩阵也有噪声时 最小二乘为何会偏}
前面的统计推导把 $H$ 当成准确已知。机器人数据中，$H$ 由位置、速度和加速度构造，它本身也受到编码器噪声、求导和时间错位影响。这叫变量含误差问题；增加样本数量可以降低随机波动，却不一定消除由此产生的系统偏差。

用一个标量模型看清楚原因。设真实关系为 $y=\beta x$，真实 $x$ 零均值、方差为 $s_x^2$，但实际记录的是 $\tilde x=x+\nu$；$\nu$ 为独立零均值误差，方差 $s_\nu^2$。用记录的 $\tilde x$ 做无截距最小二乘，很多样本时估计趋近于
\begin{equation}
 \hat\beta\approx\frac{\mathbb E[\tilde xy]}{\mathbb E[\tilde x^2]}
 =\frac{\beta\mathbb E[x^2]+\beta\mathbb E[\nu x]}{\mathbb E[x^2]+2\mathbb E[x\nu]+\mathbb E[\nu^2]}
 =\beta\frac{s_x^2}{s_x^2+s_\nu^2}.
\end{equation}
独立性使交叉项为零，但分母保留了额外噪声方差，导致估计向零衰减。这是一个明确的偏差，不是数值求逆不精确；换成更精确的线性代数求解器也不会消失。

\begin{examplebox}[算例 8F 加速度噪声可能把惯量估低]
假设只有惯性力矩 $y=J_r\ddot q$，真实惯量 $J_r=0.1$，真实加速度方差为 $1$，加速度估计误差方差为 $0.25$，且满足上述独立模型。则估计极限为
\[
 \hat J_r\approx0.1\frac1{1+0.25}=0.08\ \mathrm{kg\,m^2}.
\]
即使无限增加同类数据，仍可能低估 $20\%$。真实闭环系统的误差相关结构更复杂，偏差不一定总向零，但这个算例足以说明“数据很多”不能替代测量链质量。
\end{examplebox}

闭环采集还增加另一个相关性来源：测量噪声影响控制器输出，控制器输出又影响未来运动。因此不能简单认为力矩与位置噪声彼此独立。工具变量法尝试寻找一个与真实运动相关、但不与测量误差相关的替代信息；若工具本身来自同一个带噪信号，又没有验证独立性，就未必解决问题。

工程中应先减少可解释的误差：使用实际时间戳对齐，避免饱和样本，拟合实际周期运动，核对电流到力矩的比例，再评估是否需要更复杂估计方法。比较两个方法时，除了平均参数值，还要比较重复采集的波动、新轨迹预测和物理一致性，不能凭训练残差更小就认定偏差更小。

\section{物理一致性约束为什么比惯量正定更强}
\subsection{正惯量不一定来自真实质量分布}
惯量主值 $I_1,I_2,I_3$ 必须满足三角不等式：$I_1+I_2\ge I_3$，另外两种排列同样成立。例如 $\operatorname{diag}(1,1,3)$ 虽然正定，却不能作为任何非负质量分布的质心惯量，因为 $1+1<3$。

更统一的表达是伪惯量矩阵
\begin{equation}\label{eq:pseudoinertia_di}
 P_i=\begin{bmatrix}\Sigma_i&h_i\\h_i^{\mathsf T}&m_i\end{bmatrix},\qquad
 \Sigma_i=\tfrac12\operatorname{tr}(I_{O,i})I_3-I_{O,i}.
\end{equation}
这里 $\Sigma_i$ 是二阶质量矩，不是前面噪声协方差；本小节通过下标和定义区分两者。因为
\[
 P_i=\int\begin{bmatrix}r\\1\end{bmatrix}\begin{bmatrix}r^{\mathsf T}&1\end{bmatrix}dm,
\]
对任何向量 $z$，$z^{\mathsf T}P_iz=\int(z^{\mathsf T}[r;1])^2dm\ge0$，所以它必须半正定，并且质量为正。

对于理想零厚度杆或平面薄片，半正定矩阵可以是奇异的；严格正定对应更强的非退化空间质量分布。数值优化常施加一个小的正定裕量，但这属于建模与数值选择，不应声称所有理想物体都必须严格正定。

\begin{examplebox}[算例 8C 平行轴变换后的物理一致性]
质量 $m=2\ \mathrm{kg}$ 的均匀实心球半径 $0.1\ \mathrm m$，故 $I_C=(2mr^2/5)I_3=0.008I_3$。令其质心 $c=(0.2,0,0)^{\mathsf T}\ \mathrm m$，则
\[
 h=(0.4,0,0)^{\mathsf T},\qquad I_O=\operatorname{diag}(0.008,0.088,0.088).
\]
由式\eqref{eq:pseudoinertia_di}得 $\Sigma_i=\operatorname{diag}(0.084,0.004,0.004)$。消去右下角的质量块 $m$ 后，剩余矩阵称为 Schur 补；当 $m>0$ 时，其正定性等价于整个伪惯量的正定性。这里得到
\[
 \Sigma_i-\frac{hh^{\mathsf T}}m
 =\operatorname{diag}(0.084-0.16/2,0.004,0.004)=0.004I_3>0.
\]
因此该伪惯量严格正定。这里的减法把原点二阶矩还原成关于质心的二阶矩，物理意义十分直接。
\end{examplebox}

\subsection{为什么约束通常回到标准参数空间}
若最终要生成可用的仿真模型，应对每根连杆的标准参数施加约束。例如
\begin{equation}
 \min_{\pi}\ \|W(y-H_{\mathrm{full}}\pi)\|^2
 +\lambda\|D_0^{-1}(\pi-\pi_0)\|^2,\qquad P_i(\pi_i)\succeq0,\quad m_i>0.
\end{equation}
$\pi_0$ 是 CAD 先验，$D_0$ 用于量纲与先验尺度，$\lambda$ 控制对先验的信任。数据没有决定的参数方向，由先验和物理约束补充，不应宣称它们是被实验唯一测出来的。

质量为正、惯量符合三角不等式也不保证质心位于实际连杆包络内。若知道质量分布所在的几何范围，还需附加相应约束。检查若干姿态下 $M(q)>0$ 是有用的实现检查，却不能代替每根刚体参数的物理一致性论证。

\section{递推辨识和其他方法各自解决什么}
\subsection{递推最小二乘的一次完整更新}
当数据逐条到来，可用递推最小二乘 RLS 更新参数：
\begin{align}
 K_k&=\frac{P_{k-1}\varphi_k}{\lambda+\varphi_k^{\mathsf T}P_{k-1}\varphi_k},\\
 \hat\beta_k&=\hat\beta_{k-1}+K_k(y_k-\varphi_k^{\mathsf T}\hat\beta_{k-1}),\\
 P_k&=\lambda^{-1}(P_{k-1}-K_k\varphi_k^{\mathsf T}P_{k-1}).
\end{align}
$\varphi_k$ 是当前回归列向量，$\lambda\in(0,1]$ 是遗忘因子，$P$ 是信息矩阵的逆或其尺度版本。不能不加噪声尺度说明就总把 $P$ 当成真实参数协方差。

\begin{examplebox}[算例 8D 一次 RLS 更新]
只估一个参数，初值 $\hat\beta_0=2$，$P_0=1$，取 $\lambda=1$。新数据 $\varphi_1=2$、$y_1=6$。预测值为 $2\times2=4$，残差为 $2$；增益
\[
 K_1=\frac{1\times2}{1+2^2\times1}=0.4.
\]
所以 $\hat\beta_1=2+0.4\times2=2.8$，$P_1=1-0.4\times2=0.2$。它没有立刻变成 $6/2=3$，因为初始有限 $P_0$ 相当于保留先验。$\lambda=1$ 与批量最小二乘的等价，需要包含相同的初始先验约定。
\end{examplebox}

若 $\lambda<1$，旧数据逐渐被遗忘，可以跟踪慢变化参数，但长时间没有激励时，不确定性可能增长。不能只因为每周期都有新数据就每周期无条件更新；需要监测信息量、参数边界和应用效果。

\subsection{方法选择应对应具体误差来源}
\begin{longtable}{p{0.18\linewidth}p{0.34\linewidth}p{0.39\linewidth}}
\toprule 方法&主要用途&必须核对的假设\\\midrule\endhead
OLS 与 WLS&固定结构的线性参数估计&回归矩阵近似无误差，残差统计与权重匹配\\
正则化估计&激励不足时结合 CAD 等先验&先验会影响数据未确定的方向，需说明尺度\\
总体最小二乘 \mbox{TLS}&考虑回归矩阵与输出同时有误差&需要适合的噪声结构和量纲处理，普通 TLS 并非自动优于 OLS\\
工具变量法 IV&减少闭环与回归量噪声造成的偏差&工具变量要与真实回归量相关，并与误差不相关\\
最大似然 ML&利用明确的噪声概率模型&高斯输出噪声与已知协方差下可退化为 WLS\\
输出误差方法&通过仿真轨迹拟合位置或速度&通常是非线性优化，依赖初值、控制器模型与时间同步\\
能量或积分模型&减少对显式加速度的依赖&积分边界、摩擦功和可辨识参数集合会改变\\
残差学习&拟合已解释物理之外的重复误差&需要独立验证、输入范围约束，不保证范围外外推\\\bottomrule
\end{longtable}

\section{训练误差 验证误差与三种常见指标}
残差定义为 $e_k=\tau_k-\hat\tau_k$。RMSE 是均方根误差：
\[
 \mathrm{RMSE}=\sqrt{\frac1N\sum_{k=1}^N e_k^2}.
\]
它与力矩单位相同。所谓 RRMSE 没有唯一的跨文献定义，本书固定采用
\[
 \mathrm{RRMSE}_{\mathrm{rms}}=\frac{\sqrt{N^{-1}\sum e_k^2}}{\sqrt{N^{-1}\sum\tau_k^2}}\times100\%.
\]
如果别人用信号峰峰值、均值或去均值后的标准差作为分母，数值就不应直接比较。分母接近零时，相对误差可能巨大，应同时报告绝对 RMSE。

决定系数定义为
\[
 R^2=1-\frac{\sum e_k^2}{\sum(\tau_k-\bar\tau)^2},\qquad
 \bar\tau=N^{-1}\sum\tau_k.
\]
当分母非零时，$R^2=1$ 表示完美预测，$R^2=0$ 表示平方误差等同于始终预测这组数据的均值，$R^2<0$ 表示更差。验证集上的 $R^2$ 可以为负；常量信号的分母为零时没有通常意义下的该指标。

\begin{examplebox}[算例 8E 把三个指标都算一遍]
测量力矩 $\tau=(1,2,3,4)$，预测 $(1.1,1.8,3.2,3.9)$。误差为 $(-0.1,0.2,-0.2,0.1)$，平方和 $0.1$，因此
\[
 \mathrm{RMSE}=\sqrt{0.1/4}=0.158114\ \mathrm{N\,m}.
\]
测量力矩 RMS 是 $\sqrt{(1+4+9+16)/4}=\sqrt{7.5}=2.73861$，所以本书定义的 RRMSE 为 $5.7735\%$。均值为 $2.5$，去均值平方和为 $2.25+0.25+0.25+2.25=5$，所以 $R^2=1-0.1/5=0.98$。
\end{examplebox}

\begin{figure}[htbp]\centering
\includegraphics[width=0.91\linewidth]{figures/di_regression_validation.pdf}
\caption{合成验证轨迹中的模型预测与残差。残差与速度方向或加速度相关时，说明误差具有结构，不能仅用一个总体 RMSE 掩盖。}
\end{figure}

参数拟合、超参数选择和最终验证应使用可区分的数据。把同一周期随机拆成训练与验证点，可能因为强时间相关而高估泛化能力；完整独立轨迹、不同速度和姿态区域更能检验模型。误差是否可接受，应由最终任务决定：重力补偿看静态与低速力矩残差，快速跟踪看动态残差，外力观测还需考虑残差如何映射成假接触力。

\section{计算练习与分步解答}
\subsection*{练习 8.1 无法分开的摩擦参数}
只有恒速 $\dot q=0.2\ \mathrm{rad/s}$ 的数据，模型为 $F_c\operatorname{sgn}(\dot q)+F_v\dot q$。能否分别求 $F_c,F_v$？若测得摩擦为 $0.5\ \mathrm{N\,m}$，列举两组解。

\textbf{解答。}只能得到 $F_c+0.2F_v=0.5$，两列成比例，秩为一。例如 $(F_c,F_v)=(0.4,0.5)$ 或 $(0.3,1)$ 都成立。加入另一个不同非零速度，并确保其余条件一致，才可能把截距和斜率分开。

\subsection*{练习 8.2 白化必须同时处理输出和回归矩阵}
测量方程 $2\beta=6+\epsilon$，噪声标准差 $0.5$。写出白化后的方程。

\textbf{解答。}两边同除 $0.5$，得 $4\beta=12+\epsilon/0.5$，新噪声标准差为一。若只将输出变成 $12$ 而保留回归系数 $2$，就会错误估出 $\beta=6$，而原方程的解是 $3$。

\subsection*{练习 8.3 条件数与绝对误差}
$H=\operatorname{diag}(2,0.1)$，输出噪声为 $(0,0.01)^{\mathsf T}$。求条件数和参数误差。

\textbf{解答。}奇异值为 $2,0.1$，故条件数 $20$。参数误差 $H^{-1}\epsilon=(0,0.1)^{\mathsf T}$。若把整个 $H$ 乘十，条件数仍为 $20$，相同输出噪声导致的参数误差却变为 $(0,0.01)^{\mathsf T}$。这直接说明降低条件数与提高信号强度不是同一件事。

\subsection*{练习 8.4 主惯量检查}
质量为正，质心惯量主值分别为 $(0.01,0.02,0.04)\ \mathrm{kg\,m^2}$。它能否代表一个真实非负质量分布？

\textbf{解答。}虽然三个主惯量都为正，但 $0.01+0.02=0.03<0.04$，违反三角不等式，因此不能。不能只检查惯量矩阵可逆或正定就认定物理可行。

\section*{本章与原稿的对应}
本章展开 [A，第33--34页] 与 [B，第41--53页] 的线性参数化、基参数、最小二乘、协方差和辨识方法。原稿中针对某类机械臂的参数数量、仿真结果及经验阈值不作为普遍结论沿用；本章提供可逐项复算的小型数据集，并明确统计与物理假设。

\chapter{激励轨迹 信号处理与摩擦补偿}\label{ch:excitation_friction}
\section{实验轨迹决定哪些参数能被看见}
辨识不是先随便录一段运动，再期待求解器把所有参数找出来。一个参数对输出的影响必须在数据中足够明显，并且能够与其他参数的影响区分，才有可能估计准确。这个要求叫激励充分性。反复做同一段慢速轨迹，可以降低部分随机噪声，却未必改善惯量与摩擦之间的区分。

考虑 $\tau=J_r\ddot q+F_v\dot q$。如果全部数据都满足 $\ddot q=2\dot q$，则 $\tau=(2J_r+F_v)\dot q$，只能识别组合 $2J_r+F_v$。让速度变化更复杂，使加速度与速度不总成比例，才能把这两项分开。多关节机械臂还需通过姿态变化区分重力、耦合惯量和不同连杆的贡献。

\section{有限项傅里叶激励的公式与导数}
\subsection{阶数就是保留多少个谐波}
对第 $i$ 个关节，采用
\begin{equation}\label{eq:fourier_di}
 q_i(t)=q_{i0}+\sum_{k=1}^{L}\left[
 \frac{a_{ik}}{k\omega_0}\sin(k\omega_0t)-\frac{b_{ik}}{k\omega_0}\cos(k\omega_0t)\right].
\end{equation}
$\omega_0=2\pi/T$ 是基角频率，$T$ 是一个周期，$L$ 是最高谐波序号；第 $k$ 项频率为 $k/T$。$a_{ik},b_{ik}$ 在这个参数化下的单位是角速度，不是角度。它与直接把角度幅值放在正余弦前面的写法等价，但系数单位不同。

逐项求导，分母中的 $k\omega_0$ 恰好抵消一次：
\begin{align}
 \dot q_i(t)&=\sum_{k=1}^{L}[a_{ik}\cos(k\omega_0t)+b_{ik}\sin(k\omega_0t)],\\
 \ddot q_i(t)&=\sum_{k=1}^{L}k\omega_0[-a_{ik}\sin(k\omega_0t)+b_{ik}\cos(k\omega_0t)],\\
 q_i^{(3)}(t)&=-\sum_{k=1}^{L}(k\omega_0)^2[a_{ik}\cos(k\omega_0t)+b_{ik}\sin(k\omega_0t)].
\end{align}
第三阶导数叫加加速度或 jerk。虽然只增加一个较小高频谐波，速度变化可能不大，但加速度与 jerk 因频率乘子而明显增大。

七阶与九阶在这个定义下分别保留 $7$ 与 $9$ 个谐波，最高频率分别为 $7/T$ 与 $9/T$；不是把时间 $t$ 提到七次或九次。“五次傅里叶级数”通常是口语不严格的说法，应回到公式确认其含义。若七个关节各用 $L$ 个谐波，未施加约束前有 $2\times7L$ 个正余弦系数，另有可能需要优化的中心位置。

\subsection{周期连续不等于起止静止}
取 $t=0$ 或 $t=T$，有
\[
 q_i(0)=q_{i0}-\sum_k\frac{b_{ik}}{k\omega_0},\qquad
 \dot q_i(0)=\sum_ka_{ik},\qquad
 \ddot q_i(0)=\omega_0\sum_k kb_{ik}.
\]
整数谐波保证这些量在一个周期后重复，但它们不自动为零。若需要从静止开始，应加上
\[
 \sum_ka_{ik}=0,\qquad\sum_kkb_{ik}=0.
\]
若还要求指定初始位置，必须同时调整 $q_{i0}$。有些实验会先用平滑过渡进入稳定周期，然后只采集中间完整周期；这也是合理设计，但过渡段不能当作满足周期条件的数据直接平均。

\begin{examplebox}[算例 9A 两个谐波形成起止静止的轨迹]
令 $T=2\ \mathrm s$，所以 $\omega_0=\pi\ \mathrm{rad/s}$；取 $a_1=0.2$、$a_2=-0.2$、$b_1=b_2=0$、$q_0=0$。有
\begin{align*}
 q(t)&=\frac{0.2}{\pi}\sin\pi t-\frac{0.1}{\pi}\sin2\pi t,\\
 \dot q(t)&=0.2\cos\pi t-0.2\cos2\pi t,\\
 \ddot q(t)&=-0.2\pi\sin\pi t+0.4\pi\sin2\pi t.
\end{align*}
初始速度 $0.2-0.2=0$，初始加速度为零，两个周期端点的位置也为零。$t=0.5\ \mathrm s$ 时，$q=0.0636620\ \mathrm{rad}$，$\dot q=0.2\ \mathrm{rad/s}$，$\ddot q=-0.628319\ \mathrm{rad/s^2}$。整个周期的位置峰值约 $0.0826993\ \mathrm{rad}$，速度绝对峰值 $0.4\ \mathrm{rad/s}$，加速度绝对峰值约 $1.71896\ \mathrm{rad/s^2}$。
\end{examplebox}

\begin{figure}[htbp]\centering
\includegraphics[width=0.89\linewidth]{figures/di_fourier_trajectory.pdf}
\caption{算例 9A 的位置、速度、加速度与 jerk。起止位置、速度、加速度连续，但这些边界条件本身不等于与任意外部静止段的 jerk 连续。}
\end{figure}

\section{轨迹优化到底在迭代什么}
优化变量是傅里叶系数。给定一组系数，就能算轨迹、回归矩阵和约束；优化器据此改变系数，再重新计算。所谓“迭代 164 次”，通常表示优化器执行了 164 次更新或主要迭代，具体计数还取决于软件；它不是机械臂已经来回运动 164 次，也不是辨识模型一定更新了 164 次。

设缩放白化后的矩阵为 $H_s$，信息矩阵为 $F=H_s^{\mathsf T}H_s$。常见目标有：
\begin{align}
 \text{条件数目标：}&\quad\min\ \frac{\sigma_{\max}(H_s)}{\sigma_{\min}(H_s)},\\
 \text{D 最优：}&\quad\max\ \log\det F,\\
 \text{A 最优：}&\quad\min\ \operatorname{tr}(F^{-1}),\\
 \text{E 最优：}&\quad\max\ \lambda_{\min}(F).
\end{align}
D 最优压小参数置信椭球的体积，A 最优压小方差总和，E 最优强化最弱参数方向。它们都依赖固定的噪声和参数尺度。D 最优在固定可逆参数坐标变换下的最优轨迹不变，因为行列式只增加一个与轨迹无关的常数因子；若每次轨迹都重新选缩放，则不能据此声称不受影响。

约束至少应对应实际模型中的位置、速度、加速度、jerk 或驱动器可实现性、力矩、碰撞和工作空间。离散网格上的约束满足不保证网格之间也满足；尤其高谐波轨迹，应使用更密的复核网格或连续时间上界。对碰撞距离，仅检查关节角极值更远远不够。

把轨迹整体放慢 $s$ 倍，即 $q_{\mathrm{slow}}(t)=q(t/s)$，则速度变为原来的 $1/s$，加速度变为 $1/s^2$，jerk 变为 $1/s^3$。这可以降低动态要求，却同时减弱惯量参数的激励；慢速试运行可以核对运动，但不能据此断定正式辨识数据仍有足够信息量。

\section{采样 时间同步与微分噪声}
\subsection{时间错位会制造不存在的模型误差}
如果力矩与状态相差 $\delta t$，一阶近似给出 $\tau(t+\delta t)-\tau(t)\approx\dot\tau(t)\delta t$。仅看重力部分，$\dot g=(\partial g/\partial q)\dot q$。若某工作点的重力梯度 $40\ \mathrm{N\,m/rad}$、速度 $1\ \mathrm{rad/s}$、时间错位 $2\ \mathrm{ms}$，伪力矩误差就是 $40\times1\times0.002=0.08\ \mathrm{N\,m}$。

时间戳应明确表示采样时刻、发送时刻还是接收时刻。互相关可帮助发现大致错位，但峰值不为零也可能来自执行器动态或模型误差；不能不分析物理链条就通过任意平移数据把残差调到最小。

\subsection{求导为什么放大量化误差}
设每个位置样本独立带有均匀量化误差，量化间隔为 $\Delta$，则位置误差方差为 $\Delta^2/12$。一阶后向差分
\[
 \hat v_k=\frac{q_k-q_{k-1}}{T_s}
\]
使两个独立误差相减，方差相加，因此
\[
 \sigma_v^2=\frac{2(\Delta^2/12)}{T_s^2},\qquad
 \sigma_v=\frac{\Delta}{\sqrt6T_s}.
\]
二阶差分 $(q_{k+1}-2q_k+q_{k-1})/T_s^2$ 的三个噪声系数平方和为 $1+4+1=6$，所以
\[
 \sigma_a=\frac{\Delta}{\sqrt2T_s^2}.
\]
加速度噪声按采样周期的平方反比增长。更快采样提供更多信息，但直接使用未经处理的二阶差分，并不自动得到更好的加速度。

\begin{examplebox}[算例 9B 十七位编码器的差分误差量级]
一圈 $2^{17}$ 个计数时，$\Delta=2\pi/2^{17}=4.79369\times10^{-5}\ \mathrm{rad}$。若 $T_s=1\ \mathrm{ms}$，则 $\sigma_v\approx0.01957\ \mathrm{rad/s}$，$\sigma_a\approx33.90\ \mathrm{rad/s^2}$。若 $T_s=5\ \mathrm{ms}$，则分别约为 $0.003914\ \mathrm{rad/s}$ 和 $1.356\ \mathrm{rad/s^2}$。这只是独立量化噪声模型的估计，真实误差还可能有相关性、机械振动和编码器周期误差。
\end{examplebox}

\subsection{可以怎样估计速度和加速度}
离线可对位置做平滑后差分、Savitzky--Golay 局部多项式拟合，或在周期实验中拟合谐波再解析求导。局部多项式的本质是在一个窗口内拟合 $q(t)\approx c_0+c_1t+c_2t^2+\cdots$，窗口中心速度为 $c_1$、加速度为 $2c_2$。窗口越宽通常越降噪，却越难保留快速真实变化。

谐波拟合则使用已知频率的正余弦基函数，估计实际测量位置的系数。应拟合测得的位置，不能直接把指令轨迹当成真实运动；否则跟踪滞后和驱动器饱和会被错误归到模型参数。真实运动包含未建模谐波、漂移或非周期过渡时，过低的谐波阶数也会产生偏差。

\section{为什么力矩和回归矩阵必须同滤波}
假设 $y(t)=H(t)\beta+\epsilon(t)$，参数 $\beta$ 在这段实验中恒定。对任意线性滤波算子 $\mathcal F$，有
\begin{equation}
 \mathcal F[y]=\mathcal F[H\beta]+\mathcal F[\epsilon]
 =(\mathcal F[H])\beta+\mathcal F[\epsilon].
\end{equation}
因此可在力矩和每一列最终回归特征上应用相同滤波。这里包括惯性列、重力列以及摩擦列，后续使用同一套边缘样本掩码。

\begin{pitfallbox}[先滤位置再造矩阵 不等于对回归矩阵滤波]
一般有 $H(\mathcal F[q])\ne\mathcal F[H(q)]$，因为 $H$ 含正弦、余弦、速度乘积和符号函数等非线性。位置平滑用于得到可靠状态；构造最终回归矩阵后，对其各列和力矩同滤波，是另一件事。只把位置和力矩使用同一个滤波器，并不足以自动保证最终辨识方程两边一致。
\end{pitfallbox}

周期运动可能在回归矩阵里产生更高频率。例如 $\dot q=\sin\omega t$，则 $\dot q^2=(1-\cos2\omega t)/2$，多出二倍频；$\cos(q(t))$ 更可能包含多个谐波。故滤波截止频率不能只根据位置轨迹最高谐波机械确定。

\subsection{四阶双向低通到底是什么意思}
以四阶 Butterworth 数字低通为例，单向传递响应为 $H_f(e^{j\Omega})$。前向再反向滤波，在远离边界且近似平稳的内部区间，其有效频率响应是
\[
 H_{\mathrm{fb}}(e^{j\Omega})=H_f(e^{j\Omega})H_f(e^{-j\Omega})=|H_f(e^{j\Omega})|^2.
\]
相位相互抵消，但幅值也被再乘一次。单向设计截止频率处幅值为 $1/\sqrt2$，双向后为 $1/2$，即约 $-6.02\ \mathrm{dB}$，不是 $-3.01\ \mathrm{dB}$。说“等效八阶”仅指渐近幅度滚降，不意味着它是一个可实时运行的普通因果八阶滤波器。

\begin{examplebox}[算例 9C 二百赫兹采样下八赫兹双向滤波]
取 $f_s=200\ \mathrm{Hz}$，用四阶数字 Butterworth 设计单向截止频率 $8\ \mathrm{Hz}$。计算双向有效幅值，在 $4,8,12,16\ \mathrm{Hz}$ 处分别约为
\[
 0.99623,\quad0.50000,\quad0.035674,\quad0.0034228.
\]
因此 $8\ \mathrm{Hz}$ 的真实信号已经减半，$12\ \mathrm{Hz}$ 成分只剩约 $3.57\%$。这能显著降噪，也可能删除与惯量或柔性有关的信息。该数字说明的是一组指定设计，不是所有“八赫兹低通”的共同表现。
\end{examplebox}

\begin{figure}[htbp]\centering
\includegraphics[width=0.92\linewidth]{figures/di_filter_eight_hz.pdf}
\caption{指定采样率和设计方法下的四阶单向与双向滤波响应。双向滤波无内部相位延迟，但幅值平方、边缘效应和非因果性仍需处理。}
\end{figure}

选 $8\ \mathrm{Hz}$ 是否合适，应比较几个候选截止频率下的真实信号频谱、回归列能量、奇异值、参数稳定性和独立验证误差。若参数随截止频率大幅变化，应怀疑模型、噪声、对齐或高频真实动态。双向滤波需要未来样本，仅用于离线处理；在线控制应使用因果滤波并把其相位延迟计入闭环。

\section{不显式求加速度的积分回归}
前面的方法通过平滑得到加速度。另一条思路是先把动力学在时间窗口内积分，把加速度项变成速度差。对已扣除重力的单关节模型
\[
 \tau=J_r\dot v+F_vv+F_c\operatorname{sgn}(v),
\]
从 $t_0$ 积分到 $t_1$，有
\begin{equation}
 \int_{t_0}^{t_1}\tau\,dt
 =J_r[v(t_1)-v(t_0)]+F_v[q(t_1)-q(t_0)]
 +F_c\int_{t_0}^{t_1}\operatorname{sgn}(v)\,dt.
\end{equation}
未知参数依然线性，但回归量不再包含显式加速度。若窗口内始终正向运动，最后一个积分就是窗口时长；若发生换向，则必须根据速度符号分别计算。重力未扣除时，还需保留 $h g_0\int\cos q\,dt$，不能因为用了积分模型就把重力删掉。

\begin{examplebox}[算例 9G 用一个时间窗口核对积分动力学]
取 $J_r=0.1$、$F_v=0.2$、$F_c=0.3$，窗口时长 $0.2\ \mathrm s$，速度以常加速度由 $0.2$ 增至 $0.4\ \mathrm{rad/s}$。速度差为 $0.2$，位置增量为平均速度乘时间 $0.3\times0.2=0.06\ \mathrm{rad}$，全程正向。因此
\[
 \int\tau dt=0.1\times0.2+0.2\times0.06+0.3\times0.2=0.092\ \mathrm{N\,m\,s}.
\]
直接按瞬时力矩计算：加速度为一，力矩从 $0.1+0.2\times0.2+0.3=0.44$ 线性增至 $0.48$，平均值 $0.46$，乘以窗口长度同样得到 $0.092$。
\end{examplebox}

一个窗口只提供一条方程，不能单独估出三个参数。不同窗口还需具有独立的速度变化、位移和持续时间特征。相邻重叠窗口的噪声彼此相关，窗口太长也会把不同参数影响混在一起，并累积传感器偏置。

多关节系统不能简单把 $M(q)\ddot q$ 积成 $M\Delta\dot q$，因为 $M$ 随时间改变。正确分部积分含有
\[
 \int M\ddot q\,dt=[M\dot q]_{t_0}^{t_1}-\int\dot M\dot q\,dt.
\]
必须把额外项与原有速度项一致处理，或从已知线性回归结构推导滤波动力学、动量或能量回归。积分法降低了对直接加速度的依赖，却没有免除对模型、时间同步和参数可辨识性的要求。

\section{库仑与黏性摩擦怎样从曲线中读出来}
对非零速度，最小模型为
\begin{equation}
 \tau_f(v)=F_c\operatorname{sgn}(v)+F_vv.
\end{equation}
在正速度一侧，它是一条截距 $F_c$、斜率 $F_v$ 的直线；负速度一侧截距为 $-F_c$。库仑摩擦近似与速度大小无关，黏性摩擦随速度线性增加。这里 $\tau_f$ 放在动力学左侧，耗散要求 $\tau_fv\ge0$。

\begin{examplebox}[算例 9D 通过三种速度辨识摩擦直线]
在正速度 $0.2,0.5,1.0\ \mathrm{rad/s}$ 下，扣除其他动力学项后的平均阻力分别为 $0.32,0.35,0.40\ \mathrm{N\,m}$。任取第一与第三点，斜率
\[
 F_v=\frac{0.40-0.32}{1.0-0.2}=0.10\ \mathrm{N\,m\,s/rad}.
\]
截距 $F_c=0.32-0.10\times0.2=0.30\ \mathrm{N\,m}$。第二点预测为 $0.30+0.10\times0.5=0.35$，恰好吻合。有噪声时应对所有点一起做最小二乘，并给出重复测量离散程度。
\end{examplebox}

\subsection{正反向半和半差为什么有条件}
同一姿态、相同速度绝对值、加速度近似为零，且摩擦对称时：
\[
 \tau_+=g+F_c+F_vv,\qquad\tau_-=g-F_c-F_vv.
\]
因此 $g=(\tau_++\tau_-)/2$，$F_c+F_vv=(\tau_+-\tau_-)/2$。若 $g=2.5$，$F_c=0.3$，$F_v=0.1$，$v=0.5$，则 $\tau_+=2.85$、$\tau_-=2.15$，半和为 $2.5$、半差为 $0.35\ \mathrm{N\,m}$。

如果正反向数据来自不同角度，重力项可能不同；如果电缆力矩和摩擦不对称，半和还包含其偶对称部分；如果其他关节在运动，速度耦合也可能未消失。单姿态静止电流更不能直接视为准确重力，因为静摩擦可以承担一部分保持力矩。

\section{静摩擦不能简单等同于符号函数在零点取零}
静止时，摩擦不是固定输出一个数，而是在允许范围内恰好抵消外加的剩余驱动力。简化写作
\begin{equation}
 v=0:\quad\tau_f\in[-F_s,F_s],\qquad
 \tau_f=\tau_{\mathrm{cmd}}+\tau_{\mathrm{ext}}-g\quad\text{若静止能保持}.
\end{equation}
$F_s$ 是起动阈值。只有当所需保持摩擦超过这个区间，物体才开始滑动。把 $\operatorname{sgn}(0)$ 设为零，是某个数值函数的定义，不是静摩擦物理规律。

例如 $F_s=0.5\ \mathrm{N\,m}$、扣除重力后的驱动力为 $0.3\ \mathrm{N\,m}$，关节仍可能不动；实际静摩擦此时是 $0.3$，不是 $0.5$，也不是零。这直接解释低刚度控制中的回正死区：恢复力矩变小到静摩擦能平衡的范围，系统可能停在非零误差处。

Stribeck 曲线描述从起动到滑动过程中，摩擦随低速增加先下降的现象。一种常见滑动模型为
\begin{equation}
 \tau_f(v)=\left[F_c+(F_s-F_c)e^{-(|v|/v_s)^\alpha}\right]\operatorname{sgn}(v)+F_vv,\qquad v\ne0.
\end{equation}
$v_s$ 是过渡速度尺度，$\alpha$ 控制曲线形状。这个静态速度函数仍不能独立描述真正的预滑动记忆、迟滞和粘滑循环。

\begin{examplebox}[算例 9E 低速处摩擦并不一定更小]
取 $F_s=0.5$、$F_c=0.3\ \mathrm{N\,m}$、$F_v=0.1\ \mathrm{N\,m\,s/rad}$、$v_s=0.05\ \mathrm{rad/s}$、$\alpha=2$。在 $v=0.02$ 时
\[
 \tau_f=0.3+0.2e^{-0.16}+0.002=0.472429\ \mathrm{N\,m}.
\]
在 $v=0.1$ 时
\[
 \tau_f=0.3+0.2e^{-4}+0.01=0.313663\ \mathrm{N\,m}.
\]
速度增加五倍，摩擦反而先降低。这一段负斜率与弹性、惯性共同作用，可能形成走停交替的粘滑现象。
\end{examplebox}

\begin{figure}[htbp]\centering
\includegraphics[width=0.91\linewidth]{figures/di_friction_models.pdf}
\caption{库仑加黏性、Stribeck 与平滑补偿曲线。零速竖直线段表示静摩擦允许的集合，而非零速时只有一个固定摩擦值。}
\end{figure}

若必须描述预滑动与记忆，可引入 LuGre 等动态摩擦模型。一个常见形式是
\begin{align}
 \dot z&=v-\frac{|v|}{g_z(v)}z,\\
 \tau_f&=\sigma_0z+\sigma_1\dot z+\sigma_2v.
\end{align}
$z$ 表示微观接触的平均弹性变形，$g_z(v)$ 是由稳态摩擦曲线确定的长度尺度，$\sigma_0$ 为微接触刚度，$\sigma_1$ 为微阻尼。若稳态摩擦幅值为 $F(v)$，可令 $g_z(v)=F(v)/\sigma_0$。必须说明这一单位关系，否则不同文献中 $g(v)$ 是力还是长度会混淆。动态模型新增状态和参数，应由数据支持其必要性，不能仅因名称更复杂就直接使用。

\section{联合估计与交替估计怎样配合}
设模型 $y=Y\pi+F\gamma$，其中 $\pi$ 为惯性参数，$\gamma$ 为线性摩擦参数。联合最小二乘直接解
\[
 \min_{\pi,\gamma}\|W(y-[Y\ F][\pi;\gamma])\|^2.
\]
如果先忽略摩擦拟合惯性，摩擦残差可能投影到 $Y$ 的列空间；之后用残差估摩擦，再固定摩擦重估惯性，可以逐步减轻相互污染。但不能保证一两轮就收敛，也不能通过迭代自动创造缺失的信息。

对一个列满秩的纯线性凸二次问题，精确的交替块最小化可收敛到联合最优解；如果惯性与摩擦特征强相关，可能收敛很慢。若包含 $v_s$ 等非线性参数、数据剔除规则变化或饱和模型，甚至可能出现局部极小或不稳定的反复更新。因此，线性部分优先联合解，低速非线性部分再利用专门实验估计，并在同一验证集比较。

\section{补偿不是把拟合曲线原封不动全部加回去}
为了减少速度噪声引发的换向跳变，可用
\[
 \tau_{\mathrm{ff},f}=\gamma\left[F_c\tanh(v/\varepsilon)+F_vv\right],\qquad0\le\gamma\le1.
\]
$\varepsilon$ 是速度平滑尺度，$\gamma$ 是补偿比例。这能让零附近力矩连续，但不会准确再现静摩擦集合。$\varepsilon$ 越小，零附近斜率 $\gamma F_c/\varepsilon$ 越大，测量噪声更容易引发明显补偿变化。

若补偿基于实测速度，且模型估计过大，它会沿运动方向主动注入过多能量，表现得像负阻尼。若补偿基于期望速度，停滞时又可能与真实运动方向不一致。具体选择要结合控制结构、速度估计和目标工况；补偿后应观察力矩、速度、能量与温度变化，而不只看一次回正误差。

\begin{examplebox}[算例 9F 平滑尺度决定零附近的敏感度]
取 $F_c=0.3\ \mathrm{N\,m}$、$\gamma=0.8$、$\varepsilon=0.02\ \mathrm{rad/s}$。零附近斜率为 $0.8\times0.3/0.02=12\ \mathrm{N\,m\,s/rad}$。若把 $\varepsilon$ 缩小到 $0.002$，斜率增至 $120$。同样的 $0.001\ \mathrm{rad/s}$ 速度噪声，在局部线性近似下可引起约 $0.012$ 与 $0.12\ \mathrm{N\,m}$ 的补偿变化，差十倍。
\end{examplebox}

\section{面向七轴低阻抗的结果应该怎样记录}
辨识交付物至少应明确：几何和负载版本、关节顺序与符号、反馈力矩定义、时间对齐、温度范围、滤波器完整设计、样本剔除条件、基参数变换、噪声白化与列缩放、训练及验证轨迹、各关节绝对和相对误差。条件数应附矩阵定义，不能只留下一个“40”。

腕部关节力矩小，电流偏置、线缆拖拽和负载质心误差可能占更大比例。末端误差还要经过雅可比与刚度映射，不能仅凭“J5--J7 摩擦估计大”就认定机械摩擦异常。应把残差按速度、姿态、加速度、温度和运动方向分别作图，确认其结构，再决定增加哪一类模型。

\section{计算练习与分步解答}
\subsection*{练习 9.1 七阶与九阶的最高频率}
周期 $T=10\ \mathrm s$，分别使用七阶、九阶傅里叶级数。求基频和最高频率。

\textbf{解答。}基频 $f_0=1/T=0.1\ \mathrm{Hz}$。七阶最高 $7f_0=0.7\ \mathrm{Hz}$，九阶最高 $0.9\ \mathrm{Hz}$。如果把周期改成 $2\ \mathrm s$，两者变为 $3.5$ 和 $4.5\ \mathrm{Hz}$。阶数本身不能单独说明轨迹快慢。

\subsection*{练习 9.2 同一轨迹放慢两倍}
某轨迹速度、加速度、jerk 峰值为 $0.8$、$3$、$20$，单位依次为 $\mathrm{rad/s}$、$\mathrm{rad/s^2}$、$\mathrm{rad/s^3}$。整体放慢两倍后是多少？

\textbf{解答。}时间缩放 $s=2$，三者分别除以 $2$、$4$、$8$，得到 $0.4$、$0.75$、$2.5$。位置范围不变，动态惯性力矩却通常显著下降。

\subsection*{练习 9.3 从正反向数据分离对称摩擦}
同一姿态、速度 $\pm0.4\ \mathrm{rad/s}$ 时，稳态力矩分别为 $3.7$ 和 $2.9\ \mathrm{N\,m}$。在对称模型假设下求重力与摩擦幅值。

\textbf{解答。}半和 $(3.7+2.9)/2=3.3$，半差 $(3.7-2.9)/2=0.4\ \mathrm{N\,m}$。这次数据只确定 $F_c+0.4F_v=0.4$，不能分别确定两者；还需要其他速度。

\subsection*{练习 9.4 低刚度下的静摩擦误差范围}
单关节恢复力矩为 $-K_pe$，静摩擦阈值 $F_s=0.4\ \mathrm{N\,m}$，$K_p=8\ \mathrm{N\,m/rad}$，无其他偏置。求可能静止的误差区间。

\textbf{解答。}只要 $|K_pe|\le F_s$，静摩擦就可能平衡恢复力矩。所以 $|e|\le0.4/8=0.05\ \mathrm{rad}=2.865^\circ$。这是一个理想静态可能区间，不是预测每次一定停在其边缘。若存在重力模型残差，区间中心还会移动。

\section*{本章与原稿的对应}
本章整合 [A，第33--34页] 与 [B，第42--53页]，重点补足傅里叶边界条件、优化迭代的含义、差分噪声推导、同滤波的线性代数依据、双向滤波的有效幅值以及摩擦分离的假设。原稿中的固定截止频率倍数、固定条件数目标和固定收敛轮数均改为需要由模型和数据判断的设计选择。

\chapter{负载 关节零位 工具与相机标定}\label{ch:calibration}
\section{先明确究竟是哪一个模型不准确}
机械臂末端偏离目标，不一定是同一种误差。关节反馈角与驱动器目标角不一致，属于跟踪或执行问题；角度一致而末端实测位置不一致，可能来自几何、零位、工具变换或结构弹性；位置基本正确而低阻抗时向某方向偏，可能与重力、负载、摩擦或外力估计有关。

标定（calibration）是在一个选定模型内，通过参考测量确定参数。它与闭环控制的职责不同：控制器尽力实现自己知道的目标，标定则让“控制器知道的位置或力”与真实物理量一致。没有独立参考，只比较同一套模型输出与该模型自己的目标，不能证明绝对精度。

\begin{longtable}{p{0.20\linewidth}p{0.32\linewidth}p{0.38\linewidth}}
\toprule 问题&主要未知参数&需要的参考信息\\\midrule\endhead
负载重力补偿&负载质量及一阶质量矩&不同姿态的力矩或六维力传感器数据\\
负载动态补偿&负载质量 质心与惯量&足够变化的平动和转动加速度\\
关节零位标定&编码器零点与实际几何零点差&外部位置或姿态测量\\
TCP 标定&法兰到工具工作点的位置&同一固定点在不同法兰姿态下的约束\\
手眼标定&法兰与相机之间的刚体变换&机器人相对运动与相机相对运动\\
相机内参标定&焦距 主点及畸变参数&图像中的已知几何目标\\\bottomrule
\end{longtable}

\section{负载重力为什么只需要四个惯性组合}
\subsection{先把重力写在法兰坐标系}
设世界或基座坐标系为 $B$，法兰坐标系为 $F$。$R_{BF}$ 把法兰分量变成基座分量，重力加速度在基座中为 $a_g^B$。则
\[
 a_g^F=R_{BF}^{\mathsf T}a_g^B.
\]
负载质量为 $m_p$，质心在法兰坐标中的位置为 $c_p$，一阶质量矩为 $h_p=m_pc_p$。负载重力作用到法兰的力旋量，按本书 $[f;m]$ 顺序写为
\begin{equation}
 w_g^F=\begin{bmatrix}m_pa_g^F\\h_p\times a_g^F\end{bmatrix}
 =\underbrace{\begin{bmatrix}a_g^F&0_{3\times3}\\0_{3\times1}&-[a_g^F]_\times\end{bmatrix}}_{A_p(R_{BF})}
 \begin{bmatrix}m_p\\h_p\end{bmatrix}.
\end{equation}
$[a]_\times b=a\times b$。因为 $h\times a=-a\times h$，矩阵下半部分出现负号。四个未知量就是质量一个数、一阶矩三个数；静止数据不包含角加速度和角速度，无法据此确定全部转动惯量。

如果 $J_F$ 将关节速度映射到关于法兰点、在法兰系表达的速度旋量，负载重力对机械臂的广义作用力是 $J_F^{\mathsf T}w_g^F$，所需重力补偿为
\begin{equation}
 \Delta g(q)=-J_F(q)^{\mathsf T}w_g^F.
\end{equation}
这是“环境重力作用力”与“抵消它的驱动力矩”的区别。测力传感器给出的究竟是哪一侧作用哪一侧的力，要按其符号约定核对，不能把传感器读数直接当成补偿力矩。

\begin{examplebox}[算例 10A 两种方向就能看出一阶质量矩的作用]
假设负载 $m_p=2\ \mathrm{kg}$，$h_p=(0.2,-0.1,0.1)^{\mathsf T}\ \mathrm{kg\,m}$。第一姿态的法兰系重力为 $(0,0,-9.81)^{\mathsf T}$，则
\[
 f_1=(0,0,-19.62)^{\mathsf T}\ \mathrm N,\qquad
 m_1=(0.981,1.962,0)^{\mathsf T}\ \mathrm{N\,m}.
\]
第二姿态中重力为 $(-9.81,0,0)^{\mathsf T}$，则
\[
 f_2=(-19.62,0,0)^{\mathsf T},\qquad m_2=(0,-0.981,-0.981)^{\mathsf T}.
\]
若这些都是已去零偏、符号正确的静态载荷读数，可依次解出
\[
 m_p=-f_{1z}/9.81=2,\quad h_x=m_{1y}/9.81=0.2,
\]
\[
 h_y=-m_{1x}/9.81=-0.1,\quad h_z=-m_{2y}/9.81=0.1.
\]
再由 $c_p=h_p/m_p$ 得质心 $(0.1,-0.05,0.05)^{\mathsf T}\ \mathrm m$。第一种姿态单独不能看出沿重力方向的质心分量 $h_z$，第二种方向补足了信息。
\end{examplebox}

\begin{figure}[htbp]\centering
\includegraphics[width=0.90\linewidth]{figures/di_payload_gravity.pdf}
\caption{同一负载在不同法兰姿态下，重力在法兰坐标中的方向改变，从而暴露不同质心分量。图中力和力矩均按“负载作用于法兰”定义；几何为 $x$-$z$ 投影示意，箭头仅表示方向。}
\end{figure}

\subsection{多姿态数据怎样堆成线性方程}
将各姿态的 $A_p$ 竖向堆叠，得到 $w=A\beta_p$，其中 $\beta_p=(m_p,h_x,h_y,h_z)^{\mathsf T}$。使用上一章的 WLS 即可估计。若传感器还有未知固定零偏 $b_w$，应使用
\[
 w_i=A_{p,i}\beta_p+b_w+\epsilon_i,
\]
并把相应的常数列也加入回归。此时需要更多且方向丰富的数据，不能把一次静止读数的全部偏置都解释成质量。

使用关节电流推算负载时，同样可写 $\Delta\tau_i=-J_{F,i}^{\mathsf T}A_{p,i}\beta_p$。但负载前后差分还可能混入摩擦和线缆变化；静止关节中的静摩擦并不因两次数据相减就必然消失。可以比较正反接近同一姿态的结果，记录其差异，并通过多姿态冗余拟合检查残差。

需要多少姿态由矩阵的秩和条件数决定，不是固定的“三个”“六个”或“十个”。姿态很多却都绕同一轴小幅变化，仍可能缺少信息；两个方向在理想六维力传感器算例中足以解出四个组合，也不意味着真实系统两次测量就足够鲁棒。

\section{负载改变的力矩量级应先手算}
\begin{examplebox}[算例 10B 半公斤工具对两连杆重力的影响]
在本书两连杆末端再装质量 $0.5\ \mathrm{kg}$ 的工具，质心沿第二杆轴线超出末端 $0.1\ \mathrm m$。水平伸直时，质心到第一关节的力臂为 $0.5+0.5+0.1=1.1\ \mathrm m$，到第二关节为 $0.5+0.1=0.6\ \mathrm m$。
\[
 \Delta g_1=0.5\times9.81\times1.1=5.3955\ \mathrm{N\,m},
\]
\[
 \Delta g_2=0.5\times9.81\times0.6=2.9430\ \mathrm{N\,m}.
\]
原机械臂该姿态的重力补偿为 $(9.81,2.4525)$，加负载后为 $(15.2055,5.3955)\ \mathrm{N\,m}$。若 $q=(0,\pi/2)$，工具质心在第二关节正上方，第二关节附加重力力矩为零，第一关节则为 $0.5\times9.81\times0.5=2.4525\ \mathrm{N\,m}$。
\end{examplebox}

这个计算说明：同一个负载在不同姿态下造成的偏差不同。低阻抗下，近似平衡误差为模型残差除以刚度，因此“某个方向比较准，另一个方向回正不到位”可能来自姿态相关重力误差，而非单纯控制增益问题。

\begin{examplebox}[算例 10C 单位错误怎样放大负载力矩]
$1\ \mathrm{kg}$ 负载的水平力臂 $0.2\ \mathrm m$，重力矩为 $1.962\ \mathrm{N\,m}$。若误把 $20\ \mathrm{cm}$ 写成 $2\ \mathrm m$，就变成 $19.62\ \mathrm{N\,m}$。质量、距离、力矩三项应进行量纲和数量级检查，不能只相信程序返回了一个有限数。
\end{examplebox}

\section{把负载加入连杆之前先统一参考点}
已知负载在自身质心坐标系中的惯量 $I_{Cp}$，其质心和朝向在末连杆坐标系 $L$ 中分别为 $c_p^L$、$R_{LP}$。先旋转到连杆坐标，再平移到连杆原点：
\begin{align}
 I_{Cp}^L&=R_{LP}I_{Cp}R_{LP}^{\mathsf T},\\
 I_{Op}^L&=I_{Cp}^L+m_p\left(\|c_p^L\|^2I_3-c_p^L(c_p^L)^{\mathsf T}\right),\\
 h_p^L&=m_pc_p^L.
\end{align}
如果原末连杆标准参数已经关于同一原点、在同一坐标系，则质量、一阶矩和原点惯量可以逐项相加。若一个参数包关于法兰，另一个关于连杆原点，不能直接相加。组合后要写回质心惯量时，先算
\[
 m'=m_l+m_p,\qquad c'=\frac{h_l+h_p}{m'},
\]
再用反向平行轴定理
\[
 I_C'=I_O'-m'\left(\|c'\|^2I_3-c'(c')^{\mathsf T}\right).
\]
只关心重力时，质量和一阶矩足够；涉及快速转动和角加速度时，才需要可靠负载惯量。静态标定不应被包装为已完成全部十个负载参数辨识。

\section{关节零位与几何误差怎样线性化}
设几何参数向量 $\theta$ 包括关节零偏、连杆尺寸等，外部测得位置 $p^{\mathrm{meas}}$，模型为 $p=f(q,\theta)$。在当前参数 $\theta_0$ 附近作一阶展开：
\begin{equation}
 p^{\mathrm{meas}}-f(q,\theta_0)\approx
 J_\theta(q,\theta_0)\Delta\theta,\qquad
 J_\theta=\frac{\partial f}{\partial\theta}.
\end{equation}
把不同姿态的方程堆起来，又是最小二乘。更新参数后重新线性化，这就是常见的迭代非线性最小二乘思路。误差很大时一次线性修正可能不够，需要逐步更新而不能过度外推。

\subsection{零位误差就是雅可比的一种输入误差}
如果只估关节零偏 $\delta q$，局部有 $\Delta p\approx J(q)\delta q$。

\begin{examplebox}[算例 10D 由位置偏差估计两个关节零偏]
在本书标准姿态 $q=(0,\pi/2)$，测得相对模型的位置偏差 $\Delta p=(0.003,-0.002)^{\mathsf T}\ \mathrm m$。假设只存在关节零偏，则
\[
 \begin{bmatrix}-0.5&-0.5\\0.5&0\end{bmatrix}
 \begin{bmatrix}\delta q_1\\\delta q_2\end{bmatrix}
 =\begin{bmatrix}0.003\\-0.002\end{bmatrix}.
\]
第二行给 $\delta q_1=-0.004\ \mathrm{rad}$；第一行给 $-0.5(-0.004+\delta q_2)=0.003$，因此 $\delta q_2=-0.002\ \mathrm{rad}$。换成角度约为 $-0.2292^\circ$ 与 $-0.1146^\circ$。这个结果仅在“只有零偏”假设下成立；若工具长度和连杆长度也未知，一个姿态的两个位置分量无法区分全部误差来源。
\end{examplebox}

\subsection{长度误差与零偏可能互相补偿}
在同一标准姿态，模型末端位置为 $(l_1,l_2)$。如果只存在长度偏差，外测位置 $(0.502,0.497)\ \mathrm m$ 就对应 $\Delta l_1=2\ \mathrm{mm}$、$\Delta l_2=-3\ \mathrm{mm}$。但在另一个姿态，长度误差的作用方向会改变；多姿态数据才能检验这套解释。

还需选择“哪些参考坐标固定”。若同时允许基座坐标、末端工具坐标和部分连杆参数任意改变，可能出现多个参数组合产生相同测量结果的规范自由度。应固定参考基准，去掉不可辨识方向，而不是靠更多迭代消除数学歧义。

\begin{figure}[htbp]\centering
\includegraphics[width=0.90\linewidth]{figures/di_calibration_sensitivity.pdf}
\caption{零位偏差对不同两连杆姿态产生不同方向和大小的位置误差。右侧将误差放大显示，用多姿态数据才能区分一致的参数误差与偶发测量噪声。}
\end{figure}

\section{工具中心点标定与负载质心标定不是一回事}
TCP 是 tool center point，即工具中心点或实际工作点。它可以是针尖、夹爪接触点或相机光学中心。质心是质量分布的平均位置，两者一般不同。把 TCP 设准确可以改善位置命令，不能自动改善重力补偿。

一种常见 TCP 标定是让工具尖端始终抵住同一个固定点，改变法兰姿态。设法兰到工具尖端的未知位置为 $c_T$，固定点在基座中的位置为 $p_*$，每次法兰位姿为 $(R_i,t_i)$，则
\begin{equation}
 R_ic_T+t_i=p_*,\qquad
 \begin{bmatrix}R_i&-I_3\end{bmatrix}\begin{bmatrix}c_T\\p_*\end{bmatrix}=-t_i.
\end{equation}
堆叠多次测量即可求解。只绕一个轴转动时，该轴方向的信息可能不足；应使用足够多样的姿态，并检查残差是否来自尖端滑动或机械弹性。

\begin{examplebox}[算例 10E 二维固定点 TCP 标定]
为了便于手算，取二维旋转 $R_1=I$、$R_2=R(90^\circ)$，法兰位置 $t_1=(1,0)$、$t_2=(1.2,-0.2)\ \mathrm m$。两次工作点相同，所以
\[
 (I-R_2)c_T=t_2-t_1,\qquad
 \begin{bmatrix}1&1\\-1&1\end{bmatrix}c_T=\begin{bmatrix}0.2\\-0.2\end{bmatrix}.
\]
两行相加得 $2c_{Ty}=0$，因此 $c_{Ty}=0$；第一行给 $c_{Tx}=0.2$。固定工作点为 $p_*=t_1+c_T=(1.2,0)\ \mathrm m$。这是工具几何位置，不涉及工具质量。
\end{examplebox}

\section{手眼标定从坐标链推导到方程}
\subsection{先固定每个变换的方向}
对于眼在手上配置，相机刚性安装在法兰上。设
\[
 T_i={}^BT_{F_i},\qquad X={}^FT_C,\qquad C_i={} ^{C_i}T_{\mathrm{board}}.
\]
$T_i$ 来自机器人，$C_i$ 来自相机对固定标定板的位姿估计，$X$ 是要求的法兰到相机变换。固定标定板在基座中的位姿 $Z$ 应满足
\begin{equation}
 T_iXC_i=Z.
\end{equation}
对两次姿态 $i,j$ 有 $T_iXC_i=T_jXC_j$。左乘 $T_i^{-1}$、右乘 $C_j^{-1}$，整理得
\begin{equation}
 \underbrace{T_i^{-1}T_j}_{A_{ij}}X
 =X\underbrace{C_iC_j^{-1}}_{B_{ij}}.
\end{equation}
这就是 $AX=XB$。记住坐标链比背诵 $A$、$B$ 的文字定义更可靠，因为不同软件接口可能使用相反方向的输入。

\begin{figure}[htbp]\centering
\includegraphics[width=0.88\linewidth]{figures/di_handeye_chain.pdf}
\caption{眼在手上的坐标链。同一块静止标定板通过两次机器人与相机观测得到相同基座位姿，因此产生相对运动约束。}
\end{figure}

\subsection{先解旋转 再解平移}
把 $A,X,B$ 分别拆成旋转和平移，得到
\begin{align}
 R_AR_X&=R_XR_B,\\
 R_At_X+t_A&=R_Xt_B+t_X.
\end{align}
第一式等价于 $R_A=R_XR_BR_X^{\mathsf T}$。对不处于对数映射歧义的旋转，旋转向量满足 $a=R_Xb$，其中 $a=\operatorname{Log}(R_A)^\vee$、$b=\operatorname{Log}(R_B)^\vee$。对多组旋转向量，可求
\[
 \min_{R_X\in SO(3)}\sum_i\|a_i-R_Xb_i\|^2.
\]
令 $S=\sum_i a_ib_i^{\mathsf T}=U\Sigma V^{\mathsf T}$，则一种闭式旋转对齐解为
\[
 R_X=U\operatorname{diag}(1,1,\det(UV^{\mathsf T}))V^{\mathsf T}.
\]
对角修正确保行列式为正一，避免得到镜像反射。实际存在近 $180^\circ$ 运动或噪声时，还需处理旋转对数的分支和权重。

旋转确定后，平移方程是线性的：
\begin{equation}
 (R_A-I_3)t_X=R_Xt_B-t_A.
\end{equation}
把多组运动堆叠即可用最小二乘求 $t_X$。若全部运动只绕同一轴旋转，某些平移分量不可观察；若只有平移，$R_A-I=0$，$t_X$ 完全从方程中消失。丰富的平移可能帮助约束某些旋转信息，但不能由纯平移恢复完整手眼变换。

\begin{examplebox}[算例 10F 用具体矩阵验证 AX 与 XB]
取 $R_X=R_z(90^\circ)$，$t_X=(0.1,0.2,0.3)^{\mathsf T}\ \mathrm m$。设相机相对运动 $R_B=R_x(90^\circ)$、$t_B=(0.2,0,0)^{\mathsf T}$。由共轭关系得到 $R_A=R_y(90^\circ)$。再由平移关系得到
\[
 t_A=t_X+R_Xt_B-R_At_X=(-0.2,0.2,0.4)^{\mathsf T}.
\]
现在验证左边 $AX$ 的平移：
\[
 R_At_X+t_A=(0.3,0.2,-0.1)+(-0.2,0.2,0.4)=(0.1,0.4,0.3).
\]
右边 $XB$ 的平移为 $R_Xt_B+t_X=(0,0.2,0)+(0.1,0.2,0.3)=(0.1,0.4,0.3)$，完全相同。这个例子用于核对乘法方向；仅一组相对运动不足以唯一完成实际手眼标定。
\end{examplebox}

\section{相机内参与手眼误差如何影响机械臂定位}
针孔相机模型的核心为
\[
 u=f_x\frac XZ+c_x,\qquad v=f_y\frac YZ+c_y.
\]
$(X,Y,Z)$ 是点在相机系中的坐标，$f_x,f_y$ 以像素为单位表示焦距，$(c_x,c_y)$ 是主点，$(u,v)$ 是图像像素坐标。真实镜头还需畸变模型。相机内参不准确，标定板位姿 $C_i$ 就会带偏，随后污染手眼变换 $X$。

例如 $Z=1\ \mathrm m$、$f_x=600$ 像素，忽略深度误差时，横向像素误差 $2$ 像素对应 $\Delta X\approx Z\Delta u/f_x=0.003333\ \mathrm m$，约 $3.33\ \mathrm{mm}$。这只是局部针孔换算，不是完整三维定位误差预算；姿态、深度和标定板检测误差也会贡献。

验证手眼标定应使用没有参与求解的新姿态，检查 $T_iXC_i$ 是否在基座系中保持一致，并报告位置与旋转残差。位置残差与角度残差不能直接相加成一个没有说明尺度的数。小旋转误差还会通过工具长度放大为横向位置误差：$0.5^\circ$ 在 $0.3\ \mathrm m$ 长度上约产生 $0.3\times0.5\pi/180=2.62\ \mathrm{mm}$ 的位移。

\section{怎样把标定残差转成有用的误差预算}
标定后的一组参数不是终点，还需要知道它在哪些位置和姿态下可靠。对位置误差作局部线性化，可写
\[
 \delta p=J_q\delta q+J_\theta\delta\theta+J_T\delta\xi_T+\delta p_{\mathrm{meas}}.
\]
$\delta q$ 是关节测量或跟踪误差，$\delta\theta$ 是几何参数误差，$\delta\xi_T$ 是工具安装的小位姿误差，最后一项是外部测量误差。各个雅可比描述相应误差如何传到同一个位置坐标。所有量必须在相同参考系和相同时刻比较，否则“误差分解”本身又引入误差。

若各类误差可近似看成独立零均值随机量，位置协方差可由各项传播后相加。例如关节部分为 $\Sigma_p=J_q\Sigma_qJ_q^{\mathsf T}$。如果存在系统零偏或相关误差，就不能直接按平方和开根号缩小它们；有确定方向的偏置会持续偏向一侧。

这也解释了为何标定姿态应覆盖实际工作区。一个区域内数据丰富、参数组合可辨，另一个接近奇异或伸展极限的区域仍可能放大同样的参数误差。验证数据应包含实际使用的工具长度、负载方向和相机视角，而不只是与训练姿态相邻的一小团点。

对七轴相机云台，工具标定与手眼标定解决的是刚性安装关系。底座移动、支架弹性和镜头组件松动改变的是运行中的变换，不能靠一次静态标定永久消除。若误差随着姿态、负载或时间系统变化，先判断“固定参数有误”还是“参数确实在变化”，再选择标定、在线观测或机械加固。

\section{计算练习与分步解答}
\subsection*{练习 10.1 从一阶矩恢复质心}
估得质量 $1.5\ \mathrm{kg}$、一阶矩 $(0.12,-0.03,0.075)^{\mathsf T}\ \mathrm{kg\,m}$，求质心。

\textbf{解答。}逐分量除以质量，$c=(0.08,-0.02,0.05)^{\mathsf T}\ \mathrm m$。若质量估计误差明显，三个质心分量都会受同一个分母影响，报告不确定性时不能把它们当作独立误差。

\subsection*{练习 10.2 转换重力方向}
基座重力为 $(0,0,-9.81)$，法兰相对基座绕 $y$ 轴旋转 $90^\circ$。求法兰系重力。

\textbf{解答。}$R_{BF}=\begin{bmatrix}0&0&1\\0&1&0\\-1&0&0\end{bmatrix}$，所以
$R_{BF}^{\mathsf T}(0,0,-9.81)^{\mathsf T}=(9.81,0,0)^{\mathsf T}$。这里应乘转置；如果乘 $R_{BF}$，会得到相反的 $x$ 方向。

\subsection*{练习 10.3 工具误差与质心误差分别影响什么}
一公斤工具的质心水平距离模型少算 $0.03\ \mathrm m$，同时 TCP 的横向位置少算 $0.005\ \mathrm m$。分别计算重力矩偏差和工具几何位置偏差。

\textbf{解答。}质心误差造成重力矩偏差 $1\times9.81\times0.03=0.2943\ \mathrm{N\,m}$。TCP 几何位置偏差为 $5\ \mathrm{mm}$。修正 TCP 不会自动消除 $0.2943\ \mathrm{N\,m}$ 的重力偏差；修正质心也不自动把工作点移动到正确位置。

\subsection*{练习 10.4 纯平移为什么不能标定相机安装平移}
在 $AX=XB$ 中令 $R_A=R_B=I$，写出剩余约束。

\textbf{解答。}旋转方程变成 $R_X=R_X$，平移方程变为 $t_X+t_A=R_Xt_B+t_X$，消去 $t_X$ 后得到 $t_A=R_Xt_B$。因此相机安装平移 $t_X$ 没有出现在数据关系中，无论再采多少这种纯平移，仍不能通过它确定 $t_X$。

\section*{本章与原稿的对应}
本章扩展 [B，第44--45页、第46--50页] 的负载、运动学与手眼标定，并与 [A，第22--30页] 的坐标变换及运动学对应。重点补充了重力力旋量的符号、静态四参数回归、负载参数同坐标合并、TCP 与质心区别，以及 $AX=XB$ 的坐标链和数值验算。

\chapter{从反馈误差到关节伺服}
\label{ch:feedback}

本章用同一个教学关节把控制理论串起来。设关节转动惯量 $J_r=0.1\,\mathrm{kg\,m^2}$，黏性阻尼 $b=0.2\,\mathrm{N\,m\,s/rad}$，关节侧输入力矩为 $u$。这些数只用于计算，不是 NERO 的实测参数。重力已经由名义前馈抵消，剩余扰动力矩记作 $d$，正方向与关节角增大的方向相同：
\begin{equation}
J_r\ddot q+b\dot q=u+d.\label{eq:cf_joint}
\end{equation}
当我们说“控制得好”，至少要分别回答：能否稳定、目标变化时跟得多快、有外力时偏多少，以及为此用了多少力矩。一个响应曲线无法代替这四个问题。

\section{反馈为什么能够纠正误差}

设 $q_d$ 是目标角度，$q$ 是实际角度。本书统一定义 $e=q-q_d$。当关节比目标多转了一个正角度时，$e>0$，恢复力矩应该为负，因此比例反馈是 $u=-K_p e$，其中 $K_p>0$。如果编码器方向或电机力矩方向反了，原本的负反馈会变成正反馈，这时改变增益大小无法修正符号错误。

图中采用常见的“目标减测量”加法节点，因而控制器输入为 $q_d-q=-e$；它与本书的误差定义没有矛盾。误差的定义可以二选一，整条推导必须始终一致。
\begin{figure}[htbp]
\centering\includegraphics[width=.95\linewidth]{figures/cf_feedback.pdf}
\caption{带输入扰动和测量噪声的反馈结构。扰动 $d$ 作用在对象入口，噪声 $n$ 只进入测量通道，二者的闭环传递关系不同。}
\end{figure}

“闭环”是输出测量回到输入端形成回路；“开环”是指没有这条反馈通道，或为分析而把回路在某处断开。开环控制不一定差，例如精确的重力前馈很有用；但前馈无法凭空知道未建模负载发生了什么。反馈的价值是根据已经出现的偏差修正，而它的代价是测量噪声、延迟和未建模动态也进入闭环。

\subsection{从微分方程得到传递函数}

传递函数描述线性时不变模型在零初始状态下的输入输出关系。对式\eqref{eq:cf_joint}取拉普拉斯变换，暂令 $d=0$：
\begin{align}
J_r\big[s^2Q(s)-s q(0)-\dot q(0)\big]
+b\big[sQ(s)-q(0)\big]&=U(s),\\
q(0)=\dot q(0)=0\quad\Longrightarrow\quad
(J_rs^2+bs)Q(s)&=U(s),\\
G_q(s)=\frac{Q(s)}{U(s)}&=\frac{1}{s(J_rs+b)}.
\end{align}
这里 $s$ 的量纲是每秒，$Q(s)$、$U(s)$ 是变换后的信号，不是新的物理状态。角速度 $v=\dot q$，所以速度对象为
\begin{equation}
G_v(s)=\frac{V(s)}{U(s)}=\frac{1}{J_rs+b}
=\frac{5}{0.5s+1}.
\end{equation}
分母为零的位置称为极点。速度模型只有 $s=-2$ 这个极点，对应时间常数 $0.5\,\mathrm{s}$。位置模型还多一个零极点，因为速度必须积分才能成为位置。这个零极点意味着恒定力矩可以让位置一直增加，并不意味着电机速度会无限增加。

\begin{examplebox}[算例 11A 恒定力矩产生怎样的运动]
关节从静止开始，输入 $u=0.4\,\mathrm{N\,m}$。由 $0.1\dot v+0.2v=0.4$ 得
\[
v(t)=2(1-e^{-2t})\,\mathrm{rad/s}.
\]
在 $t=0.5\,\mathrm{s}$，$v=2(1-e^{-1})=1.2642\,\mathrm{rad/s}$。再积分：
\[
q(t)=2t-(1-e^{-2t}),\qquad q(0.5)=0.3679\,\mathrm{rad}.
\]
速度逐渐接近 $2\,\mathrm{rad/s}$，位置仍持续增加。看懂这一步，就能理解为什么位置对象和速度对象不能套用同一组增益。
\end{examplebox}

\subsection{闭环代数以及扰动的入口}

设控制器传递函数为 $C(s)$，测量噪声为 $N(s)$。对象满足 $Y=G(U+D)$，控制器满足 $U=C(R-Y-N)$。逐项代入：
\begin{align}
Y&=GC(R-Y-N)+GD,\\
(1+GC)Y&=GCR+GD-GCN,\\
Y&=T R+G S D-TN,\\
S(s)&=\frac{1}{1+GC},\qquad
T(s)=\frac{GC}{1+GC},\qquad S+T=1.
\end{align}
$S$ 称为灵敏度函数，$T$ 称为互补灵敏度函数。参考信号通过 $T$ 到输出；输入端力矩扰动通过 $GS$ 到输出；测量噪声通过 $-T$ 到输出。因此“$S$ 小就能消除所有扰动”是不完整的说法，还要知道扰动从哪里进入。

关系 $S+T=1$ 是复数函数的相加，不是说两个幅值之和处处等于 1；但由三角不等式可知，同一频率下不可能让 $|S|$ 和 $|T|$ 都任意小。设计通常希望低频 $T\approx1$、$S\approx0$，高频 $T\approx0$，由此在跟踪、抗扰和噪声之间分配频率范围。

\section{比例和微分怎样决定响应}

对静止目标采用
\begin{equation}
u=-K_p(q-q_d)-K_d\dot q.
\end{equation}
$K_p$ 的单位是 $\mathrm{N\,m/rad}$；$K_d$ 的单位是 $\mathrm{N\,m\,s/rad}$。代入对象并令 $d=0$：
\begin{align}
J_r\ddot q+(b+K_d)\dot q+K_pq&=K_pq_d,\\
\frac{Q(s)}{Q_d(s)}&=
\frac{K_p/J_r}{s^2+[(b+K_d)/J_r]s+K_p/J_r}.
\end{align}
把分母与标准二阶式 $s^2+2\zeta\omega_ns+\omega_n^2$ 对齐，得到
\begin{equation}
\omega_n=\sqrt{\frac{K_p}{J_r}},\qquad
\zeta=\frac{b+K_d}{2\sqrt{J_rK_p}}.
\label{eq:cf_pd_design}
\end{equation}
$\omega_n$ 是无阻尼自然角频率，单位 $\mathrm{rad/s}$；除以 $2\pi$ 才得到 Hz。$\zeta$ 是没有单位的阻尼比，描述能量损耗相对于振动尺度的大小。不是所有二阶系统都振荡：$0<\zeta<1$ 为欠阻尼，$\zeta=1$ 为临界阻尼，$\zeta>1$ 为过阻尼。

\begin{examplebox}[算例 11B 从指标反推 PD 增益]
希望 $\omega_n=10\,\mathrm{rad/s}$、$\zeta=0.7$。由式\eqref{eq:cf_pd_design}反解：
\[
K_p=J_r\omega_n^2=0.1\times100=10,
\]
\[
K_d=2\zeta\omega_nJ_r-b
=2\times0.7\times10\times0.1-0.2=1.2.
\]
闭环极点为 $-7\pm j\sqrt{51}=-7\pm7.1414j$。给一个 $0.1\,\mathrm{rad}$ 的位置阶跃，初始比例力矩为 $1\,\mathrm{N\,m}$。标准二阶公式给出
\[
M_p=e^{-\pi\zeta/\sqrt{1-\zeta^2}}=0.04599,
\quad t_p=\frac{\pi}{10\sqrt{0.51}}=0.4399\,\mathrm{s},
\]
\[
t_s\approx\frac{4}{\zeta\omega_n}=0.5714\,\mathrm{s}.
\]
峰值角度约为 $0.10460\,\mathrm{rad}$。调节时间这里是进入并持续留在最终值 $\pm2\%$ 带内的估算，不是精确的普适等式。额外滤波器、延迟、零点或饱和都会改变结果。
\end{examplebox}

\begin{figure}[htbp]
\centering\includegraphics[width=.95\linewidth]{figures/cf_pd_step.pdf}
\caption{左图改变阻尼比而保持自然频率不变；右图在同一阻尼下改变刚度。振动衰减和稳态抗扰是两个不同的问题。}
\end{figure}

\subsection{为什么越软越容易留偏差}

在恒定扰动 $d_0$ 下，静止时 $\dot q=\ddot q=0$，对象方程变成
\begin{equation}
0=-K_p(q-q_d)+d_0,\qquad e_\infty=\frac{d_0}{K_p}.
\end{equation}
这不是控制器没有执行，而是虚拟弹簧的恢复力矩与扰动力矩达到了平衡。

例如扰动力矩 $d_0=0.3\,\mathrm{N\,m}$，刚度 $K_p=10$，误差为 $0.03\,\mathrm{rad}$，约 $1.72^\circ$。把 $K_p$ 降为 2，误差变为 $0.15\,\mathrm{rad}$，约 $8.59^\circ$。提高刚度能压小偏差，却也增大人推动它需要的力矩。后面的辨识、重力补偿和摩擦补偿，目的正是减少 $d_0$，使较低刚度也能保持较小误差。

\section{积分消除常值误差及其代价}

定义积分状态 $z(t)=\int_0^t e(\xi)\,d\xi$，PID 控制写成
\begin{equation}
u=-K_pe-K_d\dot e-K_iz.
\end{equation}
$K_i$ 的单位为 $\mathrm{N\,m/(rad\,s)}$。积分记住了误差持续多久，即使误差很小，持续积累也会增加纠正力矩。对静止目标，闭环特征多项式为
\begin{equation}
J_rs^3+(b+K_d)s^2+K_ps+K_i=0.
\end{equation}
三阶 Routh 判据给出严格稳定条件：各系数为正，且
\begin{equation}
(b+K_d)K_p>J_rK_i.
\end{equation}
因此加入积分不只是多一个“消除误差”的按钮，它也多了一个状态和一个稳定性约束。

\begin{examplebox}[算例 11C 积分增益不是越大越好]
延续 $J_r=0.1$、$b+K_d=1.4$、$K_p=10$，得到
\[
0<K_i<\frac{1.4\times10}{0.1}=140.
\]
$K_i=20$ 满足该理想模型的稳定条件；$K_i=160$ 不满足。上界 140 不是实机可用增益，真实对象还有滤波、采样和柔性，通常需要更大的稳定余量。恒定扰动为 $0.3\,\mathrm{N\,m}$ 时，稳定平衡要求 $e=0$ 且 $-K_i z=-0.3$，因此 $z=0.015\,\mathrm{rad\,s}$。积分状态保存了抵消扰动所需的信息。
\end{examplebox}

\subsection{饱和后积分为什么会越积越多}

设计算出的力矩为 $u_{\rm raw}$，执行器只能输出 $u_{\rm sat}=\operatorname{clip}(u_{\rm raw},-u_{\max},u_{\max})$。在大误差下，控制器可能要求 10，而执行器实际只能给 1。若积分器仍按“误差还没消失”不断累加，它相当于在积累执行器无力兑现的命令。待关节到达目标附近，积分存量仍很大，关节会越过目标。

把积分贡献直接记为 $i=-K_iz$，一种反算抗饱和方法为
\begin{equation}
u_{\rm raw}=-K_pe-K_d\dot q+i+u_{\rm ff},\qquad
\dot i=-K_ie+K_{\rm aw}(u_{\rm sat}-u_{\rm raw}).
\label{eq:cf_aw}
\end{equation}
$K_{\rm aw}$ 的单位是 $\mathrm{s^{-1}}$。输出上饱和时，$u_{\rm sat}-u_{\rm raw}<0$，额外项把过大的正积分贡献拉下来。条件积分则只在积分会把饱和推得更深时停止积分；能够解除饱和的反向积分应继续进行。

\begin{examplebox}[算例 11D 手算一次抗饱和更新]
取 $T_s=0.005\,\mathrm{s}$、$e=-0.02\,\mathrm{rad}$、$K_i=10$、$K_{\rm aw}=20\,\mathrm{s^{-1}}$，当前 $u_{\rm raw}=1.6$、$u_{\rm sat}=1.0\,\mathrm{N\,m}$。本拍积分变化为
\[
\Delta i=T_s[-K_ie+K_{\rm aw}(u_{\rm sat}-u_{\rm raw})]
=0.005[0.2-12]=-0.059\,\mathrm{N\,m}.
\]
如果没有抗饱和，这一拍还会增加 $0.001\,\mathrm{N\,m}$。两者方向恰好相反，因为当前最需要的是卸掉积累，而不是继续扩大命令。
\end{examplebox}

\begin{figure}[htbp]
\centering\includegraphics[width=.95\linewidth]{figures/cf_antiwindup.pdf}
\caption{教学关节的饱和仿真。目标为 1 rad，$K_p=10$、$K_d=1.2$、$K_i=30$，输出限制为 $\pm1\,\mathrm{N\,m}$，反算增益为 $15\,\mathrm{s^{-1}}$。比较的是同一对象和同一目标。}
\end{figure}

对测量值微分可以避免目标阶跃产生微分冲击，但这时 D 项是 $-K_d\dot q$，与轨迹误差微分 $-K_d(\dot q-\dot q_d)$ 相差 $K_d\dot q_d$。移动目标若需要速度跟踪，应明确加入这项前馈。纯微分的高频增益无限大，实际通常用 $K_ds/(1+T_fs)$，并把 $T_f$ 引入的相位变化计入闭环设计。

\section{极点和频率是同一个系统的两种观察方法}

\subsection{用 Routh 表算出一个稳定区间}

考虑开环 $L(s)=K/[s(s+1)(s+2)]$。单位负反馈的闭环分母为
\[
s^3+3s^2+2s+K.
\]
Routh 表为
\begin{center}
\begin{tabular}{ccc}\toprule
幂次 & 第一列 & 第二列\\\midrule
$s^3$ & $1$ & $2$\\
$s^2$ & $3$ & $K$\\
$s^1$ & $(6-K)/3$ & $0$\\
$s^0$ & $K$ & $0$\\\bottomrule
\end{tabular}
\end{center}
第一列全正要求 $0<K<6$。$K=6$ 时，多项式可分解成 $(s+3)(s^2+2)$，其中 $\pm j\sqrt2$ 没有负实部，不能渐近衰减。根轨迹把每个 $K$ 对应的根画在复平面上，展示了这一结论。这个例子中的增益过大确实使一对极点进入右半平面，但“增加增益一定使所有系统不稳定”不是定理。

\begin{figure}[htbp]
\centering\includegraphics[width=.98\linewidth]{figures/cf_root_bode.pdf}
\caption{同一闭环问题的根轨迹与频率响应。根轨迹上虚轴交点对应 $K=6$，频域中也能找到相同的增益边界。}
\end{figure}

\subsection{Bode 图怎样读}

用正弦输入激励稳定的线性系统，稳态输出仍是同一频率，只改变幅值和相位。令 $s=j\omega$，$|G(j\omega)|$ 给幅值比例，$\arg G(j\omega)$ 给相位差。幅值常画成 $20\log_{10}|G|$，单位为 dB：幅值减半约为 $-6.02\,\mathrm{dB}$，幅值变成 $1/\sqrt2$ 约为 $-3.01\,\mathrm{dB}$。

开环幅值等于 1 的频率叫增益穿越频率 $\omega_c$。在常见的单一穿越、开环没有右半平面极点等情形下，相位裕度为
\[
\mathrm{PM}=180^\circ+\arg L(j\omega_c).
\]
相位到 $-180^\circ$ 时，增益裕度为 $\mathrm{GM}=1/|L(j\omega_{180})|$。多个穿越、开环不稳定或复杂多变量系统，应使用完整 Nyquist 或相应多变量分析，不能仅看一个正的 PM 数字。

闭环带宽常用 $|T(j\omega)|$ 相对于低频增益下降 3 dB 的频率来定义。它与开环穿越频率通常同量级，但并不相等；它更不等于控制程序每秒运行多少次。

\begin{examplebox}[算例 11E 一个可手算的速度 PI 设计]
速度对象为 $G_v=1/(0.1s+0.2)$，取 $C_v=2+8/s=2(1+4/s)$。在 $\omega=20\,\mathrm{rad/s}$：
\[
|G_v|=\frac{1}{\sqrt{0.2^2+2^2}}=0.4975,
\qquad |C_v|=\sqrt{2^2+0.4^2}=2.0396.
\]
乘积约为 1.015，因此穿越频率就在 20 附近。对象相位为 $-\arctan(10)=-84.29^\circ$，PI 相位为 $-\arctan(0.2)=-11.31^\circ$。若总等效延迟为 $7.5\,\mathrm{ms}$，额外相位约为 $-20\times0.0075\times180/\pi=-8.59^\circ$。

由此估计 PM 约为 $75.8^\circ$。数值求根得到实际穿越频率 $20.2867\,\mathrm{rad/s}=3.2287\,\mathrm{Hz}$，计入该延迟后的 PM 为 $75.76^\circ$。这仍只是当前刚体一阶模型的设计结果；若速度测量有滤波或存在柔性模态，须把它们放入同一个开环模型。
\end{examplebox}

\section{电流速度位置三环怎样分工}

电流环的输出是电压或 PWM，占空比改变电流；速度环的输出是电流或力矩目标；位置环的输出常是速度目标。只有先问清每个接口的输入输出单位，才有资格比较两个参数都叫 $K_p$ 的控制器。

若电流内环的简化对象为 $G_i=1/(Ls+R_e)$，PI 控制器为 $(K_{pi}s+K_{ii})/s$，选择 $K_{ii}/K_{pi}=R_e/L$，名义开环化成 $K_{pi}/(Ls)$。若希望其名义闭环极点为 $-\omega_i$，则
\[
K_{pi}=L\omega_i,\qquad K_{ii}=R_e\omega_i.
\]
这个推导假设电感电阻模型有效，反电动势和耦合已妥善处理，且忽略数字延迟。它说明参数的由来，不能把零极点对消当作消除实物动态的魔法。尤其不能靠精确对消隐藏不稳定极点。

当速度内环足够快时，位置外环可近似为 $\dot q=v_{\rm cmd}$。取
\begin{equation}
v_{\rm cmd}=\dot q_d-K_q(q-q_d)
\end{equation}
便有 $\dot e=-K_qe$。这里 $K_q$ 的单位是 $\mathrm{s^{-1}}$，与直接输出力矩的比例增益不同。若省略速度前馈，只用 $v_{\rm cmd}=-K_qe$，对匀速目标 $\dot q_d=v_d$，稳态误差为 $e_\infty=-v_d/K_q$。

\begin{examplebox}[算例 11F 速度前馈消除跟随落后]
外环 $K_q=5\,\mathrm{s^{-1}}$，期望匀速 $v_d=0.2\,\mathrm{rad/s}$。没有前馈时，实际关节必须落后 $0.04\,\mathrm{rad}$，比例项才能产生 $0.2\,\mathrm{rad/s}$ 的指令。加入 $\dot q_d$ 后，在理想速度内环下即使 $e=0$ 也能继续运动。

前馈改善了已知参考的跟踪，并没有增加反馈开环的积分器个数。因此不能把“加速度前馈”直接称作“系统变为二型”；系统型别描述的是反馈开环在原点的积分结构。
\end{examplebox}

\section{计算练习与逐步解答}

\textbf{练习 1}\quad $J_r=0.2$、$b=0.3$，希望 $\omega_n=8\,\mathrm{rad/s}$、$\zeta=1$。求直接力矩 PD 的 $K_p,K_d$，并求 $0.16\,\mathrm{N\,m}$ 恒定扰动造成的静态误差。

\textbf{解}\quad $K_p=0.2\times8^2=12.8\,\mathrm{N\,m/rad}$；$K_d=2\times1\times8\times0.2-0.3=2.9\,\mathrm{N\,m\,s/rad}$。静止平衡 $K_pe=d$，故 $e=0.16/12.8=0.0125\,\mathrm{rad}=0.7162^\circ$。阻尼改变瞬态，但不出现在静态误差公式中。

\textbf{练习 2}\quad 对教学关节取 $K_p=5$、$K_d=0.8$。在理想连续 PID 模型下，积分增益的稳定范围是什么？

\textbf{解}\quad 特征式为 $0.1s^3+(0.2+0.8)s^2+5s+K_i$。系数为正要求 $K_i>0$，Routh 条件要求 $1\times5>0.1K_i$，所以 $0<K_i<50$。在 $K_i=50$ 的边界处不能当作严格稳定。

\textbf{练习 3}\quad 速度前馈比例只有 80\%，控制律 $v_{\rm cmd}=0.8v_d-K_qe$。取 $K_q=4\,\mathrm{s^{-1}}$、$v_d=0.1\,\mathrm{rad/s}$，求稳态误差。

\textbf{解}\quad $\dot e=v_{\rm cmd}-v_d=-0.2v_d-K_qe$。稳态令 $\dot e=0$，得到 $e=-0.2\times0.1/4=-0.005\,\mathrm{rad}$。关节仍略落后，但比完全没有前馈时的 $-0.025\,\mathrm{rad}$ 小五倍。

\section*{本章与原稿的对应}
本章整合 [A，第5--14、34--38页] 与 [B，第59--63页]。补足了扰动入口的代数推导、增益单位、抗饱和逐拍计算，并修正“前馈改变系统型别”和“增加增益总会把极点推向不稳定”的过度概括。

\chapter{状态空间最优反馈与稳定性证明}
\label{ch:state_space}

传递函数回答“输入怎样影响输出”；状态空间进一步记录系统当前的内部处境。机械臂同时有多个关节位置和速度，如果还考虑柔性、电流、积分器或外部负载，状态会更多。本章仍从一个关节起步，把矩阵里的每一个数字算出来，再解释为什么最优控制、观测器和预测控制都建立在这套表示上。

\section{状态是预测未来所需的记忆}

对 $J_r\ddot q+b\dot q=u$，只知道当前位置无法预测下一刻的位置：关节可能正在向前，也可能正在向后。取状态
\[
x=\begin{bmatrix}q\\v\end{bmatrix},\qquad v=\dot q,
\]
便可写为
\begin{equation}
\dot x=Ax+Bu,\qquad y=Cx+Du,
\end{equation}
其中
\begin{equation}
A=\begin{bmatrix}0&1\\0&-b/J_r\end{bmatrix},\quad
B=\begin{bmatrix}0\\1/J_r\end{bmatrix},\quad
C=\begin{bmatrix}1&0\end{bmatrix},\quad D=0.
\end{equation}
状态方程第二行仍然是原来的牛顿方程，第一行只是 $\dot q=v$。矩阵并未增加新的物理假设，只是把两个一阶方程放在一起。

状态的分量可以有不同单位，因此 $A$ 的各元素未必有相同单位。例如 $A_{12}$ 把速度变成位置导数，数值为 1；$A_{22}=-b/J_r$ 的单位是每秒。直接把角度和角速度相加没有意义，而经过带正确单位的矩阵加权后可以相加。

\subsection{从状态空间回到传递函数}

零初始状态时，$sX=AX+BU$，故 $(sI-A)X=BU$，于是
\[
G(s)=C(sI-A)^{-1}B+D.
\]
对教学关节，$A=\left[\begin{smallmatrix}0&1\\0&-2\end{smallmatrix}\right]$、$B=[0,10]^\top$。直接解方程：第一行 $sX_1=X_2$，第二行 $(s+2)X_2=10U$，所以
\[
\frac{Y}{U}=\frac{10}{s(s+2)}=\frac{1}{s(0.1s+0.2)}.
\]
这和上一章完全相同。一般状态模型的内部极点是 $A$ 的特征值；只有在没有不可控、不可观消去等隐藏模态时，才会全部出现在输入输出传递函数中。因此“传递函数看起来稳定”不自动保证内部所有状态稳定。

\subsection{非线性模型怎样在工作点附近变成矩阵}

真实关节可能满足 $J_r\ddot q+b\dot q+g(q)=u$。选平衡点 $q_0$，静止时的输入为 $u_0=g(q_0)$。定义小偏差 $\delta q=q-q_0$、$\delta u=u-u_0$，使用
\[
g(q_0+\delta q)\approx g(q_0)+g'(q_0)\delta q
\]
得到
\[
\begin{bmatrix}\delta\dot q\\\delta\dot v\end{bmatrix}
=\begin{bmatrix}0&1\\-g'(q_0)/J_r&-b/J_r\end{bmatrix}
\begin{bmatrix}\delta q\\\delta v\end{bmatrix}
+\begin{bmatrix}0\\1/J_r\end{bmatrix}\delta u.
\]
$g'(q_0)$ 取决于姿态。相同机械结构在不同平衡点可能具有不同的局部稳定性。线性化结论只覆盖展开点附近，不能把它无条件推广到整条大范围运动。

\section{先问能否控制以及能否看见}

\subsection{可控性}

可控性指是否能通过输入在有限时间内改变任意状态方向。对 $n$ 维连续线性系统，可控性矩阵为
\[
\mathcal C=[B,AB,A^2B,\ldots,A^{n-1}B].
\]
若 $\operatorname{rank}\mathcal C=n$，所有状态方向都可控。秩不足意味着至少存在一个输入无法改变的组合。

\begin{examplebox}[算例 12A 教学关节的可控性和可观性]
\[
AB=\begin{bmatrix}10\\-20\end{bmatrix},\qquad
\mathcal C=\begin{bmatrix}0&10\\10&-20\end{bmatrix}.
\]
$\det\mathcal C=-100\ne0$，因此秩为 2，位置和速度都可控制。注意输入只直接出现在速度导数里；经过时间演化，改变速度也会改变位置。

如果只测位置，$C=[1,0]$，可观性矩阵为
\[
\mathcal O=\begin{bmatrix}C\\CA\end{bmatrix}
=\begin{bmatrix}1&0\\0&1\end{bmatrix},
\]
秩为 2，因此位置历史加上输入历史足以重建速度。如果只测速度，$C_v=[0,1]$，则
\[
\mathcal O_v=\begin{bmatrix}0&1\\0&-2\end{bmatrix},
\]
秩为 1。两个位置相差一个常数、速度相同的运动，会给出完全相同的速度测量，因此绝对位置不可观。
\end{examplebox}

可控和可观是模型结构的性质；噪声很大时，数学上可观的量仍可能难以精确估计。这称为“可观性很弱”的工程情形，不能只看秩是否为满秩。若某些不可控模态本来已经稳定，仍可能设计稳定反馈，条件称为可稳定性；观测器对应的较弱条件称为可检测性。

\section{用极点配置设计反馈}

对零目标取 $u=-Kx$，$K=[k_1,k_2]$，闭环为
\[
\dot x=(A-BK)x.
\]
特征多项式展开为
\begin{align}
sI-A+BK&=
\begin{bmatrix}s&-1\\k_1/J_r&s+(b+k_2)/J_r\end{bmatrix},\\
\det(sI-A+BK)&=s^2+\frac{b+k_2}{J_r}s+\frac{k_1}{J_r}.
\end{align}
对于这个二阶关节，状态反馈就是 PD；对耦合的高维系统，状态反馈矩阵可以让某个关节的状态影响另一个关节的输入，因此比逐轴独立 PD 更一般。

\begin{examplebox}[算例 12B 把闭环极点放在两个指定位置]
期望极点为 $-8$、$-12$，对应多项式
\[
(s+8)(s+12)=s^2+20s+96.
\]
系数对齐得到 $k_1=0.1\times96=9.6$，$k_2=0.1\times20-0.2=1.8$。对常值目标 $q_d$，可以直接写
\[
u=-9.6(q-q_d)-1.8v.
\]
例如 $q-q_d=0.02\,\mathrm{rad}$、$v=-0.03\,\mathrm{rad/s}$，比例项为 $-0.192$，阻尼项为 $+0.054$，合计 $u=-0.138\,\mathrm{N\,m}$。速度为负表示关节正在回程，阻尼项为正，正在减缓回程速度。
\end{examplebox}

\subsection{积分增广就是给状态多加一份记忆}

令 $z=\int e\,dt$，状态取 $x_a=[e,v,z]^\top$，常值目标时
\[
\dot x_a=
\begin{bmatrix}0&1&0\\0&-2&0\\1&0&0\end{bmatrix}x_a
+\begin{bmatrix}0\\10\\0\end{bmatrix}u.
\]
取 $u=-k_pe-k_dv-k_iz$，期望极点为 $-3,-5,-8$。将
\[
(s+3)(s+5)(s+8)=s^3+16s^2+79s+120
\]
与 $s^3+(b+k_d)s^2/J_r+k_ps/J_r+k_i/J_r$ 比较，得到
\[
k_d=1.4,\qquad k_p=7.9,\qquad k_i=12.
\]
这比“随便加一点积分”清楚得多。不过饱和时系统不再满足这条线性特征方程，积分状态仍须进行抗饱和处理。

\section{LQR 为什么会得到一组反馈增益}

LQR 是线性二次型调节器，即 Linear Quadratic Regulator。它针对线性模型，使用平方形式的代价，求把状态带回零点的反馈。这里把代价记为 $\mathcal J$，避免与雅可比矩阵 $J$ 混淆：
\begin{equation}
\mathcal J=\int_0^\infty (x^\top Qx+u^\top R_uu)\,dt,
\qquad Q\succeq0,\quad R_u\succ0.
\end{equation}
$Q$ 决定哪些状态偏差代价较高；$R_u$ 决定控制输入代价较高。$\succeq0$ 表示半正定，即任意二次型非负；$\succ0$ 表示正定。这里的 $R_u$ 是控制代价矩阵，不是旋转矩阵，也不是卡尔曼滤波的测量噪声协方差。

\subsection{从最优剩余代价推到 Riccati 方程}

设从当前状态出发的最优剩余代价为 $V(x)=x^\top Px$，其中 $P=P^\top$。在极短时间内，要付出当前代价，并使未来剩余代价发生变化。连续最优性方程可写成
\[
0=\min_u\left[x^\top Qx+u^\top R_uu+
\frac{\partial V}{\partial x}(Ax+Bu)\right].
\]
由 $\partial V/\partial x=2x^\top P$，对括号中的式子关于 $u$ 求导：
\[
2R_uu+2B^\top Px=0.
\]
于是
\begin{equation}
u^*=-R_u^{-1}B^\top Px=-Kx.
\end{equation}
再代回最优性方程并把 $x$ 的二次项合并：
\[
0=x^\top\left(A^\top P+PA-PBR_u^{-1}B^\top P+Q\right)x.
\]
要让它对所有 $x$ 成立，括号内矩阵必须为零，得到连续代数 Riccati 方程：
\begin{equation}
A^\top P+PA-PBR_u^{-1}B^\top P+Q=0.
\label{eq:cf_care}
\end{equation}
这并不是“试出了一个经验公式”，而是由代价最小条件导出。存在稳定化解需要可稳定性和相应可检测性等条件；并非任意 $A,B,Q,R_u$ 都能得到所需解。

\begin{examplebox}[算例 12C 用三个标量方程手算 LQR]
对教学关节，取 $Q=\operatorname{diag}(100,1)$、$R_u=1$，并设
\[
P=\begin{bmatrix}p_{11}&p_{12}\\p_{12}&p_{22}\end{bmatrix}.
\]
将 $A,B$ 代入式\eqref{eq:cf_care}，三个独立元素分别给出
\begin{align*}
-100p_{12}^2+100&=0,\\
p_{11}-2p_{12}-100p_{12}p_{22}&=0,\\
2p_{12}-4p_{22}-100p_{22}^2+1&=0.
\end{align*}
取稳定化的正定解 $p_{12}=1$，第三式变为 $100p_{22}^2+4p_{22}-3=0$，正根是
\[
p_{22}=\frac{-4+\sqrt{1216}}{200}=0.154356.
\]
第二式给 $p_{11}=2+100p_{22}=17.4356$。因此
\[
P=\begin{bmatrix}17.4356&1\\1&0.154356\end{bmatrix},
\qquad K=B^\top P=[10,1.54356].
\]
闭环极点是 $-8.7178\pm4.8990j$。若初始误差为 $0.1\,\mathrm{rad}$ 且速度为零，初始控制量为 $-1\,\mathrm{N\,m}$，最优剩余代价为 $0.1^2\times17.4356=0.174356$。
\end{examplebox}

\begin{figure}[htbp]
\centering\includegraphics[width=.95\linewidth]{figures/cf_lqr.pdf}
\caption{同一教学关节的 LQR 比较。只改变输入代价 $R_u$，初始误差均为 0.1 rad。越惩罚力矩，通常越愿意以较慢的误差衰减换取较小的瞬时输入。}
\end{figure}

\subsection{权重需要量纲与工程限制的解释}

一个起点是先归一化：$\tilde q=e/e_{\rm scale}$、$\tilde v=v/v_{\rm scale}$、$\tilde u=u/u_{\rm scale}$，然后惩罚 $\tilde q^2+\tilde v^2+\tilde u^2$。这等价于
\[
Q=\operatorname{diag}(1/e_{\rm scale}^2,1/v_{\rm scale}^2),
\qquad R_u=1/u_{\rm scale}^2.
\]
这些尺度不是自动满足的硬限制。即使把 $u_{\rm scale}=2\,\mathrm{N\,m}$，LQR 仍可能在误差大时要求 4 或 20。需要严格满足力矩上限时，应明确考虑饱和或带约束优化，而不能把权重当成保护边界。

LQR 的“最优”还只针对指定模型、指定代价、指定时间范围。改了负载、加入延迟、改变噪声模型或实际输出限幅，原来的最优性和鲁棒性结论不再直接适用。用观测状态替代真实状态组成 LQG 后，也不能照搬理想全状态 LQR 的特定裕度结论。

\section{用能量解释稳定而不只是看曲线}

对定点 PD 关节，误差方程为 $J_r\ddot e+(b+K_d)\dot e+K_pe=0$。取
\begin{equation}
V=\frac12J_r\dot e^2+\frac12K_pe^2.
\end{equation}
第一项是动能，第二项是虚拟弹簧势能。沿实际运动求导：
\begin{align}
\dot V&=J_r\dot e\ddot e+K_pe\dot e\\
&=\dot e[-(b+K_d)\dot e-K_pe]+K_pe\dot e\\
&=-(b+K_d)\dot e^2\le0.
\end{align}
这说明能量不会增加。仅有 $\dot V\le0$ 还不能直接宣布所有状态趋于零；还需看能量不再下降时能停在哪里。此时 $\dot e=0$，要一直保持静止还必须满足 $K_pe=0$，所以唯一不变平衡是 $e=0$。这就是 LaSalle 不变集分析在这个例子中的具体含义。

\begin{figure}[htbp]
\centering\includegraphics[width=.95\linewidth]{figures/cf_energy.pdf}
\caption{状态轨迹向原点收敛，同时能量下降。误差在某个瞬间增大，不一定代表系统不稳定；还要看速度和总能量。}
\end{figure}

正定的 $V$ 像一个最低点位于目标的碗；负的 $\dot V$ 像摩擦持续消耗能量。这种证明能揭示控制律为什么有效，而不是仅证明某一次仿真曲线没有发散。全局结论还需要状态域、函数增长性质和模型条件；对姿态空间或带关节限位的系统，不能省略这些条件。

\section{鲁棒自适应和重复学习的最小例子}

\subsection{滑模控制如何对付有界扰动}

考虑 $J_r\ddot q=u+d$，$|d|\le\bar d$，定义 $e=q-q_d$、$s=\dot e+\lambda e$。如果 $s=0$，则 $\dot e=-\lambda e$，误差沿指数曲线衰减。选择
\[
u=J_r(\ddot q_d-\lambda\dot e)-k\operatorname{sgn}(s).
\]
代入可得 $J_r\dot s=d-k\operatorname{sgn}(s)$，于是
\[
s\dot s\le-\frac{k-\bar d}{J_r}|s|.
\]
当 $k>\bar d$ 时，$|s|$ 会下降；理想不连续模型可给出有限到达时间上界
\[
t_{\rm reach}\le\frac{J_r|s(0)|}{k-\bar d}.
\]
例如 $J_r=0.1$、$|s(0)|=0.6$、$k=0.5$、$\bar d=0.2$，上界为 $0.2\,\mathrm{s}$。真实执行器无法无限快切换，采样和柔性会导致抖振，因此常把 sign 换成 $\operatorname{sat}(s/\phi)$。这会引入边界层，不能继续声称误差精确为零。

若 $d=0.2$ 为常数、$\phi=0.02$、$\lambda=4$，在边界层内稳态满足 $d-k s/\phi=0$，得到 $s=d\phi/k=0.008$，进一步 $e=s/\lambda=0.002\,\mathrm{rad}$。鲁棒性与残差精度之间的交换可以直接算出来。

\subsection{自适应控制为什么还需要激励}

设未知重力系数为 $\theta$：
\[
J_r\ddot q+b\dot q+\theta\cos q=u.
\]
定义 $v_r=\dot q_d-\lambda e$、$a_r=\ddot q_d-\lambda\dot e$、$s=\dot q-v_r$。采用
\[
u=J_ra_r+b\dot q+\hat\theta\cos q-K_s s,
\qquad \dot{\hat\theta}=-\Gamma\cos q\,s.
\]
令 $\tilde\theta=\hat\theta-\theta$，则
\[
J_r\dot s=-K_s s+\tilde\theta\cos q.
\]
候选能量为 $V=\tfrac12J_rs^2+\tfrac1{2\Gamma}\tilde\theta^2$，其导数中的交叉项为
\[
s\tilde\theta\cos q+
\frac{\tilde\theta}{\Gamma}\dot{\hat\theta}=0,
\]
最终剩下 $\dot V=-K_ss^2$。更新律的符号来自这个抵消过程，不能凭直觉任意写正号。

例如 $\cos q=0.8$、$s=0.1$、$\Gamma=0.5$，此时 $\dot{\hat\theta}=-0.04$。模型预测的重力项应逐渐减小。跟踪误差趋近零仍不意味着参数必然是真值：若长期停在 $\cos q=0$ 的姿态，数据根本看不到 $\theta$。参数收敛需要足够的持续激励，实际还需投影、死区等措施防止噪声使参数漂移。

\subsection{重复动作怎样利用上一次误差}

迭代学习控制 ILC 的“迭代”是完整动作执行次数，不是控制采样次数。用最简单的静态关系 $y_j=Gu_j$、$e_j=y_j-r$，令
\[
u_{j+1}=u_j-Le_j.
\]
代入得到 $e_{j+1}=(1-GL)e_j$。只要 $|1-GL|<1$，误差会随动作次数缩小。取 $G=0.8$、$L=1$、$r=1$、$u_0=0$，有
\[
(u_1,y_1,e_1)=(1,0.8,-0.2),
\]
\[
(u_2,y_2,e_2)=(1.2,0.96,-0.04),\qquad
(u_3,y_3,e_3)=(1.24,0.992,-0.008).
\]
真正的动态对象要处理相位、对齐、滤波和边界状态。重复扰动可以学到，随机接触或任务变化不能机械照搬上一次修正。

\section{计算练习与逐步解答}

\textbf{练习 1}\quad 教学关节希望极点为 $-4,-6$，求状态反馈 $K$。

\textbf{解}\quad 期望多项式为 $s^2+10s+24$。故 $k_1=0.1\times24=2.4$，$k_2=0.1\times10-0.2=0.8$，即 $K=[2.4,0.8]$。常值参考的控制律为 $u=-2.4(q-q_d)-0.8v$。

\textbf{练习 2}\quad 取 $e=0.03\,\mathrm{rad}$、$\dot e=-0.2\,\mathrm{rad/s}$，沿用 $J_r=0.1$、$K_p=10$、$b+K_d=1.4$。计算 $V$ 和 $\dot V$。

\textbf{解}\quad 动能为 $0.5\times0.1\times0.2^2=0.002\,\mathrm{J}$；虚拟势能为 $0.5\times10\times0.03^2=0.0045\,\mathrm{J}$，故 $V=0.0065\,\mathrm{J}$。$\dot V=-1.4\times0.2^2=-0.056\,\mathrm{W}$，能量正在减少。

\textbf{练习 3}\quad ILC 的静态增益 $G=0.5$，学习率 $L$ 的收敛范围是什么？若 $L=3$，初始误差为 0.08，三次更新后误差是多少？

\textbf{解}\quad $|1-0.5L|<1$ 等价于 $0<L<4$。取 $L=3$，误差因子为 $-0.5$，因此三次后 $e_3=(-0.5)^3\times0.08=-0.01$。误差符号每次交替，但幅值在缩小；符号交替本身不等于发散。

\section*{本章与原稿的对应}
本章整合 [A，第15--18、39--43页] 与 [B，第63--67、71--74页]。补足了 Riccati 方程的来源、标量展开、积分增广和稳定性证明，并区分局部线性稳定、理想非线性稳定、参数收敛与实际饱和后的行为。

\chapter{采样离散控制与滤波的实际影响}
\label{ch:digital}

连续控制方程描述每一瞬间，程序只在一串离散时刻执行。每个周期依次获得数据、计算、发送命令，驱动器再在两个更新之间保持上一条命令。这些步骤改变了闭环，不能把连续公式逐字翻成代码后就认为性能不变。本章以 $200\,\mathrm{Hz}$、即 $T_s=5\,\mathrm{ms}$ 的外部控制周期为主要教学例子；它只代表一种采样设置，不代表任何机械臂的固有带宽。

\section{先分清采样量化保持与延迟}

采样是只在 $t_k=kT_s$ 读信号；量化是把连续幅值舍入到有限分辨率；保持是两次更新之间维持某个输出；延迟是数据所代表的时刻与输入真正生效的时刻不同。这四件事可以同时存在，修正其中一项不会自动消除其他三项。

例如一条关节状态在 $t=1.000\,\mathrm{s}$ 被测量，主机在 $1.003\,\mathrm{s}$ 才收到，即使计算只用了 $0.2\,\mathrm{ms}$，控制器使用的状态也已经旧了 $3.2\,\mathrm{ms}$。仅打印计算函数耗时不足以描述闭环延迟，必须记录测量时间、接收时间、命令发送时间和实际生效时序。

\subsection{采样后高频会伪装成低频}

对于连续正弦 $x(t)=\sin(2\pi ft)$，采样后是 $x[k]=\sin(2\pi f k/f_s)$。把 $f$ 换成 $f+mf_s$，$m$ 为整数，采样结果不变，因为多出来的是 $2\pi mk$。这就是混叠的来源，不是软件计算误差。

\begin{examplebox}[算例 13A 170 Hz 为什么会看起来像 30 Hz]
取 $f_s=200\,\mathrm{Hz}$：
\[
\sin(2\pi170k/200)
=\sin(2\pi k-2\pi30k/200)
=-\sin(2\pi30k/200).
\]
同一串采样点既可由 170 Hz 信号产生，也可由负相位的 30 Hz 信号产生。软件拿到采样点后已经无法区分两者。应在采样前做合适的抗混叠处理，或先以更高频率采样并滤波后再降采样；不是在 200 Hz 数据上随意加一个低通就能恢复丢失的信息。
\end{examplebox}

\begin{figure}[htbp]
\centering\includegraphics[width=.95\linewidth]{figures/cf_sampling.pdf}
\caption{左图的高频信号与低频信号通过相同采样点；右图展示零阶保持的阶梯输出。二者分别是信息混叠和输出保持，不能混为一个问题。}
\end{figure}

采样定理要求信号有效带限且 $f_s>2f_{\max}$。实际还要给抗混叠滤波器的过渡带留空间。控制设计则比信号重建更严格，因为采样虽能记录某个频率，闭环仍可能没有足够相位余量去控制该频率。

\section{从连续模型精确得到一拍更新}

对 $\dot x=Ax+Bu$，在区间 $[kT_s,(k+1)T_s)$ 内保持 $u=u_k$，解微分方程得
\begin{equation}
x_{k+1}=A_dx_k+B_du_k,
\qquad A_d=e^{AT_s},\qquad
B_d=\int_0^{T_s}e^{A\xi}B\,d\xi.
\label{eq:cf_zoh_exact}
\end{equation}
这叫零阶保持等效离散化。它对区间内恒定输入在采样点的状态是精确的；不意味着区间内轨迹变成直线，也不包含额外通信延迟。

\subsection{手算教学关节的精确离散模型}

速度方程为 $\dot v=-(b/J_r)v+u/J_r$。一拍后
\[
v_{k+1}=a v_k+\frac{1-a}{b}u_k,
\qquad a=e^{-bT_s/J_r}.
\]
位置由区间内速度积分：
\begin{align}
q_{k+1}
&=q_k+\int_0^{T_s}
\left[e^{-b\xi/J_r}v_k+\frac{1-e^{-b\xi/J_r}}b u_k\right]d\xi\\
&=q_k+\frac{J_r}{b}(1-a)v_k
+\left[\frac{T_s}{b}-\frac{J_r}{b^2}(1-a)\right]u_k.
\end{align}
对 $J_r=0.1$、$b=0.2$、$T_s=0.005$，有 $a=e^{-0.01}=0.9900498$，因此
\begin{equation}
A_d=\begin{bmatrix}1&0.00497508\\0&0.99004983\end{bmatrix},
\qquad B_d=\begin{bmatrix}0.000124584\\0.04975083\end{bmatrix}.
\end{equation}

\begin{examplebox}[算例 13B 相同一拍用三种积分方法会怎样]
从 $q_0=v_0=0$ 开始，保持 $u=1\,\mathrm{N\,m}$ 一拍。精确结果是
\[
v_1=0.0497508\,\mathrm{rad/s},\qquad
q_1=0.000124584\,\mathrm{rad}.
\]
前向欧拉先用旧状态算导数，给 $v_1=0.05$、$q_1=0$；半隐式欧拉先更新速度再更新位置，给 $v_1=0.05$、$q_1=0.00025$。两种数值法都在近似同一个连续对象，但一拍的位置分别偏小和偏大。控制周期变小能减小误差，积分方法本身仍需选清楚。
\end{examplebox}

\section{连续稳定不等于离散实现稳定}

离散线性系统 $x_{k+1}=A_dx_k$ 的渐近稳定条件是全部特征值严格位于单位圆内，即 $|z_i|<1$。若用精确映射 $z=e^{sT_s}$，连续极点 $\Re(s)<0$ 对应 $|z|<1$。但前向欧拉把 $\dot x=-ax$ 变成 $x_{k+1}=(1-aT_s)x_k$，于是要求
\[
|1-aT_s|<1\quad\Longleftrightarrow\quad0<aT_s<2.
\]

\begin{examplebox}[算例 13C 一个本来稳定的衰减被欧拉算成发散]
连续系统 $\dot x=-500x$ 显然指数衰减。取 $T_s=0.005\,\mathrm{s}$，精确离散极点为 $e^{-2.5}=0.082085$，前向欧拉极点却是 $1-2.5=-1.5$。从 $x_0=1$ 开始，欧拉序列是 $1,-1.5,2.25,-3.375,\ldots$，符号交替且幅值越来越大。物理模型并未变得不稳定，失稳来自数值实现。
\end{examplebox}

\begin{figure}[htbp]
\centering\includegraphics[width=.95\linewidth]{figures/cf_discrete_stability.pdf}
\caption{连续稳定系统的精确更新与不合适的前向欧拉更新。单位圆判断针对离散极点，不能继续只看连续极点在不在左半平面。}
\end{figure}

后向欧拉的替换为 $s\leftarrow(z-1)/(T_sz)$；Tustin 双线性变换为
\[
s\leftarrow\frac{2}{T_s}\frac{z-1}{z+1}.
\]
它们对连续左半平面的映射具有更好的稳定性性质，但仍会改变频率和瞬态。后向欧拉可能把部分连续不稳定动态数值上压住，因而“看起来不发散”不能用来确认真实系统稳定。Tustin 对频率有扭曲，设计窄带陷波时通常需要预畸变。

\section{给相位建立一张预算表}

纯延迟 $T_d$ 的传递函数是 $e^{-sT_d}$。令 $s=j\omega$，幅值为 1，相位为 $-\omega T_d$。零阶保持在相对较低频率处可近似贡献半拍延迟；如果还要等下一拍输出，则常见的一种时序近似为总延迟 $1.5T_s$。这不是所有控制器的固定常数，实际值要根据时间戳和执行结构确定。

\begin{examplebox}[算例 13D 200 Hz 更新不代表能闭环控制到 200 Hz]
取 $T_s=5\,\mathrm{ms}$，假设一拍计算等待加半拍保持，总等效延迟 $7.5\,\mathrm{ms}$。相位损失为
\[
\Delta\phi=360^\circ fT_d.
\]
在 5、10、20 Hz 处分别为 $13.5^\circ$、$27^\circ$、$54^\circ$。若设计前已有 $65^\circ$ 相位裕度，想把穿越频率放在 20 Hz，单是这一项就只剩约 $11^\circ$。

再考虑 20 Hz 截止的一阶低通，在 10 Hz 处损失 $\arctan(10/20)=26.565^\circ$。若穿越仍大约在 10 Hz，延迟与滤波合计约 $53.6^\circ$。真实穿越点也可能因滤波而改变，因此最终必须用完整开环重新求穿越，而不是只做度数相减。
\end{examplebox}

“每秒算 200 次”只说明更新率。可用带宽还由力矩内环、机械柔性、测量滤波、通信、计算和保持共同决定。阻尼增益增大后看见抖动，不应立刻认定“阻尼不够”；速度噪声和相位滞后可能已经让名义阻尼变成不合适的能量输入。

\section{低通滤波的公式怎样变成代码}

一阶连续低通 $H(s)=a/(s+a)$ 对应 $\dot y=a(x-y)$，其中 $a=2\pi f_c$。若输入在一拍内保持为 $x_k$，精确递推为
\[
y_{k+1}=\alpha y_k+(1-\alpha)x_k,
\qquad\alpha=e^{-aT_s}.
\]
这里的下标明确表示：当前输入 $x_k$ 作用到下一拍状态 $y_{k+1}$。某些实现收到当前采样后立即更新并把新值称为 $y_k$，只是索引约定不同，但与其他信号同步时必须知道它代表哪个时刻。

后向欧拉常见的“收到当前输入立即更新”形式为
\[
y_k=\frac{1}{1+aT_s}y_{k-1}
+\frac{aT_s}{1+aT_s}x_k.
\]
Tustin 则给出
\[
y_k=\frac{2-aT_s}{2+aT_s}y_{k-1}
+\frac{aT_s}{2+aT_s}(x_k+x_{k-1}).
\]
三者输入时序和频率响应并不完全相同，不能只把不同的系数都塞进同一条递推式然后称作完全等价。

\begin{examplebox}[算例 13E 计算 20 Hz 一阶低通系数]
取 $f_c=20\,\mathrm{Hz}$、$T_s=0.001\,\mathrm{s}$，得到 $aT_s=0.125664$。精确更新中新输入权重为 $1-e^{-0.125664}=0.118089$；后向欧拉的新输入权重为 $0.125664/1.125664=0.111635$；Tustin 的两个输入系数均为 $0.0591174$，旧输出系数为 $0.881765$。

若使用后向欧拉，旧输出为 1，当前输入为 2，则新输出为 $1+0.111635(2-1)=1.111635$。参数越靠近采样能力的上限，不同离散方法的差异越不能忽略。
\end{examplebox}

\subsection{更高阶不自动代表更适合闭环}

Butterworth 低通追求平坦通带；Chebyshev 允许纹波换取陡峭过渡；Bessel 侧重通带内较平坦的群延迟；椭圆滤波器同时允许通带和阻带纹波以获得更陡过渡。这些是设计目标，不是优劣排名。闭环关心的是噪声衰减与相位一起变化后的结果。

高阶 IIR 常拆成二阶节，即 SOS。对一个归一化二阶节
\[
H(z)=\frac{b_0+b_1z^{-1}+b_2z^{-2}}{1+a_1z^{-1}+a_2z^{-2}},
\]
对应差分式为
\[
y_k=b_0x_k+b_1x_{k-1}+b_2x_{k-2}-a_1y_{k-1}-a_2y_{k-2}.
\]
公式右边减号来自把分母乘到左边后移项。把高阶多项式拆成二阶节通常能改善浮点数值条件，但稳定性仍需检查每节极点、系数精度和状态初始化，不能只因为用了 SOS 就保证不会出问题。

\begin{figure}[htbp]
\centering\includegraphics[width=.98\linewidth]{figures/cf_filters.pdf}
\caption{低通同时改变幅值和相位；双向滤波把幅值平方；陷波只针对一个频带。中图采用连续 Butterworth 幅频公式展示原理，实际数字滤波的精确频响应以所用系数计算。}
\end{figure}

\subsection{离线零相位不能搬进实时回路}

正向滤波再反向滤波需要未来数据，是非因果处理。对实系数滤波器，理想稳态频率响应变成
\[
H_{\rm fb}(e^{j\Omega})=H(e^{j\Omega})H(e^{-j\Omega})
=|H(e^{j\Omega})|^2.
\]
相位被抵消，幅值却不保持原样。单向截止处幅值为 $1/\sqrt2$，双向后变成 $1/2$，对应约 $-6.02\,\mathrm{dB}$，不是 $-3.01\,\mathrm{dB}$。边界填充和初值处理还会产生端点效应。

对四阶模拟 Butterworth 原型，单向幅值为 $[1+(f/f_c)^8]^{-1/2}$，双向后为 $[1+(f/f_c)^8]^{-1}$。若名义 $f_c=8\,\mathrm{Hz}$，双向结果在 8 Hz 为 0.5。实际离散滤波要考虑频率映射。第9章的辨识数据处理中还必须注意：同一个线性滤波算子应施加到回归等式两边；仅仅分别把 $q$ 和 $\tau$ 平滑一下，并不能保证非线性构造的回归矩阵仍与力矩完全同步。

\subsection{移动平均和局部多项式}

$N$ 点滑动平均是有限脉冲响应 FIR 滤波器。对称线性相位 FIR 的群延迟为 $(N-1)T_s/2$。例如 31 点、200 Hz 采样会带来 $75\,\mathrm{ms}$ 延迟；在 5 Hz 处对应 $135^\circ$ 相位。这种处理可能适合离线看曲线，却很难直接放进要求较快响应的反馈路径。

Savitzky--Golay 方法在小窗口内拟合多项式，再由拟合多项式取值或求导。五点二次拟合的一阶和二阶导数为
\begin{align}
\dot y_k&\approx\frac{-2y_{k-2}-y_{k-1}+y_{k+1}+2y_{k+2}}{10T_s},\\
\ddot y_k&\approx\frac{2y_{k-2}-y_{k-1}-2y_k-y_{k+1}+2y_{k+2}}{7T_s^2}.
\end{align}
它不是不知道未来也能算出的神奇微分器；上述居中窗口需要 $k+1,k+2$ 两个未来样本。若实时实现为延迟输出，须计入两拍以上的等待。

\begin{examplebox}[算例 13F 用五个位置样本计算速度和加速度]
取 $q(t)=t^2$，$T_s=0.1\,\mathrm{s}$，在中心时刻 $t=0.2$ 周围的五个样本为 $[0,0.01,0.04,0.09,0.16]$。代入一阶公式：
\[
\dot q=\frac{0-0.01+0.09+0.32}{1}=0.4.
\]
代入二阶公式：
\[
\ddot q=\frac{0-0.01-0.08-0.09+0.32}{0.07}=2.
\]
与解析导数 $2t=0.4$、$2$ 一致，因为窗口中真实曲线恰好是二次多项式。遇到更复杂的高频运动或噪声时，它不再精确，窗口长度仍存在分辨率与平滑程度的权衡。
\end{examplebox}

\section{量化噪声为什么会在求导后放大}

若一圈有 $N_c$ 个有效计数，位置分辨率为 $\Delta q=2\pi/N_c$。这里应使用解码后的实际计数，不要把编码器“线数”和四倍频后的“计数数”混为一谈。直接差分速度为
\[
\hat v_k=\frac{q_k-q_{k-1}}{T_s}.
\]
一个计数跳变对应 $\Delta q/T_s$ 的速度台阶。若每次量化误差独立、均匀分布在 $[-\Delta q/2,\Delta q/2]$，单个位置误差方差为 $\Delta q^2/12$，两次差值的方差相加，因此
\[
\sigma_v=\frac{\Delta q}{\sqrt6\,T_s}.
\]
这个独立均匀假设只用于估算；真实低速编码器误差常呈周期性和相关性。

\begin{examplebox}[算例 13G 17 位编码器在 200 Hz 下的速度台阶]
$N_c=2^{17}=131072$，所以 $\Delta q=4.79369\times10^{-5}\,\mathrm{rad}$。在 $T_s=0.005\,\mathrm{s}$ 下，一个计数形成 $0.0095874\,\mathrm{rad/s}$ 的速度台阶。独立量化误差模型给出 $\sigma_v\approx0.003914\,\mathrm{rad/s}$。

若把采样周期减到 1 ms 而不改变差分方法，一个计数的速度台阶扩大到 $0.0479369\,\mathrm{rad/s}$。因此更高采样率不自动让简单差分更平滑；它需要与估计窗口、模型和滤波一起设计。
\end{examplebox}

\section{一拍控制更新必须明确状态顺序}

下列是带反算抗饱和和测量微分滤波的结构性伪代码。它用本书的 $e=q-q_d$，积分变量直接保存力矩贡献。这里没有机械臂接口，目的是看清更新顺序。
\begin{lstlisting}[language=C++,basicstyle=\ttfamily\small]
// q_prev must be initialized from the first measurement.
error = q - q_des;
velocity_raw = (q - q_prev) / dt;
velocity_f += alpha * (velocity_raw - velocity_f);
u_raw = -kp * error - kd * velocity_f + i_term + u_ff;
u_sent = clamp(u_raw, u_min, u_max);
i_term += dt * (-ki * error + kaw * (u_sent - u_raw));
q_prev = q;
\end{lstlisting}

积分更新在输出之后或之前均可设计，但推导、代码和日志必须一致。发生丢帧时，不应悄悄用固定 $T_s$ 对实际间隔不同的两个样本求导。使用实测间隔可以修正差分比例，却不能消除不规则采样本身造成的闭环变化。

增量式 PID 把 $u_k-u_{k-1}$ 作为计算结果，但积分信息仍藏在历史输出中。它不会天然免疫积分饱和，也不会自动实现无扰切换。切换模式时应初始化内部状态，使新控制器当前输出与切换前一致；对式\eqref{eq:cf_aw}，可令
\[
i_0=u_{\rm old}+K_pe+K_d\dot q-u_{\rm ff}.
\]
若新的输出限制更窄，则还要考虑限幅导致的实际变化。

\section{计算练习与逐步解答}

\textbf{练习 1}\quad 控制周期为 2 ms，一拍计算延迟加零阶保持近似为 $1.5T_s$。在 20 Hz 处损失多少相位？

\textbf{解}\quad 总延迟为 $0.003\,\mathrm{s}$，相位损失为 $360^\circ\times20\times0.003=21.6^\circ$。如果额外通信还有 1 ms，新增 $7.2^\circ$，合计 $28.8^\circ$。

\textbf{练习 2}\quad $\dot x=-80x$ 用前向欧拉，求允许的采样周期范围。取 10 ms 和 30 ms 时分别会怎样？

\textbf{解}\quad 稳定要求 $0<80T_s<2$，即 $0<T_s<0.025\,\mathrm{s}$。10 ms 时离散极点为 0.2，单调衰减；30 ms 时极点为 $-1.4$，交替发散。精确离散在 30 ms 时的极点为 $e^{-2.4}=0.090718$，仍然稳定。

\textbf{练习 3}\quad 采样率 500 Hz，使用 21 点对称 FIR。求群延迟，在 10 Hz 处相当于多少相位？

\textbf{解}\quad 延迟 $(21-1)/(2\times500)=0.02\,\mathrm{s}$，相位为 $-360^\circ\times10\times0.02=-72^\circ$。这是固定群延迟，不是 CPU 算得慢产生的延迟；换更快 CPU 无法消除它。

\textbf{练习 4}\quad 单向滤波器在某频率幅值为 0.8，相位为 $-40^\circ$。理想双向滤波后的幅值与相位是什么？

\textbf{解}\quad 幅值为 $0.8^2=0.64$，相位为 $-40^\circ+40^\circ=0$。幅值对应 $20\log_{10}(0.64)=-3.876\,\mathrm{dB}$。零相位并不表示没有改变幅值，也不意味着可实时执行。

\section*{本章与原稿的对应}
本章整合 [A，第18--22页] 与 [B，第67--70、78--83页]。补足精确 ZOH 数值、索引时序、混叠推导、量化噪声方差及五点求导计算，修正了“增量式自然无扰”“双向滤波仍是原来的截止定义”等容易误用的表述。

\chapter{状态观测外力估计与模型残差}
\label{ch:observers}

编码器给出位置，却不会直接告诉我们干净的速度；电流可以换算成电机力矩，却不会直接告诉我们有人推了末端多少牛；一个新负载可能表现得很像持续外力。观测器解决的是“从已有输入输出推断看不见的量”，而不是把所有未知因素自动区分开。本章逐步推导状态观测器、卡尔曼滤波、DOB、ESO 和动量观测器，并明确它们各自估计的到底是什么。

\section{状态观测器是带测量纠正的模型}

对 $\dot x=Ax+Bu$、$y=Cx$，如果直接在计算机里运行 $\dot{\hat x}=A\hat x+Bu$，初值或模型有偏差时，估计未必能追上真实状态。加入输出残差 $y-C\hat x$：
\begin{equation}
\dot{\hat x}=A\hat x+Bu+L(y-C\hat x).
\end{equation}
$\hat x$ 表示估计，$L$ 是观测增益。定义估计误差 $\tilde x=x-\hat x$，真实方程减去估计方程：
\begin{equation}
\dot{\tilde x}=(A-LC)\tilde x.
\end{equation}
输入项抵消，因为观测器和真实系统使用了同一个输入。若观测器收到的是“期望力矩”，真实执行器却因饱和只输出一部分，输入不再相同，该差额会进入估计误差。记录实际施加的命令或可靠的实测输入因此很重要。

\begin{examplebox}[算例 14A 用位置估计关节速度]
教学关节 $A=\left[\begin{smallmatrix}0&1\\0&-2\end{smallmatrix}\right]$、$C=[1,0]$，设 $L=[l_1,l_2]^\top$，则
\[
A-LC=\begin{bmatrix}-l_1&1\\-l_2&-2\end{bmatrix}.
\]
其特征多项式为
\[
s^2+(l_1+2)s+(2l_1+l_2).
\]
希望误差极点为 $-30,-40$，期望多项式为 $s^2+70s+1200$，因此 $l_1=68$、$l_2=1200-136=1064$。

观测器两行分别为
\[
\dot{\hat q}=\hat v+68(q-\hat q),
\]
\[
\dot{\hat v}=-2\hat v+10u+1064(q-\hat q).
\]
如果位置残差是 $0.001\,\mathrm{rad}$，修正项分别为 $0.068\,\mathrm{rad/s}$ 和 $1.064\,\mathrm{rad/s^2}$。噪声也会经过同一增益放大，因此把极点设得更快并非免费。
\end{examplebox}

普通龙伯格观测器估计的是模型中已有的状态。上例只有位置和速度，没有新增“未知负载”状态，所以它不能保证独立估计并消除任意常值扰动。若需要估扰，必须增广模型或采用明确的扰动估计结构。也不能说“凡是观测器都假设扰动缓变”：缓变假设适用于某些扰动模型，不是所有状态观测问题的定义。

\subsection{分离原理的边界}

理想线性模型下，用 $u=-K\hat x$ 反馈，组合系统在适当坐标中呈块三角结构，极点由 $A-BK$ 和 $A-LC$ 两部分组成。这称为分离原理。它允许先设计反馈、再设计观测器；但不保证饱和、时延、模型失配和非线性下仍能任意独立整定。观测器“比控制器快若干倍”是常见起点，最终仍需检查噪声放大和离散极点。

\section{卡尔曼滤波把可信度放进计算}

设离散模型为
\[
x_{k+1}=F x_k+B_du_k+w_k,\qquad z_k=Hx_k+v_k.
\]
$w_k$ 是过程噪声，表示模型没有解释的状态变化；$v_k$ 是测量噪声。协方差 $Q_w=\mathbb E[w_kw_k^\top]$ 和 $R_v=\mathbb E[v_kv_k^\top]$ 描述不确定度。协方差的对角元是方差，单位是相应状态单位的平方；非对角元描述误差之间的相关性。

在常见的零均值、白噪声、不相关等假设下，预测和更新分别为
\begin{align}
\hat x_k^-&=F\hat x_{k-1}+B_du_{k-1},\\
P_k^-&=FP_{k-1}F^\top+Q_w,\\
S_k&=HP_k^-H^\top+R_v,\\
K_k&=P_k^-H^\top S_k^{-1},\\
\hat x_k&=\hat x_k^-+K_k(z_k-H\hat x_k^-).
\label{eq:cf_kalman}
\end{align}
上标 $-$ 表示还没有吸收当前测量的预测；$P$ 是估计误差协方差；$S_k$ 是新息协方差；$z_k-H\hat x_k^-$ 叫新息，即“本次测量带来的新信息”。$K_k$ 是随可信度变化的修正权重。

精确算术下可写 $P_k=(I-K_kH)P_k^-$。浮点实现常用 Joseph 形式保持对称半正定更可靠：
\begin{equation}
P_k=(I-K_kH)P_k^-(I-K_kH)^\top+K_kR_vK_k^\top.
\end{equation}
线性高斯条件下，卡尔曼滤波给出条件均值意义下的最优估计；即使噪声不是高斯，在相应二阶统计假设下仍有最佳线性估计的意义。不能把“高斯”当作能不能使用算法的简单开关，也不能把错误的噪声协方差当成最优保证。

\begin{examplebox}[算例 14B 两拍标量卡尔曼滤波]
估计一个近似不变的标量，$F=H=1$、$Q_w=0.01$、$R_v=0.25$，初值 $\hat x_0=0$、$P_0=1$。第一拍测量 $z_1=1.2$：
\[
P_1^-=1.01,\qquad K_1=\frac{1.01}{1.01+0.25}=0.801587,
\]
\[
\hat x_1=0+0.801587\times1.2=0.961905,
\qquad P_1=0.200397.
\]
第二拍测量 $z_2=0.9$：
\[
P_2^-=0.210397,\quad K_2=0.456990,
\]
\[
\hat x_2=0.961905+0.456990(0.9-0.961905)=0.933615,
\]
\[
P_2=0.114248.
\]
第一拍估计很不确定，较相信测量；吸收信息后，第二拍不再完全跟着测量走。若模型持续允许过程变化，稳态增益不会降为零。此例稳态后验方差为
\[
P=\frac{-Q_w+\sqrt{Q_w^2+4Q_wR_v}}2=0.0452494,
\]
稳态增益约为 0.180998。
\end{examplebox}

\subsection{为什么位置测量也能修正速度}

匀速模型取 $x=[q,v]^\top$，$F=\left[\begin{smallmatrix}1&T_s\\0&1\end{smallmatrix}\right]$、$H=[1,0]$。预测位置依赖速度，所以一段时间后位置误差与速度误差相关，$P$ 的非对角项不再为零。新的位置测量因此可以同时修正位置和速度。

\begin{examplebox}[算例 14C 矩阵卡尔曼的一次完整更新]
设 $T_s=0.1\,\mathrm{s}$、$\hat x=[0,1]^\top$、$P=\operatorname{diag}(0.01,0.04)$、$Q_w=\operatorname{diag}(0.0001,0.001)$、$R_v=0.0025$。这里是为了方便手算选择的教学协方差，不是编码器通用参数。

预测得到
\[
\hat x^-=[0.1,1]^\top,
\qquad P^-=\begin{bmatrix}0.0105&0.004\\0.004&0.041\end{bmatrix}.
\]
$S=0.0105+0.0025=0.013$，所以
\[
K=\frac1{0.013}\begin{bmatrix}0.0105\\0.004\end{bmatrix}
=\begin{bmatrix}0.807692\\0.307692\end{bmatrix}.
\]
测得位置 $z=0.12$，新息为 0.02。更新后
\[
\hat x=\begin{bmatrix}0.116154\\1.006154\end{bmatrix},
\quad P=\begin{bmatrix}0.00201923&0.00076923\\0.00076923&0.03976923\end{bmatrix}.
\]
测量比预测走得更远，滤波器除了增加位置估计，还小幅提高速度估计。这就是矩阵中的相关项在发挥作用。
\end{examplebox}

若把每一拍未知加速度 $a_k$ 视为方差 $\sigma_a^2$ 的随机常值，状态变化为 $[T_s^2/2,T_s]^\top a_k$，于是过程协方差为
\[
Q_w=\sigma_a^2\begin{bmatrix}T_s^4/4&T_s^3/2\\T_s^3/2&T_s^2\end{bmatrix}.
\]
若改用连续白噪声加速度驱动，积分后则会出现 $T_s^3/3$、$T_s^2/2$、$T_s$ 等系数。两式对应不同随机过程假设，不能不解释就混用。

\begin{figure}[htbp]
\centering\includegraphics[width=.95\linewidth]{figures/cf_kalman.pdf}
\caption{相同有噪位置数据的测速比较。直接差分放大噪声，低通带来滞后，卡尔曼依赖运动模型。曲线是固定随机种子的教学仿真，不表示卡尔曼对任何运动都能同时做到零噪声和零延迟。}
\end{figure}

EKF 对非线性状态方程和测量方程求局部雅可比；UKF 用代表性采样点近似传播均值和协方差；误差状态滤波在局部小误差空间工作，尤其适合旋转状态。它们的增加并不取消模型、单位、时间同步和噪声统计这些基础要求。

\section{DOB 估的是输入端总残差}

考虑 $J_r\ddot q+b\dot q=u+d$。取名义对象 $G_n$，若模型精确，$G_n^{-1}q-u=d$。用稳定低通 $Q_f(s)$ 包裹该残差：
\begin{equation}
\hat d=Q_f(s)\big[G_n^{-1}(s)q-u\big].
\end{equation}
这里扰动以正号加在输入端，因此抵消时应使用 $u=u_0-\hat d$。如果一开始把方程写成 $J_r\ddot q=u-d$，估计定义和补偿符号都要一起改变。源稿中某些公式的差异就是来自没有统一这个约定。

当 $G_n=1/(J_ns^2+b_ns)$，逆模型含有二阶微分。为了使 $Q_fG_n^{-1}$ 可实现，$Q_f$ 的相对阶数至少要补足这两阶。相对阶数是分母阶数减去分子阶数；可实现并不等于没有噪声问题，整个从位置到估计扰动的高频增益仍可能很大。

如果真实惯量与名义惯量不同，残差中还包含 $(J_n-J_r)\ddot q$；摩擦、未补偿重力和实际外力也会叠加进去。DOB 可以帮助抵消总残差，却不能仅凭这个结果宣称已经测出了“人手推力”。

\section{ESO 怎样把扰动增加成一个状态}

把单关节写成
\[
\ddot q=b_0u+f(t),
\]
其中 $b_0$ 为名义输入增益，单位是角加速度除以力矩；$f$ 把重力、摩擦、模型差异和外力等统称为角加速度形式的总扰动。增加状态 $x_3=f$，并假设其变化率在关注频段内可控制，则线性扩张状态观测器为
\begin{align}
\dot z_1&=z_2+\beta_1(q-z_1),\\
\dot z_2&=z_3+b_0u+\beta_2(q-z_1),\\
\dot z_3&=\beta_3(q-z_1).
\end{align}
忽略 $\dot f$ 时，观测误差特征式为 $s^3+\beta_1s^2+\beta_2s+\beta_3$。希望三个极点都在 $-\omega_o$，展开 $(s+\omega_o)^3$ 得
\[
\beta_1=3\omega_o,\quad\beta_2=3\omega_o^2,\quad\beta_3=\omega_o^3.
\]
因此带宽增大时，第三个增益按立方增加，这是噪声很快变严重的原因之一。

\begin{examplebox}[算例 14D ESO 补偿如何换回力矩]
取 $b_0=10$、$\omega_o=30\,\mathrm{rad/s}$，得到增益 $(90,2700,27000)$。定点目标 $q_d=0.1$，当前估计 $z_1=0.12$、$z_2=-0.05$、$z_3=1\,\mathrm{rad/s^2}$。期望误差动力学增益为 $k_p=100\,\mathrm{s^{-2}}$、$k_v=20\,\mathrm{s^{-1}}$，则
\[
u=\frac{-k_p(z_1-q_d)-k_vz_2-z_3}{b_0}
=\frac{-2+1-1}{10}=-0.2\,\mathrm{N\,m}.
\]
不能把 $z_3=1$ 直接当作 $1\,\mathrm{N\,m}$ 加到力矩上；它的单位是角加速度，必须除以 $b_0$。在 5 ms 周期下还要检查观测器离散实现，连续极点配置本身不处理实际时延和采样。
\end{examplebox}

\section{动量观测器怎样避开加速度求导}

采用全书统一动力学约定：
\begin{equation}
M(q)\ddot q+C(q,\dot q)\dot q+g(q)+\tau_f
=\tau_{\rm cmd}+\tau_{\rm ext}.
\end{equation}
$\tau_{\rm ext}$ 是环境作用于机器人的关节广义力矩。直接反解外力应为
\begin{equation}
\tau_{\rm ext}=M\ddot q+C\dot q+g+\tau_f-\tau_{\rm cmd}.
\label{eq:cf_direct_external}
\end{equation}
请注意它与“实测驱动力矩减模型所需力矩”的顺序相反。若传感器输出定义成机器人对环境的反作用力，符号会再反过来，所以必须先定义物理方向。

定义广义动量 $p=M(q)\dot q$。对时间求导：
\begin{align}
\dot p&=M\ddot q+\dot M\dot q\\
&=\tau_{\rm cmd}+\tau_{\rm ext}-g-\tau_f+(\dot M-C)\dot q.
\end{align}
对满足机器人动力学标准性质的 $C$，$\dot M-2C$ 反对称，等价地 $\dot M=C+C^\top$，因此
\begin{equation}
\dot p=\tau_{\rm cmd}+C^\top\dot q-g-\tau_f+\tau_{\rm ext}.
\end{equation}
这一步消去了显式加速度，代价是必须有质量矩阵、速度和一致的动力学模型。

定义名义动量变化项 $h_p=\tau_{\rm cmd}+\hat C^\top\dot q-\hat g-\hat\tau_f$，构造
\begin{equation}
r(t)=K_o\left[p(t)-p(0)-\int_0^t(h_p(\xi)+r(\xi))\,d\xi\right].
\end{equation}
$K_o$ 可取正对角阵，单位为 $\mathrm{s^{-1}}$。模型精确且 $K_o$ 恒定时，对上式求导：
\begin{align}
\dot r&=K_o(\dot p-h_p-r)\\
&=K_o(\tau_{\rm ext}-r).
\end{align}
所以每一轴的残差就是外力矩的一阶低通，而不是外力矩的瞬时精确测量。若初始时刻并非静止，$p(0)$ 不能擅自置零。

\begin{examplebox}[算例 14E 从观测器带宽算出检测滞后]
对一轴取 $K_o=20\,\mathrm{s^{-1}}$，时间常数 $1/K_o=0.05\,\mathrm{s}$，对应截止频率 $K_o/(2\pi)=3.183\,\mathrm{Hz}$。从零开始受到 $2\,\mathrm{N\,m}$ 阶跃外力：
\[
r(t)=2(1-e^{-20t}).
\]
50 ms 后为 $1.2642\,\mathrm{N\,m}$，150 ms 后为 $1.9004\,\mathrm{N\,m}$。若外力只持续 5 ms，峰值只到
\[
r_{\max}=2(1-e^{-0.1})=0.1903\,\mathrm{N\,m}.
\]
短冲击被严重削峰。理想一阶滤波的直流增益为 1，在完整时间区间积分、初末状态回零等条件下可以保留总面积，但这不等于峰值、接触时间或实时冲击强度都准确。
\end{examplebox}

\begin{figure}[htbp]
\centering\includegraphics[width=.95\linewidth]{figures/cf_momentum.pdf}
\caption{动量残差对阶跃和短脉冲的响应。增益提高能缩短滞后，同时也会让模型和速度测量噪声更明显。}
\end{figure}

一种实现方法是保存内部动量估计 $\hat p$，用 $r=K_o(p-\hat p)$，再积分 $\dot{\hat p}=h_p+r$。若使用前向欧拉，对理想单轴误差子系统有 $0<K_oT_s<2$ 的数值稳定要求；这只是最基本条件，不代表接近边界时具有可用性能。状态、输入和模型值还必须来自相容的时间点。

\subsection{模型误差会伪装成外力}

若模型有残差 $\Delta$，估计方程变成 $\dot r=K_o(\tau_{\rm ext}+\Delta-r)$。漏掉一个负载、摩擦方向错了、温度改变力矩常数、输出电流饱和，都可能成为 $r$ 的持续偏置。

例如没有外接触，但漏掉的重力项在某轴产生 $0.8\,\mathrm{N\,m}$，观测器最终也会报告约 $0.8\,\mathrm{N\,m}$ 的残差。静止调零能去掉当前姿态的偏置，却不能保证换姿态后仍然正确，更不能在本来有接触时把接触“调没”。

\section{关节残差怎样换成末端力}

若外力确实只作用在已知末端点，而且力旋量与雅可比在同一坐标系、同一参考点表达，则
\[
\tau_{\rm ext}=J^\top F_{\rm ext}.
\]
可用最小二乘求等效末端力
\[
\hat F=(J^\top)^+r.
\]
它回答的是“哪一个末端力最能解释这些关节残差”。如果推的是肘部、同时存在多个接触，或者残差含有零空间成分，这个结果未必是真实接触力。

\begin{examplebox}[算例 14F 两连杆末端力反算]
在标准教学位形 $q=(0,\pi/2)$，
\[
J=\begin{bmatrix}-0.5&-0.5\\0.5&0\end{bmatrix}.
\]
假设残差为 $r=[-1.5,-1]^{\top}\,\mathrm{N\,m}$。由
\[
-0.5f_x+0.5f_y=-1.5,\qquad -0.5f_x=-1
\]
先得 $f_x=2\,\mathrm{N}$，再得 $f_y=-1\,\mathrm{N}$。回代 $J^\top[2,-1]^\top=[-1.5,-1]^\top$，计算闭合。

对七轴臂应同时检查重构残差 $r-J^\top\hat F$。它明显不为零，说明“一个末端力解释全部残差”的假设不充分，而不应强行把全部残差都解释为末端接触。
\end{examplebox}

\section{从残差到事件还差哪些判断}

阈值必须基于无接触轨迹中的残差分布，并随姿态、速度、加速度和温度等因素评估。若理想独立高斯样本超过双侧 $3\sigma$ 的概率约为 0.27\%，200 Hz 下每秒期望超阈样本数约为 $200\times0.0027=0.54$。这还不是误报事件率：连续样本通常相关，去抖逻辑也会把多个样本合并成一个事件。

因此要区分偏置、瞬间尖峰和持续接触。延长连续超阈要求可以减少某类误触发，却必然增加检测时间；抬高阈值会减少误报，也会漏掉较小接触。没有一组通用阈值能绕过模型和数据质量。

\section{计算练习与逐步解答}

\textbf{练习 1}\quad 教学关节的观测器希望极点为 $-10,-20$，求 $L$。

\textbf{解}\quad 目标多项式为 $s^2+30s+200$。由 $l_1+2=30$ 得 $l_1=28$；由 $2l_1+l_2=200$ 得 $l_2=144$。比算例14A的增益小，但观测误差衰减也更慢。

\textbf{练习 2}\quad 标量先验均值为 2，方差为 0.04；测量为 2.3，测量方差为 0.09。求本次更新结果。

\textbf{解}\quad $K=0.04/(0.04+0.09)=4/13=0.307692$。后验均值为 $2+(4/13)\times0.3=2.092308$；后验方差为 $(1-4/13)0.04=0.0276923$。测量方差大于先验方差，所以更新只向测量移动约31\%。

\textbf{练习 3}\quad $K_o=40\,\mathrm{s^{-1}}$，$1\,\mathrm{N\,m}$ 阶跃外力从零开始，残差何时达到 95\%？

\textbf{解}\quad $1-e^{-40t}=0.95$，故 $t=-\ln(0.05)/40=0.074893\,\mathrm{s}$，约75 ms。把它称作40 Hz会错，角频率参数40每秒对应截止频率仅 $40/(2\pi)=6.366\,\mathrm{Hz}$。

\textbf{练习 4}\quad 单关节 $M\ddot q=0.2$、$C\dot q=0.05$、$g=1.1$、$\tau_f=0.15$、$\tau_{\rm cmd}=1.3$，单位均为力矩。求环境对机器人的外力矩。

\textbf{解}\quad 由式\eqref{eq:cf_direct_external}，$\tau_{\rm ext}=0.2+0.05+1.1+0.15-1.3=0.2\,\mathrm{N\,m}$。若反过来用驱动力矩减模型力矩，得到的是相反符号；只有明确把结果定义为机器人对环境的反作用时才能这样使用。

\section*{本章与原稿的对应}
本章整合 [A，第16--18、43、50--51页] 与 [B，第56--58、74--78页]。统一了外力符号，补足矩阵卡尔曼计算与动量观测器推导，并区分状态、总扰动、真实接触力和可解释的等效末端力。

\chapter{从模型补偿到阻抗导纳和接触力}
\label{ch:impedance}

让机械臂到达一点与让机械臂被人推时表现出指定软硬程度，是不同的控制目标。前者主要关心运动误差，后者还要规定外力与运动之间的关系。本章先解释模型补偿能替反馈承担什么，再用质量弹簧阻尼模型逐项推导柔顺控制。所有增益示例均用于教学计算；真机参数还取决于实际接口、惯量、内环、延迟和执行器限制。

\section{重力补偿在做什么以及没有做什么}

全书动力学约定为
\begin{equation}
M\ddot q+C\dot q+g+\tau_f=\tau_{\rm cmd}+\tau_{\rm ext}.
\label{eq:cf_full_dyn}
\end{equation}
采用定点控制
\begin{equation}
\tau_{\rm cmd}=\hat g(q)-K_pe-K_d\dot q,
\qquad e=q-q_d.
\label{eq:cf_gpd}
\end{equation}
$\hat g$ 提供托住模型自重的力矩；$-K_pe$ 提供虚拟弹簧；$-K_d\dot q$ 提供速度相关阻尼。阻尼在静止时为零，因此“把 $K_p$ 设零、保留一点 $K_d$ 就能托住机械臂”不成立：静止托举依赖重力补偿或其他支撑力矩，不来自阻尼。

若模型准确、外力为零、摩擦先忽略，将控制律代入得到
\[
M\ddot q+C\dot q+K_d\dot q+K_pe=0.
\]
在静止状态，若重力存在残差 $\Delta g=\hat g-g$，则 $K_pe=\Delta g-\tau_f+\tau_{\rm ext}$。恢复误差的方向可以帮助判断残差方向，但不能仅从一条末端误差曲线唯一识别所有来源。

\begin{examplebox}[算例 15A 一点重力误差怎样成为可见位置偏差]
取单轴 $K_p=10\,\mathrm{N\,m/rad}$，真实重力为 $0.55\,\mathrm{N\,m}$，补偿只给 $0.50$。忽略摩擦和外力，静止时
\[
e_\infty=\frac{0.50-0.55}{10}=-0.005\,\mathrm{rad}=-0.2865^\circ.
\]
若把刚度减为 2，偏差变为 $-0.025\,\mathrm{rad}=-1.4324^\circ$。若先把重力残差减到 $0.01\,\mathrm{N\,m}$，则在刚度2下也能回到 $0.005\,\mathrm{rad}$ 的误差尺度。较软又较准的关键工作之一是降低未补偿力矩。
\end{examplebox}

\subsection{定点稳定性的完整能量抵消}

取 $V=\tfrac12\dot q^\top M\dot q+\tfrac12 e^\top K_pe$，其中固定目标满足 $\dot e=\dot q$。求导并代入闭环：
\begin{align}
\dot V
&=\dot q^\top M\ddot q+\tfrac12\dot q^\top\dot M\dot q
+e^\top K_p\dot q\\
&=-\dot q^\top K_d\dot q
+\tfrac12\dot q^\top(\dot M-2C)\dot q\\
&=-\dot q^\top K_d\dot q\le0.
\end{align}
第二项为零是因为 $\dot M-2C$ 反对称。若未补偿摩擦是耗散的，满足 $\dot q^\top\tau_f\ge0$，则还多出一个非正项 $-\dot q^\top\tau_f$。

这个证明依赖精确重力、固定目标、理想连续力矩实现、合适的状态域及增益等条件。若加入任意模型前馈、带延迟的摩擦抵消或大幅变化的参考，必须重新分析，不能说“前馈加 PD 一定稳定”。关节限位、力矩饱和和几何约束也没有自动包括在上述证明中。

\section{计算力矩怎样把期望加速度变成实际力矩}

定义速度误差 $\dot e=\dot q-\dot q_d$，希望误差满足
\[
\ddot e+K_v^a\dot e+K_p^a e=0.
\]
上标 $a$ 提醒这组增益输出的是加速度，单位分别为 $\mathrm{s^{-1}}$ 和 $\mathrm{s^{-2}}$。先构造期望实际加速度
\[
a=\ddot q_d-K_v^a\dot e-K_p^a e,
\]
再用逆动力学给出
\begin{equation}
\tau_{\rm cmd}=\hat M(q)a+\hat C(q,\dot q)\dot q
+\hat g(q)+\hat\tau_f.
\label{eq:cf_computed}
\end{equation}
模型完全一致、外力为零时，把式\eqref{eq:cf_computed}代入式\eqref{eq:cf_full_dyn}，得到 $M\ddot q=Ma$，即 $\ddot q=a$，所要求的误差方程成立。

这里用的是\textbf{参考加速度 $\ddot q_d$}，加上实测位置与速度形成的反馈，不要求把实测位置二次差分成加速度再放回控制律。辨识可能需要估计实测加速度，计算力矩控制本身并不因此必须测量它。

\begin{examplebox}[算例 15B 计算一次模型控制力矩]
单轴取 $M=0.1$、黏性项 $b=0.2$、$\hat g=0.5$，当前 $q-q_d=0.02\,\mathrm{rad}$、$\dot q=0.05\,\mathrm{rad/s}$，目标速度和加速度均为零。取 $K_p^a=100$、$K_v^a=14$：
\[
a=-100(0.02)-14(0.05)=-2.7\,\mathrm{rad/s^2},
\]
\[
\tau=0.1(-2.7)+0.2(0.05)+0.5=0.24\,\mathrm{N\,m}.
\]
若真实重力也是 0.5，代入真实对象得到 $\ddot q=-2.7$；若真实重力为0.55，则实际加速度为 $(0.24-0.01-0.55)/0.1=-3.2$，比期望多了 $-0.5\,\mathrm{rad/s^2}$。模型误差没有消失，它成为误差动力学的输入。
\end{examplebox}

沿期望轨迹提前算 $\tau_{\rm ff}=M(q_d)\ddot q_d+C(q_d,\dot q_d)\dot q_d+g(q_d)$，再叠加物理单位的关节 PD，是另一种常用结构。它与使用当前状态的完整反馈线性化不相同：跟踪误差大时，前馈模型与实际姿态不一致，不能声称精确解耦。

\section{阻抗规定受力后怎样偏离目标}

先用一个平移方向说明。设目标位置 $x_d$，实际位置 $x$，误差 $e_x=x-x_d$。目标阻抗为
\begin{equation}
m_d\ddot e_x+D\dot e_x+Ke_x=F_{\rm ext}.
\label{eq:cf_target_impedance}
\end{equation}
$m_d$ 是期望质量，单位 kg；$D$ 是阻尼，单位 $\mathrm{N\,s/m}$；$K$ 是刚度，单位 $\mathrm{N/m}$；$F_{\rm ext}$ 是环境施加给机器人的力。恢复弹簧力为 $-Ke_x$，与位移方向相反。

以速度为端口时，机械阻抗是 $Z_v(s)=F/V=m_ds+D+K/s$；以位移为输入定义时，则常写动态刚度 $F/X=m_ds^2+Ds+K$。文献有时都称为阻抗，读公式时必须看分母是速度还是位移，不能只看名称。两者相差一个 $s$。

\begin{examplebox}[算例 15C 从软硬和质量算阻尼]
取 $m_d=2\,\mathrm{kg}$、$K=200\,\mathrm{N/m}$，希望临界阻尼 $\zeta=1$。由
\[
\omega_n=\sqrt{K/m_d}=10\,\mathrm{rad/s},\qquad
D=2\zeta\sqrt{Km_d}=40\,\mathrm{N\,s/m}.
\]
持续受到 $5\,\mathrm{N}$ 外力后，最终误差为 $e_x=5/200=0.025\,\mathrm{m}$。若误差为5 mm，恢复力只有1 N；若运动速度为0.1 m/s，阻尼力已有4 N。低刚度不等于低阻尼，也不等于拖起来处处轻。
\end{examplebox}

\begin{figure}[htbp]
\centering\includegraphics[width=.95\linewidth]{figures/cf_impedance_bias.pdf}
\caption{刚度越低，同一外力造成的位移越大，模型残差造成的偏差也越大。左图对每个刚度均采用 $D=2\sqrt{2K}$，以保持临界阻尼。}
\end{figure}

若要求某方向在未知残余力不超过 $F_b$ 时稳态误差不超过 $e_b$，单方向静态关系要求 $K\ge F_b/e_b$。例如 $F_b=1\,\mathrm{N}$、$e_b=1\,\mathrm{mm}$，要求至少 $1000\,\mathrm{N/m}$。要在 $K=100$ 时达到同样误差，残余力就必须压到约 $0.1\,\mathrm{N}$。这条关系能把“又软又准”转化成可测的补偿精度目标。

\subsection{相对速度阻尼与绝对速度阻尼}

定点时 $\dot x_d=0$，$-D\dot x$ 与 $-D(\dot x-\dot x_d)$ 相同。移动目标时不相同。例如目标以 $0.1\,\mathrm{m/s}$ 匀速运动，$D=40$、$K=200$，若采用 $F=-K(x-x_d)-D\dot x$，稳定匀速跟踪时
\[
0=-Ke_x-D(0.1)\quad\Longrightarrow\quad e_x=-0.02\,\mathrm{m}.
\]
没有任何模型误差也会落后2 cm。使用相对速度阻尼，并正确处理期望加速度，才对应式\eqref{eq:cf_target_impedance}的移动参考定义。

\section{用雅可比转置把末端恢复力变成关节力矩}

在某一任务空间中，设 $\dot x=J\dot q$，且力与速度在同一坐标系、同一点表达。由功率相等
\[
F^\top\dot x=F^\top J\dot q=(J^\top F)^\top\dot q
\]
得到 $\tau=J^\top F$。简化笛卡尔弹簧阻尼控制为
\begin{equation}
\tau_{\rm cmd}=\hat g+\hat\tau_f+
J^\top[-K_xe_x-D_x\dot e_x]+\tau_{\rm null}.
\label{eq:cf_simple_impedance}
\end{equation}
这里先只把 $x$ 看作平移任务或局部适当坐标；完整三维姿态误差、角速度与几何雅可比的关系见下一章。$\tau_{\rm null}$ 如何构造也不能只凭“投影”两个字处理。

\begin{examplebox}[算例 15D 两连杆的笛卡尔弹簧力矩]
在教学位形 $q=(0,\pi/2)$，$J=\left[\begin{smallmatrix}-0.5&-0.5\\0.5&0\end{smallmatrix}\right]$。设平移误差为 $e_x=[0.01,-0.02]^\top\,\mathrm{m}$，速度误差为 $[0.03,-0.01]^\top\,\mathrm{m/s}$，
\[
K_x=\operatorname{diag}(200,100),\quad
D_x=\operatorname{diag}(20,10).
\]
弹簧力为 $[-2,2]^\top$，阻尼力为 $[-0.6,0.1]^\top$，合力为 $F=[-2.6,2.1]^\top\,\mathrm{N}$。关节反馈力矩为
\[
J^\top F=
\begin{bmatrix}-0.5&0.5\\-0.5&0\end{bmatrix}
\begin{bmatrix}-2.6\\2.1\end{bmatrix}
=\begin{bmatrix}2.35\\1.30\end{bmatrix}\,\mathrm{N\,m}.
\]
还需叠加重力等补偿项。虚拟弹簧势能为 $\tfrac12e_x^\top K_xe_x=0.03\,\mathrm{J}$，可用作单位和符号检查。
\end{examplebox}

在无预载、小位移和固定姿态附近，$\delta x=J\delta q$，所以关节等效刚度近似为 $K_q=J^\top K_xJ$。它一般不是对角阵。关节控制器只给每轴独立的 $k_p,k_d$，通常无法任意指定末端某个空间方向的软硬。若已有接触预载，雅可比随姿态改变还会产生几何刚度项，不能把这条局部关系当成全局精确公式。

\subsection{只写弹簧阻尼不代表已经指定了虚拟质量}

式\eqref{eq:cf_simple_impedance}常用于低速柔顺，但它保留了机器人本来的惯性、耦合、驱动器动态和未补偿项。$K_x,D_x$ 的输入并不会自动把末端质量改成某个 $m_d$。如果不同方向或不同姿态的实际质量不同，同一组刚度阻尼会产生不同的自然频率和阻尼比。

对一个解耦近似方向，用有效质量 $m_{\rm eff}$ 估计 $D=2\zeta\sqrt{Km_{\rm eff}}$ 是合理起点。若任务惯量强耦合，仅把其对角元逐个套公式不保证整体模态都有指定阻尼，需考察质量与刚度组合的特征方向。

\section{惯量重塑为什么需要额外信息}

先取理想一维对象 $m\ddot x=u+F_{\rm ext}$。希望得到式\eqref{eq:cf_target_impedance}，由目标方程解出
\[
\ddot x=\ddot x_d+\frac{F_{\rm ext}-D\dot e_x-Ke_x}{m_d}.
\]
代回真实动力学所需的 $u=m\ddot x-F_{\rm ext}$，得到
\begin{equation}
u=m\ddot x_d-\frac m{m_d}(D\dot e_x+Ke_x)
+\left(\frac m{m_d}-1\right)F_{\rm ext}.
\label{eq:cf_mass_shaping}
\end{equation}
当 $m_d=m$，最后一项为零，无需测外力就能保留自然惯量。想改变惯量时，控制律通常需要外力测量或估计，或者等价的加速度反馈等附加实现条件。外力噪声和延迟会经过这条通道影响稳定性。

\begin{examplebox}[算例 15E 把2 kg的对象表现成1 kg]
取真实 $m=2$、目标 $m_d=1$、$K=100$、$D=20$，当前 $e_x=0.01\,\mathrm{m}$、$\dot e_x=0$、$\ddot x_d=0$、$F_{\rm ext}=5\,\mathrm{N}$。式\eqref{eq:cf_mass_shaping}给出
\[
u=-2\times100\times0.01+(2-1)\times5=3\,\mathrm{N}.
\]
真实加速度为 $(3+5)/2=4\,\mathrm{m/s^2}$；目标方程给出 $(5-100\times0.01)/1=4$，二者相同。控制器额外沿外力方向出力，使人感觉它更轻。这也说明质量重塑可能主动注入能量，不能仅用“弹簧刚度为正”保证整个实现无源。
\end{examplebox}

\section{导纳把外力变成一个可以执行的位置目标}

当机器人已有较强位置内环，外层可以运行一个虚拟模型：
\begin{equation}
M_d\ddot x_c+D_d\dot x_c+K_d(x_c-x_d)=F_{\rm ext}.
\end{equation}
这里 $x_c$ 是生成给位置控制器的目标，不是真实末端位置。只有内环能足够准确、足够及时地跟随 $x_c$，实际机器人才能接近这个导纳。外层目标、内环动态、力滤波和环境共同组成闭环，任何一项都不能省略。

对固定 $x_d$，半隐式欧拉更新为
\begin{align}
a_k&=M_d^{-1}[F_k-D_dv_k-K_d(x_k-x_d)],\\
v_{k+1}&=v_k+T_sa_k,\\
x_{k+1}&=x_k+T_sv_{k+1}.
\end{align}
这只是一个离散实现选项，稳定性受虚拟参数、周期和真实位置环影响。若 $x_d$ 本身运动，应明确是直接使用绝对指令动力学，还是对相对偏移状态积分并叠加参考运动。

\begin{examplebox}[算例 15F 手算两拍零刚度导纳]
取 $M_d=2\,\mathrm{kg}$、$D_d=40\,\mathrm{N\,s/m}$、$K_d=0$、$T_s=0.005\,\mathrm{s}$，初始 $x_0=v_0=0$，持续推力为4 N。

第一拍 $a_0=4/2=2$，故 $v_1=0.01\,\mathrm{m/s}$、$x_1=0.00005\,\mathrm{m}$。第二拍 $a_1=(4-40\times0.01)/2=1.8$，故 $v_2=0.019$、$x_2=0.000145\,\mathrm{m}$。

连续理想稳态速度为 $F/D=0.1\,\mathrm{m/s}$。零刚度时持续施力会持续移动，而不是最终回到原位。若传感器存在0.5 N未补偿偏置，即使没人推也会趋于 $0.5/40=0.0125\,\mathrm{m/s}$ 的漂移速度；加入 $K=100$ 后，同一静态偏置转为5 mm的位置偏差。
\end{examplebox}

\begin{figure}[htbp]
\centering\includegraphics[width=.95\linewidth]{figures/cf_admittance.pdf}
\caption{零刚度导纳在持续外力下趋于恒定速度；加上位置弹簧后才会趋于有限偏移。真实位置内环和力传感器动态尚未包括在这张理想模型图中。}
\end{figure}

阻抗与导纳主要区别在因果实现：前者根据运动产生力矩，后者根据外力产生运动指令。不能给出“导纳在硬环境一定更稳定”或“阻抗碰硬物就一定稳定”的通用结论。驱动接口、回驱摩擦、位置内环、力测量和通信时延共同决定可实现范围。

\section{静态接触力和撞击峰值必须分开计算}

设真实墙面在 $x=0$，进入墙内的正位移为 $x>0$，环境弹簧给机器人的力为 $F_{\rm ext}=-k_ex$。虚拟目标位于墙内 $x_d>0$。静止时
\[
0=-K(x-x_d)-k_ex,
\]
因此
\begin{equation}
x=\frac{K}{K+k_e}x_d,
\qquad F_{\rm contact}=k_ex=\frac{Kk_e}{K+k_e}x_d.
\end{equation}
$F_{\rm contact}>0$ 在这里表示压向墙的力大小，与环境对机器人的有符号力 $F_{\rm ext}$ 相反。

\begin{examplebox}[算例 15G 两个弹簧串联的接触力]
取 $K=200\,\mathrm{N/m}$、$k_e=20000\,\mathrm{N/m}$、$x_d=0.01\,\mathrm{m}$，则
\[
x=\frac{200}{20200}\times0.01=0.00009901\,\mathrm{m},
\]
\[
F_{\rm contact}=20000x=1.9802\,\mathrm{N}.
\]
机器人实际只把墙压入约0.099 mm，大部分目标偏移被虚拟弹簧吸收。若直接假设理想刚性位置源把真实墙压入全部10 mm，得到的会是200 N，这对应完全不同的控制和结构假设。
\end{examplebox}

冲击峰值还取决于接触瞬间动能。对无阻尼的单质量弹性撞击模型，初速度 $v_0$ 的动能全部存入环境弹簧时，
\[
\tfrac12mv_0^2=\tfrac12k_e\delta_{\max}^2,
\quad F_{\max}=k_e\delta_{\max}=v_0\sqrt{mk_e}.
\]
取 $m=2\,\mathrm{kg}$、$k_e=20000\,\mathrm{N/m}$、$v_0=0.1\,\mathrm{m/s}$，有 $F_{\max}=20\,\mathrm{N}$。这不是通用碰撞上界，而是展示即使稳态目标力很小，已有动能仍可造成较大的瞬时接触力。阻尼、接触材料、执行器动作和结构柔性都会改变实际峰值。

\begin{figure}[htbp]
\centering\includegraphics[width=.95\linewidth]{figures/cf_contact.pdf}
\caption{左图是无阻尼弹性碰撞的第一接触段，右图是两个弹簧的静态平衡。速度影响冲击，刚度与目标偏移影响最终压力，两类指标不能互相代替。}
\end{figure}

\section{直接力控制与力位混合}

如果已知接触方向，令 $f=k_ex$ 表示正向压缩力大小，并用
\[
\dot x_{\rm cmd}=k_f(f_d-f)
\]
生成位置指令。在理想位置跟随和始终接触的条件下，
\[
\dot f=k_e\dot x=k_ek_f(f_d-f),
\]
即时间常数为 $1/(k_ek_f)$ 的一阶力跟踪。$k_f$ 的单位是 $\mathrm{m/(N\,s)}$。环境越硬，相同增益下外力闭环越快，可能超出真实位置环和通信所能支撑的范围。

\begin{examplebox}[算例 15H 环境刚度改变力环速度]
取 $k_e=20000\,\mathrm{N/m}$、$k_f=0.0005\,\mathrm{m/(N\,s)}$，则 $k_ek_f=10\,\mathrm{s^{-1}}$，时间常数0.1 s。对5 ms前向欧拉更新，误差因子为 $1-k_ek_fT_s=0.95$。

若环境刚度扩大十倍，理想时间常数缩为0.01 s，离散误差因子变成0.5。虽然这个简化离散模型仍稳定，真实位置环若没有足够带宽，就不能实现该结果。只检查 $0<k_ek_fT_s<2$ 不足以判断真实接触稳定。
\end{examplebox}

力位混合控制在任务坐标系中分配方向。例如沿平面切向 $x,y$ 跟踪位置，沿法向 $z$ 跟踪压力，可用 $S_p=\operatorname{diag}(1,1,0)$、$S_f=I-S_p$ 表示分工。$S_pF_p+S_fF_f$ 是各通道输出在共同力坐标中的组合，不是随意把位置量和力量相加。表面法向变化时，选择矩阵必须随坐标变换；完整动态解耦还与操作空间惯量和约束几何有关。

某些力控制结构加积分可以消除常值偏差，但“直接力控必须含积分”不是普适定理：理想直接力矩源加正确前馈也能产生目标力。是否需要积分要由对象、传感器偏置和控制结构决定。积分进入接触过程后，同样需要限幅、抗饱和和切换时状态一致性。

\section{无源性和变刚度的能量账}

对固定目标、常数质量和刚度，目标阻抗的储能为
\[
E=\tfrac12\dot e_x^\top M_d\dot e_x+
\tfrac12e_x^\top K_de_x.
\]
代入目标动力学得
\[
\dot E=F_{\rm ext}^\top\dot e_x-
\dot e_x^\top D_d\dot e_x.
\]
外界输入的机械功率一部分存储、一部分被阻尼耗散，没有凭空产生的能量，这就是该端口模型的无源性直觉。与另一个无源系统互联时，可以利用储能关系讨论稳定与有界性；若还要得到渐近收敛，需要严格耗散或可检测性等额外条件。不能把“两个无源系统连接”直接简化为一切运动都会迅速停止。

若刚度随时间变化，即使目标固定，还多出
\[
\tfrac12e_x^\top\dot K_de_x
\]
这一项。误差尚未消失时突然提高刚度，会增加虚拟势能。

\begin{examplebox}[算例 15I 刚度变化需要多少能量]
一维误差为2 cm，刚度从100提高到400 N/m，在位置未改变的瞬间，势能增加
\[
\Delta E=\tfrac12(0.02)^2(400-100)=0.06\,\mathrm{J}.
\]
如果控制设计只允许用一个剩余可用能量0.04 J的虚拟能量罐支付这次变化，那么该变化不能原样执行；应减小增量、等待阻尼积累能量或采用明确的其他能量来源。仅限制 $\dot K$ 可以让变化较平滑，却不自动证明总能量条件满足。
\end{examplebox}

数字采样、延迟、力矩估计误差、摩擦过补偿及零空间耦合都会改变实际能量收支。某些单自由度触觉装置能推导出采样刚度和物理阻尼的标量条件，但它们依赖特定采样、速度估计和控制结构，不能把某个标量公式当成所有七轴机械臂的通用刚度上限。

\section{计算练习与逐步解答}

\textbf{练习 1}\quad 目标有效质量3 kg，刚度300 N/m，希望阻尼比0.8。求阻尼、自然频率和6 N恒力下的偏移。

\textbf{解}\quad $\omega_n=\sqrt{300/3}=10\,\mathrm{rad/s}$；$D=2\times0.8\sqrt{300\times3}=48\,\mathrm{N\,s/m}$；静态偏移 $6/300=0.02\,\mathrm{m}$。阻尼改变动态过程，不改变此静态偏移。

\textbf{练习 2}\quad 零刚度导纳 $M=1.5$、$D=30$，受到3 N恒力。求最终速度和速度时间常数。

\textbf{解}\quad 方程为 $1.5\dot v+30v=3$。稳态令 $\dot v=0$，得到 $v_\infty=0.1\,\mathrm{m/s}$。时间常数为 $M/D=0.05\,\mathrm{s}$。从静止开始：
\[
v(0.05)=0.1(1-e^{-1})=0.06321\,\mathrm{m/s}.
\]

\textbf{练习 3}\quad 实际质量4 kg，希望表现为2 kg；静止目标，$K=100$、$D=20$，当前误差0.02 m、速度误差0.1 m/s、外力8 N。求质量重塑控制力。

\textbf{解}\quad 恢复项 $D\dot e+Ke=20\times0.1+100\times0.02=4\,\mathrm{N}$。由式\eqref{eq:cf_mass_shaping}，$u=-(4/2)4+(4/2-1)8=0$。真实加速度为 $(0+8)/4=2$；目标方程给 $(8-4)/2=2$，一致。

\textbf{练习 4}\quad 环境刚度10000 N/m、控制刚度500 N/m，目标在表面内5 mm。求稳态压力和真实压入量。

\textbf{解}\quad 真实压入量为 $x=500\times0.005/10500=0.000238095\,\mathrm{m}$；压力 $f=10000x=2.38095\,\mathrm{N}$。也可用串联等效刚度 $500\times10000/10500=476.1905\,\mathrm{N/m}$ 乘5 mm得到相同结果。

\section*{本章与原稿的对应}
本章整合 [A，第39--52页] 与 [B，第54--58、83--87页]。补足了静态偏差、移动参考、质量重塑、导纳两拍更新、接触冲击和能量变化的计算，并对阻抗与导纳的适用边界、实测加速度需求和稳定性前提作了明确限定。

\chapter{操作空间控制与七轴主任务和零空间}
\label{ch:operational_space}

七个关节不等于末端有七个独立空间运动方向。完整末端位姿通常包含三维位置和三维姿态，共六维；在雅可比满行秩时，七轴机械臂还剩一维瞬时自运动。理解这条维数关系，才能正确讨论“末端柔顺”和“肘部柔顺”。本章从关节动力学推导末端动力学，再用可手算反例区分运动学零空间与力矩零空间。

\section{从关节加速度推到任务加速度}

设所选任务有 $m$ 个独立速度分量，关节有 $n$ 维，$J\in\mathbb R^{m\times n}$。先写
\begin{equation}
v_x=J(q)\dot q,
\qquad a_x=\dot v_x=J\ddot q+\dot J\dot q.
\end{equation}
对于平移任务，$v_x=\dot x$、$a_x=\ddot x$。对于完整几何任务，本书用 $v_x=[v^\top,\omega^\top]^\top$，其中 $\omega$ 是角速度，不能把它直接当成欧拉角导数。

令 $h=C\dot q+g+\tau_f$，并假设外部作用只来自已知任务端口，$\tau_{\rm ext}=J^\top F_{\rm ext}$。由动力学得到
\begin{align}
\ddot q&=M^{-1}(\tau_{\rm cmd}-h+J^\top F_{\rm ext}),\\
a_x&=JM^{-1}\tau_{\rm cmd}-JM^{-1}h
+\dot J\dot q+JM^{-1}J^\top F_{\rm ext}.
\label{eq:cf_task_acc}
\end{align}
定义
\begin{equation}
A_x=JM^{-1}J^\top,\qquad
\Lambda=A_x^{-1}.
\label{eq:cf_lambda}
\end{equation}
当 $M\succ0$ 且 $J$ 满行秩时，$A_x\succ0$，因此逆存在。$\Lambda$ 称为操作空间惯量矩阵。平移块的单位类似质量；旋转块类似转动惯量；混合块含有不同单位，因此不能不做尺度处理就比较全部元素大小。

若取 $\tau_{\rm cmd}=J^\top F+\tau_{\rm null}$，并且 $JM^{-1}\tau_{\rm null}=0$，将式\eqref{eq:cf_task_acc}左乘 $\Lambda$，得到
\begin{equation}
\Lambda a_x+\eta=F+F_{\rm ext},
\qquad \eta=\Lambda JM^{-1}h-\Lambda\dot J\dot q.
\end{equation}
若把 $h$ 分开，可写成科氏速度项 $\mu$、重力项 $p_g$ 和摩擦项：
\begin{align}
\mu&=\Lambda(JM^{-1}C\dot q-\dot J\dot q),\\
p_g&=\Lambda JM^{-1}g,\qquad
p_f=\Lambda JM^{-1}\tau_f.
\end{align}
因此期望任务加速度 $a_d$ 对应控制力
\[
F=\Lambda a_d+\mu+p_g+p_f.
\]
$\dot J\dot q$ 表示雅可比本身变化造成的速度映射变化，快速运动时不能随意删掉。操作空间方法的历史来源可参阅 Khatib 的原始工作\cite{khatib}；这里的符号与数值计算均按本书约定重新展开。

\section{完整算出两连杆的任务惯量}

取统一教学两连杆：$l_1=l_2=0.5\,\mathrm{m}$，$m_1=m_2=1\,\mathrm{kg}$，质心均在连杆中点，质心转动惯量均为 $1/48\,\mathrm{kg\,m^2}$。在 $q_1=0,q_2=\pi/2$、速度为零时，第6章动力学给出
\[
M=\begin{bmatrix}5/12&1/12\\1/12&1/12\end{bmatrix},\quad
J=\begin{bmatrix}-1/2&-1/2\\1/2&0\end{bmatrix},\quad
 g=\begin{bmatrix}7.3575\\0\end{bmatrix}.
\]
先求 $M$ 的逆。行列式为 $(5-1)/144=1/36$，因此
\[
M^{-1}=\begin{bmatrix}3&-3\\-3&15\end{bmatrix}.
\]
再乘雅可比：
\[
JM^{-1}=\begin{bmatrix}0&-6\\1.5&-1.5\end{bmatrix},\qquad
JM^{-1}J^\top=\begin{bmatrix}3&0\\0&0.75\end{bmatrix}.
\]
因此
\begin{equation}
\Lambda=\begin{bmatrix}1/3&0\\0&4/3\end{bmatrix}\,\mathrm{kg}.
\end{equation}
同一位形在两个平移方向的有效质量相差四倍；不能因两根连杆总质量是2 kg，就把所有方向的操作空间质量都写成2 kg。

\begin{examplebox}[算例 16A 从期望末端加速度算到关节力矩再验算]
速度为零，故 $\mu=0$、$\dot J\dot q=0$，先忽略摩擦。重力任务力为
\[
p_g=\Lambda JM^{-1}g=
\begin{bmatrix}0\\14.715\end{bmatrix}\,\mathrm{N}.
\]
要求 $a_d=[1,2]^\top\,\mathrm{m/s^2}$，得到
\[
F=\Lambda a_d+p_g=
\begin{bmatrix}0.333333\\17.381667\end{bmatrix}\,\mathrm{N}.
\]
映射到关节：
\[
\tau=J^\top F=
\begin{bmatrix}8.524167\\-0.166667\end{bmatrix}\,\mathrm{N\,m}.
\]
最后不信任自己的推导，直接回代真实关节方程：
\[
\ddot q=M^{-1}(\tau-g)
=\begin{bmatrix}3&-3\\-3&15\end{bmatrix}
\begin{bmatrix}1.166667\\-0.166667\end{bmatrix}
=\begin{bmatrix}4\\-6\end{bmatrix}.
\]
于是 $J\ddot q=[1,2]^\top$，确实得到所需末端加速度。这种“先正向设计，再回代原方程”的检查比只看公式外观可靠。
\end{examplebox}

接近奇异位形时，$J$ 某个奇异值变小，对应可实现运动能力下降，$\Lambda$ 在相关方向可能变得很大或不可定义。奇异不只是数值求逆麻烦，还表示机构本身暂时缺少某个独立运动方向；加正则化可以缓解数值爆炸，但不能创造实际不存在的自由度。

\section{运动学零空间和力矩零空间不是同一个投影}

\subsection{速度投影需要满足什么}

动力学一致广义逆定义为
\begin{equation}
\bar J=M^{-1}J^\top\Lambda.
\end{equation}
在满行秩条件下，$J\bar J=I_m$。定义速度零空间投影
\begin{equation}
N=I_n-\bar J J.
\end{equation}
于是
\[
JN=J-J\bar JJ=0,\qquad N^2=N.
\]
$N\dot q_0$ 是不会瞬时改变主任务速度的关节速度分量。$N$ 一般不对称；它是在质量度量下构造的投影，而不是欧氏正交投影。

\subsection{力矩投影需要满足什么}

力矩首先经过 $M^{-1}$ 变成加速度，所以次任务不能影响主任务加速度的条件是
\[
JM^{-1}\tau_{\rm null}=0,
\]
不是简单的 $J\tau_{\rm null}=0$。采用
\begin{equation}
\tau_{\rm null}=N^\top\tau_0
=(I-J^\top\bar J^\top)\tau_0
\label{eq:cf_null_torque}
\end{equation}
即可验证：
\begin{align}
JM^{-1}N^\top
&=JM^{-1}-JM^{-1}J^\top\bar J^\top\\
&=JM^{-1}-\Lambda^{-1}\Lambda JM^{-1}=0.
\end{align}
上式用到了 $M,\Lambda$ 对称，因此 $\bar J^\top=\Lambda JM^{-1}$。这就是为什么速度形式用 $N$，力矩形式用 $N^\top$。如果某篇资料把力矩投影本身命名为 $N$，它也可以写 $N\tau_0$；关键是定义，不能只比较字母。

\begin{examplebox}[算例 16B 普通投影为什么会在主任务上漏出加速度]
用两个独立滑块构造最小冗余例子，广义坐标单位为米，质量矩阵 $M=\operatorname{diag}(1,4)\,\mathrm{kg}$，主任务为 $x=q_1+q_2$，故 $J=[1,1]$。此例不是两转动关节模型，专门用于看清投影。

普通伪逆为 $J^+=[0.5,0.5]^\top$，对应欧氏投影
\[
N_k=I-J^+J=\begin{bmatrix}0.5&-0.5\\-0.5&0.5\end{bmatrix}.
\]
给原始次任务力 $\tau_0=[1,0]^\top\,\mathrm{N}$，投影后为 $[0.5,-0.5]^\top$。若只检查两力之和，会看到零；但真实加速度为
\[
M^{-1}N_k^\top\tau_0=
\begin{bmatrix}0.5\\-0.125\end{bmatrix},
\]
所以主任务加速度为 $0.5-0.125=0.375\,\mathrm{m/s^2}$，并不为零。

动力学一致计算则给
\[
\Lambda=(1+1/4)^{-1}=0.8,
\quad\bar J=\begin{bmatrix}0.8\\0.2\end{bmatrix},
\quad N=\begin{bmatrix}0.2&-0.8\\-0.2&0.8\end{bmatrix}.
\]
正确力矩投影为
\[
N^\top\tau_0=\begin{bmatrix}0.2\\-0.8\end{bmatrix}.
\]
真实加速度为 $[0.2,-0.2]^\top$，主任务加速度恰好为零。区别来自质量不同，而不是矩阵符号的排版习惯。
\end{examplebox}

\begin{figure}[htbp]
\centering\includegraphics[width=.95\linewidth]{figures/cf_nullspace.pdf}
\caption{相同原始次任务力经过不同投影，主任务加速度不同。右图显示加速度方向；零主任务要求两坐标加速度之和为零。}
\end{figure}

\section{为什么只做任务重力补偿可能仍让肘部动}

任务重力项 $p_g=\Lambda JM^{-1}g$ 映射回关节后是 $J^\top p_g$。对非冗余可逆映射，这可能覆盖全部 $g$；对冗余系统，它通常只能提供主任务可见的那一部分。利用式\eqref{eq:cf_null_torque}的定义，恒等式为
\begin{equation}
g=J^\top\Lambda JM^{-1}g+N^\top g
=J^\top p_g+N^\top g.
\end{equation}
因此只用 $J^\top p_g$，可能使末端瞬时不加速，而零空间构型仍受未补偿重力推动。

\begin{examplebox}[算例 16C 末端没动但冗余坐标在漂]
沿用两个滑块，设广义重力向量 $g=[1,0]^\top\,\mathrm{N}$。任务重力为 $p_g=0.8$，任务补偿给关节的力为 $J^\top p_g=[0.8,0.8]^\top$。

真实净力是 $[0.8,0.8]^\top-[1,0]^\top=[-0.2,0.8]^\top$，加速度为 $[-0.2,0.2]^\top$。主任务 $x=q_1+q_2$ 的加速度为零，两个内部坐标却在反向移动。

缺少的部分恰好是 $N^\top g=[0.2,-0.8]^\top$，把它加回去才得到完整 $g=[1,0]^\top$。这说明“只先补笛卡尔重力”需要明确目标：只保持选定任务，还是同时托住所有冗余自由度。
\end{examplebox}

一种容易避免重复补偿的结构是先在关节层补全部名义 $h$：
\begin{equation}
\tau_{\rm cmd}=\hat h+
J^\top\Lambda(a_d-\dot J\dot q)+N^\top\tau_{\rm post}.
\label{eq:cf_full_joint_comp}
\end{equation}
在精确模型条件下，它实现 $a_x=a_d+A_xF_{\rm ext}$，次任务瞬时不进入主任务加速度。另一种结构是在任务力中放入 $\mu+p_g+p_f$，同时用零空间项补齐剩余 $N^\top h$。两种账法可以等价，但不能既加完整 $g$，又在任务力中再次加完整投影重力而不做扣除。

低速笛卡尔弹簧常采用
\[
\tau_{\rm cmd}=\hat g+\hat\tau_f+
J^\top(-K_xe_x-D_x\tilde v_x)+N^\top\tau_{\rm post}.
\]
它与式\eqref{eq:cf_full_joint_comp}不是同一个精确加速度控制器：前者主要指定恢复力和阻尼，保留实际惯量；后者显式使用模型控制任务加速度。选择哪一条，应由目标手感、模型可信度和实现条件决定。

\section{阻尼伪逆会牺牲严格解耦}

为缓解近奇异问题，可定义正则化广义逆
\[
\bar J_\lambda=M^{-1}J^\top(A_x+\lambda^2I)^{-1},
\qquad N_\lambda=I-\bar J_\lambda J.
\]
在位置与姿态混合任务中，正则项必须按统一的任务尺度或度量定义，$\lambda^2I$ 不能无视各行单位直接照抄。

由代数展开：
\begin{align}
JM^{-1}N_\lambda^\top
&=[I-A_x(A_x+\lambda^2I)^{-1}]JM^{-1}\\
&=\lambda^2(A_x+\lambda^2I)^{-1}JM^{-1}.
\end{align}
$\lambda\ne0$ 时右边一般不为零。因此正则化使逆解更平滑，同时放弃严格零空间性质。也通常有 $N_\lambda^2\ne N_\lambda$，它不再是严格幂等投影。

\begin{examplebox}[算例 16D 算出正则化带来的主任务泄漏]
沿用滑块例子，$A_x=1.25$，取 $\lambda^2=0.25$。则
\[
\bar J_\lambda=\frac1{1.5}\begin{bmatrix}1\\0.25\end{bmatrix}
=\begin{bmatrix}2/3\\1/6\end{bmatrix}.
\]
对 $\tau_0=[1,0]^\top$，
\[
N_\lambda^\top\tau_0=\begin{bmatrix}1/3\\-2/3\end{bmatrix},
\quad JM^{-1}N_\lambda^\top\tau_0=\frac13-\frac16=\frac16.
\]
额外任务加速度为 $0.166667\,\mathrm{m/s^2}$。调大阻尼能减少奇异放大，但不能同时声称末端完全不受次任务影响。
\end{examplebox}

\section{七轴末端与肘部柔顺能独立到什么程度}

对 $J\in\mathbb R^{6\times7}$，若秩为6，零空间维数为 $7-6=1$。完整末端位姿已经占用了六个独立瞬时方向，肘部可以沿剩下的一维自运动方向调整，不能在保持末端完整位姿的同时再任意指定肘部三维位置。

常见姿态项为
\[
\tau_{\rm post}=-K_n(q-q_0)-D_n\dot q,
\qquad \tau_{\rm null}=N^\top\tau_{\rm post}.
\]
$K_n$ 小表示对参考构型的恢复力较弱，$D_n$ 控制自运动中的速度阻尼。$K_n=0$ 并不违法，它可以服务于自由拖动；但这时构型可能没有恢复目标。是否需要位置型零空间项应由任务决定，而不是“七轴一定必须加弹簧”。

需要同时明确四个限制：第一，投影保证的是当前模型和当前位形下的瞬时关系；第二，$N(q)$ 随构型变化；第三，饱和后实际力矩不再等于投影结果；第四，普通投影姿态弹簧并不自动继承一个全局无源性证明。完整稳定性需要连同主任务、补偿、时变投影和约束一起分析。

例如低阻抗主任务与低阻抗零空间同时启用时，主要托举必须由重力补偿承担，不能指望两个很软的弹簧纠正很大的模型残差。若机械摩擦和死区较大，投影结构正确也可能出现“肘部推不动”；这不是零空间公式自动失效，而是可用外力尚未克服执行器非理想项。

\section{姿态误差必须与角速度处在同一坐标约定}

设 $R,R_d$ 都把工具坐标分量转到基座坐标。定义基座表达的相对旋转
\[
E=RR_d^\top,\qquad e_R=\operatorname{Log}(E)^\vee.
\]
$e_R$ 是当前姿态相对目标的旋转向量，单位为弧度，其正负与本书“实际减目标”的误差方向一致。$\vee$ 把反对称矩阵取回三维向量。若使用 $\operatorname{Log}(R_d^\top R)^\vee$，结果在目标坐标表达，必须再转换后才能与基座雅可比配合。

\subsection{旋转向量导数不总等于角速度差}

基座角速度满足 $\dot R=[\omega]_\times R$、$\dot R_d=[\omega_d]_\times R_d$，于是
\begin{align}
\dot E&=[\omega]_\times E-E[\omega_d]_\times\\
&=[\omega-E\omega_d]_\times E.
\end{align}
令 $J_l(e_R)$ 为 SO(3) 左雅可比，即把旋转向量导数映射到相对旋转的空间角速度的矩阵，则
\begin{equation}
\dot e_R=J_l(e_R)^{-1}(\omega-E\omega_d).
\end{equation}
小误差下 $E\approx I$、$J_l^{-1}\approx I$，才有 $\dot e_R\approx\omega-\omega_d$。因此把姿态三维误差和角速度误差直接拼进欧氏六维 PD，是一种常用局部设计，需要知道近似范围。

对固定目标和各向同性旋转刚度 $k_RI$，$-k_Re_R$ 与距离势能 $\tfrac12k_R\|e_R\|^2$ 相容；对各向异性 $K_R$，若严格采用势能 $\tfrac12e_R^\top K_Re_R$，基座力矩梯度为
\[
m_R=-J_l(e_R)^{-\top}K_Re_R,
\]
而 $-K_Re_R$ 是小误差近似。接近 $\pi$ 的旋转还涉及对数分支与轴方向不唯一，不能声称用一个光滑公式能在所有姿态上无条件全局收敛。

\begin{examplebox}[算例 16E 十度误差不能把数字10直接送进弧度刚度]
目标为 $R_d=I$，实际为绕基座 $z$ 轴正转 $10^\circ$，故 $e_R=[0,0,\pi/18]^\top$。取 $k_R=10\,\mathrm{N\,m/rad}$，恢复力矩为
\[
m_z=-10\frac{\pi}{18}=-1.74533\,\mathrm{N\,m}.
\]
如果把10度中的数字10直接代入，得到 $-100\,\mathrm{N\,m}$，相差约57.3倍。单位错误比高级算法选择更容易彻底破坏控制。
\end{examplebox}

\section{工具坐标的软硬怎样转换到基座}

在同一参考点，仅改变坐标表达时，平移力和位移都按旋转变换。若工具系刚度为 $K_t$，基座系刚度应为
\[
K_b=R_{bt}K_tR_{bt}^\top.
\]
\begin{examplebox}[算例 16F 工具转45度后刚度矩阵出现耦合项]
二维工具刚度为 $K_t=\operatorname{diag}(500,50)\,\mathrm{N/m}$，工具相对基座转 $45^\circ$：
\[
R=\frac1{\sqrt2}\begin{bmatrix}1&-1\\1&1\end{bmatrix}.
\]
逐项相乘得到
\[
K_b=\begin{bmatrix}275&225\\225&275\end{bmatrix}\,\mathrm{N/m}.
\]
若基座位移为 $[0.01,0]^\top$，恢复力为 $[-2.75,-2.25]^\top\,\mathrm{N}$。出现 $y$ 向恢复力是旋转后的各向异性弹簧应有行为，不一定是“坐标耦合错误”。如果希望基座 $x,y$ 互不耦合，就应该直接在基座定义对角刚度。
\end{examplebox}

\begin{figure}[htbp]
\centering\includegraphics[width=.95\linewidth]{figures/cf_stiffness_frame.pdf}
\caption{左图为旋转后的各向异性弹簧等能量椭圆，窄轴方向更硬；右图展示姿态角度显示与内部弧度运算的关系。}
\end{figure}

如果改变力旋量的参考点，而不只是旋转坐标，还必须加入力臂引起的力矩。例如力 $f$ 在距新参考点 $r$ 的位置作用，力矩项会增加 $r\times f$。相应速度变换也包含旋转引起的平移速度，二者通过功率不变联系。不能只旋转六维向量的前三项然后保留后三项不变。

\section{把七轴回正问题分成可以测量的几件事}

首先记录关节目标、实际角度、速度、发送力矩、实际力矩或电流、限幅状态、末端目标和模型末端位姿，并统一时间戳。若末端位姿来自同一套 FK，它只能说明控制器与自身模型的一致性；要判断绝对几何精度，还需外部测量。

对静止或足够慢的回正过程，可把关节残差近似写成
\[
K_pe\approx\hat g-g-\tau_f+\tau_{\rm ext}
\]
或在任务空间解释恢复力 $-K_xe_x$。对于七轴，零空间弹簧与投影也参与平衡，不能只用 $J^\top K_xe_x$ 就唯一反推出所有模型误差。

从相反两个方向接近同一目标，若停止偏差有明显方向相关，摩擦、回差或不同冗余构型是候选原因；若换姿态后偏差按重力方向改变，质量质心和安装姿态是候选原因；若只有某些位形出现严重响应变弱，应检查雅可比秩与任务尺度；若增益提高只让命令更常限幅，应先检查可实现力矩与轨迹要求。

这些现象是排查证据，不是单凭曲线就能确诊。例如“左右方向主要由第一轴完成”只在具体机构、坐标和位形下可能近似成立，不能未读 URDF、未计算雅可比就断定 NERO 的所有方向响应原因。

区分两层误差也很重要。若 $q_{\rm meas}$ 明显不同于 $q_d$，主要是关节跟踪或接口执行问题；若 $q_{\rm meas}\approx q_d$，但外部测得末端与 FK 不符，则应检查几何标定、零位、工具、底座和柔性。冗余臂比较时优先统一用 FK，不要把外部位姿做一次不受分支约束的 IK 后强行比较关节角。

\section{计算练习与逐步解答}

\textbf{练习 1}\quad 对滑块模型 $J=[1,1]$、$M=\operatorname{diag}(1,4)$，给次任务力 $\tau_0=[0,1]^\top$。计算动力学一致零空间力与加速度。

\textbf{解}\quad $N^\top=\left[\begin{smallmatrix}0.2&-0.2\\-0.8&0.8\end{smallmatrix}\right]$，所以 $\tau_{\rm null}=[-0.2,0.8]^\top$。加速度为 $[-0.2,0.2]^\top$，两者之和为零，主任务不产生额外瞬时加速度。

\textbf{练习 2}\quad 教学2R在 $q=(0,\pi/2)$、零速度、重力已完整补偿时，希望末端加速度为 $[0,3]^\top$。求额外任务力和关节力矩。

\textbf{解}\quad $\Lambda=\operatorname{diag}(1/3,4/3)$，所以 $F=\Lambda a=[0,4]^\top\,\mathrm{N}$。$\tau=J^\top F=[2,0]^\top\,\mathrm{N\,m}$。回代 $M^{-1}\tau=[6,-6]^\top$，再乘 $J$ 得 $[0,3]^\top$，与要求一致。

\textbf{练习 3}\quad 七轴臂只控制三维末端位置，位置雅可比秩为3。局部零空间维数是多少？如果再加入两个独立朝向约束且总体秩为5，剩多少？

\textbf{解}\quad 秩零化度定理给 $\dim\mathcal N(J)=n-\operatorname{rank}J$。第一种为 $7-3=4$，第二种为 $7-5=2$。维数取决于独立约束的秩，不能仅按“发了几项命令”计数。限位、碰撞和奇异还会限制实际可用自运动。

\textbf{练习 4}\quad 旋转刚度 $k_R=8\,\mathrm{N\,m/rad}$，绕一个轴的实际误差为 $-5^\circ$，求恢复力矩。

\textbf{解}\quad 弧度误差为 $-5\pi/180=-0.0872665$，恢复力矩为 $-8(-0.0872665)=+0.698132\,\mathrm{N\,m}$。恢复方向与误差相反。该计算针对单轴旋转，避免了多轴大角度下旋转误差导数的额外几何关系。

\section*{本章与原稿的对应}
本章整合 [A，第29--30、44页] 与 [B，第55、83--86、122--123页]，并参考操作空间控制原始论文\cite{khatib}。补足了 $\Lambda$、$\bar J$、$N^\top$ 的推导和数值回代，纠正严格解耦、任务重力覆盖范围、正则化零空间、姿态误差坐标以及七轴冗余维数等容易影响真机实现的细节。

\chapter{从目标点到连续轨迹}
\label{ch:trajectory}

把目标写成“关节走到 $1.2\,\mathrm{rad}$”，只说明了终点。电机还需要知道每一时刻该在哪里、该多快、该如何加减速。轨迹生成的任务，就是把一个终点变成连续的参考信号 $q_d(t),\dot q_d(t),\ddot q_d(t)$，并检查它是否超过机构、驱动器和任务允许的范围。

\section{路径 轨迹 插补和控制的分工}

\textbf{路径（path）}是几何曲线 $q(s)$，$s\in[0,1]$ 只表示进度，没有秒的单位。\textbf{时间律（time scaling）}是 $s(t)$，告诉我们第几秒走到路径的几成。二者复合为\textbf{轨迹（trajectory）} $q(t)=q(s(t))$。\textbf{插补（interpolation）}是在离散控制时刻求出连续轨迹的值。\textbf{反馈控制}根据参考和实际测量之间的差，计算位置、速度或力矩命令。

例如，同一条 $0$ 到 $1\,\mathrm{rad}$ 的关节路径可以用 $1\,\mathrm{s}$ 或 $5\,\mathrm{s}$ 走完；它们的几何完全相同，但所需加速度、惯性力矩和发热不同。避障规划找到一条路，不能据此断言这条路可以直接执行。

\begin{table}[htbp]\centering\small
\begin{tabularx}{\linewidth}{p{25mm}XX}\toprule
层次 & 输入和输出 & 需要回答的问题\\\midrule
任务 & 起点 终点 工艺要求 & 相机应到哪里 面向哪里\\
路径规划 & 环境与目标到 $q(s)$ & 连杆及工具是否碰撞\\
时间参数化 & 路径与限值到 $s(t)$ & 速度 加速度 力矩是否可行\\
周期插补 & 当前时间到参考状态 & 这一拍的 $q_d,\dot q_d,\ddot q_d$ 是多少\\
反馈控制 & 参考和测量到执行器命令 & 如何克服重力 摩擦和扰动\\\bottomrule
\end{tabularx}
\caption{轨迹与控制的接口必须按物理量区分。}
\end{table}

MoveJ 通常描述关节空间的运动指令；MoveL 通常描述末端直线；MoveC 通常描述圆弧。MoveP 在不同厂家中可能指点到点位姿运动，也可能指连续工艺运动，不能仅凭名字判断轨迹。所谓 MIT 模式常见的是关节位置、速度、刚度、阻尼和前馈力矩组合接口，它属于执行层接口，不是一种路径算法。具体字段、单位和控制模式必须以该驱动器协议为准。

\section{先亲手解出三次和五次多项式}

\subsection{三次多项式为什么有四个系数}

设运动从 $q_0$ 到 $q_f$，总时间 $T>0$，起止速度均为零。选
\begin{equation}
q(t)=a_0+a_1t+a_2t^2+a_3t^3.
\end{equation}
四个未知系数需要四个独立条件。由 $q(0)=q_0$ 和 $\dot q(0)=0$ 得 $a_0=q_0,a_1=0$。令 $\Delta=q_f-q_0$，另两个条件给出
\begin{align}
a_2T^2+a_3T^3&=\Delta,\\
2a_2T+3a_3T^2&=0.
\end{align}
第二式除以 $T$ 后得 $a_2=-3a_3T/2$，代回第一式，得到 $a_3=-2\Delta/T^3$，于是
\begin{equation}
q(t)=q_0+\Delta(3s^2-2s^3),\qquad s=t/T.
\end{equation}
注意 $s$ 是无量纲数。对 $t$ 求导要再乘 $ds/dt=1/T$：
\begin{align}
\dot q&=\frac{\Delta}{T}(6s-6s^2),\\
\ddot q&=\frac{\Delta}{T^2}(6-12s).
\end{align}
速度峰值在 $s=1/2$，为 $1.5|\Delta|/T$。如果在运动前后保持静止，外部加速度为零，而运动刚开始的加速度是 $6\Delta/T^2$，因此拼接处加速度跳变。一个区间内的多项式很光滑，不等于它与前后区间的拼接也很光滑。

\subsection{五次多项式的六个边界条件}

要让起止加速度也为零，再增加两个条件，并把次数增加到五次。采用无量纲写法
\begin{equation}
q(t)=q_0+\Delta h(s),\quad
h(s)=b_0+b_1s+b_2s^2+b_3s^3+b_4s^4+b_5s^5.
\end{equation}
$h(0)=h'(0)=h''(0)=0$ 立即给出 $b_0=b_1=b_2=0$。终点条件构成
\begin{equation}
\begin{bmatrix}1&1&1\\3&4&5\\6&12&20\end{bmatrix}
\begin{bmatrix}b_3\\b_4\\b_5\end{bmatrix}
=\begin{bmatrix}1\\0\\0\end{bmatrix}.
\end{equation}
第二行减第一行的三倍：$b_4+2b_5=-3$；第三行减第二行的两倍：$4b_4+10b_5=0$。联立得 $b_5=6,b_4=-15,b_3=10$。
\begin{align}
h(s)&=10s^3-15s^4+6s^5,\\
\dot q&=\frac{\Delta}{T}(30s^2-60s^3+30s^4),\\
\ddot q&=\frac{\Delta}{T^2}(60s-180s^2+120s^3),\\
\dddot q&=\frac{\Delta}{T^3}(60-360s+360s^2).
\end{align}
$\dddot q$ 叫\textbf{加加速度（jerk）}，角运动单位为 $\mathrm{rad/s^3}$。它描述加速度改变得多快，关系到力矩变化和柔性结构的激励。五次曲线的加速度可以与静止段连续拼接，但 jerk 在边界仍从零跳到非零。jerk 有界与 jerk 连续是两件事。

\begin{examplebox}[完整算例 从零走到一点二弧度]
取 $q_0=0,q_f=1.2\,\mathrm{rad},T=2\,\mathrm{s}$。
\[
q(t)=1.5t^3-1.125t^4+0.225t^5.
\]
在 $t=1\,\mathrm{s}$，$s=0.5$，所以 $h=0.5$，
\[
q=0.6\,\mathrm{rad},\quad \dot q=1.125\,\mathrm{rad/s},\quad \ddot q=0.
\]
加速度极值满足 $\dddot q=0$，即 $6s^2-6s+1=0$，得 $s=(3\pm\sqrt3)/6$。代回后
\[
|\ddot q|_{\max}=\frac{10\sqrt3}{3}\frac{|\Delta|}{T^2}
=1.7321\,\mathrm{rad/s^2}.
\]
最大 jerk 在区间端点，为 $60|\Delta|/T^3=9\,\mathrm{rad/s^3}$。若速度上限只有 $1\,\mathrm{rad/s}$，这条两秒轨迹虽然终点正确，仍然不可执行。
\end{examplebox}

\begin{figure}[htbp]\centering
\includegraphics[width=0.94\linewidth]{figures/rt_polynomials.pdf}
\caption{相同位移和总时间下的三次与五次轨迹。五次曲线使起止加速度为零，但中间峰值速度更高，不能把“更平滑”误读成“所有峰值都更小”。}
\end{figure}

\subsection{从三个限值反算运动时间}

对零速度、零加速度起止的五次曲线，速度、加速度、jerk 限制共同要求
\begin{equation}
T\ge\max\left\{
\frac{1.875|\Delta|}{v_{\max}},\quad
\sqrt{\frac{(10\sqrt3/3)|\Delta|}{a_{\max}}},\quad
\sqrt[3]{\frac{60|\Delta|}{j_{\max}}}\right\}.
\end{equation}
取 $\Delta=1.2,v_{\max}=1,a_{\max}=2,j_{\max}=10$，三项分别为 $2.25,1.8612,1.9310\,\mathrm{s}$，故取 $T=2.25\,\mathrm{s}$。这只是\textbf{所选五次曲线家族内}满足这些限值的最小时间，不是所有可能轨迹中的时间最优解；也尚未检查力矩和碰撞。

五次曲线常被称为“最小 jerk 轨迹”，这个称呼有前提：固定总时间与两端位置、速度、加速度，最小化 $\int_0^T\dddot q^2dt$。对 $q+\epsilon\eta$ 做变分，并三次分部积分，端点的 $\eta,\dot\eta,\ddot\eta$ 为零，得到 $\delta J=-2\int q^{(6)}\eta dt=0$，因此 $q^{(6)}=0$。这给出五次多项式。若加入中间点、障碍或力矩约束，不能继续照搬这个无约束结论。

\subsection{起止速度不为零时怎样计算}

令 $q(t)=\sum_{k=0}^{5}a_kt^k$，已知两端 $(q_0,v_0,\alpha_0)$ 和 $(q_f,v_f,\alpha_f)$。先得 $a_0=q_0,a_1=v_0,a_2=\alpha_0/2$，再求
\begin{equation}
\begin{bmatrix}T^3&T^4&T^5\\3T^2&4T^3&5T^4\\6T&12T^2&20T^3\end{bmatrix}
\begin{bmatrix}a_3\\a_4\\a_5\end{bmatrix}
=\begin{bmatrix}
q_f-q_0-v_0T-\alpha_0T^2/2\\
v_f-v_0-\alpha_0T\\
\alpha_f-\alpha_0
\end{bmatrix}.
\end{equation}
在程序中解线性方程即可，不要显式求逆。取 $T=2,q_0=0,v_0=0.2,\alpha_0=0,q_f=1,v_f=\alpha_f=0$，右端为 $(0.6,-0.2,0)^\mathsf T$，解为
\[
a_3=0.95,\qquad a_4=-0.7375,\qquad a_5=0.15.
\]
验算 $q(2)=0.4+7.6-11.8+4.8=1$；$\dot q(2)=0.2+11.4-23.6+12=0$；$\ddot q(2)=11.4-35.4+24=0$。如果 $T$ 极大或极小，应先用 $s=t/T$ 归一化，避免 $T^3,T^4,T^5$ 数量级悬殊。

\section{梯形速度与七段 S 曲线}

\subsection{梯形曲线由面积求出来}

设位移绝对值为 $d$，加速度上限为 $a_m$，速度上限为 $v_m$。加速时间 $t_a=v_m/a_m$；加速与减速两块三角形的总面积为 $v_m^2/a_m$。因此
\begin{equation}
T=\begin{cases}
d/v_m+v_m/a_m,&d\ge v_m^2/a_m,\\
2\sqrt{d/a_m},&d<v_m^2/a_m.
\end{cases}
\end{equation}
短距离情形没有匀速段，实际峰值速度为 $\sqrt{da_m}$。例如 $d=0.1\,\mathrm{rad},a_m=2,v_m=1$，因为 $0.1<0.5$，得到 $T=0.4472\,\mathrm{s}$，峰值速度 $0.4472\,\mathrm{rad/s}$。不能先算出负的匀速时间，再把负值强行置零而不重算其它段。

\subsection{七段曲线的适用条件}

令 jerk 上限为 $j_m$。当速度和加速度上限都能达到，且有非负匀速段时，前半段为：以 $+j_m$ 增加加速度、保持 $a_m$、以 $-j_m$ 降回零，再匀速。后三段对称减速。
\begin{align}
t_j&=a_m/j_m,\\
t_a&=v_m/a_m-t_j,\\
t_c&=d/v_m-(2t_j+t_a),\\
T&=4t_j+2t_a+t_c=d/v_m+v_m/a_m+a_m/j_m.
\end{align}
这里 $t_a$ 是“恒加速度段”的时间，不包含两端 jerk 段。这个公式只在 $t_a\ge0,t_c\ge0$ 时成立。

\begin{examplebox}[七段曲线的逐段时间]
取 $d=3\,\mathrm{rad},v_m=1,a_m=2,j_m=8$。先算 $t_j=0.25\,\mathrm{s}$，$t_a=0.25\,\mathrm{s}$，$t_c=2.25\,\mathrm{s}$，总时间 $T=3.75\,\mathrm{s}$。七段持续时间是
\[
(0.25,0.25,0.25,2.25,0.25,0.25,0.25)\,\mathrm{s}.
\]
任意恒 jerk 段从 $(q_s,v_s,a_s)$ 起步，经过 $\delta t$ 后都有
\[
a=a_s+j\delta t,\quad v=v_s+a_s\delta t+\tfrac12j\delta t^2,
\]
\[
q=q_s+v_s\delta t+\tfrac12a_s\delta t^2+\tfrac16j\delta t^3.
\]
第一段末 $a=2,v=0.25,q=0.020833$。把这三个值作为第二段初值，逐段累积即可得到整条曲线，不需要猜曲线形状。
\end{examplebox}

\begin{figure}[htbp]\centering
\includegraphics[width=0.93\linewidth]{figures/rt_scurve.pdf}
\caption{七段 S 曲线的速度、加速度和 jerk。jerk 分段常数，加速度连续；这些曲线来自逐段积分，位移总和为 $3\,\mathrm{rad}$。}
\end{figure}

极短距离时，加速度上限和速度上限可能都达不到。若只有四个等长 jerk 段，符号依次为 $+,-,-,+$，每段时间为 $t_j$，则 $a_{\mathrm{peak}}=j_mt_j$，$v_{\mathrm{peak}}=j_mt_j^2$，对称面积给出 $d=2j_mt_j^3$。于是 $t_j=(d/(2j_m))^{1/3}$。求完还要验证两个峰值小于相应限值。这说明 S 曲线必须分情况讨论；在线运动生成器把这些分支系统地封装起来。Ruckig 的官方文档明确区分状态到状态运动和含中间点的情况，不能把它概括为“任何避障任务都能直接时间最优”。\cite{ruckig}


\subsection{两种中间分支也要会判断}

如果达到加速度上限但没有匀速段，先令 $t_j=a_m/j_m$，由位移面积
\[
d=a_m(t_a+t_j)(t_a+2t_j)
\]
解得
\[
t_a=\frac{-3t_j+\sqrt{t_j^2+4d/a_m}}{2},\qquad t_c=0.
\]
必须验证 $t_a\ge0$，实际峰值速度 $a_m(t_a+t_j)\le v_m$。例如 $d=1,a_m=2,j_m=8,v_m=2$，$t_j=0.25$，$t_a\approx0.34307$，峰值速度约 $1.18614$，总时间 $4t_j+2t_a\approx1.68614\,\mathrm s$。

另一种情况是先达到速度上限、却达不到加速度上限。此时
\[
t_j=\sqrt{v_m/j_m},\quad t_a=0,\quad
 t_c=d/v_m-2t_j.
\]
验证 $j_mt_j\le a_m$ 和 $t_c\ge0$。例如 $d=1,v_m=0.5,j_m=8,a_m=4$，$t_j=0.25$，峰值加速度 $2$，匀速时间 $1.5$，总时间 $2.5\,\mathrm s$。至此，零起止速度和加速度的对称点到点运动有了完整的四类情形。非零边界速度或加速度会引入更多分支，需按新的边界状态重新求解。

\section{多轴同步与时间缩放}

七个关节分别得到最小可用时间后，在\textbf{零起止速度和加速度、同一曲线形状}的条件下，可以取最大的时间作为公共时间 $T$，再令 $q_i=q_{i0}+\Delta_i h(t/T)$。所有关节同时开始、同时结束。位移为零的轴可以保持不动，无需除以位移。

把已有轨迹放慢 $\gamma>1$ 倍，令 $q_{\mathrm{new}}(t)=q_{\mathrm{old}}(t/\gamma)$，链式法则给出
\begin{equation}
\dot q_{\mathrm{new}}=\gamma^{-1}\dot q_{\mathrm{old}},\quad
\ddot q_{\mathrm{new}}=\gamma^{-2}\ddot q_{\mathrm{old}},\quad
\dddot q_{\mathrm{new}}=\gamma^{-3}\dddot q_{\mathrm{old}}.
\end{equation}
减速通常迅速降低惯性需求，但重力力矩不会随速度减小而消失。若初末速度必须保持某个非零数值，直接拉长时间会改变边界速度，必须重新求轨迹。

\section{末端直线和圆弧怎样落到关节上}

\subsection{关节直线一般不是末端直线}

二维两连杆长度都为 $0.5\,\mathrm m$。起点 $q_A=(0,90^\circ)$，末端为 $(0.5,0.5)$；终点 $q_B=(60^\circ,60^\circ)$，末端为 $(0,\sqrt3/2)$。关节插值到一半，$q=(30^\circ,75^\circ)$，由正运动学
\[
p=\begin{bmatrix}0.5\cos30^\circ+0.5\cos105^\circ\\0.5\sin30^\circ+0.5\sin105^\circ\end{bmatrix}
=\begin{bmatrix}0.30360\\0.73296\end{bmatrix}\mathrm m.
\]
而末端直线的中点是 $(0.25,0.68301)\,\mathrm m$，两者相差约 $0.0733\,\mathrm m$。即使关节跟踪误差为零，MoveJ 的末端仍不会自动走直线。

真正的末端直线先规定 $p(s)=p_A+s(p_B-p_A)$，再逐点求逆解。在上述中点，
\[
\cos q_2=\frac{p_x^2+p_y^2-l_1^2-l_2^2}{2l_1l_2}\approx0.05801.
\]
选择与起点连续的 $q_2>0$ 分支，$q_2\approx86.674^\circ$，$q_1\approx26.56^\circ$。它与关节插值的中点明显不同。

\begin{figure}[htbp]\centering
\includegraphics[width=0.93\linewidth]{figures/rt_cartesian.pdf}
\caption{同一对关节起止状态对应的两种路径。左图看末端几何，右图看关节随进度的变化。路径选在哪个空间，决定中间状态。}
\end{figure}

\subsection{姿态和圆弧的参数化}

姿态使用同一个进度 $s(t)$。给定两端旋转矩阵，可写
\begin{equation}
R(s)=R_A\exp\!\left(s\log(R_A^\mathsf TR_B)\right).
\end{equation}
若用四元数，先确保相邻四元数点积非负，再用球面插值 SLERP；在相对转角接近 $\pi$ 时须明确选定的轴和分支。欧拉角逐分量插值在特定小角度、远离奇异的场景可以使用，但它不等于最短旋转路径，不应当作一般姿态插值算法。

圆弧可用圆心 $c$、正交单位基 $u,v$、半径 $r$ 和角度 $\theta(s)$ 表示：
\begin{equation}
p(s)=c+r\big(u\cos\theta(s)+v\sin\theta(s)\big).
\end{equation}
例如平面圆弧 $c=(0.4,0.4),r=0.1,\theta=\pi s/2$。在 $s=0.5$，$p=(0.47071,0.47071)\,\mathrm m$。若匀速转角 $\dot\theta=1\,\mathrm{rad/s}$，切向速度为 $r\dot\theta=0.1\,\mathrm{m/s}$，法向加速度为 $r\dot\theta^2=0.1\,\mathrm{m/s^2}$。所以“沿路径匀速”不代表空间加速度为零。

\section{给定路径后的时间最优参数化}

设几何路径已经无碰撞。对 $q(s(t))$ 连续求两次导数：
\begin{equation}
\dot q=q'(s)\dot s,\qquad
\ddot q=q'(s)\ddot s+q''(s)\dot s^2.
\end{equation}
撇号表示对 $s$ 求导，点表示对时间求导，不能互换。令 $x=\dot s^2\ge0$，则在正向运动 $\dot s>0$ 时 $dx/ds=2\ddot s$。每个关节的加速度约束变成
\begin{equation}
-a_{i,\max}\le q_i'(s)\ddot s+q_i''(s)x\le a_{i,\max}.
\end{equation}
在固定的 $s$ 上，这是关于 $(x,\ddot s)$ 的线性不等式。刚体惯性和科氏项也可整理成 $a(s)\ddot s+b(s)x+g(s)$，但速度相关摩擦、jerk 和接触切换需要另行处理，不能全部无条件塞进同一个线性形式。

\begin{examplebox}[一段路径上的可行加速度区间]
某一点有 $q_1'=2,q_1''=0.5$，关节加速度限值为 $3$，当前 $\dot s=1$，因此 $x=1$。约束为 $-3\le2\ddot s+0.5\le3$，故 $-1.75\le\ddot s\le1.25$。
若另一个关节给出 $-0.8\le\ddot s\le0.9$，共同可行区间是 $[-0.8,0.9]$。规划器不能只采用第一关节允许的 $1.25$，必须取所有约束的交集。若交集为空，就需降低路径速度，或重新规划几何路径。
\end{examplebox}

在最简单的双积分器、固定路径、仅速度和加速度约束问题中，加速到上限再尽可能晚地制动会产生时间最优解。一般路径含曲率和多个约束时，仍需系统地求可行速度包络与切换点。TOPP-RA 用可达性分析求时间参数化；它保留给定路径，作用与在线状态到状态轨迹生成不同。\cite{toppra}

\section{构型空间与避障为什么要检查整段}

\textbf{构型空间（configuration space）}把一整组关节角看作一个点。二维两关节机械臂的构型空间是一张 $(q_1,q_2)$ 图；七轴机械臂则是七维空间。障碍物在工作空间中只是一个圆或方块，在构型空间中却可能变成复杂区域，因为任何一段连杆和工具都不能碰撞。

\begin{figure}[htbp]\centering
\includegraphics[width=0.95\linewidth]{figures/rt_cspace.pdf}
\caption{圆障碍物对应的两关节构型空间禁区。两端姿态都无碰撞，连接它们的关节直线仍可穿过禁区。右图由逐点检查两段连杆到圆心的距离计算，未包含连杆厚度以外的额外安全余量。}
\end{figure}

RRT 的基本过程是随机取构型 $q_{\rm rand}$，在现有树上找最近节点 $q_{\rm near}$，沿二者方向走一小步得到 $q_{\rm new}$，仅当整段连接有效才加入树。RRT-Connect 从起终点两端生长并尝试连接。它们适合高维查询，但输出一般仍需路径简化和时间参数化。\textbf{概率完备}表示在相应假设下，若存在具有足够间隙的可行路径，采样数趋于无穷时找到解的概率趋于一；有限时间内没有解，不能据此证明问题不可行。\cite{lavalle}

CHOMP、STOMP、TrajOpt 等把离散路径或轨迹作为优化变量，用平滑代价和碰撞代价改进它，通常存在局部最优。图搜索和采样搜索寻找路径，轨迹优化改良路径，时间参数化赋予时间，它们可以组合。所有后处理都可能改变几何，必须重新检查碰撞。特别是拐角处\textbf{圆滑过渡（blending）}通常主动偏离原航点；如果航点必须精确通过，就要降低速度或在该点停下。

\section{计算练习与解答}

\subsection*{题一 五次轨迹的时间}
一个关节移动 $0.8\,\mathrm{rad}$，限值 $v_m=0.5,a_m=1,j_m=4$。求零起止五次曲线的可用最短时间。

解：三个下界为 $1.875\times0.8/0.5=3.0\,\mathrm{s}$，$\sqrt{5.7735\times0.8}=2.1491\,\mathrm{s}$，$\sqrt[3]{60\times0.8/4}=2.2894\,\mathrm{s}$。故取 $T=3.0\,\mathrm{s}$。随后还需核查力矩；重力上限并不会因拉长到三秒自动满足。

\subsection*{题二 变速回放}
轨迹原峰值速度 $0.6$、加速度 $1.2$、jerk $3$，将总时间缩短到原来的 $0.8$，求新峰值。

解：$\gamma=0.8$，故速度 $0.6/0.8=0.75$，加速度 $1.2/0.8^2=1.875$，jerk $3/0.8^3=5.8594$。加速回放时，jerk 的增长比速度快得多。

\subsection*{题三 可行区间}
若 $q'=1,q''=-2,\dot s=0.5$，关节加速度范围为 $[-1,1]$，求 $\ddot s$ 范围。

解：$\dot s^2=0.25$，故 $\ddot q=\ddot s-0.5$。由 $-1\le\ddot s-0.5\le1$ 得 $-0.5\le\ddot s\le1.5$。若速度继续增加，曲率项也改变，不能把此区间当成整条路径的常数。

\section*{本章与原稿的对应}
本章整合 [A，第52--60页] 与 [B，第87--93页]。补出多项式系数、峰值和 S 曲线的分支条件，修正“轨迹平滑就自动无碰撞”“所有 S 曲线都适用同一总时间公式”等容易产生的误读。

\chapter{拖动示教 回放与运动学习}
\label{ch:teaching}

示教的本质是把人的意图转换成一条可以复现的参考轨迹。拖动时的力矩控制、数据录制、轨迹修整和回放时的跟踪控制，是四个不同问题。把它们拆开，才容易判断“推起来不顺”“回放不平滑”和“回放位置不准”分别出在哪里。

\section{拖动时机器人到底应该做什么}

对力矩可控机械臂，最简单的教学模型是
\begin{equation}
\tau_{\rm cmd}=\hat g(q)+\hat\tau_f(\dot q)-D_q\dot q.
\end{equation}
$\hat g$ 估计并补偿重力；$\hat\tau_f$ 补偿部分摩擦；$D_q\succeq0$ 提供阻尼。理想补偿且没有外力时，剩余运动方程为 $M\ddot q+C\dot q+D_q\dot q=0$。半正定阻尼只保证不增加机械能，未阻尼的运动模态可能持续存在。若 $D_q\succ0$，并满足轨迹有界等常规条件，才可用能量与不变集分析得到速度趋于零；系统没有位置弹簧，不会自动回到某个固定构型。

\textbf{阻尼不能在静止时托住重力。}因为 $\dot q=0$ 时 $-D_q\dot q=0$。如果重力补偿少了 $1\,\mathrm{N\,m}$，加一点阻尼只能限制下落速度，不能使任意姿态静止平衡。以一维近似 $J\ddot q+D\dot q=-1$ 为例，稳态速度为 $-1/D$。取 $D=2\,\mathrm{N\,m\,s/rad}$，仍会以 $-0.5\,\mathrm{rad/s}$ 漂移。原稿中把小阻尼描述成“托住”的说法需要这样修正。

若采用导纳拖动，先由测得或估计的外力生成参考位移，再交给位置内环：
\begin{equation}
M_v\ddot x_r+D_v\dot x_r+K_v(x_r-x_0)=F_{\rm ext}.
\end{equation}
$K_v=0$ 时松手后速度会衰减，但停留位置由此前运动决定；$K_v>0$ 时会回到 $x_0$。关节拖动、末端某几个方向自由、肘部自由，可以用不同任务约束实现，不能把一个标量“软硬”旋钮当成全部设计。

\section{记录哪些量才能在以后解释结果}

每条样本至少包括单调时间戳、按固定关节名排序的实测 $q$、控制模式以及数据有效标记。建议同时记录实测或估计速度、指令力矩、电流换算力矩、末端位姿和所用模型版本。关节命令不能替代关节测量；电流换算力矩也不能在没有校准时当成独立高精度力矩传感器。

\begin{table}[htbp]\centering\small
\begin{tabularx}{\linewidth}{p{31mm}XX}\toprule
字段 & 用途 & 常见错误\\\midrule
$t_k$ 与序号 & 重建真实时间和丢帧 & 用数组下标假定每帧等间隔\\
关节名与 $q_k$ & 保留实际构型 & 左右臂或关节顺序串位\\
参考与测量 & 分开路径问题和跟踪问题 & 只保存指令画出“完美跟踪”\\
工具与基座变换 & 重建实际任务坐标 & 换工具后继续用旧 TCP\\
质量与质心版本 & 解释重力和观测器残差 & 只记录文件名没有参数版本\\
模式与限幅状态 & 找切换和饱和问题 & 把被截断的输出当控制器原输出\\\bottomrule
\end{tabularx}
\caption{一条轨迹文件应该同时保存轨迹本身及解释它所需的上下文。}
\end{table}

假定录制名义频率为 $100\,\mathrm{Hz}$，真实样本时间是 $0,0.010,0.023\,\mathrm{s}$，对应位置 $0,0.002,0.006\,\mathrm{rad}$。在 $0.020\,\mathrm{s}$ 重采样，线性插值为
\[
q=0.002+\frac{0.020-0.010}{0.023-0.010}(0.006-0.002)
=0.0050769\,\mathrm{rad}.
\]
若误把第三帧当成 $0.020\,\mathrm{s}$，就直接取 $0.006$，人为引入约 $0.0009231\,\mathrm{rad}$ 差异。以后再对它求速度，时间错误还会继续放大。

\section{把一百赫兹记录交给两百赫兹控制器}

插补到 $200\,\mathrm{Hz}$ 不会创造新的测量信息，只是在已有样本之间构造参考。步骤是：去掉重复或倒退时间戳，按关节名映射，检查跳变，选择拟合或滤波方法，再在控制器自己的时钟上查询轨迹。它可以提高参考的连续性，但不能把原始 $100\,\mathrm{Hz}$ 信号的可恢复频率变成 $100\,\mathrm{Hz}$。

\subsection{平滑和保形不是同一个要求}

三次样条通常具有连续的二阶导数，但可能在稀疏点之间过冲。PCHIP 是分段三次 Hermite 保形插值，在标量单调数据区间可以保持单调性，通常只有一阶导数连续。把它说成“同时没有过冲且加速度处处连续”是不正确的。Akima 也是局部三次方法，但不能一般性保证单调和无过冲。

若关节必须满足限位，应该检查连续插值曲线的极值；只检查记录点不够。如果必须让加速度连续，可采用带边界条件的平滑样条或优化求解，并把限位作为约束。对姿态则需要在旋转空间插值；先统一四元数符号，再进行 SLERP 或适当的姿态样条。

\begin{figure}[htbp]\centering
\includegraphics[width=0.94\linewidth]{figures/rt_teach.pdf}
\caption{示教轨迹的采样、噪声和重建示例。左图是位置，右图比较差分速度与平滑后速度。平滑会改变原数据，必须同时看位置偏差与动态峰值，不能只看曲线是否好看。}
\end{figure}

离线双向滤波可以避免净相位滞后，但必须使用未来数据，并会改变幅值与边界。处理前后各保留一份数据，裁边范围和滤波参数写入轨迹元数据。如果要保留关键时刻的接触动作，不应盲目把所有高频成分都删掉。

\subsection{回放起点和时间缩放}

回放前，当前位置 $q_{
m now}$ 通常不同于轨迹第一点 $q_0$。需要先生成一段从 $q_{
m now}$ 到 $q_0$ 的过渡轨迹，完成后再启用录制时间轴。把第一帧直接发给内环，会产生一个不在示教数据里的跳变。

若把两秒示教压到一秒，位置路径不变，速度加倍、加速度四倍、jerk 八倍。假定原峰值分别为 $0.4\,\mathrm{rad/s}$、$0.8\,\mathrm{rad/s^2}$、$2\,\mathrm{rad/s^3}$，压缩后就是 $0.8,3.2,16$。这也是为什么“人手推得动”不能作为自动高速回放的可行性证明。

\section{关节回放和笛卡尔回放如何选择}

关节回放直接复用 $q_d(t)$，保留示教时的肘部构型，避免重新选择逆解分支。但实际轨迹仍可能因控制误差、弹性、障碍变化和工具变化而偏离；它不保证绝对定位精度，也不自动保证无碰撞。

笛卡尔回放保存 $T_{BE}(t)$，方便把整个动作变换到新的工件坐标系。例如原工作坐标系是 $W$，工具轨迹为 $T_{WE}(t)$；新工作坐标系在基座中的位置是 $T_{BW}'$，那么新目标是
\begin{equation}
T_{BE}'(t)=T_{BW}'T_{WE}(t).
\end{equation}
这一步是坐标变换，不是把每个坐标数直接加上一个偏移。之后仍要做连续逆解、冗余姿态处理和完整可行性检查。

\begin{examplebox}[把运镜轨迹绕基座旋转]
原轨迹一点在工作系中为 $p_W=(0.2,0.1,0)\,\mathrm m$。新工作系相对基座绕 $z$ 旋转 $90^\circ$，原点为 $(0.5,0,0)$。于是
\[
p_B=\begin{bmatrix}0&-1&0\\1&0&0\\0&0&1\end{bmatrix}
\begin{bmatrix}0.2\\0.1\\0\end{bmatrix}
+\begin{bmatrix}0.5\\0\\0\end{bmatrix}
=\begin{bmatrix}0.4\\0.2\\0\end{bmatrix}\mathrm m.
\]
如果只把工作系原点加上去，就会错误地得到 $(0.7,0.1,0)$。工具姿态也必须左乘同一个旋转矩阵。
\end{examplebox}

\section{动态运动基元从稳定系统加形状项得到}

\textbf{DMP（dynamic movement primitive，动态运动基元）}把示教形状写进一个动力系统。这里用 $y$ 表示一个标量运动坐标，$g$ 表示目标，$\rho$ 表示从一衰减到零的相位，$T_d>0$ 是时间尺度，取 $\alpha_z>0,\beta_z>0,\alpha_\rho>0,\rho(0)=1$：
\begin{align}
T_d\dot z&=\alpha_z\big(\beta_z(g-y)-z\big)+f(\rho),\\
T_d\dot y&=z,\\
T_d\dot\rho&=-\alpha_\rho\rho.
\end{align}
消去 $z$ 得
\begin{equation}
T_d^2\ddot y+\alpha_zT_d\dot y+\alpha_z\beta_z(y-g)=f(\rho).
\end{equation}
当 $f=0$，它就是一个稳定的弹簧阻尼系统；取 $\beta_z=\alpha_z/4$ 时，对应临界阻尼。相位解为 $\rho(t)=\exp(-\alpha_\rho t/T_d)$。形状项采用高斯基函数
\begin{equation}
\psi_i(\rho)=\exp[-h_i(\rho-c_i)^2],\qquad
f(\rho)=\frac{\sum_i\psi_i(\rho)w_i}{\sum_i\psi_i(\rho)}\rho a_s.
\end{equation}
$c_i$ 是中心，$h_i>0$ 控制宽度，$w_i$ 是要学习的权重，$a_s$ 是空间尺度。只要权重有界、分母有效、目标固定，且没有破坏稳定性的额外耦合，$\rho\to0$ 会让形状项消失，最终回到稳定目标系统。加入避障耦合、时变目标或真实执行器之后，需要重新分析，不能再用一句“DMP一定稳定”代替证明。

\subsection{权重怎样从示教数据算出来}

已有示教 $y_d,\dot y_d,\ddot y_d$ 时，先把它代入方程反求所需形状力：
\begin{equation}
f_d=T_d^2\ddot y_d-\alpha_z\big(\beta_z(g-y_d)-T_d\dot y_d\big).
\end{equation}
定义 $\Phi_{ki}=\rho_k a_s\psi_i(\rho_k)/\sum_j\psi_j(\rho_k)$，于是 $f_d\approx\Phi w$，求解 $\min_w\|\Phi w-f_d\|^2+\lambda\|w\|^2$。这正是前面学过的正则化最小二乘，DMP 没有绕开辨识数学。

\begin{examplebox}[一个 DMP 样本怎样产生回归目标]
取 $T_d=1,\alpha_z=4,\beta_z=1,g=1$。某时刻示教为 $y_d=0.5,\dot y_d=1,\ddot y_d=0$，则
\[
f_d=0-4[(1-0.5)-1]=2.
\]
若只用一个基函数，此处 $\rho=0.5,a_s=1$，归一化后 $\Phi=0.5$，因此这个样本要求 $0.5w=2$，即 $w=4$。多个时刻通常不能被同一个权重精确满足，所以要堆叠成回归问题，而不是每个时刻单独改一次权重。
\end{examplebox}

常见写法令 $a_s=g-y_0$。如果任务起点等于终点，尺度变成零，整个形状项会丢失。可以使用独立尺度、分段动作或周期型 DMP，不能只在除法里加一个很小的数就声称问题已解决。

\begin{figure}[htbp]\centering
\includegraphics[width=0.94\linewidth]{figures/rt_dmp.pdf}
\caption{同一组 DMP 权重的重建与时间缩放示例。改变时间尺度会改变速度；改变目标或尺度后仍应重新检查关节和碰撞限制。曲线由方程数值积分得到。}
\end{figure}

\section{概率运动基元与一次条件化计算}

\textbf{ProMP（probabilistic movement primitive，概率运动基元）}不只学一条曲线，而是学权重分布。设
\begin{equation}
y(t)=\phi(t)^\mathsf Tw,\qquad w\sim\mathcal N(\mu,\Sigma).
\end{equation}
那么 $y(t)$ 的均值为 $\phi^\mathsf T\mu$，方差为 $\phi^\mathsf T\Sigma\phi$。不同示教在某段很一致，通常使该段方差较小；但这也可能只是示教样本不足，不能直接把方差当成任务重要程度的唯一证据。

要让时刻 $t_*$ 经过观测值 $y_*$，并假设观测噪声方差为 $r$，高斯条件公式为
\begin{align}
K&=\frac{\Sigma\phi_*}{\phi_*^\mathsf T\Sigma\phi_*+r},\\
\mu'&=\mu+K(y_*-\phi_*^\mathsf T\mu),\\
\Sigma'&=\Sigma-K\phi_*^\mathsf T\Sigma.
\end{align}
它与卡尔曼滤波的测量更新形式相同，但这里估计的是轨迹权重。

取 $\phi(t)=(1,t)^\mathsf T,\mu=(0,1)^\mathsf T,\Sigma=\operatorname{diag}(0.04,0.09)$，在 $t_*=0.5$ 希望经过 $0.8$，令 $r=0.01$。原预测均值为 $0.5$、方差为 $0.0625$，因此
\[
K=\frac{(0.04,0.045)^\mathsf T}{0.0725}
=(0.551724,0.620690)^\mathsf T,
\]
\[
\mu'=(0.165517,1.186207)^\mathsf T.
\]
新轨迹在 $0.5$ 时的均值为 $0.758621$，没有精确到 $0.8$，因为我们允许观测有噪声。取 $r\to0$ 才趋向硬约束，但多条相互矛盾的硬约束可能无法同时满足。

\section{动态时间规整与重复学习}

\subsection{DTW 先对齐动作进度}

\textbf{DTW（dynamic time warping，动态时间规整）}寻找两条序列的单调对应关系。给 $a_i,b_j$，局部代价 $d_{ij}=(a_i-b_j)^2$，累计代价满足
\begin{equation}
D_{ij}=d_{ij}+\min(D_{i-1,j},D_{i,j-1},D_{i-1,j-1}).
\end{equation}
边界设 $D_{00}=0$，其它零行零列为正无穷，再从 $(1,1)$ 计算。对 $a=(0,1,2)$、$b=(0,0,1,2)$，可以将第一个零与两个零对应，总代价为零。按数组下标硬比则会把速度差误判成形状差。实际应限制匹配带宽，防止把明显错误动作过度拉伸成“匹配”；DTW 代价也不是以秒定义的时序跟踪误差。

\subsection{ILC 在重复任务之间修正前馈}

\textbf{迭代学习控制（iterative learning control，ILC）}使用第 $r$ 次完整执行的误差，更新第 $r+1$ 次的输入。对离散整段任务写 $y_r=Pu_r+d$，其中 $P$ 是输入到输出的线性映射，$d$ 是每次相同的扰动。令 $e_r=y_d-y_r$，用
\begin{equation}
u_{r+1}=u_r+Le_r,
\end{equation}
则
\begin{equation}
e_{r+1}=(I-PL)e_r.
\end{equation}
标量例中 $P=0.8,L=1$，误差每次乘以 $0.2$。若首次为 $0.1\,\mathrm{rad}$，后两次为 $0.02$ 和 $0.004\,\mathrm{rad}$。但如果 $L=3$，乘子为 $-1.4$，误差会变号并放大。ILC 是跨次学习，不是实时积分；初始条件和接触情况变化过大时，上一轮学习可能失效。

\section{怎样评价回放}

同一时间轴上的轨迹跟踪误差可用
\begin{equation}
E_{q,\rm RMS}=\sqrt{\frac1N\sum_{k=1}^N\|q_k-q_{d,k}\|^2},\quad
E_{p,\rm RMS}=\sqrt{\frac1N\sum_{k=1}^N\|p_k-p_{d,k}\|^2}.
\end{equation}
不同单位不要混进同一个未加权范数。\textbf{重复性}描述多次到同一目标时的离散程度；\textbf{准确度}描述均值与真实目标之间的偏差。每次都偏 $5\,\mathrm{mm}$、散布只有 $0.2\,\mathrm{mm}$，代表重复性好但存在系统偏差。仅用同一个错误运动学模型计算参考和实际 FK，发现不了真实末端的安装和标定误差。

\section{计算练习与解答}

\subsection*{题一 时间偏差折成位置偏差}
某关节近似匀速 $0.2\,\mathrm{rad/s}$，参考和测量错位 $10\,\mathrm{ms}$。求伪误差及 $0.5\,\mathrm m$ 力臂上的近似位移。

解：$\delta q=0.2\times0.01=0.002\,\mathrm{rad}$，小角度位移 $\delta p\approx0.5\times0.002=0.001\,\mathrm m=1\,\mathrm{mm}$。毫米级比较先要求毫秒级时间对齐。

\subsection*{题二 DMP 的时间尺度}
保持同一相位轨迹，把 $T_d$ 从 $2$ 改为 $4\,\mathrm s$。原峰值速度为 $0.8$，加速度为 $1.6$，求新值。

解：整体时间增加两倍，速度为 $0.4$，加速度为 $0.4$。对实际机器人还应检查是否在新速度区间进入更明显的静摩擦或爬行状态。

\subsection*{题三 ILC 的收敛范围}
标量对象 $P=0.8$，求使误差收敛的学习增益 $L$。

解：要求 $|1-0.8L|<1$。展开得 $-1<1-0.8L<1$，因此 $0<L<2.5$。端点 $L=0$ 不学习，$L=2.5$ 误差等幅变号，均不渐近收敛。

\par\smallskip
{\small\noindent\textbf{来源与扩展}\quad 本章对应 [B，第93--97页]，承接 [A，第52--60页]；补充示教时序、插值、DMP/ProMP与ILC计算，并明确回放精度与重力补偿的条件。\par}

\chapter{遥操作 视觉伺服与通信延迟}
\label{ch:teleoperation}

遥操作是一条闭环信息链：人给出运动意图，映射层转换坐标和尺度，机器人控制器执行，再把图像、状态或力反馈给人。即使每一部分单独正确，坐标混用、时间错位或模式切换仍会破坏整条链。本章先把运动映射算清楚，再讨论为什么有力反馈后延迟尤其重要。

\section{主从同构和笛卡尔映射}

同构主从臂常用关节映射
\begin{equation}
q_s=q_{s0}+S_q(q_m-q_{m0}),
\end{equation}
其中 $m$ 表示主端，$s$ 表示从端；$S_q$ 包含比例和方向修正。关节数相同不代表关节轴顺序、零位和符号相同，必须逐项定义。若主从结构不同，用末端位姿映射通常更自然。

对于位置增量，设 $R_{BM}$ 把主端参考系 $M$ 中的向量转到机器人基座系 $B$，按下离合时记录两端参考位置，随后用
\begin{equation}
p_{Bd}(t)=p_{Bd0}+s_pR_{BM}\big[p_{MH}(t)-p_{MH0}\big].
\end{equation}
$s_p$ 是无量纲位移缩放。取 $R_{BM}=R_z(90^\circ)$，$s_p=0.3$，手柄向自身参考系 $x$ 方向移动 $0.2\,\mathrm m$，从端增量为
\[
0.3\begin{bmatrix}0&-1&0\\1&0&0\\0&0&1\end{bmatrix}
\begin{bmatrix}0.2\\0\\0\end{bmatrix}
=\begin{bmatrix}0\\0.06\\0\end{bmatrix}\mathrm m.
\]
只乘缩放而漏掉旋转，会让手柄和机器人运动方向相差 $90^\circ$。

\subsection{姿态增量不能按三个欧拉角相减}

令 $R_{MH}(t)$ 是手柄姿态。以主端固定参考系表达的相对转动为
\begin{equation}
\Delta R_M=R_{MH}(t)R_{MH0}^{\mathsf T}.
\end{equation}
将它换到机器人基座系：$\Delta R_B=R_{BM}\Delta R_MR_{BM}^{\mathsf T}$。若不缩放角度，目标姿态为 $R_{Bd}=\Delta R_BR_{Bd0}$。若将旋转角缩小到 $s_R$ 倍，可以在选定对数分支后使用
\begin{equation}
R_{Bd}=\exp\!\big(s_R\log\Delta R_B\big)R_{Bd0}.
\end{equation}
右乘增量则表示在工具局部轴上的转动，是另一种操作语义。二者都可以合理，但必须与用户看到的坐标轴一致。

\subsection{离合为什么需要重新锁定参考}

\textbf{离合（clutch）}允许操作者暂时松开映射，重新摆放手柄。松开期间机器人保持当前目标或按任务设定停止；再次按下时，把此刻手柄和机器人目标都作为新的零增量参考。若只更新手柄零点而保留旧机器人原点，可能发生目标跳变；若把收到的第一帧直接当绝对目标，也可能突然跨越很大距离。

\section{从目标位姿到连续关节命令}

设几何任务误差用当前减目标表示为 $e_x$，局部闭环差分逆运动学可写成
\begin{equation}
\dot q=J^\#(v_d-K_xe_x)+N\dot q_0.
\end{equation}
$J^\#$ 是选定的伪逆或阻尼逆，$N$ 是相应冗余处理。对位置任务，$\dot e_x=J\dot q-v_d$；理想满行秩且使用精确右逆时，主任务满足 $\dot e_x=-K_xe_x$。姿态任务的对数误差还要处理误差雅可比，不能把三维位置推导直接当成任意大姿态误差的精确公式。

只做 $q_{k+1}=q_k+T_sJ^\#v_d$ 的开环积分，模型误差和跟踪误差会累积。每拍根据实测 $q$ 重算当前位姿并反馈误差，才形成真正闭环。若用阻尼逆，主任务本身会有近似误差；若关节速度限幅被触发，还要重新判断实际能够实现的任务速度。

MoveIt Servo 在 Humble 文档中提供带时间戳的关节点动和末端速度输入，并包含奇异和碰撞检查。它能帮助实现反应式运动，但具体接口取决于版本；这些软件检查也不等于对真实环境的无限可靠保证。\cite{moveit}

\section{延迟怎样变成相位损失和空间误差}

固定延迟 $T_d$ 的频率响应为
\begin{equation}
G_d(j\omega)=e^{-j\omega T_d},\qquad |G_d|=1,\qquad
\angle G_d=-\omega T_d.
\end{equation}
延迟不衰减幅值，却使反馈“看见过去”。在 $5\,\mathrm{Hz}$ 上，$40\,\mathrm{ms}$ 延迟造成 $-2\pi\times5\times0.04=-1.2566\,\mathrm{rad}=-72^\circ$ 相位；如果一个闭环经过两次这种延迟，总计就可能损失 $144^\circ$。真实链路还包含滤波、采样和执行器，不应只测网络 ping。

\begin{figure}[htbp]\centering
\includegraphics[width=0.92\linewidth]{figures/rt_latency.pdf}
\caption{同一正弦运动在不同延迟下的相位差，以及延迟造成的相位损失。频率越高，同样毫秒数的影响越大。}
\end{figure}

在低速匀速近似下，位置落后量为 $\delta p\approx vT_d$。以 $v=0.2\,\mathrm{m/s},T_d=0.1\,\mathrm s$ 为例，落后为 $0.02\,\mathrm m$。假定停止指令到达后还能以 $1\,\mathrm{m/s^2}$ 恒定制动，制动距离为 $v^2/(2a)=0.02\,\mathrm m$，总计约 $4\,\mathrm{cm}$。这是给定假设下的计算，不是任何设备的统一停止距离。

\section{双边力反馈需要同时看运动和能量}

无力反馈时，人通过图像判断远端状态。加入力反馈后，远端接触力又作用到人的手，形成一个包含主端、从端、操作者、环境和通信的闭环。位置误差型反作用力和直接环境力反射不是同一种反馈：前者连自由空间中的跟踪滞后也会反映成“手感很重”。

对位置比例关系 $\delta p_s=s_pR_{BM}\delta p_m$，相同虚功要求反射到手上的环境力为
\begin{equation}
F_{m,\rm display}=s_pR_{BM}^{\mathsf T}F_{s,\rm env}.
\end{equation}
因为 $F_{m,\rm display}^{\mathsf T}\delta p_m=F_{s,\rm env}^{\mathsf T}\delta p_s$。其中 $F_{s,\rm env}$ 指环境施加在从端上的力，反射方向应让操作者感到阻碍，不能误把机器人施加给环境的反向力直接照搬。人为改变力缩放当然可以，但此时不能再声称这个映射天然保持相同功率。

\subsection{波变量为什么能处理理想恒定延迟}

\textbf{无源（passive）}系统不会凭空向外输出超过初始储能和此前输入的能量。对一个通信端口，令 $F$ 为流入通道的力，$v$ 为按端口约定的速度，选择正的\textbf{波阻抗} $Z_w$，单位 $\mathrm{N\,s/m}$。定义
\begin{equation}
a^+=\frac{F+Z_wv}{\sqrt{2Z_w}},\qquad
a^-=\frac{F-Z_wv}{\sqrt{2Z_w}}.
\end{equation}
两式平方相减，分子中 $F^2$ 和 $Z_w^2v^2$ 抵消，留下
\begin{equation}
Fv=\tfrac12\big[(a^+)^2-(a^-)^2\big].
\end{equation}
功率变成输入波能量与输出波能量之差。理想、恒定且非负延迟只把已经进入的波延后输出，可以把未到达部分解释成通道中的储能。因此，在相应端口连接和无源终端假设下，通信通道能保持无源。时变延迟、丢包处理、采样保持以及非无源控制器还需额外设计。它并不保证“无论怎样写软件，任何延迟都稳定”。\cite{niemeyer}

\begin{examplebox}[验证功率恒等式]
取 $F=4\,\mathrm N,v=0.2\,\mathrm{m/s},Z_w=20\,\mathrm{N\,s/m}$。
$Z_wv=4\,\mathrm N$，于是 $a^+=8/\sqrt{40}=1.26491$，$a^-=0$。
波能量差为 $1.26491^2/2=0.8\,\mathrm W$，恰好等于 $Fv=4\times0.2=0.8\,\mathrm W$。
\end{examplebox}

\subsection{时域无源控制的符号必须讲清}

用采样周期 $T_s$ 累积端口流入能量 $E_n=E_0+\sum_{k=0}^nF_kv_kT_s$。如果某一步试算 $E_n<0$，表示当前离散实现输出了账面上没有的能量。对 $v_n\ne0$，可添加耗散系数
\begin{equation}
b_{\rm pc}=\frac{-E_n}{T_sv_n^2}.
\end{equation}
把通道端口力改为 $F_n+b_{\rm pc}v_n$，对机器人或人产生的反作用阻尼则是 $-b_{\rm pc}v_n$。两个符号属于相反的端口，不能混用。

例如 $E_n=-0.002\,\mathrm J,T_s=0.001\,\mathrm s,v_n=0.2\,\mathrm{m/s}$，需 $b_{\rm pc}=50\,\mathrm{N\,s/m}$，对应反作用力 $-10\,\mathrm N$，该步耗散 $10\times0.2\times0.001=0.002\,\mathrm J$。当速度接近零时，这个除法会产生很大增益；真实实现需要一致的离散能量账本、限幅和近零速度处理。若力已被限幅，不能仍记账为“所需能量已全部耗散”。

\section{视觉伺服把图像误差变成运动}

\textbf{视觉伺服（visual servoing）}使用相机观测闭环控制运动。\textbf{PBVS}先估计三维位姿再控制；\textbf{IBVS}直接控制图像特征。这里推导一个点特征的最小例子，说明它与雅可比的关系。\cite{corke}

设固定场景点在相机坐标中为 $P=(X,Y,Z)$，$Z>0$。归一化成像坐标 $u=X/Z,v=Y/Z$；像素坐标可由焦距和主点换算得到。若相机平移速度为 $(v_x,v_y,v_z)$、角速度为 $(\omega_x,\omega_y,\omega_z)$，都表达在相机系，固定点相对相机满足
\begin{equation}
\dot P=-v_c-\omega_c\times P.
\end{equation}
先对 $u=X/Z$ 用商法则：$\dot u=(\dot XZ-X\dot Z)/Z^2$；$v$ 同理。逐项代入叉乘后得到
\begin{equation}
\begin{bmatrix}\dot u\\\dot v\end{bmatrix}
=L\begin{bmatrix}v_x\\v_y\\v_z\\\omega_x\\\omega_y\\\omega_z\end{bmatrix},\quad
L=\begin{bmatrix}
-1/Z&0&u/Z&uv&-(1+u^2)&v\\
0&-1/Z&v/Z&1+v^2&-uv&-u
\end{bmatrix}.
\end{equation}
$L$ 叫\textbf{交互矩阵（interaction matrix）}，又称图像雅可比。它把相机速度映射到图像特征速度。一个点只有两条观测方程，无法唯一决定六维相机速度，需要多个特征或约束部分自由度。

\begin{examplebox}[只沿相机横向移动]
假定只允许 $v_x$，目标 $u_d=0$，当前 $u=0.1$，深度 $Z=1\,\mathrm m$，取误差衰减率 $\lambda=2\,\mathrm{s^{-1}}$。要求 $\dot u=-\lambda(u-u_d)=-0.2\,\mathrm{s^{-1}}$。
因为 $\dot u=-v_x/Z$，所以 $v_x=0.2\,\mathrm{m/s}$。
相机向右移，静止目标在图像中向左移，符号由几何决定。
若深度实际为 $2\,\mathrm m$，同一速度只产生 $-0.1\,\mathrm{s^{-1}}$ 的特征速度，说明深度误差会改变闭环响应。
\end{examplebox}

要控制机械臂，还需把相机速度映射到关节：$\dot s_{
m img}=LJ_c\dot q$，其中 $J_c$ 必须是相机参考点、相机坐标轴、线速度在前角速度在后的雅可比。相机和工具存在偏置时，不能直接拿工具原点雅可比代替。视觉通常还存在曝光、图像传输和特征计算延迟，末端快速稳定不能只依赖低频图像反馈。

\section{计算练习与解答}

\subsection*{题一 延迟相位}
闭环关键频率为 $8\,\mathrm{Hz}$，额外延迟为 $15\,\mathrm{ms}$。求相位损失。

解：$\omega T_d=2\pi\times8\times0.015=0.75398\,\mathrm{rad}$，换成角度为 $43.2^\circ$。这是额外损失，不是整个系统的相位裕度。

\subsection*{题二 坐标和缩放}
主端位移为 $(0.1,-0.2,0)\,\mathrm m$，$R_{BM}=R_z(90^\circ),s_p=0.5$。求从端位移。

解：旋转后为 $(0.2,0.1,0)$，缩放后为 $(0.1,0.05,0)\,\mathrm m$。旋转保留长度，缩放把长度减半。

\subsection*{题三 图像误差}
归一化横坐标误差为 $0.05$，深度为 $0.8\,\mathrm m$，只控制横移，$\lambda=1.5\,\mathrm{s^{-1}}$。求横向相机速度。

解：$\dot u=-1.5\times0.05=-0.075\,\mathrm{s^{-1}}$。由 $\dot u=-v_x/Z$，得 $v_x=0.8\times0.075=0.06\,\mathrm{m/s}$。这是相机坐标系速度，仍需通过相机雅可比转成关节命令。

\section*{本章与原稿的对应}
本章对应 [B，第97--100页]，并延伸其相机云台应用。波变量和无源控制补出端口方向及功率推导；视觉伺服为整合新增内容，用于把坐标变换、雅可比和相机任务连接起来。

\chapter{把算法接到 ROS 2 驱动器和实时循环}
\label{ch:software}

控制公式只规定理想的输入输出关系，工程实现还必须规定数据属于哪个时刻、哪个关节、哪个坐标系，以及谁拥有写命令的权利。本章以固定基座机械臂为例，把数学变量和软件接口逐个对应起来。软件接口说明按 ROS 2 Humble 体系组织；不同发行版和驱动器的实现细节不能直接混用。

\section{先定义接口 再选择框架}

一个可复用控制器至少需要三类结构：测量、参考和输出。测量是 $\{q,\dot q,\tau_{\rm meas},t_{\rm sample},\mathrm{valid}\}$；参考是 $\{q_d,\dot q_d,\ddot q_d,T_d,\mathrm{mode}\}$ 中本控制器实际需要的量；输出是明确单位的关节命令。不是每个机器人都能提供真实关节力矩，也不是每个控制器都需要全部字段。

\begin{table}[htbp]\centering\small
\begin{tabularx}{\linewidth}{p{28mm}XX}\toprule
边界 & 必须写清楚的约定 & 数学意义\\\midrule
编码器到状态 & 零位 方向 角度单位 时间戳 & $q$ 不是未经转换的计数值\\
动力学到力矩 & 电机侧或关节侧 重力方向 & $\tau$ 与 $M\ddot q$ 同单位\\
任务到关节 & 末端参考点 坐标轴 行顺序 & $\dot x=J\dot q$ 必须同约定\\
控制器到驱动 & 接口模式 内置增益 补偿所有权 & 避免同一 PD 或重力补偿做两次\\
线程之间 & 快照一致性 数据年龄 & 防止把不同拍关节拼成一个状态\\\bottomrule
\end{tabularx}
\caption{这些约定是控制算法的一部分，不能只写在口头说明里。}
\end{table}

例如关节消息顺序为 $[j_3,j_1,j_2]$，值为 $[0.3,0.1,0.2]$，模型顺序为 $[j_1,j_2,j_3]$，正确输入应是 $[0.1,0.2,0.3]$。直接把消息数组复制到模型向量不会报矩阵维度错误，却会让正运动学和重力补偿全部算错。

\section{ROS 2 的节点 消息和坐标树}

\textbf{节点（node）}是逻辑计算单元，不等于操作系统进程；多个节点可以组合在同一进程中。\textbf{话题（topic）}提供发布订阅流；\textbf{服务（service）}适合请求与响应；\textbf{动作（action）}适合有持续时间、进度反馈和取消语义的任务，例如执行一条轨迹。机器人模型与坐标变换不是实时运动命令。

\textbf{URDF}描述连杆、关节、几何、惯量和部分限位；\textbf{xacro}帮助复用模型描述；\textbf{tf2}维护带时间信息的坐标变换树。某个变换数值正确但时间不匹配，仍可能把移动相机的图像放到错误位置。仿真时相关节点应使用一致的时间基准；实际时钟和仿真时钟不能随意相减。

\textbf{QoS（quality of service，服务质量）}规定可靠性、历史队列和持久性等通信策略。可靠传输关注是否重传，实时控制还关注这份数据是否已经过期；队列越长不一定越好。订阅端所请求的 QoS 与发布端提供的 QoS 不兼容时，可能无法连接。可以检查端点信息和事件，而不是只看话题名字存在。

下列命令用于只读检查接口；其中话题名来自具体系统，不能当成所有 ROS 2 机器人的标准控制接口：
\begin{lstlisting}[language=bash]
ros2 node list
ros2 topic list -t
ros2 topic info /control/move_j --verbose
ros2 topic type /control/move_j
ros2 interface show sensor_msgs/msg/JointState
ros2 topic echo /feedback/joint_states --once
\end{lstlisting}
发布运动之前，要先查真实消息类型、关节名、订阅者和接口语义。消息叫 \texttt{JointState} 并不证明它只能用作状态，也不证明任意填一个 \texttt{position} 数组就能执行；有些厂家把标准消息用于私有命令约定，解释权在驱动实现。

\section{ros2\_control 的一个周期}

ros2\_control 把硬件通信与控制算法分离。硬件提供\textbf{状态接口}和\textbf{命令接口}；控制器读取状态并写命令；控制器管理器组织生命周期和周期更新。典型循环为读取硬件、计算控制、写出命令。官方 Humble 文档把 \texttt{update\_rate} 定义为这条更新循环的频率，同时强调需要控制抖动。\cite{ros2control}

硬件插件不只是两个孤立的函数。它还要完成初始化、资源分配、激活、停用、错误处理和命令模式切换。控制器激活时应把位置目标与当前状态对齐，并按控制律处理积分器和滤波器状态；仅把所有内部量清零不一定无扰。

\subsection{轨迹控制器到底把什么交给驱动器}

Humble 的关节轨迹控制器按照时间对航点插补。位置命令接口通常转发插补后的位置；速度或力矩接口可通过配置的 PID 把位置和速度跟踪误差映射为命令。可用接口组合和状态反馈要求应以对应版本文档为准。\cite{jointtrajectory}

因此，同样一条关节轨迹，在位置接口和力矩接口下的闭环结构不同。位置接口下面通常还有驱动器的内部伺服；力矩接口则要求上层承担更多动态控制职责。把上层输出的 $-K_pe-K_d\dot e$ 再送进另一个含同样 PD 的接口，可能叠加刚度和阻尼。

以常见组合式关节接口为例，若驱动器实现
\begin{equation}
\tau=K_{p,\rm drv}(q_r-q)+K_{d,\rm drv}(v_r-\dot q)+\tau_{\rm ff},
\end{equation}
则上位机可以选择让驱动器承担局部 PD，自己计算前馈；也可以在协议允许且行为已确认时，使用合适设置把主要力矩律放在上位机。必须明确两个层级各自做什么。协议中的增益数值范围、编码比例、饱和和电流限幅，不能由这个公式反推出来。

\section{MoveIt 与动力学库的分工}

MoveIt 的模型、规划场景和规划管线负责运动规划相关工作；ros2\_control 组织执行接口；Pinocchio 这类动力学库负责计算运动学、惯量、非线性项和雅可比。库能很快计算公式，但不会自动决定你的控制器应该用哪一种误差、哪一种摩擦补偿或哪一种安全停机策略。

\begin{table}[htbp]\centering\small
\begin{tabularx}{\linewidth}{p{34mm}XX}\toprule
配置或数据 & 主要用途 & 与其它部分的边界\\\midrule
URDF 及碰撞几何 & 机构 轴向 惯量 模型限位 & CAD 惯量不等于已完成辨识\\
SRDF & 规划组 命名姿态 允许碰撞对 & 禁止滥用允许碰撞掩盖几何错误\\
运动学配置 & IK 插件 超时 搜索策略 & 七轴不意味着某个插件必然不能用\\
规划限值 & 速度 加速度等 & 与驱动器硬限值保持一致和余量\\
执行配置 & 控制器名字 动作接口 关节集合 & 规划成功不等于执行已成功\\
工具附件 & 相机或抓持物的碰撞形状 & 只标 TCP 不会自动添加碰撞体\\\bottomrule
\end{tabularx}
\caption{软件配置表达的是算法前提，出现问题时要回到其物理含义。}
\end{table}

选择 IK 插件时，应按实际模型、限位、目标分布、初值和超时，比较成功率及耗时。原稿“七轴不要用 KDL”的说法过于绝对；适用性应由具体任务与求解表现判断。

\subsection{模型配置维度和速度维度}

配置向量的维度记为 $n_q$，切空间速度维度记为 $n_v$。普通有界转动关节各占一个配置量和一个速度量；若某库用正余弦表示无限转动关节，配置量可能有两个。浮动基座常用位置加单位四元数存七个数，速度则是三维线速度加三维角速度，共六个数。

因此一般的配置更新是 $q_{k+1}=\operatorname{integrate}(q_k,v_kT_s)$，不是不加区分地执行 $q+vT_s$。对本书固定基座、普通角度坐标的 2R 模型，普通加法可以成立；推广时需重新确认库的表示约定。

同样，\texttt{LOCAL}、\texttt{WORLD} 和\texttt{LOCAL\_WORLD\_ALIGNED} 的区别不仅是坐标轴旋转，还涉及空间速度参考点。控制末端点的线速度时，需要末端参考点的雅可比；绕世界原点表达的空间旋量，其前三项未必等于末端点的实际线速度。第3章的旋量变换给出了这一区别的数学原因。

\section{实时性用时间预算解释}

设控制周期为 $T_s=5\,\mathrm{ms}$，对应 $200\,\mathrm{Hz}$。一个教学预算为读取 $0.7\,\mathrm{ms}$、计算 $0.9\,\mathrm{ms}$、写出 $0.5\,\mathrm{ms}$、其它固定开销 $0.2\,\mathrm{ms}$，总计 $2.3\,\mathrm{ms}$，余量为 $2.7\,\mathrm{ms}$。若偶发一次文件写入耗时 $8\,\mathrm{ms}$，平均耗时依旧可能很好看，但这一拍已经错过截止时间。

\begin{figure}[htbp]\centering
\includegraphics[width=0.93\linewidth]{figures/rt_timing.pdf}
\caption{五毫秒周期的时间预算示例。绝对时间调度可以避免把每次计算时间累积进下一次周期，但不能消除过载、抢占和通信阻塞。}
\end{figure}

\textbf{周期}是理想更新间隔；\textbf{唤醒抖动}是实际开始时刻偏离计划开始时刻的量；\textbf{执行耗时}是本轮代码实际用了多久；\textbf{端到端延迟}从传感器采样一直算到执行器产生作用。它们相关但不相同。

若每轮先计算 $0.8\,\mathrm{ms}$ 再固定睡眠 $5\,\mathrm{ms}$，实际周期至少为 $5.8\,\mathrm{ms}$，频率约为 $172.4\,\mathrm{Hz}$。采用绝对时刻 $t_k=t_0+kT_s$ 作为下一次唤醒目标，可以避免这类漂移；某轮超时后则必须记录并按设计处理，不能用“快速补跑多次旧命令”伪装没超时。

热路径中应避免未经界定的动态分配、可能阻塞的输入输出和无上界等待。采用预分配存储、一致的数据快照、有界队列和经过验证的线程交互。无锁本身也不是正确性的保证：内存顺序和读写一致性仍需证明。普通 ROS 2 话题并不自动提供硬实时保证，但能否满足具体周期取决于体系结构与测量，不能用“所有 ROS 消息永远不能进入实时系统”替代分析。

\section{CAN 带宽先算账再谈频率}

只做一个明确假设的估算：每个关节每周期一帧命令、一帧反馈，每帧连同开销平均按 $130\,\mathrm{bit}$ 估计，总线速率 $1\,\mathrm{Mbit/s}$。七关节、$200\,\mathrm{Hz}$ 时
\begin{equation}
B=7\times2\times200\times130=364000\,\mathrm{bit/s},
\end{equation}
平均占用约 $36.4\%$。改为 $1000\,\mathrm{Hz}$ 时，需求为 $1.82\,\mathrm{Mbit/s}$，已超过单路总线能力。这个算例不代表某款机械臂的实际协议：位填充、帧格式、打包、多路总线、额外状态报文和错误重发都会改变结果。

平均带宽还不等于最坏延迟。仲裁会让低优先级帧等待；同一拍七个关节反馈不一定同时到达。需要记录每轴反馈年龄，必要时做有说明的状态外推，并给异常数据设置明确的处理行为。总线速率高，不代表控制闭环带宽也相同。

\section{编码器 电气接口和力矩单位}

编码器是传感器，不会自己驱动机械臂。增量式编码器的 A、B 两路有相位差，用边沿先后判断方向；Z 通常是每转一次的参考标记。若每路每转 $N=2500$ 个周期，采用四倍频计数，一转有 $10000$ 个计数，对应
\begin{equation}
\Delta q=2\pi/10000=0.00062832\,\mathrm{rad}=0.036^\circ.
\end{equation}
接在电机端时，还需按减速比换算到关节端，并考虑零位与传动弹性。标称分辨率不等于实际准确度。

A+/A-、B+/B-、Z+/Z- 通常表示差分信号对。DB9 只是连接器形状，把导线焊到 DB9 上不会自动变成 RS232 串行数据。差分接收、正交计数、电源和逻辑电平匹配是不同功能；接口类型必须根据设备输出电路确定。普通串口适配器不会仅凭插头形状就读出角度。

力矩换算也要保留层级。电机轴 $\tau_m=K_ti_q$，理想减速器输出约为 $\tau_j=\eta N\tau_m$。如果驱动反馈已经是关节侧力矩，再乘一次 $N$ 就会放大几十或上百倍。厂家的 $k,b,c$ 等系数名称没有跨厂家的统一含义，要把每个系数的单位、标定依据和已经包含的变换写下来。

\section{控制状态机与命令所有权}

一个教学状态机可以包含未激活、保持、轨迹、拖动、受控停止和故障。状态转移同时改变参考生成器、控制律和内部状态，不能只把字符串变量换掉。每个关节在同一时刻只能有明确的一套最终命令来源；“规划器还在发位置，同时拖动程序又发 MIT 力矩”不是控制融合，而是竞争写入。

停止也不是把所有数值写成零。重力存在时，零力矩可能导致下落；位置置零可能发起大幅回零运动。正常减速停止、失联处理、硬件紧急停止和断电保持属于不同机制，其具体动作由机械结构、抱闸和驱动器设计决定。教材不规定某款设备通用的使能延时或抱闸时序。

\section{计算练习与解答}

\subsection*{题一 总线预算}
采用上述每关节收发各一帧、每帧 $130\,\mathrm{bit}$ 假设，七关节运行 $500\,\mathrm{Hz}$，占用多少带宽？

解：$7\times2\times500\times130=910000\,\mathrm{bit/s}$，即 $91\%$。这尚未包含额外报文和重发，不能只因为小于 $100\%$ 就认定有充分余量。

\subsection*{题二 电机侧与关节侧}
已知 $K_t=0.08\,\mathrm{N\,m/A},i_q=3\,\mathrm A,N=50,\eta=0.9$。求电机和关节侧力矩。

解：$\tau_m=0.08\times3=0.24\,\mathrm{N\,m}$，$\tau_j=0.9\times50\times0.24=10.8\,\mathrm{N\,m}$。这个简化静态换算不含方向相关摩擦和动态损耗。

\subsection*{题三 周期与数据年龄}
控制周期为 $5\,\mathrm{ms}$，最新反馈采样时间距当前已过去 $12\,\mathrm{ms}$。若关节可能以 $0.5\,\mathrm{rad/s}$ 运动，忽略加速度，位置可能相差多少？

解：$0.5\times0.012=0.006\,\mathrm{rad}\approx0.344^\circ$。它跨越了 $2.4$ 个周期，不能只检查“这一拍是否收到一条消息”，还要检查消息对应的真实采样时刻。

\section*{本章与原稿的对应}
本章整合 [A，第67--74页] 与 [B，第100--111、116--120页]，并补充编码器接口与带宽计算。纠正节点等于进程、ROS 通信天然实时、七轴必须排除某个 IK 插件、停机等于写零以及不同 MoveIt 版本接口混用等问题。

\chapter{仿真模型 数值积分与真实误差}
\label{ch:simulation}

仿真应回答具体问题：公式是否写对，模型误差会造成多大偏差，传感器和执行器的非理想因素是否能解释真实现象。只有动画动起来，无法回答这些问题。本章把“对象真值”“控制器模型”“测量”和“评价量”分开，让同一套仿真既可以做正确性核验，也可以研究模型失配。

\section{正动力学如何推进一个时间步}

忽略接触时，动力学为
\begin{equation}
M(q)\ddot q=\tau_{\rm act}+\tau_{\rm ext}-C(q,\dot q)\dot q-g(q)-\tau_f.
\end{equation}
先解线性方程求加速度，再用数值积分更新速度和位置。\textbf{数值积分器}不是控制器，它只是近似连续微分方程的演化。即使控制律相同，不同积分器和步长也可能产生不同的数值误差。

取一维关节 $J=0.1\,\mathrm{kg\,m^2}$，当前驱动力矩为 $1$，黏性摩擦 $b\dot q=0.2\times0.5=0.1$，此处重力矩为 $0.4\,\mathrm{N\,m}$，则
\[
\ddot q=(1-0.1-0.4)/0.1=5\,\mathrm{rad/s^2}.
\]
用半隐式欧拉、步长 $h=0.001\,\mathrm s$，从 $q=0.2,\dot q=0.5$ 出发：
\[
\dot q^+=0.5+0.001\times5=0.505,
\quad q^+=0.2+0.001\times0.505=0.200505.
\]
若在这一小步把加速度当常数，精确运动学为 $q^+=0.2+0.5h+5h^2/2=0.2005025$。差值 $2.5\times10^{-6}\,\mathrm{rad}$ 是离散积分误差，不是控制器突然变差。

\section{仿真周期与控制周期为什么要分开}

设控制器 $T_s=5\,\mathrm{ms}$ 更新一次，物理仿真步长 $h=0.5\,\mathrm{ms}$，则每十个物理步计算一次新命令，其间保持上一命令，这就是\textbf{零阶保持}。传感器采样、网络到达和执行器生效还可以各有自己的时间。

较小物理步长有助于解析执行器和接触的快动态，但并不是“必须永远小于控制周期”的数学定理。恰当的离散模型可以直接按控制周期仿真；关键是明确离散化和保持关系。若为了计算更精细而每个物理小步都更新控制器，就把本来 $200\,\mathrm{Hz}$ 的控制器悄悄变成了 $2000\,\mathrm{Hz}$，结论会偏乐观。

\begin{figure}[htbp]\centering
\includegraphics[width=0.94\linewidth]{figures/rt_zoh_actuator.pdf}
\caption{控制命令零阶保持与一阶执行器的实际力矩。命令更新、传输延迟和力矩建立是不同过程；给仿真增加一个“延迟参数”时要说明它放在哪一段。}
\end{figure}

一个与具体仿真器无关的调度结构如下。这里的代码是结构示意，传感器和执行器封装应有实际实现。
\begin{lstlisting}[language=Python]
for physics_step in range(n_steps):
    if physics_step % control_stride == 0:
        measurement = sensors.sample_and_delay(plant_state)
        reference = trajectory.sample(control_time)
        command = controller(measurement, reference, nominal_model)
    applied_torque = actuator.update(command, physics_dt)
    plant_state = plant.step(applied_torque, physics_dt)
    evaluator.record(plant_state, measurement, command)
\end{lstlisting}
控制器只拿到约定测量，评价器才读取真值；两条数据通道不能在不知情时混合。

\section{执行器一阶滞后可以精确离散}

用 $T_a>0$ 表示力矩建立时间常数：
\begin{equation}
T_a\dot\tau_a+\tau_a=\tau_c.
\end{equation}
若一小步内命令 $\tau_c$ 恒定，令误差 $z=\tau_a-\tau_c$，得到 $\dot z=-z/T_a$，所以
\begin{equation}
\tau_a[k+1]=e^{-h/T_a}\tau_a[k]+(1-e^{-h/T_a})\tau_c[k].
\end{equation}
取 $T_a=3\,\mathrm{ms},h=0.5\,\mathrm{ms}$，系数 $e^{-1/6}=0.846482$。命令从零跳到 $10\,\mathrm{N\,m}$，第一步实际力矩只有 $1.53518\,\mathrm{N\,m}$；经过 $5\,\mathrm{ms}$ 达到 $10(1-e^{-5/3})=8.11124\,\mathrm{N\,m}$。若还有一拍通信延迟，就先把命令放入延迟队列，再送进这个动态方程。

力矩限幅、力矩变化率限幅、速度相关降额和温度降额是不同约束。把它们都简化成同一个常数截断，只能作为第一版模型，不能声称模拟了驱动器全部行为。

\section{测量噪声为什么让加速度特别难看}

对于量化间隔 $\Delta q$，若量化误差近似独立均匀分布于 $[-\Delta q/2,\Delta q/2]$，其方差为 $\Delta q^2/12$。二阶差分
\begin{equation}
\hat{\ddot q}_k=\frac{q_{k+1}-2q_k+q_{k-1}}{T_s^2}
\end{equation}
的噪声方差为
\begin{equation}
\sigma_{\ddot q}^2=\frac{(1^2+(-2)^2+1^2)\Delta q^2}{12T_s^4}
=\frac{\Delta q^2}{2T_s^4}.
\end{equation}
十七位每转编码器的 $\Delta q=2\pi/2^{17}=4.79369\times10^{-5}\,\mathrm{rad}$。在 $T_s=1\,\mathrm{ms}$ 时，标准差约 $33.90\,\mathrm{rad/s^2}$；在 $T_s=5\,\mathrm{ms}$ 时约 $1.356\,\mathrm{rad/s^2}$。这些数依赖独立量化误差假设，真实静止编码器可能表现为相关或零变化误差，但足以说明“采样越快，直接差分就越准”不成立。

离线辨识可以用轨迹拟合、平滑微分或积分形式回归；实时速度估计可以用滤波器或状态观测器。把未来点参与的中心差分直接写进实时控制器，则需要明确它等价于引入等待延迟。

\section{MuJoCo 中的模型和状态}

MuJoCo 使用 \texttt{mjModel} 保存编译后的模型，\texttt{mjData} 保存状态和计算工作区；\texttt{mj\_step} 推进一步动力学，\texttt{mj\_forward} 在当前状态计算派生量，\texttt{mj\_inverse} 处理逆动力学。模型和数据对象应相互对应。\cite{mujocosim}

原生 MJCF 通过嵌套 body 表达机构树，关节定义该 body 相对父 body 的运动自由度。\texttt{inertial} 给出质量、质心和惯量；\texttt{armature}、\texttt{damping}、\texttt{frictionloss} 分别用于附加关节惯量、黏性阻尼和干摩擦相关约束。它们的求解语义需要分开，不应把所有字段都理解成可随意相加的同一种力矩。\cite{mujocomodel}

下面是用于本书平面 2R 思考实验的模型核心。角度用弧度，重力沿 $-y$；几何用于显示，动力学惯量显式指定。沿杆轴的小惯量 $I_{xx}$ 用于构造合法三维刚体，绕运动平面法向的转动采用 $I_{zz}=1/48\,\mathrm{kg\,m^2}$。
\begin{lstlisting}[language=XML]
<mujoco model="teaching_2r">
  <compiler angle="radian"/>
  <option timestep="0.0005" gravity="0 -9.81 0"/>
  <worldbody>
    <body name="link1">
      <joint name="q1" type="hinge" axis="0 0 1"/>
      <inertial pos="0.25 0 0" mass="1"
          diaginertia="0.0001 0.0208333333 0.0208333333"/>
      <geom type="capsule" fromto="0 0 0 0.5 0 0"
          size="0.015" contype="0" conaffinity="0"/>
      <body name="link2" pos="0.5 0 0">
        <joint name="q2" type="hinge" axis="0 0 1"/>
        <inertial pos="0.25 0 0" mass="1"
            diaginertia="0.0001 0.0208333333 0.0208333333"/>
        <geom type="capsule" fromto="0 0 0 0.5 0 0"
            size="0.015" contype="0" conaffinity="0"/>
        <site name="tip" pos="0.5 0 0"/>
      </body>
    </body>
  </worldbody>
  <actuator>
    <motor joint="q1" gear="1"/>
    <motor joint="q2" gear="1"/>
  </actuator>
</mujoco>
\end{lstlisting}
在这个模型中关闭了接触，目的是先对账自由空间动力学。它不包含真实电机、线缆、抱闸或传动弹性。控制输入在上述关节与 \texttt{gear=1} 约定下可解释为关节力矩；一般执行器的 \texttt{ctrl} 含义由传动和增益共同决定，不能对所有模型都直接赋值 $\tau$。

在 $q=(0,\pi/2)$、速度零时，本书的手算结果为
\begin{equation}
M=\begin{bmatrix}5/12&1/12\\1/12&1/12\end{bmatrix},\qquad
g=\begin{bmatrix}7.3575\\0\end{bmatrix}.
\end{equation}
模型核验阶段可以读取仿真器的质量矩阵和偏置项来比较；在评价未知参数辨识或模型失配的控制器时，则应禁止控制器读取这些真值。核验阶段知道答案是合理的，实验阶段把答案偷偷传给待评价算法才是问题。

\section{接触参数不是简单的材料常数复制}

MuJoCo 的接触约束采用带正则化的软接触模型。\texttt{solref} 和 \texttt{solimp} 共同影响约束参考动力学和阻抗，且受接触几何、有效质量及求解设置影响。直接参数形式中的“刚度”和“阻尼”不能不经核对就认作材料的 $\mathrm{N/m}$ 与 $\mathrm{N\,s/m}$ 常数。\cite{mujocomputation}

若需要研究一个力控算法，先用明确的教学弹簧阻尼接触模型建立预期。例如法向压入量 $\delta>0$ 时 $F_n=k_e\delta+d_e\dot\delta$，必要时施加单向非拉力约束；取 $k_e=10000\,\mathrm{N/m}$、静态压入 $1\,\mathrm{mm}$，弹簧力为 $10\,\mathrm N$。然后再对照仿真器的接触力和压入量标定参数。不要因为动画中“已经碰到墙”，就默认法向力和接触刚度正确。

真实力传感器对应的仿真测量，可以来自接触力或传感器接口，并添加相应带宽、偏置和噪声。若现实中没有该传感器，却把无噪声的接触真值送进控制器，就改变了问题的可观测条件。反过来，现实有力传感器时，完全禁止仿真控制器使用任何接触测量也不合理；关键是传感器能力一致。

\section{真值模型与控制器模型的分离}

建立三套概念上独立的参数：名义参数 $\theta_{\rm nom}$，仿真对象真值 $\theta_{\rm true}$，以及数据辨识结果 $\hat\theta$。真值可以按实验设计扰动，但惯量必须保持物理一致，不能给每个惯量元素独立乘随机数后就认定仍是合法刚体。

最简单的模型误差算例是常值重力残差。考虑
\[
J\ddot q+b\dot q+g_0=-K_p(q-q_d)-K_d\dot q+\hat g_0.
\]
在静态平衡，
\begin{equation}
e_\infty=q-q_d=(\hat g_0-g_0)/K_p.
\end{equation}
取 $g_0=2,\hat g_0=1.8\,\mathrm{N\,m},K_p=8\,\mathrm{N\,m/rad}$，稳态误差为 $-0.025\,\mathrm{rad}$。若错误地把真值 $g_0$ 给了控制器，这个偏差就消失，仿真会告诉你“模型补偿完美”，但没有评价你真正关心的误差。

\begin{figure}[htbp]\centering
\includegraphics[width=0.92\linewidth]{figures/rt_model_mismatch.pdf}
\caption{同一关节控制器在不同重力补偿误差下的响应。图为本教材的一维教学模型数值积分，曲线不代表 NERO 实测。末端稳态偏差与残余力矩除以刚度一致。}
\end{figure}

\section{如何证明结论不是步长造成的}

固定控制周期、模型和输入，只将物理积分步长从 $h$ 改为 $h/2$，比较轨迹误差、接触峰值和能量残差。如果关键结论大幅改变，说明还未达到足够的数值分辨率，或接触/控制存在快动态。不能在缩小物理步长时顺手加快控制器，再把改进全归功于积分器。

对无阻尼、无输入的保守系统，比较总机械能 $E=\tfrac12\dot q^\mathsf TM\dot q+U(q)$。对有力矩和摩擦的系统，比较能量收支残差
\begin{equation}
r_E(t)=E(t)-E(0)-\int_0^t\dot q^\mathsf T(\tau_{\rm act}+\tau_{\rm ext}-\tau_f)\,dt.
\end{equation}
能量不恒定不一定是错误：执行器做功和阻尼耗能本来就会改变能量。自由摆动仿真中“能量完全不变”也不能证明所有模型参数正确，它只检查了某些数值与结构性质。

\section{计算练习与解答}

\subsection*{题一 多速率仿真}
控制频率 $200\,\mathrm{Hz}$，物理步长 $0.25\,\mathrm{ms}$。每多少物理步更新控制？十秒有多少次控制更新？

解：控制周期 $5\,\mathrm{ms}$，故步数比为 $5/0.25=20$。十秒共有 $10/0.005=2000$ 个完整控制周期；起点是否额外记录一帧属于记录约定，不应算成额外持续时间。

\subsection*{题二 量化和差分}
若编码器分辨率增加一位，采样周期不变，二阶差分量化噪声标准差怎样变化？

解：一转计数翻倍，$\Delta q$ 减半。由 $\sigma_{\ddot q}=\Delta q/(\sqrt2T_s^2)$，标准差也减半，方差变为四分之一。若同时把周期减半，标准差反而变为原来的两倍。

\subsection*{题三 补偿残差}
静态误差 $-0.04\,\mathrm{rad}$，刚度 $5\,\mathrm{N\,m/rad}$，忽略摩擦及外力，估计重力补偿残差。

解：$\hat g-g=K_pe=5\times(-0.04)=-0.2\,\mathrm{N\,m}$，表示补偿不足。若有静摩擦或外接触，这一计算只能得到合成残差，不能唯一指向重力模型。

\section*{本章与原稿的对应}
本章对应 [A，第72--73页] 与 [B，第111--116、125--127页]。补出传感器差分噪声、执行器离散化、多速率循环与模型分离的计算；修正接触参数直接等于材料常数、控制周期必须等于仿真步长等不够准确的表述。

\chapter{约束优化 MPC 全身控制与学习方法}
\label{ch:advanced}

前面的控制律多数写成一条显式公式。遇到力矩上限、关节边界、接触不滑动和多个任务同时存在时，可以把“希望做什么”和“必须满足什么”写进同一个优化问题。本章不把新名词当成自动升级，而是从能手算的二次问题出发，说明这些方法增加了什么能力以及需要承担什么代价。

\section{二次规划从一个变量开始}

\textbf{优化变量}是要求解的数；\textbf{目标函数}衡量好坏；\textbf{约束}规定可接受的范围。\textbf{二次规划（quadratic programming，QP）}常写成
\begin{equation}
\min_z\ \tfrac12z^\mathsf THz+c^\mathsf Tz,
\qquad Az=b,\quad Gz\le h.
\end{equation}
当 $H\succeq0$ 时目标是凸的，可行集合由仿射等式和线性不等式给出，所以局部最优也是全局最优；要保证解唯一，还需要在可行方向上有足够的正定性。凸不代表永远有可行解，也不代表求解时间总能满足控制周期。

先看 $\min_z\tfrac12(z-2)^2$，约束 $-1\le z\le1$。无约束驻点满足 $z-2=0$，得到 $z=2$，但不可行。可行区间内最优解为 $z=1$。这个约束在最优点成为等号，叫\textbf{活跃约束}。把输出先算成二再截成一，在这个简单标量问题中恰好得到约束最优解；在多维、有耦合和未来动态的情形中，逐分量截断通常不等价于重新优化。

\subsection{拉格朗日乘子和约束反作用}

约束为 $a^\mathsf Tz=b$ 时，引入乘子 $\lambda$：
\[
\mathcal L=\tfrac12z^\mathsf THz+c^\mathsf Tz+\lambda(a^\mathsf Tz-b).
\]
对 $z$ 和 $\lambda$ 求导，可写成
\begin{equation}
\begin{bmatrix}H&a\\a^\mathsf T&0\end{bmatrix}
\begin{bmatrix}z\\\lambda\end{bmatrix}
=\begin{bmatrix}-c\\b\end{bmatrix}.
\end{equation}
对不等式还需非负乘子与互补条件。力学中的约束反力和优化中的拉格朗日乘子有密切联系：二者都在保证约束成立时抵消不允许的变化，但具体单位和符号由约束写法决定。

\section{MPC 为什么只执行第一步}

\textbf{模型预测控制（model predictive control，MPC）}每个周期用当前测量作为起点，预测未来 $N$ 步，优化输入序列，再只执行第一步。下一次新测量到来后重新求解，因此它仍然是反馈控制。

对线性离散模型 $x_{k+1}=Ax_k+Bu_k$，可构造
\begin{align}
\min_{u_0,\ldots,u_{N-1}}\quad
&\sum_{k=0}^{N-1}(x_k^\mathsf TQx_k+u_k^\mathsf TRu_k)+x_N^\mathsf TPx_N,\\
\text{满足}\quad
&x_{k+1}=Ax_k+Bu_k,\quad x_0=x_{\rm measured},\\
&u_{\min}\le u_k\le u_{\max},\quad G_xx_k\le h_x.
\end{align}
这里可把 $x$ 定义为相对恒定参考的状态误差；移动参考时需要一致地扩展参考和动力学。$Q,R,P$ 是权重矩阵，不是系统物理惯量。

将预测递推展开：$x_1=Ax_0+Bu_0$，$x_2=A^2x_0+ABu_0+Bu_1$。堆叠以后 $X=\mathcal Ax_0+\mathcal BU$。代入二次代价，得到输入序列 $U$ 的 QP，这说明“预测”怎样变成普通矩阵计算。

\subsection{两步 MPC 的完整手算}

取一个离散位置误差模型 $e_{k+1}=e_k+u_k$，其中 $u$ 表示每拍的参考增量，单位也是弧度；它不是电机力矩。初始 $e_0=1$，两步预测，代价
\begin{equation}
J=e_1^2+e_2^2+u_0^2+u_1^2.
\end{equation}
代入 $e_1=e_0+u_0,e_2=e_0+u_0+u_1$，展开得
\begin{equation}
J=2e_0^2+4e_0u_0+2e_0u_1+3u_0^2+2u_0u_1+2u_1^2.
\end{equation}
两个偏导置零，除以二后为
\[
3u_0+u_1=-2e_0,\qquad u_0+2u_1=-e_0.
\]
因此 $u_0=-0.6e_0,u_1=-0.2e_0$，在 $e_0=1$ 时为 $(-0.6,-0.2)$，预测误差为 $(0.4,0.2)$，总代价为 $0.6$。

现在加每拍约束 $|u_k|\le0.3$。第一个输入不能取 $-0.6$。在 $u_0=-0.3$ 的边界上，对 $u_1$ 求导，得 $u_1=-(e_0+u_0)/2=-0.35$，仍超限，因此最优输入为 $(-0.3,-0.3)$。预测误差为 $(0.7,0.4)$，总代价为
\[
0.7^2+0.4^2+0.3^2+0.3^2=0.83.
\]
实际只执行 $u_0=-0.3$。若下一拍测得 $e=0.7$，重新优化后第一步仍为 $-0.3$，第二步变为 $-0.2$。若有扰动使实测变成 $0.8$，优化起点就改用 $0.8$，而不是盲目沿用原计划。

\begin{figure}[htbp]\centering
\includegraphics[width=0.94\linewidth]{figures/rt_mpc.pdf}
\caption{上述两步 MPC 的实际闭环过程和初始预测。输入限值已进入优化；每次只执行第一步，使后续计划可以随测量修正。}
\end{figure}

\subsection{什么时候 MPC 不只是截断 LQR}

一维无状态约束的例子看起来像限幅反馈。更一般时，预测器会考虑\textbf{未来刹车距离}。例如关节离上限仅 $0.1\,\mathrm{rad}$，当前速度 $1\,\mathrm{rad/s}$，最大减速度 $2\,\mathrm{rad/s^2}$，最短制动距离为 $1^2/(2\times2)=0.25\,\mathrm{rad}$，已经无法在限位前停住。正确设计应更早限制速度，不能等到位置超限才截断目标。

无约束、时不变线性二次问题可得到状态反馈形式；无限时域下与 LQR 联系起来。在有限时域、终端条件合适时，也可能恰好得到相同反馈。MPC 的价值不只在有无解析解，而在预测、约束和每拍更新的组织方式。约束问题也可以在特定低维条件下预先求成分段仿射的显式 MPC，因此“有约束就绝不可能有显式解”也过于绝对。

\section{稳定 可行与按时求解是三个要求}

有一个可行的优化解，不自动证明闭环稳定。终端代价、终端集合、递归可行性和模型误差都会影响理论结论。\textbf{递归可行性}是指这次可行并执行第一步后，下次仍能找到可行解。\textbf{软约束}引入松弛变量并惩罚违约量，能提高数值可行性，但把碰撞或硬件极限随意软化会改变任务要求。

若求解超时，上一周期计划的下一步只有在仍有效、状态误差可接受且符合约束时才能使用。切回某个 PD 也必须处理参考、积分状态和切换瞬间力矩，不能把“有兜底”当作稳定性证明。需要分别记录求解状态、迭代次数、最坏耗时、约束残差和实际跟踪偏差。

\section{WBC 把任务和物理约束写在一起}

\textbf{全身控制（whole-body control，WBC）}是一类多任务控制框架。固定基座机械臂、双臂和浮动基座人形都可以采用类似思路，但驱动和接触结构不同。常见优化变量是关节广义加速度 $\ddot q$、驱动力矩 $\tau$ 和接触力 $f$：
\begin{align}
\min\quad &\sum_i\|J_i\ddot q+\dot J_i\dot q-a_i^*\|_{W_i}^2
+\|\tau\|_{R_\tau}^2,\\
\text{满足}\quad &M\ddot q+h=S^\mathsf T\tau+J_c^\mathsf Tf,\\
&J_c\ddot q+\dot J_c\dot q=0,\\
&\tau_{\min}\le\tau\le\tau_{\max},\qquad f\in\mathcal K.
\end{align}
$h=C\dot q+g$，需要时还包括明确定义的其它已知项；$S$ 选择实际驱动的自由度；$J_c$ 是保持接触的雅可比。$\|z\|_W^2=z^\mathsf TWz$。任务加速度可由参考加速度加误差反馈得到。

在\textbf{当前构型、当前接触模式固定}时，动力学和加速度约束对这些优化变量是线性的。若把摩擦锥用线性不等式近似，就得到常见的凸 QP。若同时优化未来构型、接触时序或采用精确二阶摩擦锥，则问题类别会改变，不能再笼统称整件事为一个线性约束 QP。

\begin{examplebox}[接触力为什么会改变关节所需力矩]
在本书 2R 姿态 $q=(0,\pi/2)$，
\[
J=\begin{bmatrix}-0.5&-0.5\\0.5&0\end{bmatrix},\qquad
g=(7.3575,0)^\mathsf T.
\]
静态时环境向上托住末端，$f=(0,4)^\mathsf T\,\mathrm N$。
由 $g=\tau+J^\mathsf Tf$，
\[
J^\mathsf Tf=(2,0)^\mathsf T\,\mathrm{N\,m},\quad
\tau=(5.3575,0)^\mathsf T\,\mathrm{N\,m}.
\]
环境帮助承担了一部分重力，驱动力矩应减小。若把环境对机器人的力和机器人对环境的力混用，结果会错误地增加到 $9.3575\,\mathrm{N\,m}$。
\end{examplebox}

\section{任务权重不等于严格优先级}

同一变量 $z$，任务一希望 $z=1$，任务二希望 $z=-1$。加权目标
\begin{equation}
J=(z-1)^2+w(z+1)^2
\end{equation}
的导数为 $2(z-1)+2w(z+1)$，最优解
\begin{equation}
z^*=(1-w)/(1+w).
\end{equation}
取 $w=4$，$z^*=-0.6$，两个任务都没被精确满足。把任务一的权重加得很大，结果会接近一，但有限权重仍是一种折中。

严格分层则先确定第一任务的最优可行集合，再在这个集合中优化第二任务。在本例第一层要求 $z=1$，已经没有剩余自由度，第二层只能接受误差二。这与七轴机械臂用剩余零空间处理肘部任务的逻辑一致，只是带不等式约束时需要更一般的求解结构。

\begin{figure}[htbp]\centering
\includegraphics[width=0.87\linewidth]{figures/rt_task_weights.pdf}
\caption{两个互相冲突任务的加权最优解。权重连续改变的是折中位置，并不会突然产生严格任务优先级。}
\end{figure}

\section{浮动基座与摩擦锥}

人形机器人基座的六个运动自由度没有直接电机驱动，所以不能把固定基座机械臂的所有自由度都当作可独立施加力矩。地面反力通过接触约束影响基座运动。

对接触法向 $f_z\ge0$，各向同性库仑摩擦锥为
\begin{equation}
\sqrt{f_x^2+f_y^2}\le\mu f_z.
\end{equation}
取 $f_z=100\,\mathrm N,\mu=0.5$，切向 $(40,30)\,\mathrm N$ 的长度为 $50\,\mathrm N$，恰在边界。若只分别写 $|f_x|\le50,|f_y|\le50$，就会错误允许 $(50,50)$，其长度为 $70.71$，已超出真实摩擦锥。采用 $|f_x|+|f_y|\le\mu f_z$ 是一种保守的内接菱形近似，却会排除部分真实可行力。使用哪一种近似，需要连同保守性一起说明。

\section{学习方法放在哪一层}

\textbf{强化学习（reinforcement learning，RL）}用策略 $\pi_\theta$ 从观测产生动作，训练目标常是折扣累计奖励
\begin{equation}
J(\theta)=\mathbb E\left[\sum_{k=0}^{T}\gamma^k r_k\right],\qquad0\le\gamma<1.
\end{equation}
奖励是人为定义的训练目标，不等于物理能量或稳定性函数。三步奖励都为一、$\gamma=0.9$ 时，回报为 $1+0.9+0.81=2.71$；无限重复才趋于 $1/(1-0.9)=10$。\cite{sutton}

机器人策略可以输出关节目标、力矩残差、阻抗参数或上层动作。它们对应不同的闭环。若输出关节目标，底层 PD 有助于执行参考，但不能保证任意策略目标序列都稳定、无碰撞或不超限。若输出残差 $\tau=\tau_{\rm base}+\Delta\tau_\theta$，即便残差有幅值限幅，也可能持续注入能量，需要研究频率、方向和状态依赖。

对数据驱动辨识，神经网络或高斯过程可以拟合物理模型残差，例如
\begin{equation}
\tau=Y(q,\dot q,\ddot q)\beta+f_\theta(q,\dot q,\text{温度}).
\end{equation}
先保留明确的刚体结构，再让学习模型表达难建模部分，便于检查重力、摩擦和外推行为。但如果训练数据只覆盖某个姿态与温度，不能据此保证全部工作空间。评价至少要分开训练轨迹、验证轨迹和新的运行条件；低训练误差不是因果辨识或闭环稳定的证明。

\section{计算练习与解答}

\subsection*{题一 无约束两步 MPC}
沿用本章 $J=e_1^2+e_2^2+u_0^2+u_1^2$，若 $e_0=0.5$，求输入和预测误差。

解：$u_0=-0.6e_0=-0.3,u_1=-0.2e_0=-0.1$；$e_1=0.2,e_2=0.1$。总代价 $0.04+0.01+0.09+0.01=0.15$，与初始误差平方缩放规律一致。

\subsection*{题二 权重折中}
要求两个目标 $1$ 和 $-1$ 的加权最优位置为 $0.8$，求 $w$。

解：$0.8=(1-w)/(1+w)$，得 $0.8+0.8w=1-w$，所以 $w=1/9$。第一任务权重是第二的九倍，但第一任务仍有 $0.2$ 的误差。

\subsection*{题三 摩擦可行性}
$f_z=80\,\mathrm N,\mu=0.4$，切向力 $(20,20)\,\mathrm N$ 是否在真实摩擦锥内？是否满足本章的内接菱形近似？

解：切向长度 $\sqrt{800}=28.284\,\mathrm N$，上限 $32\,\mathrm N$，所以真实摩擦锥可行。菱形近似要求 $|20|+|20|=40\le32$，不满足，说明该线性近似更保守。

\section*{本章与原稿的对应}
本章整合 [A，第60--67页] 与 [B，第65--68、92--93页] 的优化和学习内容。补出两步 MPC 的约束手算、任务权重反例、接触力矩和平面摩擦锥计算；修正“底层 PD 自动给学习策略提供稳定保证”等不成立的泛化。

\chapter{贯穿项目 从两连杆到七轴相机云台}
\label{ch:integration}

这一章把前面的工具放到同一条链上。先在模型明确的两连杆对象上核对运动学、动力学、辨识和跟踪，再讨论七轴云台的底座、相机、柔顺与冗余。原稿记录的 NERO 回正问题作为诊断情景使用；没有真机日志支持的结论，不能在教材里写成已定位的原因。

\section{两连杆项目的统一参数和边界}

两连杆长度均为 $0.5\,\mathrm m$，名义质量均为 $1\,\mathrm{kg}$，质心位于半长处，绕运动平面法向的质心惯量为 $1/48\,\mathrm{kg\,m^2}$；重力沿 $-y$，第一关节角 $q_1$ 从水平正方向起算，$q_2$ 是相对第一杆的角度。起点 $q_A=(0,\pi/2)$，终点 $q_B=(\pi/3,\pi/3)$，三秒五次轨迹，之后保持一秒。

本章实际生成的仿真采用 $200\,\mathrm{Hz}$ 控制、$0.5\,\mathrm{ms}$ 数值积分、一个控制周期的命令延迟和 $4\,\mathrm{ms}$ 执行器时间常数。对象质量设为 $(1.1,0.9)\,\mathrm{kg}$，杆惯量按质量同比缩放；摩擦为平滑库仑加黏性模型，参数与辨识数据一起记录。所有这些是\textbf{可复现的教学设定，不是 NERO 规格或测量}。

\subsection{先得到不依赖任何库的核对点}

对名义模型，在起点，末端 $p_A=(0.5,0.5)\,\mathrm m$，
\begin{equation}
J_A=\begin{bmatrix}-0.5&-0.5\\0.5&0\end{bmatrix},\quad
M_A=\begin{bmatrix}5/12&1/12\\1/12&1/12\end{bmatrix},\quad
g_A=(7.3575,0)^\mathsf T.
\end{equation}
若要瞬时末端速度 $(0,0.1)\,\mathrm{m/s}$，第二行给 $0.5\dot q_1=0.1$，所以 $\dot q_1=0.2$；第一行要求 $\dot q_1+\dot q_2=0$，故 $\dot q_2=-0.2\,\mathrm{rad/s}$。

若施加环境力 $f=(2,-3)\,\mathrm N$，其关节外力矩为
\[
J_A^\mathsf Tf=(-2.5,-1)^\mathsf T\,\mathrm{N\,m}.
\]
虚功核对：$\dot q^\mathsf TJ^\mathsf Tf=0.2(-2.5)+(-0.2)(-1)=-0.3\,\mathrm W$，同时 $v^\mathsf Tf=(0,0.1)\cdot(2,-3)=-0.3\,\mathrm W$。两边相等，可以同时检查关节顺序、转置和力方向。

\section{如何让辨识结果进入控制器}

两连杆刚体动力学可写成五个合成参数：
\begin{align}
\beta_1&=I_1+I_2+m_1c_1^2+m_2(l_1^2+c_2^2),\\
\beta_2&=m_2l_1c_2,\qquad \beta_3=I_2+m_2c_2^2,\\
\beta_4&=m_1c_1+m_2l_1,\qquad \beta_5=m_2c_2.
\end{align}
前三个有惯量单位，后两个为质量一阶矩。它们不是七轴机器人“43个基参数”的缩小版计数规则，而是当前 2R 方程的一种线性参数化。由于本例质心沿杆轴、$l_1$ 已知，存在 $\beta_2=l_1\beta_5$；若再采用已知的匀质杆关系 $I_i=m_il_i^2/12$ 和 $c_i=l_i/2$，五个组合全部由两个质量决定。因此本章实际实验只辨识两个质量，不将五个组合当成五个独立物理未知量。一般机器人完整辨识仍需采用第8章的秩分析与物理约束。

例如将第一行回归特征写为 $Y_1=(y_{11},y_{12},y_{13},y_{14},y_{15})$，逐项为
\begin{align}
y_{11}&=\ddot q_1,\qquad y_{13}=\ddot q_2,\\
y_{12}&=2\cos q_2\ddot q_1+\cos q_2\ddot q_2
-\sin q_2(2\dot q_1\dot q_2+\dot q_2^2),\\
y_{14}&=g_0\cos q_1,\qquad y_{15}=g_0\cos(q_1+q_2).
\end{align}
这里的余弦和正弦均显式写出，避免与质心距离 $c_2$ 混用。

教材配套计算先用一条解析可微的合成激励轨迹产生位置、速度和加速度，再由对象真值生成带固定随机种子噪声的力矩记录，利用已知匀质杆几何拟合两个质量和四个摩擦系数，再重建上述合成参数；另一条轨迹用于验证。控制器只使用拟合结果或名义模型，仿真对象单独保留真值。这里为了隔离辨识结构问题，位置与导数取解析值；真实编码器的噪声处理在第9章与第21章讨论，不能把本实验误认为完整复现了真机采集链。

\section{四种控制器的实际数值比较}

使用相同的关节反馈增益 $K_p=\operatorname{diag}(30,12)$、$K_d=\operatorname{diag}(8,3)$，分别计算：纯 PD；名义重力补偿加 PD；名义逆动力学轨迹前馈加 PD；辨识模型及摩擦前馈加 PD。这里的反馈增益分别是力矩每弧度和力矩每角速度的量纲。

逆动力学前馈使用参考状态：
\begin{equation}
\tau_{\rm ff}=\hat M(q_d)\ddot q_d+\hat C(q_d,\dot q_d)\dot q_d+\hat g(q_d)+\hat\tau_f(\dot q_d),
\end{equation}
再加 $-K_p(q-q_d)-K_d(\dot q-\dot q_d)$。这与把反馈加速度也乘以实测状态惯量的计算力矩控制有区别。前馈中没有使用实测加速度，也不需要通过二阶差分去制造一条高噪声信号。

\begin{figure}[htbp]\centering
\includegraphics[width=0.97\linewidth]{figures/rt_integrated_2r.pdf}
\caption{四种控制器在同一教学对象、同一轨迹、相同延迟和反馈增益下的数值结果。末端误差按真实对象的几何计算；不存在几何参数误差，因此本图只比较动态跟踪链。}
\end{figure}

\begin{table}[htbp]\centering\small
\begin{tabularx}{\linewidth}{Xrrrr}\toprule
控制器 & RMS / mm & 最大 / mm & 4 s末值 / mm & 饱和拍数\\\midrule
纯 PD & 140.190 & 212.088 & 56.603 & 0 \\
名义重力补偿加 PD & 3.052 & 5.745 & 3.861 & 0 \\
名义动力学前馈加 PD & 5.729 & 7.640 & 6.890 & 0 \\
辨识模型与摩擦前馈加 PD & 1.414 & 3.952 & 0.117 & 0 \\
\bottomrule\end{tabularx}
\caption{统一两连杆实验的实际计算结果。误差在200 Hz的0至4 s共801个记录点上统计，末值为4 s、到达后保持1 s的误差。饱和拍数统计4 s之前800次控制计算中至少一轴指令被裁剪的次数。}\end{table}

RMS为末端二维误差向量模长的均方根；最大值也是上述记录点上的最大误差，不是逐坐标百分比误差。

训练轨迹包含 3000 个时刻，每时刻两条力矩方程。独立验证的两轴力矩 RMSE 分别为 $0.02032$ 和 $0.01992\,\mathrm{N\,m}$；验证力矩含标准差 $0.02\,\mathrm{N\,m}$ 的独立合成噪声。

本实验将匀质杆几何和摩擦平滑速度尺度 $v_s=0.02\,\mathrm{rad/s}$ 作为已知设定，只辨识两个质量与四个摩擦系数，共六个参数。估计值为
\[
\begin{aligned}
(\hat m_1,\hat m_2)&=(1.09992,0.89993)\,\mathrm{kg},\\
(\hat F_{c1},\hat F_{c2})&=(0.14999,0.10173)\,\mathrm{N\,m},\\
(\hat b_1,\hat b_2)&=(0.07985,0.04637)\,\mathrm{N\,m\,s/rad}.
\end{aligned}
\]

关节力矩上限为 $(15,6)\,\mathrm{N\,m}$；四组实验的延迟队列与执行器力矩从零初始化，所以起始短暂下沉也计入误差。该初始化用于比较延迟影响，不代表推荐的真实硬件激活流程。



表中数据由配套脚本实际计算。注意，名义动力学前馈的 RMS 为 $5.729\,\mathrm{mm}$，反而大于名义重力补偿的 $3.052\,\mathrm{mm}$。两者并不是“在同一个控制律上准确地多加一项”：重力补偿在实测 $q$ 计算，而轨迹前馈在参考 $q_d$ 计算，模型质量有偏差，摩擦又未补偿，时延也使前馈与对象状态错位。增加模型项不自动保证误差下降。

不能拿某一列单独证明某算法普遍更好。改变轨迹、摩擦模型、滤波、饱和或时延，排序可能改变；本实验最有价值的是把“补了什么”和“误差改变多少”对应起来。所有曲线都有相同评价区间和单位，才能公平比较。

\subsection{补偿的求值点为什么改变局部刚度}

在保持目标 $q_d$ 附近，令 $G=\partial g/\partial q|_{q_d}$，$\hat G=\partial\hat g/\partial q|_{q_d}$。忽略高阶项，将重力展开为 $g(q_d+e)\approx g(q_d)+Ge$。若补偿固定为 $\hat g(q_d)$，恢复项近似为 $(K_p+G)e$；若补偿随实测位置变化为 $\hat g(q)$，则近似为 $(K_p+G-\hat G)e$。两种控制器对相同误差的局部恢复作用已经不同。

在本实验终点，这两个矩阵约为
\[
K_p+G=\begin{bmatrix}21.9291&-1.9115\\-1.9115&10.0885\end{bmatrix},\qquad
K_p+G-\hat G=\begin{bmatrix}30.4248&0.2124\\0.2124&12.2124\end{bmatrix}.
\]
单位为 $\mathrm{N\,m/rad}$。这只是保持阶段的一阶解释；完整轨迹还受惯量误差、摩擦残差和时延影响。若想单独评估“多加动力学项”的价值，应固定重力求值方式，逐项增加补偿并重新比较，避免把同时改变的机制归为一种原因。

\section{七轴云台首先是一个带柔性的支撑系统}

相机在世界中的位姿为
\begin{equation}
T_{WC}=T_{WB}T_{BE}(q)T_{EC}.
\end{equation}
$T_{WB}$ 是基座在世界中的位姿，将基座分量转换到世界系；$T_{BE}(q)$ 是工具在基座中的位姿，来自机械臂正运动学；$T_{EC}$ 是相机在工具中的固定安装位姿。关节编码器主要观测中间一段；底座移动、支架变形或相机安装松动，不会自动出现在关节位置误差里。

小角度基座扰动对相机位置的一阶影响包含
\begin{equation}
\delta p_{WC}\approx\delta p_{WB}
-R_{WB}[p_{BC}]_\times\delta\theta_B
+R_{WB}J_{p,C}\delta q.
\end{equation}
本式中 $\delta p_{WB}$ 在世界系表达，$\delta\theta_B$ 为在基座系表达的右扰动小角，$J_{p,C}$ 是相机参考点在基座系中的位置雅可比。因此机械臂关节保持得很好，也可能拍出抖动画面。若相机距基座 $0.8\,\mathrm m$，基座角度晃动 $0.1^\circ=0.0017453\,\mathrm{rad}$，垂直于力臂方向的位移约为 $0.8\times0.0017453=1.40\,\mathrm{mm}$。若相机等效焦距为 $1000$ 像素，对图像中心附近的特征，纯小旋转可造成约 $1.75$ 像素的角向画面位移；有限深度下的相机平移还会产生另一项图像运动。

\begin{figure}[htbp]\centering
\includegraphics[width=0.94\linewidth]{figures/rt_gimbal_base.pdf}
\caption{基座角度误差通过力臂放大成相机位移。右图为小角度一阶计算，不代表某种支架的实测刚度。控制器没有测到的基座运动，不能仅靠关节误差消除。}
\end{figure}

\subsection{从允许角度误差反推支撑刚度}

用一维转动弹簧近似底座：$\tau_b=k_\theta\theta_b$。若任务下变化力矩为 $20\,\mathrm{N\,m}$，允许静态角度变化为 $0.05^\circ$，则
\[
k_\theta\ge\frac{20}{0.05\pi/180}=22918\,\mathrm{N\,m/rad}.
\]
这是线性静态刚度估算。动态还需要等效转动惯量 $I_b$ 和阻尼 $d_b$，模型为 $I_b\ddot\theta_b+d_b\dot\theta_b+k_\theta\theta_b=\tau_b$，自然频率 $\omega_n=\sqrt{k_\theta/I_b}$。提高控制增益不一定能消除底座共振，可能反而把更多力矩变化传给支架。

\section{没到点至少有三种不同误差}

定义目标末端位置 $p_d$，规划出的关节目标 $q_d$，实测关节 $q_m$，名义 FK 为 $\hat p(q)$，真实位置为 $p_{\rm true}(q)$。严格加减同一项可得
\begin{align}
p_{\rm true}(q_m)-p_d={}&
\underbrace{p_{\rm true}(q_m)-\hat p(q_m)}_{\text{几何 安装和柔性误差}}\\
&+\underbrace{\hat p(q_m)-\hat p(q_d)}_{\text{关节跟踪引起的任务误差}}\\
&+\underbrace{\hat p(q_d)-p_d}_{\text{规划或逆解残差}}.
\end{align}
这三项应在相同坐标系和相同时刻比较。外部测量不准还会增加测量误差。把真实末端位置再次求一次 IK 后与 $q_d$ 比较，会受到冗余和多解影响，通常不是清楚的分解办法。

\begin{examplebox}[关节误差怎样变成毫米偏差]
在 $q_d=(0,\pi/2)$ 附近，测得 $\delta q=(0.01,-0.006)\,\mathrm{rad}$。
一阶近似给出
\[
\delta p\approx J_A\delta q=(-0.002,0.005)\,\mathrm m.
\]
用精确 FK 计算为约 $(-0.002025,0.004996)\,\mathrm m$，与一阶近似很接近。若外部相机测到偏差约 $(0.008,0.005)\,\mathrm m$，多出来的约 $10\,\mathrm{mm}$ 横向差不能再全部归因于这组关节跟踪误差，应检查几何、基座、TCP 或外部测量。
\end{examplebox}

\section{低阻抗回正偏差怎样从残差反推}

一维静态阻抗的核心关系为 $Ke=F_{\rm residual}$，其中残差的符号由统一动力学约定决定，可能包含重力估计误差、摩擦、线缆和外力。在 $K=100\,\mathrm{N/m}$ 时，$1\,\mathrm N$ 残差就对应 $1\,\mathrm{cm}$ 偏差。把允许偏差收紧到 $1\,\mathrm{mm}$，而刚度仍保持 $100\,\mathrm{N/m}$，则合成残差需要约束到 $0.1\,\mathrm N$ 量级。

这说明“软”和“准”可以通过更准的模型、更低摩擦或适当外层修正同时改善，但不能只把刚度调小然后要求相同扰动下误差不变。若增加积分，要重新考虑接触时积分积累和模式切换；若增加摩擦补偿，要考虑零速符号和过补偿能量注入。

从相反方向接近同一目标，可帮助区分方向相关和方向无关部分。例如两次静态位置误差为 $e_+=12\,\mathrm{mm}$、$e_-=4\,\mathrm{mm}$，同一方向刚度 $K=100\,\mathrm{N/m}$，对应残差 $1.2$ 和 $0.4\,\mathrm N$。平均值 $0.8\,\mathrm N$ 是方向无关部分的候选，半差 $0.4\,\mathrm N$ 是方向相关部分的候选。只有姿态、温度和其它条件近似相同，且摩擦近似对称时，这种分解才有解释力；平均值不能唯一证明是重力，半差也不能唯一证明是库仑摩擦。

\section{为什么不能凭方向说只是哪一个关节的问题}

原稿用“左右主要靠 J1，所以左右准”解释一个现象，这可以作为待验证假设，但不能作为任意构型下的结论。某方向末端运动需要哪些关节，由当前雅可比和约束共同决定；不同构型、姿态任务和零空间策略会改变分配。

更有用的做法是记录当前 $J$、关节力矩残差、任务误差、饱和状态及缩放后的奇异值。给定小位移方向 $d$，求受限的 $\delta q$，检查实际参与的关节。再看这些关节的重力残差、低速摩擦和力矩余量，而不是先按方向给关节定责。

\section{七轴的剩余自由度要服务于任务}

对满秩六维末端任务，七关节局部有一维运动学零空间；若只控制三维位置，局部可剩四维。剩余维度是 $n-\operatorname{rank}(J)$，不是永远固定为一。零空间可以帮助保持肘姿、远离限位、减少相机遮挡、降低重力力矩或避开线缆扭转。

“让肘部自由拖动”和“让肘部自动回到固定姿态”对应不同的零空间势能。较小或零的零空间刚度让肘部更自由，适当阻尼耗散运动；非零刚度会产生回正倾向。加上动态一致投影还要满足模型和秩条件，实际力矩限幅或阻尼伪逆会产生主任务耦合。第16章给出了这些条件下的完整公式和数值反例。

对相机云台，可以把任务分为过渡定位、离线运镜、在线跟随和拖动录制。它们共享模型与状态，但参考生成器和控制目标不同。画面平稳要同时检查工具姿态误差、角速度谱、相机安装、底座振动和时间同步，不能只看 TCP 平移的 RMS。

\section{一条能解释问题的学习与实施路线}

第一阶段只用两连杆完成 FK、IK、雅可比、动力学和功率核对。第二阶段做本章统一轨迹，区分纯反馈、模型前馈和执行器延迟。第三阶段通过独立训练与验证数据完成辨识，观察误差如何影响控制。第四阶段加入笛卡尔阻抗和外力观测，比较静态柔顺与动态稳定。第五阶段把同样的数据结构、单位和诊断量迁移到七轴模型，再分析冗余和真实工具。

每一阶段的交付证据应与问题对应：运动学需要位置和雅可比核对；辨识需要训练与独立验证残差；控制需要参考、测量、力矩和饱和曲线；实时实现需要数据年龄和耗时分布。不要把原稿里某个教学项目的“误差小于百分之五”“一毫秒循环”等数字直接当作所有真实机械臂的通用验收标准。具体指标应来自任务要求、传感器能力和硬件限制。

\section{计算练习与解答}

\subsection*{题一 底座误差}
相机到支撑转轴距离 $0.6\,\mathrm m$，允许支撑导致的位移误差 $0.5\,\mathrm{mm}$，小角度近似下允许角度是多少？

解：$\theta\le0.0005/0.6=0.00083333\,\mathrm{rad}=0.04775^\circ$。这是由支撑单独占用全部误差预算时的上界；其它误差还需分配预算。

\subsection*{题二 低刚度残差预算}
要求 $K=80\,\mathrm{N/m}$ 下静态误差不超过 $5\,\mathrm{mm}$，残余力应低于多少？

解：$|F_r|\le K|e|=80\times0.005=0.4\,\mathrm N$。这包括所有未被消除的等效残差，不仅是摩擦模型误差。

\subsection*{题三 七轴冗余维度}
七轴机械臂，某姿态下三维位置雅可比秩为三，六维几何雅可比秩为五。两种任务分别有多少维零空间？

解：位置任务为 $7-3=4$ 维；六维几何任务为 $7-5=2$ 维。但后者丢失了一个任务方向，不能把奇异点多出来的零空间当成获得了额外完整任务能力。

\section*{本章与原稿的对应}
本章整合 [A，第74--79页]、[B，第120--127页]，增加实际执行的统一 2R 数值实验、世界到相机误差分解及底座刚度预算。原稿中的 NERO 现场描述保留为诊断背景，所有新数值都明确为教学计算。

\appendix
\titleformat{\chapter}[display]{\sffamily\bfseries\color{black}}{\large 附录\thechapter}{7pt}{\LARGE}
\renewcommand{\chaptermark}[1]{\markboth{附录\thechapter\quad #1}{}}
\chapter{原稿内容对应与关键修订}\label{app:sources}

本附录帮助读者在两份原稿和本书之间定位同一主题。A、B 的全名见前言，以下页码均为上传 PDF 的物理页序；对应关系按主题排列，不表示逐页逐句重写。本书把重复知识合并，把缺少的中间计算补齐，并用统一模型重新计算例子。\textbf{补充条件、改变教学模型和纠正错误是三种不同的处理}，不能因为新旧数字不同就认定原稿有误。

\section{主题与阅读位置}

\begingroup
\small
\setstretch{1.10}
\renewcommand{\arraystretch}{1.20}
\begin{longtable}{>{\raggedright\arraybackslash}p{30mm}>{\raggedright\arraybackslash}p{54mm}>{\raggedright\arraybackslash}p{70mm}}
\toprule
主题 & 原稿对应页码 & 本书处理与阅读位置\\
\midrule
\endfirsthead
\toprule 主题 & 原稿对应页码 & 本书处理与阅读位置\\\midrule
\endhead
\midrule\multicolumn{3}{r}{接下页}\\\endfoot
\bottomrule\endlastfoot
线性代数与概率 & A：2--5；B：3、8--17 & 第1章。补充完整 SVD、条件数、QR、Cholesky、协方差传播和最小二乘数字。\\
微积分、方程与变换 & A：3--5；B：17--24 & 第2章。加入初值、单位、变换条件、积分器稳定性和变分法边界项。\\
三维位姿与李群 & A：22--26；B：3--8、24--27 & 第3章。固定坐标变换方向、四元数顺序、姿态误差符号与旋量参考点。\\
运动学与雅可比 & A：26--30；B：14--17、26--32 & 第4--5章。串起解析逆解、局部迭代、奇异性、速度零空间与功率映射。\\
动力学与执行器 & A：30--38；B：32--40、116--120 & 第6--7章。采用均匀杆模型，逐项求导，并补充传动、柔性、接触和电机约定。\\
辨识、激励与摩擦 & A：33--34；B：41--53 & 第8--9章。明确可辨识组合、噪声假设、激励尺度、滤波和摩擦分离条件。\\
负载、工具与相机标定 & A：22--30；B：44--50 & 第10章。展开静态负载回归、零偏、TCP 和手眼坐标链。\\
反馈、状态空间与数字控制 & A：5--22、34--43；B：59--74、78--83 & 第11--13章。增加频域、LQR、精确离散化、饱和与逐拍更新的计算。\\
观测、阻抗与操作空间 & A：16--18、29--30、39--52；B：54--58、74--87、122--123 & 第14--16章。区分扰动与接触力，推导质量重塑、任务惯量和力矩零空间。\\
轨迹与示教 & A：52--60；B：87--97 & 第17--18章。补出边界系数、峰值、S 曲线分支、DMP 和 ProMP 算例。\\
遥操作与视觉任务 & B：97--100 & 第19章。明确功率端口、位置和力缩放，并新增与相机坐标相连的视觉伺服算例。\\
软件、接口与仿真 & A：67--74；B：100--120、125--127 & 第20--21章。区分模型、接口、控制时序与物理积分；软件内容结合已核对的一手文档。\\
优化与学习方法 & A：60--67；B：65--68、92--93 & 第22章。用小规模 MPC、任务权重和接触约束解释适用边界。\\
综合项目与七轴云台 & A：74--79；B：120--127 & 第23章。用可复算的 2R 教学实验连接全书，分开几何、跟踪、底座和低阻抗误差。\\
\end{longtable}
\endgroup

\section{需要特别留意的修订}

以下条目只指出会改变计算或判断的内容。原稿中的经验建议若没有充分的适用条件，本书将它改写为需要验证的设计选择，而不把它一律判作错误。

\begin{enumerate}
\item \textbf{特征值、奇异值和可对角化不是同一件事。}对于对称不定矩阵，奇异值是特征值的绝对值；一般矩阵也未必有完整特征向量基。第1章将这些结论分别说明，并给出数值例子。[A，第2--3页；B，第8--17页]

\item \textbf{条件数要连同量纲和尺度一起报告。}本书保留奇异值判断局部病态的用途，同时说明米与毫米、位置与姿态的权重会改变数值。条件数较小不等于激励绝对强度足够，某个固定阈值也不是所有模型的通用合格线。[A，第3、33--34页；B，第14--17、41--53页]

\item \textbf{几何角速度不是欧拉角导数。}第3、5章分开位置导数、角速度、欧拉角速率，以及空间旋量的平移部分。关于 $\log(R_d^{\mathsf T}R)$，关键是把误差的定义、正负号和表达坐标共同写出；不能只凭“左误差”或“右误差”的名称判断代码。[A，第22--30页；B，第24--32页]

\item \textbf{速度耦合项不能笼统说成不做功。}标准刚体模型中，$\dot q^{\mathsf T}C\dot q=\tfrac12\dot q^{\mathsf T}\dot M\dot q$，来自 $\dot M-2C$ 的反对称性，而左侧不必为零。第6章给出 $0.00025\,\mathrm W$ 的直接计算。[A，第33页；B，第32--40页]

\item \textbf{均匀杆和杆端点质量的数字不同，是模型不同。}本书主模型保留质心转动惯量；第6章另设练习恢复点质量模型的结果。不能拿一组模型的 $M$ 搭配另一组模型的质心和重力项。[A，第30--34页；B，第32--36页]

\item \textbf{外力方向必须统一后再做减法。}本书规定环境对机器人的作用为 $\tau_{\rm ext}$，所以外力矩等于模型所需力矩减去执行器力矩。如果传感器报告的是机器人对环境的反作用，需先改变符号。第14章据此统一观测器与残差解释，完整公式见附录C。[A，第50--51页；B，第56--58、74--78页]

\item \textbf{辨识权重、基参数数量和物理约束各有条件。}白化乘标准差的倒数，平方代价的权重为方差倒数；可辨识参数数目来自具体回归矩阵的秩。惯量矩阵正定还不足以保证刚体质量分布物理一致。第8--9章不把原稿中的某个参数数量、截止频率倍数或迭代轮数作为普遍规则。[A，第33--34页；B，第41--53页]

\item \textbf{前馈、离散实现和滤波不能借用过强结论。}添加加速度前馈不会自动增加反馈系统型别；增量式 PID 仍需处理状态初始化和饱和；双向滤波的有效幅值响应与单向滤波不同。连续系统的稳定性也不能直接替代离散闭环分析。[A，第5--22页；B，第59--70、78--83页]

\item \textbf{速度零空间和力矩零空间必须分别检查。}前者要求 $JN=0$，后者要检查 $JM^{-1}N^{\mathsf T}=0$，并使用相应的动态一致定义。阻尼伪逆一般破坏严格零空间性质；仅补偿任务方向的重力也未必托住全部冗余方向。[A，第29--30、44页；B，第55、83--86、122--123页]

\item \textbf{小阻尼、平滑轨迹和仿真曲线都不自动完成其它目标。}零速度下黏性阻尼不提供静态托举力；平滑轨迹仍需单独满足限位与碰撞约束；教学仿真不能证明真机已经达到某项性能。第17--23章把结论分别限定到模型、输入、时序和验证范围。[A，第52--60、72--79页；B，第87--100、111--127页]
\end{enumerate}

\chapter{术语 符号与单位速查}\label{app:glossary}

同一个字母在不同学科中可能重复使用。例如 $R$ 可表示旋转矩阵，也可能是优化权重；$C$ 可表示动力学矩阵，也可能是状态空间输出矩阵。\textbf{本表按定义查量，不按字母猜量。}表中转动单位保留 rad，虽然 SI 量纲分析通常把弧度视为无量纲。平移与转动组成的混合向量没有一个统一的物理单位。

\begingroup
\small
\renewcommand{\arraystretch}{1.18}
\begin{longtable}{>{\raggedright\arraybackslash}p{33mm}>{\raggedright\arraybackslash}p{39mm}>{\raggedright\arraybackslash}p{82mm}}
\toprule 术语与常用英文 & 符号、单位或维度 & 使用时要知道的含义\\\midrule
\endfirsthead
\toprule 术语与常用英文 & 符号、单位或维度 & 使用时要知道的含义\\\midrule
\endhead
\midrule\multicolumn{3}{r}{接下页}\\\endfoot
\bottomrule\endlastfoot
\multicolumn{3}{l}{\textbf{数学与统计基础}}\\\midrule
% TERM 01
标量与列向量\newline scalar, vector & $a\in\mathbb R$；$x\in\mathbb R^n$ & 一个数与有序的一列数；向量的分量顺序必须固定。\\
% TERM 02
转置\newline transpose & $A^{\mathsf T}$；$n\times m$ & 对 $A\in\mathbb R^{m\times n}$ 交换行列；$(AB)^{\mathsf T}=B^{\mathsf T}A^{\mathsf T}$。\\
% TERM 03
秩\newline rank & $r=\operatorname{rank}A$ & 独立列或独立行的数量；不等于矩阵非零元素的数量。\\
% TERM 04
核与零空间\newline kernel, null space & $\ker A$；维数 $n-r$ & 满足 $Az=0$ 的全部输入方向；这里 $A$ 有 $n$ 列。\\
% TERM 05
特征值与特征向量\newline eigenpair & $Av=\lambda v$ & 沿 $v$ 作用后方向不变；实矩阵也可能有复特征值。\\
% TERM 06
奇异值分解\newline SVD & $A=U\Sigma V^{\mathsf T}$ & 奇异值 $\sigma_i\ge0$；显示输入方向经线性映射后的伸缩。\\
% TERM 07
条件数\newline condition number & $\kappa_2=\sigma_{\max}/\sigma_{\min}$ & 此比值用于满秩相应空间；需要同时说明单位与缩放，秩亏时通常记为无穷。\\
% TERM 08
摩尔彭若斯伪逆\newline pseudoinverse & $A^+$；$n\times m$ & 给出最小二乘的最小欧氏范数解；不能自行处理关节限位。\\
% TERM 09
二范数\newline Euclidean norm & $\|x\|_2=\sqrt{x^{\mathsf T}x}$ & 是所选坐标中的长度；混合单位分量应先定义权重或特征尺度。\\
% TERM 10
梯度与海森矩阵\newline gradient, Hessian & $\nabla f$；$\nabla^2 f$ & 分别描述一阶变化与局部曲率；单位随自变量而变。\\
% TERM 11
最小二乘\newline least squares & $\min_\beta\|H\beta-y\|^2$ & 最小化残差平方和；小训练残差不自动说明参数正确。\\
% TERM 12
协方差\newline covariance & $\Sigma=E[\delta x\delta x^{\mathsf T}]$ & 第 $i,j$ 项单位是 $x_i$ 单位乘 $x_j$ 单位，非对角项记录相关性。\\
\midrule\multicolumn{3}{l}{\textbf{坐标 位姿与空间量}}\\\midrule
% TERM 13
坐标系\newline frame & $\{A\},\{B\}$ & 由原点和有序坐标轴组成；点的位置还取决于原点。\\
% TERM 14
位置\newline position & $p\in\mathbb R^3$；m & 指定参考点相对某原点的位置分量；不能省掉所属坐标系。\\
% TERM 15
旋转矩阵\newline rotation matrix & $R_{AB}\in\mathbb R^{3\times3}$ & 把 B 系分量变为 A 系分量；$R^{\mathsf T}R=I$、$\det R=1$。\\
% TERM 16
齐次变换\newline rigid transform & $T_{AB}\in\mathbb R^{4\times4}$ & 同时包含旋转与平移；$T_{AC}=T_{AB}T_{BC}$。\\
% TERM 17
叉乘矩阵\newline skew matrix & $[a]_\times b=a\times b$ & 三阶反对称矩阵；把叉乘写成矩阵乘法，不能交换乘法顺序。\\
% TERM 18
旋转群与旋转代数 & $SO(3)$；$\mathfrak{so}(3)$ & 前者是旋转矩阵集合，后者是其单位元处的反对称切向量集合。\\
% TERM 19
刚体运动群与代数 & $SE(3)$；$\mathfrak{se}(3)$ & 组合平移和旋转；指数与对数连接局部增量和有限位姿。\\
% TERM 20
旋转向量\newline rotation vector & $\phi=\theta u$；rad & 方向是旋转轴，长度是转角；接近 $\pi$ 时需关注分支。\\
% TERM 21
单位四元数\newline quaternion & $\boldsymbol q=(w,v)$；四个分量 & 本书标量在前；$\boldsymbol q$ 与 $-\boldsymbol q$ 表示同一旋转，勿与关节 $q$ 混淆。\\
% TERM 22
角速度\newline angular velocity & $\omega\in\mathbb R^3$；rad/s & 表示瞬时旋转轴和速率；一般不等于三个欧拉角的导数。\\
% TERM 23
六维几何速度 & $V=[v;\omega]$ & 线速度在前，角速度在后；线速度对应哪个点必须说明。\\
% TERM 24
力旋量\newline wrench & $F=[f;m]$；N 与 N m & 力矩关于指定参考点；与速度同点同轴时 $F^{\mathsf T}V$ 才是配对功率。\\
% TERM 25
几何雅可比\newline geometric Jacobian & $J\in\mathbb R^{m\times n}$ & 满足任务速度 $v_x=J\dot q$；所选任务决定行数和各行单位。\\
\midrule\multicolumn{3}{l}{\textbf{运动学与路径}}\\\midrule
% TERM 26
关节状态 & $q,\dot q,\ddot q$ & 本书转动关节分别用 rad、rad/s、rad/s$^2$；移动关节需换为长度单位。\\
% TERM 27
正解与逆解\newline FK, IK & $T=f(q)$；求 $q$ & 正解从关节到任务；逆解可能多解、无解或仅能局部迭代求得。\\
% TERM 28
DH 参数 & $a_i,\alpha_i,d_i,\theta_i$ & 相邻坐标系的一种建模规则；标准与改进 DH 不能混搭。\\
% TERM 29
指数积\newline PoE & $\prod_i\exp([S_i]q_i)M_0$ & 由关节旋量和参考位姿构造正解；必须注明旋量及乘法约定。\\
% TERM 30
运动学奇异 & $\operatorname{rank}J$ 降低 & 机构在该位形失去某个独立瞬时任务方向；数值正则化不能创造运动自由度。\\
% TERM 31
工作空间\newline workspace & 可达位置或位姿集合 & 需区分位置可达、姿态可达和受限位碰撞约束后的可达。\\
% TERM 32
工具中心点\newline TCP & $p_{\rm TCP}$；m & 工具的任务参考点；它通常不同于负载质心。\\
% TERM 33
阻尼最小二乘\newline DLS & $J^{\mathsf T}(JJ^{\mathsf T}+\lambda^2I)^{-1}$ & 用有限误差换取较小关节速度；$\lambda$ 需与所用尺度配套。\\
\midrule\multicolumn{3}{l}{\textbf{动力学与执行器}}\\\midrule
% TERM 34
质量与质心\newline mass, center of mass & $m$：kg；$c$：m & 质心坐标必须指定表达系和参考原点；不是随意选定的几何中心。\\
% TERM 35
质心转动惯量 & $I_C$；kg m$^2$ & 三维为对称 $3\times3$ 矩阵；本书 2R 只用绕平面法向的分量。\\
% TERM 36
关节惯量矩阵 & $M(q)\in\mathbb R^{n\times n}$ & 全转动关节常用 kg m$^2$；$\tfrac12\dot q^{\mathsf T}M\dot q$ 是动能。\\
% TERM 37
科氏与离心项 & $c(q,\dot q)=C\dot q$；N m & 由速度耦合产生；$c$ 固定时，满足该乘积的 $C$ 表达不唯一。\\
% TERM 38
重力项与重力加速度 & $g(q)$：N m；$g_0$：m/s$^2$ & $g(q)=\partial V/\partial q$；本书表示用于抵消重力的左端项。\\
% TERM 39
执行器与外力矩 & $\tau_{\rm cmd},\tau_{\rm ext}$；N m & 均按关节侧定义；后者是环境施加给机器人，而非其反作用。\\
% TERM 40
摩擦项 & $\tau_f$；N m & 本书放在动力学左侧；耗散摩擦满足 $\dot q^{\mathsf T}\tau_f\ge0$。\\
% TERM 41
机械能与功率 & $E$：J；$P$：W & 能量是状态储存量，功率是其变化率；$1\,\mathrm W=1\,\mathrm{J/s}$。\\
% TERM 42
电机力矩常数 & $K_t$；N m/A & 电流的峰值、有效值或 $dq$ 定义必须与力矩常数一致。\\
% TERM 43
传动比与效率 & $N$、$\eta$；无量纲 & 先注明 $N$ 是哪两个角速度的比；电机转动惯量折算涉及 $N^2$。\\
% TERM 44
刚度\newline stiffness & $K$；N/m 或 N m/rad & 把位移误差映射为恢复力或恢复力矩；混合任务矩阵有不同单位块。\\
% TERM 45
阻尼\newline damping & $D$；N s/m 或 N m s/rad & 把相对速度映射为耗散作用；零速度时黏性阻尼力为零。\\
% TERM 46
正动力学与逆动力学 & 求 $\ddot q$；求 $\tau$ & 前者由力矩预测加速度，后者由指定运动计算所需力矩。\\
\midrule\multicolumn{3}{l}{\textbf{辨识 标定与数据}}\\\midrule
% TERM 47
待辨识参数 & $\beta\in\mathbb R^p$ & 可由惯性、重力或摩擦参数组合而成；每个分量的单位可能不同。\\
% TERM 48
回归矩阵\newline regressor & $H\in\mathbb R^{N_y\times p}$ & 满足 $y\approx H\beta$；列的单位需与对应参数配对后得到测量单位。\\
% TERM 49
残差\newline residual & $r=y-H\hat\beta$ & 用来检查模型与数据的差；不能仅靠其大小唯一定位误差来源。\\
% TERM 50
白化与加权 & $W\Sigma W^{\mathsf T}=I$ & 采用白化矩阵 $W$ 时，代价为 $\|Wr\|^2$，其二次权重是 $W^{\mathsf T}W$。\\
% TERM 51
参数协方差 & $P_\beta$；$p\times p$ & 单位为参数单位的两两乘积；解释依赖噪声模型和辨识假设。\\
% TERM 52
持续激励\newline persistent excitation & 回归信息随时间积累 & 要使目标参数方向获得独立信息；运动幅度大不等于全部方向都可辨识。\\
% TERM 53
基参数\newline base parameters & 由秩决定的参数组合 & 是数据可区分的最小组合；数量依赖结构、约束和观测，不能只看关节数。\\
% TERM 54
均方根与均方根误差 & RMS、RMSE；与被评估量同单位 & 要写清评价时间、各关节或各任务分量，以及是否经过归一化。\\
% TERM 55
递推最小二乘 & RLS；遗忘因子 $\lambda_f$ & 用新数据更新参数和信息；遗忘加快适应，也减少旧数据权重。\\
\midrule\multicolumn{3}{l}{\textbf{控制与估计}}\\\midrule
% TERM 56
跟踪误差 & $e=q-q_d$；rad & 本书默认实际减期望；恢复反馈因此写为 $-K_pe-K_d\dot e$。\\
% TERM 57
PD 与 PID & 比例、微分、积分反馈 & 增益单位由被控对象和输出决定；积分状态需考虑限幅和模式切换。\\
% TERM 58
自然频率与阻尼比 & $\omega_n$：rad/s；$\zeta$：无量纲 & 用于标准二阶线性模型；不能把多关节系统逐轴任意视为独立二阶。\\
% TERM 59
状态、输入与输出 & $x,u,y$；维数由模型定义 & 状态含预测未来所需的信息；输出未必包含全部状态。\\
% TERM 60
状态空间矩阵 & $A,B,C_s,D_s$ & $\dot x=Ax+Bu$，$y=C_sx+D_su$；这里 $C_s$ 与动力学 $C$ 无关。\\
% TERM 61
极点\newline pole & 连续 $s$；离散 $z$ & 连续稳定要求极点实部为负；离散稳定要求位于单位圆内部。\\
% TERM 62
可控与可观 & controllable, observable & 是否能由输入影响全部状态、由输出重建全部状态；需结合具体模型判断。\\
% TERM 63
线性二次调节器 & LQR；$Q,R_u,P$ & 权重带有状态与输入尺度；最优性只针对指定模型和代价。\\
% TERM 64
卡尔曼滤波 & $K_f,Q_w,R_v$ & 增益由预测与测量不确定性决定；协方差单位应跟随对应状态和测量。\\
% TERM 65
扰动与扩张状态观测 & DOB、ESO & 通常估计模型未解释的总扰动；该量不自动等于纯接触力。\\
% TERM 66
阻抗与导纳 & impedance, admittance & 分别描述运动到力、力到运动的目标关系；实现取决于底层接口和传感信息。\\
% TERM 67
操作空间惯量 & $\Lambda=(JM^{-1}J^{\mathsf T})^{-1}$ & 需 $M\succ0$ 且 $J$ 满行秩；平移块类似质量，转动块类似惯量。\\
% TERM 68
动态一致零空间 & $\bar J=M^{-1}J^{\mathsf T}\Lambda$；$N=I-\bar JJ$ & 速度用 $N$，力矩用 $N^{\mathsf T}$；严格解耦还依赖模型、秩和实际执行。\\
\midrule\multicolumn{3}{l}{\textbf{规划 人机输入与工程时序}}\\\midrule
% TERM 69
控制周期与积分步长 & $T_s,h$；s & 一个是控制器更新间隔，一个是物理数值积分间隔；两者可以不同。\\
% TERM 70
采样率与奈奎斯特频率 & $f_s=1/T_s$；$f_s/2$，Hz & 带限无混叠条件涉及采样前信号；采样后滤波不能恢复已经混叠的信息。\\
% TERM 71
延迟与抖动 & delay, jitter；s & 延迟是信息或作用的滞后，抖动是时序波动；均值不能代替最坏情况。\\
% TERM 72
时间戳与数据年龄 & $t_k$；$t_{\rm now}-t_k$，s & 区分采样时刻、接收时刻和执行时刻；同步误差会污染速度和力矩配对。\\
% TERM 73
零阶保持 & ZOH & 两次更新之间保持输入不变；精确离散化公式依赖这个输入假设。\\
% TERM 74
加加速度\newline jerk & $q^{(3)}$；rad/s$^3$ & 是加速度的变化率；限制它可以控制加速度变化，不能替代所有动力学约束。\\
% TERM 75
五次轨迹与 S 曲线 & quintic, S-curve & 前者满足六个端点条件；后者需根据速度、加速度和 jerk 限制选择分支。\\
% TERM 76
运动基元 & DMP、ProMP & 前者用受迫动态系统表达运动，后者用权重分布表达轨迹及条件化。\\
% TERM 77
预测控制与全身控制 & MPC、WBC & 前者滚动优化未来，后者协调多任务与约束；权重和严格优先级不同。\\
% TERM 78
无源性\newline passivity & 功率 $F^{\mathsf T}V$；储能 J & 通过能量账约束系统净产生的能量；必须定义端口方向和储能函数。\\
% TERM 79
实时性\newline real time & 截止期、最坏耗时、抖动 & 关注是否按时完成，而不只是平均运算快；具体要求来自任务。\\
% TERM 80
模型描述文件 & URDF、MJCF & 分别按各自格式表达机器人或仿真模型；导入成功仍需核对轴向、惯量和执行器定义。\\
\end{longtable}
\endgroup

\chapter{手算复核顺序与公式适用条件}\label{app:recheck}

本附录供实现或复习时查用。正文负责推导来源，这里保留一套可反复回代的数字和检查顺序。后续代码复现说明另列；这里的计算不依赖某个机器人库，也不把某种仿真器的运行当作已经完成的证据。

\section{先固定维度 坐标与符号}

对于 $n$ 个转动关节和 $m$ 维选定任务，$q,\dot q,\tau$ 均为 $n\times1$，$M$ 为 $n\times n$，$J$ 为 $m\times n$。完整六维任务用 $V=[v;\omega]$、$F=[f;m]$；二维平移任务只用 $v\in\mathbb R^2$ 和 $f\in\mathbb R^2$。不要把二维、三维和六维数组仅通过补零混成同一个物理定义。

坐标变换约定是 $p_A=R_{AB}p_B+p_{AB}$，其中 $p_{AB}$ 是 B 系原点在 A 系中的位置。把雅可比的速度参考点从法兰改到 TCP 时，线速度块随力臂改变；力旋量中的力矩也要关于相同 TCP。关节顺序、角正方向和传动侧在整个计算链中保持一致。

全书默认动力学与误差为
\begin{align}
M\ddot q+c+g+\tau_f&=\tau_{\rm cmd}+\tau_{\rm ext},\\
e&=q-q_d,\qquad c=C\dot q.
\end{align}
如果控制器接口接受的是电流或电机侧力矩，应在动力学计算之外显式加入转换；不能把这两种输入直接当作关节侧 $\tau_{\rm cmd}$。

\section{统一两连杆模型的完整速查}

\subsection{参数与运动学}

模型是固定基座、竖直平面、两个转动关节和两根均匀细杆：
\[
l_1=l_2=0.5\,\mathrm m,\quad
m_1=m_2=1\,\mathrm{kg},\quad
c_1=c_2=0.25\,\mathrm m,\quad
I_1=I_2=\frac1{48}\,\mathrm{kg\,m^2}.
\]
重力为 $-y$ 方向，$g_0=9.81\,\mathrm{m/s^2}$。$q_1$ 从 $+x$ 轴逆时针起算，$q_2$ 是第二杆相对于第一杆的角。以下 $q_{12}=q_1+q_2$、$\dot q_{12}=\dot q_1+\dot q_2$，末端是第二杆端点。
\begin{align}
p(q)&=\frac12\begin{bmatrix}\cos q_1+\cos q_{12}\\\sin q_1+\sin q_{12}\end{bmatrix},\\
J(q)&=\frac12\begin{bmatrix}
-\sin q_1-\sin q_{12}&-\sin q_{12}\\
\cos q_1+\cos q_{12}&\cos q_{12}
\end{bmatrix},\\
\dot p&=J\dot q,\qquad
\ddot p=J\ddot q+\dot J\dot q,\\
\dot J\dot q&=-\frac12\begin{bmatrix}
\cos q_1\dot q_1^2+\cos q_{12}\dot q_{12}^2\\
\sin q_1\dot q_1^2+\sin q_{12}\dot q_{12}^2
\end{bmatrix}.
\end{align}
这里 $\det J=\tfrac14\sin q_2$。完整二维位置任务的瞬时奇异位形为 $q_2=k\pi$；若任务只选一个方向，应重新检查选中行的秩。

\subsection{动力学与能量}

将统一参数代入第6章的推导，得到
\begin{align}
M(q)&=\begin{bmatrix}
\frac5{12}+\frac14\cos q_2&\frac1{12}+\frac18\cos q_2\\
\frac1{12}+\frac18\cos q_2&\frac1{12}
\end{bmatrix},\\
c(q,\dot q)&=\frac18\sin q_2
\begin{bmatrix}-(2\dot q_1\dot q_2+\dot q_2^2)\\\dot q_1^2\end{bmatrix},\\
g(q)&=9.81\begin{bmatrix}
\frac34\cos q_1+\frac14\cos q_{12}\\
\frac14\cos q_{12}
\end{bmatrix},\\
T&=\tfrac12\dot q^{\mathsf T}M\dot q,\qquad
V=9.81\left(\tfrac34\sin q_1+\tfrac14\sin q_{12}\right).
\end{align}
$M$ 的单位为 kg m$^2$，$c,g$ 为 N m，$T,V$ 为 J。不要把这里的质心距离 $c_i$ 与速度耦合向量 $c(q,\dot q)$ 混用。

若末端只受到二维环境力 $f_{\rm ext}$，且无其它外部接触，
\begin{align}
\tau_{\rm ext}&=J^{\mathsf T}f_{\rm ext},\\
\tau_{\rm cmd}&=M\ddot q+c+g+\tau_f-J^{\mathsf T}f_{\rm ext},\\
\ddot q&=\operatorname{solve}\!\left(M,\,
\tau_{\rm cmd}+J^{\mathsf T}f_{\rm ext}-c-g-\tau_f\right).
\end{align}
最后一式表示解线性方程，不要求显式构造 $M^{-1}$。如果环境还施加末端纯力矩，要同时加入对应角速度雅可比的转置映射；若接触发生在连杆中部，要使用该接触点的雅可比。

\section{一组从位置到力矩再回代的核对数字}

取第6章已经使用的运动状态，再显式加入教学外力和摩擦：
\begin{gather}
q=(0,\pi/2)^{\mathsf T},\quad
\dot q=(0.2,-0.1)^{\mathsf T},\quad
\ddot q=(0.5,-0.2)^{\mathsf T},\\
f_{\rm ext}=(2,-3)^{\mathsf T}\,\mathrm N,\qquad
\tau_f=(0.04,-0.03)^{\mathsf T}\,\mathrm{N\,m}.
\end{gather}
这里的摩擦向量只是该运动状态下的指定值；两关节分别正向、反向运动，且 $\dot q^{\mathsf T}\tau_f=0.011\,\mathrm W\ge0$，符合左端耗散约定。

\begin{enumerate}
\item \textbf{先检查几何和速度。}逐项代入三角函数后应得到
\[
p=\begin{bmatrix}0.5\\0.5\end{bmatrix},\quad
J=\begin{bmatrix}-0.5&-0.5\\0.5&0\end{bmatrix},\quad
\dot p=\begin{bmatrix}-0.05\\0.10\end{bmatrix}\,\mathrm{m/s}.
\]
再算 $J\ddot q=(-0.15,0.25)^{\mathsf T}$ 和
$\dot J\dot q=(-0.02,-0.005)^{\mathsf T}$，所以
$\ddot p=(-0.17,0.245)^{\mathsf T}\,\mathrm{m/s^2}$。这里删去 $\dot J\dot q$ 会得到不同结果。

\item \textbf{分别计算惯性、速度和重力。}
\[
M=\begin{bmatrix}5/12&1/12\\1/12&1/12\end{bmatrix},\quad
M\ddot q=\begin{bmatrix}0.1916667\\0.025\end{bmatrix},\quad
c=\begin{bmatrix}0.00375\\0.005\end{bmatrix},\quad
g=\begin{bmatrix}7.3575\\0\end{bmatrix}.
\]
$M_{11}>0$ 且 $\det M=1/36>0$，可先确认该两阶惯量矩阵正定；此处的 $c$ 和 $g$ 都是力矩。

\item \textbf{用功率检查雅可比转置。}
\[
\tau_{\rm ext}=J^{\mathsf T}f_{\rm ext}
=\begin{bmatrix}-2.5\\-1\end{bmatrix}\,\mathrm{N\,m}.
\]
任务侧功率是 $2(-0.05)-3(0.1)=-0.4\,\mathrm W$；关节侧是
$(-2.5)(0.2)+(-1)(-0.1)=-0.4\,\mathrm W$。这一步若不相等，先查坐标、作用点、分量顺序和力方向。

\item \textbf{按统一符号合成命令。}无摩擦无外力所需力矩为
$(7.5529167,0.03)^{\mathsf T}\,\mathrm{N\,m}$；加入摩擦，再减去环境力矩，得到
\[
\tau_{\rm cmd}
=\begin{bmatrix}7.5529167\\0.03\end{bmatrix}
+\begin{bmatrix}0.04\\-0.03\end{bmatrix}
-\begin{bmatrix}-2.5\\-1\end{bmatrix}
=\begin{bmatrix}10.0929167\\1.00\end{bmatrix}\,\mathrm{N\,m}.
\]
环境端口在此吸收机械能，执行器还需承担相应作用；不能仅看到 $\tau_{\rm ext}$ 的负号就再次反向。

\item \textbf{回代原方程，不只检查中间公式。}先求净加速力矩
\[
\tau_{\rm cmd}+\tau_{\rm ext}-c-g-\tau_f
=\begin{bmatrix}0.1916667\\0.025\end{bmatrix}.
\]
用 $M$ 解线性方程，应回到
$\ddot q=(0.5,-0.2)^{\mathsf T}\,\mathrm{rad/s^2}$；再通过 $J\ddot q+\dot J\dot q$ 回到前面的末端加速度。

\item \textbf{最后核对总能量。}本状态满足
\[
\dot T=0.0360833\,\mathrm W,\qquad
\dot V=\dot q^{\mathsf T}g=1.4715\,\mathrm W.
\]
两者相加为 $1.5075833\,\mathrm W$。另一条独立路径为
\[
\dot q^{\mathsf T}\tau_{\rm cmd}
+\dot q^{\mathsf T}\tau_{\rm ext}
-\dot q^{\mathsf T}\tau_f
=1.9185833-0.4-0.011
=1.5075833\,\mathrm W.
\]
这项相等检查同时约束惯量、重力、速度项与摩擦的符号。
\end{enumerate}

\section{公式能用到哪一步}

下表只列最容易在实现中遗漏的前提。正则化、限幅、时延和状态估计一旦加入，实际闭环已不再等同于原来的理想代数模型，需要用相应模型重新判断。

\begingroup
\small
\begin{longtable}{>{\raggedright\arraybackslash}p{54mm}>{\raggedright\arraybackslash}p{100mm}}
\toprule 公式或判断 & 适用条件与检查点\\\midrule
\endfirsthead
\toprule 公式或判断 & 适用条件与检查点\\\midrule
\endhead
\midrule\multicolumn{2}{r}{接下页}\\\endfoot
\bottomrule\endlastfoot
$\hat\beta=(H^{\mathsf T}H)^{-1}H^{\mathsf T}y$ & 该写法要求 $H$ 列满秩。数值实现优先解最小二乘问题；无偏与协方差结论还要求相应噪声假设。回归矩阵自身有噪声时，普通最小二乘可能有偏。\\
$J^+=J^{\mathsf T}(JJ^{\mathsf T})^{-1}$ & 只适用于满行秩的右伪逆公式。关节与任务尺度改变时，最小范数的物理含义也改变。对一般秩使用 SVD 等合适方法。\\
$\dot q=J^+v+N\dot q_0$ & 欧氏正交投影用 $N=I-J^+J$。这是瞬时速度关系；有限步积分会产生漂移，且它本身不满足限位、碰撞和执行器约束。\\
阻尼伪逆产生有限输出 & 正则化参数为正时改善病态求解，但一般有 $J(I-J_\lambda^\#J)\ne0$；不能再声称次任务完全不影响主任务。\\
$\Lambda=(JM^{-1}J^{\mathsf T})^{-1}$ & 要求 $M$ 正定、所选 $J$ 满行秩。任务接近奇异时有效惯量会变大；六维混合单位需先定尺度。\\
$\tau_{\rm null}=N^{\mathsf T}\tau_0$ & 这里 $N=I-\bar JJ$、$\bar J=M^{-1}J^{\mathsf T}\Lambda$。严格瞬时解耦的代数条件为 $JM^{-1}N^{\mathsf T}=0$；模型误差与力矩限幅会改变实际结果。\\
$D=2\sqrt{MK}$ & 对应常系数、单自由度、$M>0,K>0$ 的标准二阶临界阻尼条件。一般矩阵不能逐元素开平方就得到多轴临界阻尼。\\
$\dot V\le0$ 推导稳定性 & 储能不增并不自动意味着误差或速度趋零。渐近结论还需正定性、有界性及不变集等条件；半正定阻尼允许存在未耗散方向。\\
$A_d=e^{AT_s}$ 与 ZOH 输入积分 & 对固定线性模型、固定采样周期和区间内输入保持不变是精确关系。对非线性或时变模型局部冻结矩阵时，只是相应近似。\\
$f_s>2f_{\max}$ & 是理想带限信号的无混叠采样条件，实际还涉及模拟抗混叠和带外噪声。它不是“闭环采样率只需两倍带宽”的设计准则。\\
五次最小 jerk 轨迹 & 固定时长、给定位置速度加速度的端点条件下，最小化平方 jerk 积分可得到五次多项式；关节界限、障碍和力矩限制须另外检查。\\
S 曲线时间公式 & 必须先判断是否达到加速度上限、速度上限及是否存在巡航段。不同分支的持续时间公式不能混用，负的分段时间表示分支选择不成立。\\
残差换算为末端外力 & 需要已知接触点与力的结构、合适的雅可比秩，以及足够准确的动力学残差。摩擦、标定和未建模力会与真实接触作用混在一起。\\
刚度静态偏差 $e=K^{-1}F_r$ & 需静态平衡、明确的残差方向和可逆任务刚度。它不能直接预测撞击峰值，也不能证明带延迟接触闭环稳定。\\
\end{longtable}
\endgroup

若检查失败，先保留一组状态、输入、单位和预期数字，让正解、雅可比、力矩、功率与回代相互一致，再扩大到整条轨迹。

\chapter{配套程序与复现说明}
\label{app:reproduce}

本书同时提供可编辑的 \LaTeX{} 正文、全部图形的生成脚本、数值校核脚本和第23章合成数据。PDF 本身包含推导、算例、练习和解答；配套程序用于把数字与曲线重新算出来，并在改变条件后观察结果。所有程序都只进行本地数学计算或作图，不连接机器人。

\section{文件怎样对应正文}

解压配套压缩包后，进入 \texttt{robotics\_textbook} 目录。正文入口为 \texttt{textbook.tex}，23章位于 \texttt{chapters/}，矢量图与 PNG 位于 \texttt{figures/}。下表按学习用途列出程序。

\begin{longtable}{>{\raggedright\arraybackslash}p{61mm}>{\raggedright\arraybackslash}p{90mm}}
\toprule 文件 & 内容与输出\\\midrule\endhead
\texttt{math\_checks.py} & 第1--5章的矩阵、几何、运动学与数值方法检查；输出 \texttt{math\_check\_results.json}\\
\texttt{dynamics\_checks.py} & 第6--10章的动力学、辨识和标定计算；输出 \texttt{dynamics\_numbers.json}\\
\texttt{control\_checks.py} & 第11--16章的反馈、离散模型、估计与操作空间核对；输出 \texttt{control\_numbers.json}\\
\texttt{math\_figures.py}\newline\texttt{dynamics\_figures.py}\newline\texttt{control\_figures.py} & 生成前三组章节的精确几何图与数值曲线\\
\texttt{root\_figures.py} & 第17--23章的轨迹、示教、时序、MPC和云台图；还输出 \texttt{root\_numeric\_checks.json}\\
\texttt{integrated\_experiment.py} & 第23章完整2R辨识及控制实验；生成数据、结果、图和教材中的指标表\\
\texttt{integration\_data.npz} & 合成训练和验证数据，不是真机日志\\
\texttt{integration\_results.json} & 参数、仿真设定、验证残差和四控制器指标\\
\texttt{README.txt} & 运行顺序、依赖和源文件说明\\
\bottomrule
\end{longtable}

\section{最小运行顺序}

数值程序使用 Python 3，以及 NumPy、SciPy、SymPy、Matplotlib。中文作图需要 Noto CJK 字体；若系统缺少该字体，可安装字体或在脚本中改成已有的中文字体。编译正文需要支持中文字体的 XeLaTeX 和完整的常用宏包。配套的 PDF 图已经生成，因此只编译正文时无需重新运行全部作图程序。

\begin{lstlisting}[language=bash]
python math_checks.py
python dynamics_checks.py
python control_checks.py
python integrated_experiment.py

# Optional: regenerate the remaining figures.
python math_figures.py
python dynamics_figures.py
python control_figures.py
python root_figures.py

# Compile twice for the contents, labels and references.
xelatex textbook.tex
xelatex textbook.tex
\end{lstlisting}

不同平台的浮点库和字体可能造成极小数字差异或分页差异。判断计算一致性时，比较合理数值容差、物理关系和曲线形状；不要用 PDF 文件字节是否完全相同来判断算法是否正确。

\section{第23章实验实际做了什么}

实验对象是二维竖直平面两连杆，使用第6章推导的显式刚体方程。程序以四阶 Runge--Kutta 推进关节状态，并在每个积分阶段使用执行器一阶方程的解析力矩值；控制命令每隔多个物理步更新一次，另保留一个控制周期的命令延迟。

\begin{enumerate}
\item 用解析激励轨迹产生 $q,\dot q,\ddot q$；以独立的对象真值生成合成力矩，再加入固定随机种子的噪声。
\item 按两个质量和四个摩擦系数构造六列回归矩阵；用最小二乘估计，不把真值直接传给辨识控制器。
\item 使用另一条轨迹和另一组噪声计算验证残差。报告包含测量噪声的误差，同时在结果文件中保留对无噪声对象的误差。
\item 对相同的起终点、三秒五次轨迹、一秒保持、反馈增益和执行器设定，分别运行四个控制器。
\item 从关节状态计算末端实际位置，统计零至四秒的误差。RMS 是二维误差向量模长的均方根，终点误差是四秒时的误差。
\end{enumerate}

本实验没有调用 MuJoCo，也没有调用 NERO 驱动。第21章的 MJCF 是用于解释结构的模型示例；软件接口的讲解和本次实际执行的仿真要分别理解。这里还没有模拟编码器噪声、真实静摩擦集合、齿隙、连杆柔性、网络随机抖动或接触。加上这些因素后，性能指标应重新计算，而不能沿用第23章的数字。

\section{怎样自己改一个条件并解释结果}

推荐一次只改一个明确因素。例如先保持模型和增益不变，把控制周期从 $5\,\mathrm{ms}$ 改为 $10\,\mathrm{ms}$，同时确保物理步仍为 $0.5\,\mathrm{ms}$，则更新步数应从10变为20。若仍保留“一控制周期延迟”，通信延迟也同时从5毫秒变为10毫秒；这实际上改变了两个物理量。为了只研究采样率，应把延迟单独固定，并明确新的队列实现。

再如只把重力补偿的质量改错，并比较静态误差与 $K^{-1}F_r$ 的近似；只加大工具质量，检查力矩余量和轨迹误差；只改变零空间增益，观察主任务是否因限幅或正则化发生耦合。每次先写预测，再运行，再用方程解释差异。

最后，为每张新增曲线写下对象参数、控制律、初始状态、采样周期、积分步长、时延、限幅、评价区间和单位。这样的记录能让一次仿真成为可以讨论和复核的证据。

\section{一个可独立运行的2R核对片段}

下面的最小代码复现名义模型在 $q=(0,\pi/2)$ 的 $J,M,g$，并验证速度和力映射的功率相等。它不依赖任何机器人工具箱。

\begingroup
\setstretch{1.0}
\begin{lstlisting}[language=Python]
import numpy as np

q1, q2 = 0.0, np.pi / 2
l1 = l2 = 0.5
c1 = c2 = 0.25
m1 = m2 = 1.0
I1 = I2 = 1.0 / 48.0
g0 = 9.81

J = np.array([
    [-l1*np.sin(q1)-l2*np.sin(q1+q2),
     -l2*np.sin(q1+q2)],
    [ l1*np.cos(q1)+l2*np.cos(q1+q2),
      l2*np.cos(q1+q2)]
])
a = I1+I2+m1*c1**2+m2*(l1**2+c2**2)
b = m2*l1*c2
d = I2+m2*c2**2
M = np.array([
    [a+2*b*np.cos(q2), d+b*np.cos(q2)],
    [d+b*np.cos(q2), d]
])
g = g0*np.array([
    (m1*c1+m2*l1)*np.cos(q1)
    + m2*c2*np.cos(q1+q2),
    m2*c2*np.cos(q1+q2)
])

v = np.array([0.0, 0.1])
f = np.array([2.0, -3.0])
qd = np.linalg.solve(J, v)
tau_ext = J.T @ f
assert np.allclose(qd, [0.2, -0.2])
assert np.allclose(tau_ext, [-2.5, -1.0])
assert np.allclose(qd @ tau_ext, v @ f)
assert np.all(np.linalg.eigvalsh(M) > 0)
print(J, M, g, qd @ tau_ext, sep="\n")
\end{lstlisting}
\endgroup

这一检查只证明这些量在本模型和本姿态下自洽。它不会替代完整参数核验，也不能证明控制器在任意轨迹、延迟和接触条件下稳定。应把每个检查的结论限定在实际覆盖的范围。

\backmatter
% 已核对的参考资料；网页访问日期统一为 2026-11-23。
% 使用 \url 而非不可断行的 \texttt，便于在正文宽度内排版。
\begingroup
\renewcommand{\bibname}{参考资料与进一步阅读}
\begin{thebibliography}{99}
\addcontentsline{toc}{chapter}{参考资料与进一步阅读}
\small
\setlength{\itemsep}{9pt}
\setlength{\parskip}{2pt}
\raggedright
\urlstyle{same}

\bibitem[A]{sourceA}
《机械臂与机器人控制算法：从零开始的基础巩固手册》。用户提供的 PDF，共 90 页；本书简称原稿 A。各章所列的 A 页码按上传 PDF 的页面顺序计。原文件未给出可核实的正式作者与出版信息，本文不补写。

\bibitem[B]{sourceB}
《机械臂与仿真工程教材（数学及格也能懂版）》。用户提供的 PDF，共 129 页；本书简称原稿 B。各章所列的 B 页码按上传 PDF 的页面顺序计。原文件未给出可核实的正式作者与出版信息，本文不补写。

\bibitem{khatib}
Oussama Khatib.
A Unified Approach for Motion and Force Control of Robot Manipulators: The Operational Space Formulation.
IEEE Journal of Robotics and Automation, RA-3(1):43--53, 1987。
作者所在研究组提供的论文原文：
\url{https://khatib.stanford.edu/publications/pdfs/Khatib_1987_RA.pdf}。
访问日期：2026-11-23。用于核对操作空间动力学、任务惯量与冗余控制的来源。

\bibitem{ruckig}
Lars Berscheid and Torsten Kröger.
Jerk-limited Real-time Trajectory Generation with Arbitrary Target States.
Proceedings of Robotics: Science and Systems XVII, 2021。
DOI：\url{https://doi.org/10.15607/RSS.2021.XVII.015}。
会议官方条目：\url{https://www.roboticsproceedings.org/rss17/p015.html}。
Ruckig 官方文档：\url{https://docs.ruckig.com/}。
访问日期：2026-11-23。正文中的软件能力范围以该文档所区分的状态到状态运动和中间航点情形为准。

\bibitem{toppra}
Hung Pham and Quang-Cuong Pham.
A New Approach to Time-Optimal Path Parameterization Based on Reachability Analysis.
IEEE Transactions on Robotics, 34(3):645--659, 2018。
DOI：\url{https://doi.org/10.1109/TRO.2018.2819195}。
作者维护的 TOPP-RA 文档：\url{https://hungpham2511.github.io/toppra/}。
访问日期：2026-11-23。用于核对固定几何路径的时间参数化与可达性分析。

\bibitem{lavalle}
Steven M. LaValle.
Planning Algorithms.
Cambridge University Press, 2006。
作者提供的完整电子版及勘误：\url{https://lavalle.pl/planning/}。
访问日期：2026-11-23。第 4--5 章可继续学习构型空间、采样规划与概率完备性；第 14 章进一步讨论微分约束下的规划。

\bibitem{niemeyer}
Günter Niemeyer and Jean-Jacques E. Slotine.
Stable Adaptive Teleoperation.
IEEE Journal of Oceanic Engineering, 16(1):152--162, 1991。
DOI：\url{https://doi.org/10.1109/48.64895}。
出版方条目：\url{https://ieeexplore.ieee.org/document/64895/}。
访问日期：2026-11-23。用于核对波变量、时延通道与无源性的理论来源。

\bibitem{corke}
Peter Corke.
Robotics, Vision and Control: Fundamental Algorithms in MATLAB.
第 2 版，Springer, 2017。
出版方书目信息：\url{https://doi.org/10.1007/978-3-319-54413-7}。
作者的视觉伺服配套资源：
\url{https://petercorke.com/rvc2/v-visual-servoing/}。
访问日期：2026-11-23。第 15--16 章对应图像与位姿视觉伺服。

\bibitem{moveit}
MoveIt 项目。
Realtime Arm Servoing.
MoveIt 官方文档，ROS 2 Humble 分支。
\url{https://moveit.picknik.ai/humble/doc/examples/realtime_servo/realtime_servo_tutorial.html}。
访问日期：2026-11-23。该链接对应正文讨论的 Humble 版本；其他版本的接口应查看其各自文档。

\bibitem{ros2control}
ros2\_control 项目。
Controller Manager.
ros2\_control 官方文档，Humble 分支。
\url{https://control.ros.org/humble/doc/ros2_control/controller_manager/doc/userdoc.html}。
访问日期：2026-11-23。用于核对控制循环、生命周期、更新率与抖动要求。

\bibitem{jointtrajectory}
ros2\_controllers 项目。
joint\_trajectory\_controller.
官方文档，Humble 分支。
\url{https://control.ros.org/humble/doc/ros2_controllers/joint_trajectory_controller/doc/userdoc.html}。
访问日期：2026-11-23。用于核对轨迹插补、命令与状态接口，以及各接口下的 PID 作用方式。

\bibitem{mujocosim}
MuJoCo 项目。
Simulation.
MuJoCo 官方文档，stable 分支。
\url{https://mujoco.readthedocs.io/en/stable/programming/simulation.html}。
访问日期：2026-11-23。用于核对模型与数据对象、正向计算、逆动力学和仿真步进的关系。

\bibitem{mujocomodel}
MuJoCo 项目。
Modeling；XML Reference.
MuJoCo 官方文档，stable 分支。
\url{https://mujoco.readthedocs.io/en/stable/modeling.html}；
\url{https://mujoco.readthedocs.io/en/stable/XMLreference.html}。
访问日期：2026-11-23。模型指南解释机构树和建模方式；XML 参考手册定义具体字段的含义。

\bibitem{mujocomputation}
MuJoCo 项目。
Computation.
MuJoCo 官方文档，stable 分支。
\url{https://mujoco.readthedocs.io/en/stable/computation/index.html}。
另见 Modeling 的 Solver parameters 小节：
\url{https://mujoco.readthedocs.io/en/stable/modeling.html#solver-parameters}。
访问日期：2026-11-23。用于核对软接触、约束求解以及 \texttt{solref}、\texttt{solimp} 的参数语义。

\bibitem{sutton}
Richard S. Sutton and Andrew G. Barto.
Reinforcement Learning: An Introduction.
第 2 版，The MIT Press, 2018。ISBN：9780262039246。
出版社官方书目及资源入口：
\url{https://mitpress.mit.edu/9780262039246/reinforcement-learning/}。
访问日期：2026-11-23。用于核对回报、折扣、策略与价值函数的基础定义；本书的机器人学习例子并不替代该书对学习算法的系统展开。

\bibitem{modernrobotics}
Kevin M. Lynch and Frank C. Park.
Modern Robotics: Mechanics, Planning, and Control.
Cambridge University Press, 2017。ISBN：9781107156302。
作者官方教材主页：
\url{https://hades.mech.northwestern.edu/index.php/Modern_Robotics}。
作者提供的更新版第一版预印本（2019 年 12 月）：
\url{https://hades.mech.northwestern.edu/images/2/25/MR-v2.pdf}。
官方章节视频：\url{https://modernrobotics.northwestern.edu/}。
访问日期：2026-11-23。
这是贯通本书数学、刚体运动、运动学、动力学、轨迹与控制的进一步阅读；其中第 3、5、8、9、11 章分别对应刚体运动、速度与力、动力学、轨迹生成及机器人控制。预印本文件名中的 v2 表示更新版本，不应误写成正式的第 2 版教材。

\end{thebibliography}
\endgroup

\end{document}
